%\RequirePackage[notext,nomath]{stix}
\documentclass[12pt]{amsart}
\usepackage{amsmath,amsthm,amssymb}
\newcommand\hmmax{0}
\newcommand\bmmax{0}
\usepackage{dsfont}
\usepackage{mathrsfs} %use \mathscr
\usepackage[all,cmtip]{xy}
\usepackage{patchcmd}
\input{diagxy}
\usepackage{hyperref}
\usepackage{bm}
\usepackage{mathtools}
\usepackage{enumerate}
\usepackage{stmaryrd}
%\usepackage{color}
\usepackage[usenames,dvipsnames,svgnames,table]{xcolor}
\usepackage[normalem]{ulem}
\usepackage[top=2.3cm, bottom=2.3cm, left=2cm, right=2cm]{geometry}
%\usepackage{showkeys}
\allowdisplaybreaks

%MYlogic

\newcommand{\memph}[1]{{\color{magenta}\emph{#1}}}

%table of contents------------------------------------------------------------------
\makeatletter
\patchcommand\@starttoc{\begin{quote}}{\end{quote}}
\def\@tocline#1#2#3#4#5#6#7{\relax
  \ifnum #1>\c@tocdepth % then omit
  \else
    \par \addpenalty\@secpenalty\addvspace{#2}%
    \begingroup \hyphenpenalty\@M
    \@ifempty{#4}{%
      \@tempdima\csname r@tocindent\number#1\endcsname\relax
    }{%
      \@tempdima#4\relax
    }%
    \parindent\z@ \leftskip#3\relax \advance\leftskip\@tempdima\relax
    \rightskip\@pnumwidth plus4em \parfillskip-\@pnumwidth
    #5\leavevmode\hskip-\@tempdima
      \ifcase #1
       \or\or \hskip 1em \or \hskip 2em \else \hskip 3em \fi%
      #6\nobreak\relax
    \dotfill\hbox to\@pnumwidth{\@tocpagenum{#7}}\par
    \nobreak
    \endgroup
  \fi}
\makeatother

\newcommand{\quotnum}[1]{ \vspace{.1in} \refstepcounter{equation}%
\noindent (\thesection.\arabic{equation}) \hspace{.6cm} \pbox{0.85\textwidth}{#1} \vspace{.1in}}

% THEOREM Environments ---------------------------------------------------
 \theoremstyle{plain}
 \newtheorem{thm}{Theorem}[section]
 \newtheorem{cor}[thm]{Corollary}
 \newtheorem{lem}[thm]{Lemma}
 \newtheorem{fact}[thm]{Fact}
 \newtheorem{prop}[thm]{Proposition}
 \newtheorem{cla}[thm]{Claim}
 \newtheorem{conj}[thm]{Conjecture}
\theoremstyle{definition}
 \newtheorem{defn}[thm]{Definition}
 \newtheorem{test}[thm]{Test}
 \newtheorem{hyp}[thm]{Hypothesis}
 \newtheorem{crit}[thm]{Criterion}
 \newtheorem{axio}[thm]{Axiom}
\theoremstyle{remark}
 \newtheorem{rem}[thm]{Remark}
 \newtheorem{ques}[thm]{Question}
 \newtheorem{quess}[thm]{Questions}
 \newtheorem{ter}[thm]{Terminology}
 \newtheorem{nota}[thm]{Notation}
 \newtheorem{conv}[thm]{Convention}
 \newtheorem{exam}[thm]{Example}

 \numberwithin{equation}{section}
 \renewcommand{\theequation}{\thesection.\arabic{equation}}

\newtheorem*{incla}{Claim}
\newtheorem*{sta}{\textbf{\textup{Claim}} $(\star)$}

\theoremstyle{plain}
\newtheorem*{theo}{Main Theorem}
\newtheorem*{thm*}{Theorem}

% MATH objects ---------------------------------------------------
\newcommand{\VF}{\textup{VF}}
\newcommand{\ACVF}{\textup{ACVF}}
\newcommand{\RV}{\textup{RV}}
\newcommand{\DC}{\textup{DC}}
\newcommand{\MM}{\mdl M}
\newcommand{\LL}{\mathcal{L}}
\newcommand{\OO}{\mdl O}
\newcommand{\II}{\mathcal{I}}
\newcommand{\CC}{\mathcal{C}}
\newcommand{\UU}{\mathfrak{U}}
% MATH operators---------------------------------------------------
\DeclareMathOperator{\rk}{rk}
\DeclareMathOperator{\otp}{otp}
 \DeclareMathOperator{\ann}{Ann}
 \DeclareMathOperator{\ran}{ran}
 \DeclareMathOperator{\dom}{dom}
 \DeclareMathOperator{\cod}{cod}
 \DeclareMathOperator{\id}{id}
 \DeclareMathOperator{\ad}{AD}
 \DeclareMathOperator{\ed}{ED}
 \DeclareMathOperator{\cd}{CD}
 \DeclareMathOperator{\Th}{Th}
 \DeclareMathOperator{\tp}{tp}
 \DeclareMathOperator{\Hom}{Hom}
 \DeclareMathOperator{\im}{im}
 \DeclareMathOperator{\lh}{lh}
 \DeclareMathOperator{\lcm}{lcm}
 \DeclareMathOperator{\spec}{Spec}
 \DeclareMathOperator{\ck}{coker}
 \DeclareMathOperator{\aut}{Aut}
 \DeclareMathOperator{\ext}{Ext}
 \DeclareMathOperator{\maxlim}{max\,lim}
 \DeclareMathOperator{\td}{tr\,deg}
 \DeclareMathOperator{\tc}{tr\,cl}
 \DeclareMathOperator{\ord}{Ord}
 \DeclareMathOperator{\cha}{char}
 \DeclareMathOperator{\ac}{\overline{ac}}
 \DeclareMathOperator{\supp}{supp}
 \DeclareMathOperator{\quot}{Quot}
 \DeclareMathOperator{\acl}{acl}
 \DeclareMathOperator{\dcl}{dcl}
 \DeclareMathOperator{\pr}{pr}
 \DeclareMathOperator{\alg}{ac}
 \DeclareMathOperator{\avg}{avg}
 \DeclareMathOperator{\tr}{tr}
 \DeclareMathOperator{\mgl}{GL}
\DeclareMathOperator{\msl}{SL}
\DeclareMathOperator{\jcb}{Jcb}
\DeclareMathOperator{\K}{\mathds{k}}
\DeclareMathOperator{\dvn}{div}
\DeclareMathOperator{\Tr}{Tr}
\DeclareMathOperator{\Jac}{Jac}
\DeclareMathOperator{\GL}{GL}
\DeclareMathOperator{\res}{res}  % map into the residue field

%---integrals-----------------------------------------------------------
\def\Xint#1{\mathchoice
{\XXint\displaystyle\textstyle{#1}}%
{\XXint\textstyle\scriptstyle{#1}}%
{\XXint\scriptstyle\scriptscriptstyle{#1}}%
{\XXint\scriptscriptstyle\scriptscriptstyle{#1}}%
\!\int}
\def\XXint#1#2#3{{\setbox0=\hbox{$#1{#2#3}{\int}$}
\vcenter{\hbox{$#2#3$}}\kern-.5\wd0}}
\def\ddashint{\Xint=}
\def\dashint{\Xint-}

% shorthand for fonts and compounds-------------------------------------

\newcommand{\Z}{\mathds{Z}}
\newcommand{\A}{\mathds{A}}
\newcommand{\C}{\mathds{C}}
\newcommand{\D}{\mathds{D}}
\newcommand{\F}{\mathds{F}}
\newcommand{\G}{\mathds{G}}
\newcommand{\I}{\mathds{I}}
\newcommand{\KK}{\mathds{K}}
\newcommand{\Q}{\mathds{Q}}
\newcommand{\N}{\mathds{N}}
\newcommand{\M}{\mathds{M}}
\newcommand{\R}{\mathds{R}}
\newcommand{\dS}{\mathds{S}}
\newcommand{\dO}{\mathds{O}}
\newcommand{\PP}{\mathds{P}}
\newcommand{\p}{$p$\nobreakdash}
\newcommand{\omin}{$o$\nobreakdash}
\newcommand{\bmin}{$b$\nobreakdash}
\newcommand{\cmin}{$C$\nobreakdash}
\newcommand{\vmin}{$V$\nobreakdash}
\newcommand{\T}{$T$\nobreakdash}
%\newcommand{\M}{$M$\nobreakdash}
\newcommand{\NN}{$N$\nobreakdash}
\newcommand{\aldash}{$\alpha$\nobreakdash}
\newcommand{\bedash}{$\beta$\nobreakdash}
\newcommand{\omdash}{$\omega$\nobreakdash}
\newcommand{\gadash}{$\gamma$\nobreakdash}
\newcommand{\lexp}[2]{{\vphantom{#2}}^{#1}{#2}}
\newcommand{\dand}{\quad \text{and} \quad}
\newcommand{\tand}{~\text{and}~}

%----fonts---------------------------------------------------------------

\newcommand{\gP}{\mathfrak{P}}
\newcommand{\gR}{\mathfrak{R}}
\newcommand{\gA}{\mathfrak{A}}
\newcommand{\gB}{\mathfrak{B}}
\newcommand{\gC}{\mathfrak{C}}
\newcommand{\gD}{\mathfrak{D}}
\newcommand{\gE}{\mathfrak{E}}
\newcommand{\gG}{\mathfrak{G}}
\newcommand{\gH}{\mathfrak{H}}
\newcommand{\gI}{\mathfrak{I}}
\newcommand{\gJ}{\mathfrak{J}}
\newcommand{\gM}{\mathfrak{M}}
\newcommand{\gN}{\mathfrak{N}}
\newcommand{\gO}{\mathfrak{O}}
\newcommand{\gZ}{\mathfrak{Z}}
\newcommand{\ga}{\mathfrak{a}}
\newcommand{\gb}{\mathfrak{b}}
\newcommand{\gc}{\mathfrak{c}}
\newcommand{\gd}{\mathfrak{d}}
\newcommand{\gee}{\mathfrak{e}}
\newcommand{\gf}{\mathfrak{f}}
\newcommand{\frakg}{\mathfrak{g}}
\newcommand{\gh}{\mathfrak{h}}
\newcommand{\gi}{\mathfrak{i}}
\newcommand{\gj}{\mathfrak{j}}
\newcommand{\gk}{\mathfrak{k}}
\newcommand{\gl}{\mathfrak{l}}
\newcommand{\gm}{\mathfrak{m}}
\newcommand{\gn}{\mathfrak{n}}
\newcommand{\go}{\mathfrak{o}}
\newcommand{\gp}{\mathfrak{p}}
\newcommand{\gq}{\mathfrak{q}}
\newcommand{\gr}{\mathfrak{r}}
\newcommand{\gs}{\mathfrak{s}}
\newcommand{\gt}{\mathfrak{t}}
\newcommand{\gu}{\mathfrak{u}}
\newcommand{\gv}{\mathfrak{v}}
\newcommand{\gw}{\mathfrak{w}}
\newcommand{\gx}{\mathfrak{x}}
\newcommand{\gy}{\mathfrak{y}}
\newcommand{\gz}{\mathfrak{z}}

\newcommand{\0}{\emptyset}

%-------------------------------

\DeclareMathAlphabet{\mathpzc}{OT1}{pzc}{m}{it}

\newcommand{\disa}{\mathpzc a}
\newcommand{\disb}{\mathpzc b}
\newcommand{\disc}{\mathpzc c}
\newcommand{\disd}{\mathpzc d}
\newcommand{\dise}{\mathpzc e}
\newcommand{\dish}{\mathpzc h}
\newcommand{\diso}{\mathpzc o}
\newcommand{\disr}{\mathpzc r}
\newcommand{\diss}{\mathpzc s}
\newcommand{\pM}{\mathpzc M}
\newcommand{\pF}{\mathpzc F}
\newcommand{\pL}{\mathpzc L}

% double bars etc.-----------------------
\newcommand{\dbra}[1]{% \llrrbracket{..}
  \,[\mkern-6.8mu[\, #1 \,]\mkern-6.8mu]}
\newcommand{\dpar}[1]{% \llrrparen{..}
   (\mkern-4mu (#1 )\mkern-4mu )}

% double bars etc.-----------------------
 \DeclarePairedDelimiter\abs{\lvert}{\rvert}
%\newcommand{\abs}[1]{\left\vert#1\right\vert}
 \newcommand{\ceil}[1]{\left\lceil#1\right\rceil}
% \newcommand{\set}[1]{\left\{#1\right\}}
\newcommand{\seq}[1]{\left<#1\right>}
 \DeclarePairedDelimiter\norm{\lVert}{\rVert}
 \newcommand{\wt}[1]{\widetilde{#1}}
 \newcommand{\wh}[1]{\widehat{#1}}
 \newcommand{\usub}[2]{#1_{\textup{#2}}}
 \newcommand{\dhat}[1]{\Hat{\Hat{#1}}}
 \newcommand{\dlbr}{( \! (}
 \newcommand{\drbr}{) \! )}
 \newcommand{\btd}{\bigtriangledown}
 \newcommand{\btu}{\bigtriangleup}
 \newcommand{\lan}[3]{\mathcal{L}_{#1 \textup{#2} #3}}

% special fonts for theory names, functions, axioms, etc.
\newcommand{\mdl}[1]{\mathcal{#1}}  % model; e.g. M
\newcommand{\bo}[1]{\mathbf{#1}}    % boldface;
\newcommand{\bb}[1]{\mathbb{#1}}
\newcommand{\comp}{\thicksim}


% logical connectives
\newcommand{\limplies}{\rightarrow}
\newcommand{\lfollows}{\rightarrow}
\newcommand{\liff}{\leftrightarrow}
\newcommand{\ex}[1]{\exists #1 \;} % exists x ...
\newcommand{\fa}[1]{\forall #1 \;} % forall x ...
\newcommand{\lImplies}{\Rightarrow}
\newcommand{\lFollows}{\Leftarrow}
\newcommand{\lIff}{\Longleftrightarrow}

\newcommand{\inot}{\mathord{\sim}}
\newcommand{\lam}[1]{\lambda #1 \;} % lambda x ...

% useful symbols
\newcommand{\proves}{\vdash}
\newcommand{\forces}{\Vdash}
\newcommand{\nforces}{\nVdash}

\newcommand{\wtbo}[1]{\widetilde{\mathbf{#1}}}
\newcommand{\rest}{\upharpoonright}
\newcommand{\lbar}{\vec}
\newcommand{\ulin}{\underline}
\newcommand{\fun}{\longrightarrow}
\newcommand{\efun}{\longmapsto}
\newcommand{\sub}{\subseteq}
\newcommand{\esub}{\preceq}
\newcommand{\mi}{\smallsetminus}
\newcommand{\xfun}{\xrightarrow}
\newcommand{\dbar}[1]{\bar{\bar{#1}}}
\newcommand{\ddlim}[1]{\underset{#1}{\varprojlim} \,\,}
\newcommand{\dclim}[1]{\underset{#1}{\varinjlim} \,\,}
\newcommand{\dlim}{\varprojlim}  %inverse limit or projective limit or limit
\newcommand{\clim}{\varinjlim}  %direct limit or colimit
\newcommand{\dia}{\diamond}
\newcommand{\colim}[1]{\underset{#1}{\text{colim}}\,}  %direct limit or colimit

\newcommand{\dash}{\mathalpha{\mbox{-}}} % e.g. for Sigma^1_1-AC
\newcommand{\goedel}[1]{\ulcorner #1 \urcorner}

\newcommand{\la}{\langle}
\newcommand{\ra}{\rangle}
\newcommand{\half}{\frac{1}{2}}
%\newcommand{\st}{ \; | \; } % x such that ...
%\newcommand{\res}{\upharpoonright}
\newcommand{\concat}{\mathord{\hat{\;}}}

%set-------------------------------------------------
% just to make sure it exists
\providecommand\given{}
% can be useful to refer to this outside \set
\newcommand\SetSymbol[1][]{%
\nonscript \: #1 \vert
\allowbreak
\nonscript\:
\mathopen{}}
\DeclarePairedDelimiterX\set[1]\{\}{%
\renewcommand\given{\SetSymbol[\delimsize]}
#1
}
%-------------------------------------------------------------------
%+++ The following code is (essentially) from Sam Buss +++
% It computes the height of the ulcorner and urcorner correctly.

\newbox\gnBoxA
\newdimen\gnCornerHgt
\setbox\gnBoxA=\hbox{$\ulcorner$}
\global\gnCornerHgt=\ht\gnBoxA
\newdimen\gnArgHgt

\def\code #1{%
        \setbox\gnBoxA=\hbox{$#1$}%
        \gnArgHgt=\ht\gnBoxA%
        \ifnum \gnArgHgt<\gnCornerHgt
                \gnArgHgt=0pt%
        \else
                \advance \gnArgHgt by -\gnCornerHgt%
        \fi
        \raise\gnArgHgt\hbox{$\ulcorner$} \box\gnBoxA %
                \raise\gnArgHgt\hbox{$\urcorner$}}
%+++ End{Sam Buss} +++

%-------------------------------------------------------------------
\newcommand{\mVF}{{\mu}{\VF}}
%\DeclareMathOperator{\vVF}{_1 \negthickspace \VF}
\newcommand{\mgVF}{{\mu} {\VF}}
\newcommand{\mRV}{{\mu}{\RV}}
\newcommand{\mgRV}{{\mu}{\RV}}
\newcommand{\mG}{{\mu}{\Gamma}}
\newcommand{\RES}{\textup{RES}}
\newcommand{\mgRES}{{\mu}{\RES}}
\newcommand{\mRES}{{\mu}{\RES}}
\newcommand{\isp}{\textup{I}_{\textup{sp}}}
\newcommand{\dsp}{\textup{D}_{\textup{sp}}}
\newcommand{\misp}{\mu \isp}
\newcommand{\ispt}{\isp^{\times}}
\newcommand{\mgisp}{\mu \isp}
\newcommand{\mispdb}{\mu \dsp}
\newcommand{\mgispdb}{\mu_{\Gamma} \dsp}
\newcommand{\mgD}{{\mu}{\bb D}}
\newcommand{\mgE}{{\mu}{\bb E}}
\newcommand{\mgL}{{\mu}{\bb L}}
\newcommand{\mL}{{\mu}{\bb L}}
\newcommand{\vgRES}{{\vol_{\Gamma}}{\RES}}
\newcommand{\maRV}[1]{{\mu_{#1}}{\RV}}
\newcommand{\maG}[1]{{\mu_{#1}}{\Gamma}}
\newcommand{\maRES}[1]{{\vol_{#1}}{\RES}}

%-------------------------------------------------------------------

\DeclareMathOperator{\rv}{rv}
\DeclareMathOperator{\sn}{sn}
\DeclareMathOperator{\stan}{st}
\DeclareMathOperator{\csn}{csn}
\DeclareMathOperator{\rcsn}{\overline {csn}}
\DeclareMathOperator{\vv}{val}
\DeclareMathOperator{\orb}{orb}
\DeclareMathOperator{\vs}{vs}
\DeclareMathOperator{\mor}{Mor}
\DeclareMathOperator{\gsk}{\mathbf{K}_+}
\DeclareMathOperator{\ggk}{\mathbf{K}}
\DeclareMathOperator{\sgsk}{!\mathbf{K}_+}
\DeclareMathOperator{\sggk}{!\mathbf{K}}
\DeclareMathOperator{\ob}{Ob}
\DeclareMathOperator{\fn}{FN}
\DeclareMathOperator{\fib}{fib}
\DeclareMathOperator{\fin}{fin}
\newcommand{\vol}{\textup{vol}}
\DeclareMathOperator{\mR}{\mathsf{R}}
\DeclareMathOperator{\mC}{\mathsf{C}}
\DeclareMathOperator{\hen}{HEN}
\DeclareMathOperator{\rad}{rad}
\DeclareMathOperator{\vcr}{vcr}
\DeclareMathOperator{\vrv}{vrv}
\DeclareMathOperator{\RVH}{RVH}
\DeclareMathOperator{\MON}{MON}
\DeclareMathOperator{\can}{\mathbf{c}}
\DeclareMathOperator{\wgt}{\bm w}
\DeclareMathOperator{\pvf}{pr_{\VF}}
\DeclareMathOperator{\prv}{pr_{\RV}}
\DeclareMathOperator{\der}{d}
\DeclareMathOperator{\mds}{mod}
\DeclareMathOperator{\gal}{Gal}
\DeclareMathOperator{\ito}{int}
\DeclareMathOperator{\cl}{cl}
\DeclareMathOperator{\sym}{Sym}
\newcommand{\RCVF}{\textup{RCVF}}
\newcommand{\RCF}{\textup{RCF}}
\DeclareMathOperator{\db}{d \! b}
\DeclareMathOperator{\bdd}{b \! d}
\newcommand{\var}{\textup{Var}}
\DeclareMathOperator{\TVF}{TVF}
\DeclareMathOperator{\TRV}{TRV}
\DeclareMathOperator{\TRES}{TRES}
\DeclareMathOperator{\TCVF}{TCVF}
\DeclareMathOperator{\RVV}{RV_\infty}
%local-------------------------------------------------------------------
\newcommand{\ol}[1]{\overline{#1}}
\DeclareMathOperator{\lant}{\lan{T}{}{}}
\DeclareMathOperator{\dist}{dist}
\newcommand{\LT}{$\lan{T}{}{}$\nobreakdash}
\DeclareMathOperator{\len}{length}
\DeclareMathOperator{\gra}{Gr}
\newcommand{\pard}[1]{\tfrac{\partial}{\partial x_{#1}}}
\newcommand{\pardy}[1]{\tfrac{\partial}{\partial y_{#1}}}
\newcommand{\ddx}{\tfrac{d}{d x}}
\DeclareMathOperator{\bfk}{\bf k}
\DeclareMathOperator{\BSA}{BSA_{\R}}
\DeclareMathOperator{\puR}{\tilde \R}
\DeclareMathOperator{\puC}{\tilde \C}
\DeclareMathOperator{\abv}{\abs{\vv}}
\newcommand{\ticvf}[1]{\text{T}_{#1}\text{CVF}}
\DeclareMathOperator{\ssin}{\Yright}
\newcommand{\mzG}{{\mu_0}{\Gamma}}
\newcommand{\mzRES}{{\mu_0}{\RES}}
\newcommand{\mzRV}{{\mu_0}{\RV}}
\newcommand{\loc}{\textup{loc}}
\DeclareMathOperator{\tor}{Tor}
\DeclareMathOperator{\dsh}{\sharp\sharp}
\DeclareMathOperator{\Vol}{Vol}
\newcommand{\Var}{\textup{Var}}
\DeclareMathOperator{\Def}{Def}
\DeclareMathOperator{\sgn}{sgn}
\DeclareMathOperator{\tbk}{tbk}
\newcommand{\gdv}{\ggk^{\hat \delta}\var_{\R}}
\newcommand{\gmv}{\ggk^{\hat \mu}\var_{\C}}
\newcommand{\kuq}{{\K^{\times}} \cup \Q}
\newcommand{\ccirc}{\circ \circ}
%local-------------------------------------------------------------------


\DeclareMathOperator{\Cl}{Cl}


\begin{document}

\title[Bounded integral and motivic Milnor fiber]{Bounded integral and motivic Milnor fiber}

\author[Arthur Forey]{Arthur Forey}
\address[A. Forey]{D-MATH, ETH Z\"urich, R\"amistrasse 101, CH-8092 Z\"urich, Switzerland}
\email{arthur.forey@math.ethz.ch}

\author[Yimu Yin]{Yimu Yin}
\address[Y. Yin]{Pasadena, California}
\email{yimu.yin@hotmail.com}

\thanks{Many thanks to Ehud Hrushovski and Fran\c cois Loeser for their encouragement and enlightening discussions during the preparation of this manuscript. Thanks also to Raf Cluckers and Martin Hils for inspiring discussions and comments. A.F.\ is supported by a DFG-SNF lead agency program grant
  (grant 200020L\_175755). }

%\keywords{Motivic integration}
%\subjclass[2000]{}


\begin{abstract}
We construct a new Hrushovski-Kazhdan style motivic integral, the so-called bounded integral, that interpolates the integral with volume forms and that without volume forms. This is done within the framework of model theory of algebraically closed valued fields of equicharacteristic zero. As the main application, we rectify and refine some results of Hrushovski and Loeser on how to construct the motivic Milnor fiber without resolution of singularities.
\end{abstract}


\maketitle
\tableofcontents


%\section*{Notes}
%Many things in the notes are cut-pasted from other sources, so there are  some notational idiosyncrasies, all recorded in \S~\ref{sec:notation}, though.
%
%\S~\ref{sec:HK:main} is of course largely redundant (we want to establish an analog of \cite{hru:loe:lef} for the real algebraic case, the excessive details are intended for that paper), but just in case there are differences as to what the categories are.
%
%The main issue is the failure of (8.1.4) and (8.1.7) in \cite{hru:loe:lef}. So the homomorphism $h_m$ needs to be constructed in a different way. This is in \S~\ref{sec:spec:pui}.
%
%To make \S~7 work, however, we actually need a variation of \cite[Theorem~10.11]{hrushovski:kazhdan:integration:vf}. This is \S~\ref{sec:retract}.
%
%I also found issues concerning (8.5.4) and (8.5.5) in \cite{hru:loe:lef}. Firstly $\Upsilon$ does not vanish on $\isp$ and hence the expression in \cite[Corollary~8.5.3]{hru:loe:lef} is not well-defined. We can force it by modding out $[\A]$, but this is detrimental to our purpose. Secondly the integrand in \cite[Corollary~8.5.3]{hru:loe:lef} does not a priori produce a class in the 2bdd $\RV$-category, hence \cite[Proposition~8.5.1]{hru:loe:lef} cannot be applied directly. The modified definition of $\Upsilon$ now appears towards the end of \S~\ref{sec:retract}.  To solve the second issue, a new integral for the $\VF$-category of invariant objects and covariant morphisms is constructed. This is \S~\ref{sec:inv:cov}.
%
%The work in \S~\ref{sec:inv:cov} is a variation on \cite[Proposition~3.2.2]{hru:loe:lef}. The theorem there can be established, but the length of the construction will  not be desirable. Maybe there are other routes that are more direct for our purpose here? See the last few remarks in \S~\ref{sec:milnor}.

\pagebreak

\section{Introduction}
Recent years have seen significant development in applying the Hrushovski-Kazhdan integration theory to the  study of a variety of topics in geometry and topology, especially those around Denef-Loeser's motivic Milnor fiber. Such an approach  was first articulated in  \cite{hru:loe:lef} for the purpose of  finding a more conceptual (read ``resolution-free'') construction of the complex  motivic Milnor fiber, among other things; unfortunately, it contains a substantial technical oversight. The main goal of this paper is to present a new Hrushovski-Kazhdan style integral,  the so-called \memph{bounded integral}, that interpolates (or refines) the integral with volume forms and that without volume forms, and thereby fully rectify the results in  \cite{hru:loe:lef}.
The overhaul taken up here constitutes a comprehensive expansion of the theoretical framework of the Hrushovski-Kazhdan style integration that should prove amenable to new applications in the future.


We now explain our construction in more detail.
%The Grothendieck semigroup $\gsk \mdl C$ of a category $\mdl C$ endowed with a binary operation ``$\mi$'' and a binary relation ``$\subseteq$'' (subobject) is the free semigroup generated by the isomorphism classes of $\mdl C$, subject to the usual scissor relation
%$[A \mi B] + [B] = [A]$, when $B\subseteq A$,
%where $[A]$, $[B]$ denote the isomorphism classes of the objects $A$, $B$, and ``$\mi$'' and ``$\subseteq$'' are compatible with isomorphisms between objects, usually they are just set subtraction and set inclusion. Sometimes $\mdl C$ is also equipped with a binary operation (for example, cartesian product) that induces multiplication in $\gsk \mdl C$, in which case $\gsk \mdl C$ becomes a (commutative) semiring. The formal groupification of  $\gsk \mdl C$, which is then a ring, is denoted by $\ggk \mdl C$. If a group $G$ acts on the objects of $\mdl C$ and the morphisms of $\mdl C$ are $G$-equivariant, that is, they commute with $G$-actions, then the corresponding $G$-equivariant Grothendieck ring is denoted by $\ggk^G \mdl C$. If $G = \lim_n G_n$ is profinite then we shall always impose the condition that a $G$-action factor through some $G_n$-action.
In this introduction, for concreteness, we work in  the field $\puC = \bigcup_{m \in \Z^+} \C \dpar{ t^{1 / m} }$ of complex Puiseux series. This field is the algebraic closure of the field $\C \dpar t$ of complex Laurent series. A typical element  takes the form $x = \sum_{n \in \Z} a_n t^{n/m}$ for some $m \in \Z^+$ such that its support $\supp(x) = \set{n/m \in \Q \given a_n \neq 0 }$ is well-ordered, in other words, there is a $q \in \Q$ such that $a_n = 0$ for all $n/m < q$. We think of $\bfk \coloneqq \C$ as a subfield of $\puC$ via the embedding $a \efun at^0$. The map $\vv: \tilde \C^\times \fun \Q$ given by $x \efun \min\supp(x)$ is indeed a valuation. With the convention $\vv(0) = \infty$, its valuation ring $\OO$ consists of those series $x$ with $\vv(x) \geq 0$ and its maximal ideal $\MM$ of those series $x$ with $\vv(x) > 0$. Its residue field $\K$ admits a section onto $\bfk$ and hence is isomorphic to $\C$. For a nonzero series $x = \sum_{n \in \Z} a_n t^{n/m} \in \puC$ with $\vv(x) = p/m$, let $\rv(x) = a_p t^{p/m}$, which is called the \memph{leading term} of $x$.


It is well-known that $(\puC, \OO)$ is an algebraically closed valued field of equicharacteristic zero. Thus, we may  think of it as an $\lan{}{RV}{}$-model of the first-order theory $\ACVF(0,0)$ of such fields, with the parameter space (the ``ground field'') $\bb S = \C \dpar t$.  Here  the two sorts $\VF$, $\RV$ of the language $\lan{}{RV}{}$ are  interpreted, respectively, as $\puC$, $\tilde \C^\times / (1 + \MM)$ (or, equivalently, the set of leading terms) and the cross-sort function $\rv : \VF^\times \fun \RV$ as the quotient map (or the leading term map described above). The obvious epimorphism from $\RV$ onto the value group $\Gamma = \Q$, also referred to as the $\Gamma$-sort, with the kernel $\K^\times$, is denoted by $\vrv$. See \cite[\S~2]{Yin:special:trans} for the precise definitions. The Galois group $\gal(\puC / \C \dpar t)$, which  is also the group of automorphisms over $\C \dpar t$ of $\puC$ as a model of $\ACVF(0,0)$, may be identified with the procyclic group $\hat \mu$ of roots of unity (the limit of the inverse system of groups $\mu_n$ of $n$th roots of unity).


The category $\VF_*$ essentially consists of the definable subsets of $\VF^n$, $n \geq 0$, as objects and the definable bijections between them as morphisms. The objects of the category $\RV[k]$ are the pairs $(U, f)$, where $U$ is a definable set in a finite power of $\RV$ and $f : U \fun \RV^k$ is a definable finite-to-one function, and a morphism $(U, f) \fun (V, g)$ is a definable bijection $U \fun V$. The category  $\RV[*]$  is the coproduct of $\RV[k]$, $k \geq 0$, and hence is graded by $k$, the rank of the target of $f$ (not by the ambient dimension of $U$). The actual definitions will be recalled in \S~\ref{sec:HK:main}.


To such categories one associates Grothendieck rings $\ggk\VF_*$ and $\ggk\RV[*]$, roughly constructed as the free group generated by isomorphism classes of objects, subject to the usual scissor relation $[A \mi B] + [B] = [A]$ for any objects $B \sub A$, and the product being given by cartesian product. See  Definition~\ref{def:Kgroup} for details.


Pulling back from $\RV$ to $\VF$ along the map $\rv: \VF^\times \fun \RV$ yields a ring homomorphism between the   Grothendieck rings
\begin{equation}\label{lifting:hom}
\bb L : \ggk\RV[*] \fun \ggk\VF_*
\end{equation}
One of the  main results of \cite{hrushovski:kazhdan:integration:vf} states that $\bb L$ is surjective with the kernel $(\bm P-1)$, where $\bm P$ stands for the element $[\{1\}] - [\rv(\MM \mi 0)]$  in $\ggk \RV[1]$ (so the principal ideal $(\bm P - 1)$ is not homogeneous), and hence, inverting, we get a canonical isomorphism
\[
\int : \ggk\VF_* \fun \ggk\RV[*]/(\bm P-1).
\]

The structure of $\ggk \RV[*]$ can be significantly elucidated. To wit, it is a quotient of a tensor product of two other Grothendieck rings $\ggk \RES[*]$ and $\ggk \Gamma[*]$ by explicit coefficient-translation relations (Proposition~\ref{red:D:iso}), where  $\RES[*]$ is the full subcategory of $\RV[*]$ of twisted constructible sets in the residue field $\K$, with the same grading, and $\Gamma[*]$ is the category of definable sets in the value group $\Gamma$ (as an \omin-minimal group), graded by ambient dimension. The objects of $\RES[*]$ are twisted because the short exact sequence
\[
1 \fun \K^\times \fun \RV \fun \Gamma \fun 0
\]
does not admit a natural splitting, and $\ggk \Gamma[*]$ is not the Grothendieck ring of \omin-minimal groups because not all definable bijections are admitted as morphisms. Anyway, we have two homomorphisms $\bb E_{b}$, $\bb E_{g}$ on $\ggk \RV[*]$ that vanish on $(\bm P-1)$, reflecting the fact that there are two Euler characteristics in the $\Gamma$-sort, the bounded one and the geometric one (\S~\ref{subs:uni:retr:novol}). The first is a retraction onto the quotient $\sggk \RES$ of $\ggk \RES$ of Notation~\ref{res:note} (the gradation is forgotten). The second requires inverting the class $[\A]$ of the affine line and takes values in the degree-zero part of $\sggk \RES[*][[\A]^{-1}]$; for instance, after forgetting the gradation, $\bb E_{b}([1]) = 1$ whereas $\bb E_{g}([1]) = [\A]^{-1}$ (Remark~\ref{eb:no:eg}).
%We shall only need $\bb E_{b}$; see the left column of the diagram (\ref{diag-diamond-interpol}).


The objects of $\VF_*$, $\RV[*]$ may be equipped with $\Gamma$-volume forms, and the resulting categories are denoted by $\mgVF[*]$, $\mgRV[*]$. There is again a  canonical isomorphism
\[
\int^\mu : \ggk\mgVF[*] \fun \ggk\mgRV[*]/(\bm P),
\]
where the principal ideal $(\bm P)$ is now homogeneous. The ring $\ggk\mgRV[*]$ is likewise a quotient of a tensor product of  $\ggk \mgRES[*]$ and $\ggk \mG[*]$, which induces two graded homomorphisms $\mgE_ {b}$, $\mgE_{g}$ from $\ggk\mgRV[*]$ to $\sggk \mgRES[*]/([\A])$ that vanish on $(\bm P)$ and restrict to the quotient map on $\ggk \mgRES[*]$ (Proposition~\ref{prop:eu:retr:k}). The construction forces $[\A] = 0$, which is detrimental to our purpose; this is avoided by the doubly bounded categories below.


We think of $\int$ as the universal additive invariant (or generalized Euler characteristic) of definable sets and $\int^\mu$ as a motivic integral. They are not directly related, though: in  $\ggk\VF_*$, two objects are identified if there is a definable bijection between them, whereas in $\ggk\mgVF[*]$, two objects are identified if there is a measure-preserving bijection between two dense subobjects.


The main result of this paper is  that the two isomorphisms $\int$, $\int^\mu$ can be interpolated by a third one, the bounded integral (Theorem \ref{main:prop}):
\[
\int^\diamond : \ggk \mgVF^\diamond[*] \fun \ggk  \mgRV^{\db}[*] / J_{\mathrm{var}}.
\]
The category $\mgVF^\diamond[*]$ consists of those \memph{proper invariant} objects of $\mgVF[*]$, which roughly means a bounded definable set that is a union of open boxes of valuative radius $\gamma$, $\gamma \in \Gamma$. The category $\mgRV^{\db}[*]$ consists of those \memph{doubly bounded} objects of $\mgRV[*]$, in the sense that its image in the value group is doubly bounded. See \S~\ref{sec:inv:cov} for the precise definitions. Note that $\mgVF^{\diamond}[*]$  is also graded since, as in classical measure theory, gradation by ambient dimensions is a necessity in the presence of volume forms (a curve has no volume if considered as a subset of a surface). The homogeneous ideal $J_{\mathrm{var}}$ is generated by the relative annular relations of Notation~\ref{relative:annular:ideal}, that is, by the differences $[\bm U^{\tau, j}] - [\bm U]$ between an object $\bm U$ of $\mgRV^{\db}[*]$ and its elementary annular blowups $\bm U^{\tau, j}$ in a coordinate $j$ (Definition~\ref{defn:blowup:coa}), whose aperture $\tau$ is a definable function on $\bm U$ with values in $\RV$ and with nonnegative bounded valuation. Constant apertures give the ideal $(\bm P_\Gamma)$ of \S~\ref{sec:retract}, which is therefore contained in $J_{\mathrm{var}}$. By Lemma~\ref{relative:annular:compatibility}, the bounded retraction $\bb E^\diamond : \ggk \mgRV^{\db}[*] \fun \sggk \RES$ of \S~\ref{sec:retract} vanishes on $J_{\mathrm{var}}$; unlike $\mgE_{b}$ and $\mgE_{g}$, it does not require $[\A] = 0$.


The situation is summarized in the  commutative diagram (\ref{diag-diamond-interpol}). More precisely speaking, the bounded integral $\int^\diamond$ is a common refinement of $\int$ and $\int^\mu$: it is finer than $\int$ because it carries volume forms, and it is finer than $\int^\mu$ because it only admits those measure-preserving morphisms that are defined everywhere (as opposed to on dense subobjects).
\begin{figure}[htb]
\begin{equation}\label{diag-diamond-interpol}
\bfig
  \hSquares(0,0)/<-`->`->`->`->`<-`->/<400>[{\ggk \VF_*}`{\ggk \mgVF^\diamond[*]}`{\ggk \mgVF[*]}`{\ggk \RV[*] / (\bm P - 1)}`{\ggk \mgRV^{\db}[*] / J_{\mathrm{var}}}`{\ggk \mgRV[*] / (\bm P)}; ``\int`\int^{\diamond}`\int^\mu``]
\efig
\end{equation}
%\begin{equation}\label{diag-diamond-interpol}
%\bfig
%  \hSquares(0,0)/<-`->`->`->`->`<-`->/<400>[{\ggk \VF_*}`{\ggk \mgVF^\diamond[*]}`{\ggk \mgVF[*]}`{\ggk \RV[*] / (\bm P - 1)}`{\ggk \mgRV^{\db}[*] / J_{\mathrm{var}}}`{\ggk \mgRV[*] / (\bm P)}; ``\int`\int^{\diamond}`\int^\mu``]
%  \hSquares(0,-400)/<-`->`->`->`->`<-`./<400>[{\ggk \RV[*] / (\bm P - 1)}`{\ggk \mgRV^{\db}[*] / J_{\mathrm{var}}}`{\ggk \mgRV[*] / (\bm P)}`{\sggk \RES}`{\sggk \RES[*]}`{\sggk \mgRES[*] / ([\A])}; ``\bb E_b`\mgE^{\db}`\mgE_{b}`\phi`]
%\efig
%\end{equation}
\end{figure}
%with $\bb E^{\diamond} = \phi \circ \mgE^{\db}$,  where $\phi$ is the forgetful homomorphism, and the dotted line means ``no meaningful relation'' insofar as the results in this paper are concerned.


% \bfig
% \hSquares(0,0)/->`->`->`->`->`->`->/<400>[{\ggk \mgVF^\diamond[*]}`{\ggk  \mgRV^{\db}[*] / J_{\mathrm{var}}}`{\sggk \RES}`{\ggk \VF_*}`{\ggk \RV[*] / (\bm P - 1)}`{\sggk \RES}; \int^{\diamond}`\bb E^{\diamond}```\id`\int`\bb E_b]

%
%
%\bfig
%\node a(0,800)[\ggk {\mgVF^\diamond[*]}]
%\node b(1000,800)[\ggk {\mgRV^{\db}[*]} / J_{\mathrm{var}}]
%\node c(0,0)[\ggk \VF_*]
%\node d(1000,0)[\ggk {\RV[*]} / (\bm P - 1)]
%%\node e(2000,0)[{\ggk^{\hat \mu} \var_{\C}[[\A]^{-1}]}]
%\node f(700,400)[\ggk {\mgVF[*]}]
%\node g(1700,400)[\ggk {\mgRV[*]} / (\bm P)]
%%\node h(2700,400)[{\ggk^{\hat \mu} \var_{\C}[[\A]^{-1}][Y]_\dag}]
%\arrow|a|/->/[a`b;{\int^{\diamond}}]
%\arrow|l|/->/[a`c;]
%\arrow|l|/->/[b`d;]
%\arrow|l|/->/[b`g;]
%\arrow|l|/->/[a`f;]
%\arrow|l|/->/[c`d;{\int}]
%\arrow|l|/->/[f`g;{\int^\mu}]
%%\arrow|l|/->/[g`h;Z]
%%\arrow|l|/->/[d`e;\Theta \circ \bb E_b]
%%\arrow|b|/->/[h`e;{-\lim}]
% \efig


The general strategy for the construction of the bounded integral $\int^\diamond$ follows roughly that of the original Hrushovski-Kazhdan integral. One first defines a lifting homomorphism from $\ggk  \mgRV^{\db}[*]$ to $\ggk \mgVF^\diamond[*]$,  shows that it is surjective, then studies its kernel.
On the technical side, in order to stay within the realm of proper invariant sets at each step of the construction, one is led to the notion of a \memph{proper covariant} bijection, first formulated in \cite{hru:loe:lef}. It goes as follows (see Definition~\ref{prop:pseu} for more details). For $\alpha \in \Gamma$, let $\pi_\alpha : \VF\fun \VF/\MM_\alpha$ be the projection map;  if $\alpha\in \Gamma^n$ then $\pi_{\alpha}$ denotes the product of the maps $\pi_{\alpha_i}$. A definable bijection $f$ from $A\sub \VF^n$ onto  $B \sub \VF^m$ is said to be proper covariant if $A$ and $B$ are bounded and there are $\alpha \in \Gamma^n$, $\beta \in \Gamma^m$, and a function $f_\downarrow : \pi_{\alpha}(A) \fun \pi_{\beta}(B)$ such that $A$ is $\alpha$-invariant, $B$ is $\beta$-invariant, and  $\pi_\beta \circ f = f_\downarrow \circ \pi_\alpha$. In order to show the existence of such proper covariant bijections of a very special kind, we establish in \S~\ref{sec:cont:fib:prop} some general properties of definable continuous functions, reminiscent of those of definable continuous functions in \omin-minimal theories (see, for instance,  \cite[Chapter~6]{dries:1998}).


Next, we turn to the construction of the motivic Milnor fiber attached to a nonconstant polynomial function $f : (\C^d, 0) \fun (\C, 0)$,  where $0$ is assumed to be a singular point, that is, $\nabla f (0) = 0$.

By an (algebraic) variety over a field $\bfk$, we mean a reduced separated $\bfk$-scheme of finite type. We denote by $\Var_{\bfk}$ the category of  varieties over $\bfk$.
Let $\ggk \var_{\C}$ be the Grothendieck ring of complex algebraic varieties. An action of $\hat \mu$ on a variety is \memph{good} if it factors through some $\mu_n$ and every orbit is contained in an affine open subset. Let $\ggk^{\hat \mu} \var_{\C}$ be the Grothendieck ring of varieties with good $\hat \mu$-actions modulo the additional relation $[X \times (\A^r, \rho)] = [X \times \A^r]$, for every variety $X$ with a good $\hat \mu$-action and every linear action $\rho$ of some $\mu_n$ on $\A^r$, where the action on the factor $\A^r$ on the right is trivial (see Notation~\ref{bipo:gdv} and \cite[\S~2.4]{denefloeser:arc}).


Let $\mathscr L$ be the space of formal arcs on $\C^d$ at $0$. So each element in $\mathscr L$ is of the form $\gamma(t) = (\gamma_1(t), \ldots, \gamma_d(t))$, where $\gamma_i(t)$ is a complex formal power series and $\gamma_i(0) = 0$. For each $m \in \Z^+$, let $\mathscr L_m$ be the space of such arcs modulo $t^{m+1}$ (also referred to as ``truncated arcs''). Consider the following locally closed subset of $\mathscr L_m$:
\begin{equation*}
  \mdl X_{f,m}  = \set{\gamma(t) \in \mathscr L_m \given f(\gamma(t)) = t^m \mod t^{m+1} }.
\end{equation*}
It may be viewed in a natural way as the set of closed points of an algebraic variety over $\C$ and, as such, carries a natural $\mu_{m}$-action. The \memph{motivic zeta function} attached to $f$ is then the generating series whose coefficients are in effect the ``$\hat \mu$-equivariant motivic volumes'' of the sets of truncated arcs above:
\[
 Z_f(T) = \sum_{m \in \Z^+} [\mdl X_{f,m}] [\A]^{-md}T^m \in \gmv[[\A]^{-1}] \dbra{T}.
\]
It is  shown in \cite{denefloeser:arc, den:loe:2002} that $Z_f(T)$ is rational and belongs to $\gmv[[\A]^{-1}][T]_\dag$, where the subscript ``$\dag$'' indicates that the ring is localized by  a certain multiplicative family of elements. The \memph{motivic Milnor fiber}
\[
\mathscr S_{f} = - \lim_{T \limplies \infty} Z_{f}(T)
\]
is then extracted from this rational expression via a formal process of sending the variable $T$ to infinity (this process is also summarized in \cite[\S~8.4]{hru:loe:lef}). Of course, to justify calling $\mathscr S_{f}$ a ``Milnor so-and-so'' one needs to show, at the very least, that invariants of the \memph{topological Milnor fiber} $F_{f}$ attached to $f$ can be recovered from it.  This is indeed the case for, say, the Euler characteristic and the Hodge characteristic.

Originally,  the proofs that $Z_f(T)$ is rational and  the Euler (or Hodge) characteristics coincide both rely on resolution of singularities.  More recently, in \cite{hru:loe:lef}, these results are established by way of a more conceptual construction. To begin with, $Z_{f}(T)$ is expressed as
\begin{equation*}\label{intro:zeta:mot}
\sum_{m \in \Z^+} H_m([\mdl X_f]) T^m,
\end{equation*}
where the coefficients $H_m([\mdl X_f])$ are certain integrals of definable sets that take values in the ring $\gmv [[\A]^{-1}]$,  and the so-called \memph{nonarchimedean Milnor fiber} of $f$ is the definable set
\[
\mdl X_f = \set{x \in \MM^d \given \rv(f(x)) = \rv(t) }.
\]
Formulated in this way, the rationality of $Z_{f}(T)$ essentially follows from  certain computation rules of (convergent) geometric series. That the Euler characteristics of $\mathscr S_{f}$ and $F_{f}$ coincide follows from the fact that we can express both the Euler characteristic of each coefficient of $Z_{f}(T)$ and that of $F_{f}$ in terms of traces of the monodromy action on the cohomology groups of $F_{f}$, where the first expression relies on  the resolution-free proofs of the A'Campo-Denef-Loeser formula (this is the main point of \cite{hru:loe:lef}) and quasi-unipotence of local monodromy (see \cite[Remark~8.5.5]{hru:loe:lef}). However, technical difficulties, arising from the fact that none of the morphisms (8.1.4), (8.1.6), and (8.1.7) in \cite{hru:loe:lef} is surjective, contrary to what is claimed in \cite[Proposition~10.10(2)]{hrushovski:kazhdan:integration:vf} (see Remark~\ref{rem:HK:fail} for an explicit object outside their images), make their construction of the integral $H_m$ incomplete.


To remedy this, as an application of our main result, we complete their construction by recasting it in the formalism of the bounded integral $\int^\diamond$. More specifically, recall from \cite[\S~4.3]{hru:loe:lef} the canonical isomorphism
\[
\Theta : \sggk \RES \fun \gmv,
\]
which is constructed in Notation~\ref{bipo:gdv}.  Let $\ggk^{\hat \mu} \var_{\C}[[\A]^{-1}][T]_{\dag}$ be the localization of $\gmv[[\A]^{-1}][T]$ by the multiplicative family generated by the elements $1 - [\A]^{a}T^b$, where $a \in \Z$ and $b \in \Z^+$; it may be regarded as a subring of  $\ggk^{\hat \mu} \var_{\C}[[\A]^{-1}] \dbra T$. Then we have a  commutative diagram (\ref{diag-diamond-zeta-function}).
\begin{figure}[htb]
\begin{equation}\label{diag-diamond-zeta-function}
\bfig
  \square(0,800)/<-`->`->`<-/<1500,400>[{\ggk \VF_*}`{\ggk^{\Z} \mgVF^\diamond[*]}`{\ggk \RV[*] / (\bm P - 1)}`{\ggk^{\Z} \mgRV^{\db}[*] / J_{\mathrm{var}}}; `\int`\int^{\diamond}`]
  \square(0,400)/`->`->`/<1500,400>[{\ggk \RV[*] / (\bm P - 1)}`{\ggk^{\Z} \mgRV^{\db}[*] / J_{\mathrm{var}}}`\gmv`\gmv{[[\A]^{-1}][T]_{\dag}^{\leq 0}}; `\Theta \circ \bb E_b`Z`]
  \square(0,0)|aamb|/`->`->`=/<1500,400>[\gmv`\gmv{[[\A]^{-1}][T]_{\dag}^{\leq 0}}`\gmv{[[\A]^{-1}]}`\gmv{[[\A]^{-1}]};
  ``-\lim`]
%  \hSquares(0,400)/<-`->`->`->`->`<-`->/<400>[{\ggk \VF_*}`{\ggk^\natural \mgVF^\diamond[*]}`{\ggk^\natural \mgVF[*]}`{\ggk \RV[*] / (\bm P - 1)}`{\ggk^\natural \mgRV^{\db}[*] / J_{\mathrm{var}}}`{\ggk^\natural \mgRV^{\bdd}[*] / (\bm P)}; ``\int`\int^{\diamond}`\int^\mu``]
%  \hSquares(0,0)/<-`->`->`->`->``->/<400>[{\ggk \RV[*] / (\bm P - 1)}`{\ggk^\natural \mgRV^{\db}[*] / J_{\mathrm{var}}}`{\ggk^\natural \mgRV^{\bdd}[*] / (\bm P)}`\gmv{[[\A]^{-1}]}`\gmv{[[\A]^{-1}][T]_{\dag}}`\gmv{[[\A]^{-1}]_{\loc}[T]_{\dag}}; ``\Theta \circ \bb E_b``Z``]
%  \morphism(0,0)/<-/<2240,0>[\gmv{[[\A]^{-1}]}_{\loc}`\gmv{[[\A]^{-1}]}_{\loc};-\lim]
 \efig
\end{equation}
\end{figure}
The superscript ``$\Z$'' indicates the subrings generated by the classes of the objects whose volume form is a constant integer (Remark~\ref{dia:to:db:dag}); in the middle row it stands for the image of that subring in the quotient by $J_{\mathrm{var}}$. The integral $\int^\diamond$ maps $\ggk^{\Z} \mgVF^\diamond[*]$ into $\ggk^{\Z} \mgRV^{\db}[*] / J_{\mathrm{var}}$. The map $Z$ assigns, to each class $[\bm U] \in \ggk^{\Z} \mgRV^{\db}[*]$, a power series
\[
Z([\bm U])(T) = \sum_{m \in \Z^+} \bm H_m([\bm U]) T^m \in \ggk^{\hat \mu} \var_{\C}[[\A]^{-1}] \dbra T,
\]
that is, the motivic zeta function of $\bm U$, where $\bm H_m = \bm \eta \circ \bm h_m$ is the rectified version of $H_m$ (\S~\ref{piece:res}). Since $\bm H_m$ vanishes on the elements of $J_{\mathrm{var}}$ in its domain (Lemma~\ref{PG:van}), $Z$ is defined on the quotient. Note that $Z$ is a ring homomorphism with respect to the Hadamard product in the target. It turns out that $Z([\bm U])(T)$ is rational: it belongs to $\gmv[[\A]^{-1}][T]_{\dag}$ and is a quotient $P/Q$ with $\deg P \leq \deg Q$ (see the proof of Theorem~\ref{direct:mil}). The superscript ``$\leq 0$'' denotes the subgroup of these elements of degree at most $0$. On it, the formal process $T \limplies \infty$ mentioned above is defined and additive, and lands us in $\gmv[[\A]^{-1}]$ (\S~\ref{sec:milnor}); this is the map $-\lim$. The lower part of the diagram commutes by (the proof of) Theorem~\ref{direct:mil}: for every $x \in \ggk^{\Z} \mgRV^{\db}[*]/J_{\mathrm{var}}$, $-\lim Z(x)$ is the image of $(\Theta \circ \bb E^\diamond)(x)$, and $\bb E^{\diamond}$ is $\bb E_b$ composed with the forgetful map, by (\ref{edb:eb:com}). If we take $[\bm U] / J_{\mathrm{var}} = \int^{\diamond}([\mdl X_f])$, which is possible since $\mdl X_f$ carries the constant volume form $0$, then the end result is the motivic Milnor fiber $\mathscr S_{f}$.


Since $\mdl X_f$ also lives in $\VF_*$, the class $(\Theta \circ \bb E_b \circ \int)([\mdl X_f])$ is defined in $\gmv$, without even inverting $[\A]$, and its image in $\gmv[[\A]^{-1}]$ is $\mathscr S_{f}$ (Theorem~\ref{direct:mil} and Remark~\ref{rem-comp-milnor-fib}). This is indeed a salient point: it is a canonical lift of $\mathscr S_{f}$ to $\gmv$. The lift is not determined by $\mathscr S_f$ alone, since the natural homomorphism $\gmv \fun \gmv[[\A]^{-1}]$ is not injective (Remark~\ref{rem-comp-milnor-fib}). However, for comparison with the Denef-Loeser construction, we have to go down to  $\gmv[[\A]^{-1}]$, since that is where the coefficients of the motivic zeta function live.


The composite  $\Theta \circ \bb E_b \circ \int$ is referred to as the \memph{motivic volume operator}, usually denoted by $\Vol$, and is already used implicitly in \cite{hrushovski:kazhdan:integration:vf}. Since its reemergence in \cite{hru:loe:lef}, it has been featured in a growing list of applications, see \cite{thuong-int-conj, nicaise_ps_tropical, Nic:Pay:trop:fub, forey_virtual_2017, NS-deg-stab-rat}.

The nonarchimedean Milnor fiber $\mdl X_f$ is closely related to the analytic Milnor fiber $\mathcal{F}^{\mathrm{an}}_{f}$ introduced in \cite{nicaise:sebag:motivic:serre}. The $\puC$-rational points of $\mathcal{F}^{\mathrm{an}}_{f}$ may be understood as an object of  $\VF_*$ (but not $\mgVF^\diamond[*]$) and hence, as such, admits a motivic volume $\Vol(\mathcal{F}^{\mathrm{an}}_{f})$. By Remark~\ref{rem-comp-milnor-fib-bis}, $\Vol(\mathcal{F}^{\mathrm{an}}_{f}) = \Vol([\mdl X_f])$; this is the same canonical lift, so its image in $\gmv[[\A]^{-1}]$ is the motivic Milnor fiber $\mathscr S_{f}$. This fact is already proven  in \cite{nicaise_ps_tropical}, but without taking the $\hat \mu$-actions into account, and also in \cite{Nic:Pay:trop:fub, forey_virtual_2017}. The arguments there all rely on resolution of singularities.

The bounded integral $\int^\diamond$ is also used in \cite{FY:Mot-Int-Mil-Fib} to obtain similar results on real motivic Milnor fiber.


The paper is organized as follows. In \S~\ref{sec-tech} we recall and extend certain facts on definable sets in $\ACVF(0,0)$ that are needed later on. In particular,  some new results about continuous definable functions in this environment are obtained, which could be of independent interest, such as fusion of continuous fibers (Proposition~\ref{prop:continuity:fiberwise}) and curve selection (Lemma~\ref{curve:select}), which closely resemble their counterparts in the \omin-minimal setting. The short \S~\ref{sec:HK:main} introduces notation and results from the original Hrushovski-Kazhdan paper \cite{hrushovski:kazhdan:integration:vf}. In \S~\ref{sec:retract} we show how to replace classes of polytopes with numbers  and thereby construct a retraction map from the Grothendieck ring in $\RV$ to that of varieties. In \S~\ref{csndag} we recall from \cite{FY:Fubini} the monomial decomposition of definable sets in $\RV$ and the isomorphism $\Theta$, which are used in \S~\ref{sec:spec:pui}. Invariance and covariance are defined and investigated  in \S~\ref{sec-pscb}, where the surjectivity of the lifting homomorphism from $\ggk \mgRV^{\db}[*]$ to $\ggk \mgVF^\diamond[*]$ is also established (Corollary~\ref{L:sur:c}). Then in \S~\ref{section:int} we study its kernel and finish the construction of the bounded integral $\int^\diamond$. Finally, \S~\ref{sec:spec:pui} is devoted to the construction of the motivic zeta function and the motivic Milnor fiber, rectifying the result in \cite{hru:loe:lef}.


\section{Technicalities}
\label{sec-tech}

\subsection{Setting and notation}\label{sec:notation}

\begin{nota}[Coordinate projections]\label{indexing}
For each $n \in \N$, let $[n]$ denote the set $\{1, \ldots, n\}$. For any $E \sub [n]$, we write $\pr_E(A)$, or even $A_E$, when there is no danger of confusion, for the projection of $A$ into the coordinates contained in $E$. It is often more convenient to use simple standard descriptions as subscripts. For example, if $E$ is a singleton $\{i\}$ then we shall always write $E$ as $i$ and $\tilde E \coloneqq [n] \mi E$ as $\tilde i$; similarly, if $E = [i]$, $\set{k \given  i \leq k \leq j}$, $\set{k \given  i < k < j}$, $\{\text{all the coordinates in the sort $S$}\}$, etc., then we may write $\pr_{\leq i}$, $\pr_{[i, j]}$, $A_{(i, j)}$, $A_{S}$, etc.; in particular, we shall frequently write $A_{\VF}$ and $A_{\RV}$ for the projections of $A$ into the $\VF$-sort and $\RV$-sort coordinates.


Unless otherwise specified, by writing $a \in A$ we shall mean that $a$ is a finite tuple of elements (or ``points'') of $A$, whose length, denoted by $\lh(a)$, is not always indicated. If $a = (a_1, \ldots, a_n)$ then, for $1 \leq i < j \leq n$, following the notational scheme above, $a_i$, $a_{\tilde i}$, $a_{\leq i}$, $a_{[i, j]}$, $a_{\VF}$, etc., are shorthand for the corresponding subtuples of $a$.

We shall write $\{t\} \times A$, $\{t\} \cup A$, $A \mi \{t\}$, etc., simply as $t \times A$, $t \cup A$, $A \mi t$, etc., when it is clearly understood that $t$ is an element and hence must be interpreted as a singleton in these expressions.

For $a \in A_{\tilde E}$, the fiber $\{b : ( b, a) \in A \} \sub A_E$ over $a$ is often denoted by $A_a$. Note that, in the discussion below, the distinction between the two sets $A_a$ and $A_a \times \{a\}$ is usually immaterial and hence they may and  shall be tacitly identified. In particular, given a function $f : A \fun B$ and $b \in B$, the pullback $f^{-1}(b)$ is sometimes written as $A_b$ as well. This is a special case since functions are identified with their graphs. This notational scheme is especially useful when the function $f$ has been clearly understood in the context and hence there is no need to spell it out all the time.
\end{nota}


From here on we work in a sufficiently saturated model $\bb U$ of $\ACVF(0,0)$, together with a parameter  set $\mathbb{S}$, which is supposed to be a substructure of $\bb U$. So by a definable set  we shall mean an $\bb S$-definable set, unless indicated otherwise. For the rest of this paper, we assume that $\Gamma(\mathbb S)$ is nontrivial.

\begin{nota}
There is a special element $\infty = \rv(0)$ in the $\RV$-sort. For simplicity, we shall write $\RV$ to mean the $\RV$-sort without the element $\infty$, and $\RV_\infty$ otherwise --- that is, $\RV = \rv(\VF^\times)$ and $\RV_\infty = \rv(\VF)$ --- although the difference rarely matters (when it does we will of course provide further clarification).
\end{nota}

\begin{ter}[Sets and subsets]\label{nota:sub}
By a definable set in $\VF$ we mean a definable subset in the sort $\VF$, by which we just mean a subset of $\VF^n$ for some $n$, unless indicated otherwise; similarly for other sorts or structures  in place of $\VF$ that have been clearly understood in the context, such as $\RV_\infty$, $\K$, $\MM$, or any substructure $\bb M$ of $\bb U$. In particular,  a definable set without further qualification means a definable set in $\bb U$, that is, a  definable subset of $\VF^n \times \RV^m_\infty$ for some $n, m \in \N$.
\end{ter}

\begin{ter}[Relations as functions]
We may think of a relation $C$ between two sets $A$ and $B$, that is, a subset $C \sub A \times B$, as a function from $A$ into the powerset $\mdl P(B)$. For instance, if $A \sub \VF^n \times \RV^m$ is a definable set then it may be thought of as a \memph{definable} function from $A_{\VF}$ into  $\mdl P(\RV^m)$. Thus a definable function $f : A \fun \mdl P(B)$ is ontologically the same as the definable set $\bigcup_{a \in A} a \times f(a)$. We say that $f$ is \memph{finitary} if $f(a)$ is finite for every $a \in A$.
\end{ter}

The default topology in $\bb U$ is the valuation topology.

\begin{nota}\label{defn:disc}
A set $\gp \sub \VF^n \times \RV_\infty^m$ of the form $\prod_{i \leq n} \gb_i \times t$ is called an (\memph{open, closed, $\RV$-}) \memph{polydisc} if each $\gb_i$ is an (open, closed, $\RV$-) disc. The \memph{radius} $\rad(\gp)$ of the polydisc $\gp$ is the tuple $(\rad(\gb_1), \ldots, \rad(\gb_n))$. The open polydisc centered at $(a,t) \in \VF^n \times \RV_\infty^m$ with radius $\gamma \in \Gamma^n$ is denoted by $\go((a,t), \gamma)$; the closed polydisc with radius $\gamma \in \Gamma_\infty^n$ is denoted by $\gc((a,t), \gamma)$. Radius $\infty$ is used only for closed discs, with $\gc(a,\infty)=\{a\}$; open discs have finite valuative radius.

The subdiscs $\go(0, \gamma)$, $\gc(0, \gamma)$ of $\VF$ are more suggestively denoted by $\MM_\gamma$, $\OO_\gamma$, respectively, and $\MM_0$, $\OO_0$ are even more suggestively denoted by $\MM$, $\OO$. Write $\RV^{\ccirc}_\infty = \rv(\MM)$ and $\RV^{\ccirc} = \RV^{\ccirc}_\infty \mi \infty$. More generally, for $\gamma \in \Gamma$, denote by $\RV^{\ccirc}_\gamma$ the set $\rv(\MM_{\gamma} \mi 0)$.

An \memph{$\RV$-disc} is a fiber $\rv^{-1}(t)$ for $t\in\RV_\infty$, including $\rv^{-1}(\infty)=\{0\}$. In a mixed ambient space, an $\RV$-polydisc keeps the existing $\RV_\infty$-coordinates fixed. The \memph{$\RV$-hull} of a set $A$, denoted by $\RVH(A)$, is the union of all the $\RV$-polydiscs whose intersections with $A$ are nonempty. If $A$ equals $\RVH(A)$ then $A$ is called an \memph{$\RV$-pullback}.
\end{nota}


\begin{ter}\label{rvfiber}
A \memph{$\VF$-fiber} of a set $A$ is a set of the form $A_t$, where $t \in A_{\RV}$; in particular,  a $\VF$-fiber of a function $f : A \fun B$ is a set of the form $f_t$ for some  $t \in f_{\RV}$ (here $f$ also stands for its own graph) which is indeed (the graph of) a function.  We say that $A$ is open if every one of its $\VF$-fibers is, $f$ is continuous if every one of its $\VF$-fibers is, and so on. Composition is understood fiberwise, retaining the intermediate $\RV$-coordinates: for every fixed input, intermediate, and output $\RV$-tuple, the corresponding $\VF$-fiber maps compose continuously on their natural domain. Thus, if $f$, $g$ are continuous functions with $\ran(f)\sub\dom(g)$, the composite with these intermediate labels retained is continuous in this fiberwise sense.
\end{ter}


\begin{nota}\label{gamwhat}
Semantically we shall treat the value group $\Gamma$ as a definable sort (the $\Gamma$-sort). However, syntactically any reference to $\Gamma$ may be eliminated in the usual way and we can still work with $\lan{}{RV}{}$-formulas for the same purpose. This implies that $\gamma\in \Gamma$ is definable means that $\vv^{-1}(\gamma)$ is definable.

We shall write $\gamma^\sharp$, $\gamma \in \Gamma$, when we want to emphasize that it is the set $\vrv^{-1}(\gamma) \sub \RV$  that is being considered. More generally, if $I$ is a set in $\Gamma$ then we write $I^\sharp = \bigcup\set{\gamma^\sharp \given \gamma \in I}$. Similarly, if $U$ is a set in $\RV$ then $U^\sharp$ stands for $\bigcup\set{\rv^{-1}(t) \given t \in U}$. In particular, $\gamma^{\sharp \sharp}=\set{ x \in \VF \given \vv(x)=\gamma}$.
\end{nota}


\begin{nota}\label{jcbG:op}
Let $U \sub \RV^n \times \Gamma^m$, $V \sub \RV^{n'} \times \Gamma^{m'}$, and $C \sub U \times V$. The \memph{$\Gamma$-Jacobian} of $C$ at $((u, \alpha), (v, \beta)) \in C$, written as $\jcb_{\Gamma} ((u, \alpha), (v, \beta))$, is the element
\[
\jcb_{\Gamma} ((u, \alpha), (v, \beta)) =  - \Sigma (\vrv(u), \alpha) + \Sigma (\vrv(v), \beta),
\]
where $\Sigma (\gamma_1, \ldots, \gamma_n) = \gamma_1 + \ldots + \gamma_n$. If $C$ is the graph of a function then we just write $\jcb_{\Gamma}C(u, \alpha)$ instead of $\jcb_\Gamma((u, \alpha), C(u, \alpha))$.
\end{nota}


%\begin{conv}\label{how:gam}
%
%\end{conv}

\begin{conv}\label{topterm}
We may and shall assume that in any $\lan{}{RV}{}$-formula, every $\VF$-sort term (polynomial) occurs in the scope of an instance of the function symbol $\rv$. For example, if $f(x)$, $g(x)$ are polynomials then the formula $f(x) = g(x)$ is equivalent to $\rv(f(x) - g(x)) = \infty$. The polynomial $f(x)$ in $\rv(f(x))$ is referred to as a \memph{top term} (of the $\lan{}{RV}{}$-formula in question).
\end{conv}


\begin{nota}[The definable sort $\DC$ of discs]\label{disc:exp}
At times it will be more convenient to work in the traditional expansion $\bb U^{\textup{eq}}$ of $\bb U$ by all definable sorts. However, for our purpose, a much simpler expansion $\bb U^{\bullet}$ suffices. This expansion has only one additional sort $\DC$ that contains, as elements, all the open and closed discs. Since each point in $\VF$ may be regarded as a closed disc of valuative radius $\infty$, for convenience, we may and occasionally do think of $\VF$ as a subset of $\DC$. Heuristically, we may think of a disc that is properly contained in an $\RV$-disc as a ``thickened'' point of certain stature in $\VF$. There are two uniform cross-sort maps $\VF\times\Gamma\fun\DC$ in $\bb U^{\bullet}$, sending $(a,\gamma)$ to the open disc $\go(a,\gamma)$ and to the closed disc $\gc(a,\gamma)$, respectively. The closed-disc map extends to $\VF\times\Gamma_\infty$ by $(a,\infty)\mapsto\{a\}$. The radius is an argument of these maps, rather than an index for separately named function symbols.

The expansion $\bb U^{\bullet}$ can help reduce the technical complexity of our discussion. However, as is the case with the definable $\Gamma$-sort, it is conceptually inessential since, for the purpose of this paper, all allusions to discs as (imaginary) elements may be eliminated in favor of objects already definable in $\bb U$.

Whether parameters in $\DC$ are used or not shall be indicated explicitly, if it is necessary, by means of  introducing parameters of the form $\code \ga$ or the phrase ``in $\bb U^{\bullet}$,'' in which case the substructure $\mathbb{S}$ may contain names for discs that may or may not be definable from $\VF(\mathbb{S}) \cup \RV(\mathbb{S})$.

A disc centered at $0$ may be identified with its valuative radius together with its open/closed designation ($\Gamma$ is already treated as a definable sort). At a fixed finite radius the open and closed discs are distinct; the radius $\infty$ designates only the closed singleton disc.

For a disc $\ga \sub \VF$, the corresponding imaginary element in $\DC$ is denoted by $\code{\ga}$ when notational distinction makes the discussion more streamlined; $\code{\ga}$ may be heuristically thought of as the ``name'' of $\ga$. Conversely, a set $D \sub \DC$ is often identified with the set $\set{\ga \given \code \ga \in D}$, in which case $\bigcup D$ denotes a subset of $\VF$.
\end{nota}


\subsection{Of points, discs, and atoms}
The substructure $\bb S$ is \memph{$\VF$-generated} if it is generated solely by elements in $\VF$, similarly for being $(\VF, \RV)$-generated, and so on. If $\bb S$ is $\VF$-generated then the map $\rv$ is surjective in $\bb S$, and the standing assumption makes $\Gamma(\mathbb S)$ nontrivial. Then the model-theoretic algebraic closure $\acl \mathbb{S}$ of $\mathbb{S}$ is a model of $\ACVF(0,0)$; see, for instance, \cite[Corollary~4.18]{Yin:QE:ACVF:min}. In that case, we also know what the definable closure $\dcl \mathbb{S}$ of $\mathbb{S}$ is:

\begin{lem}\label{cut:to:hensel:substru}
Suppose that $\mathbb{S}$ is $\VF$-generated and $\Gamma(\mathbb{S})$ is nontrivial. Then $\dcl \mathbb{S} = \mathbb{S}$ if and only if the underlying valued field $(\VF(\mathbb{S}), \OO(\mathbb{S}))$ of $\mathbb{S}$ is henselian.
\end{lem}
\begin{proof}
We have $\VF(\acl \mathbb{S}) = \VF(\mathbb{S})^{\alg}$, where the latter is the field-theoretic algebraic closure of $\VF(\mathbb{S})$. Since the valued field automorphisms of $(\VF(\mathbb{S})^{\alg}, \OO(\acl \mathbb{S}))$ over $(\VF(\mathbb{S}), \OO(\mathbb{S}))$ are in one-to-one correspondence with the $\lan{}{RV}{}$-automorphisms of $\acl \mathbb{S}$ over $\mathbb{S}$ and $\VF(\mathbb{S})$ is the fixed field with respect to these automorphisms if and only if $(\VF(\mathbb{S}), \OO(\mathbb{S}))$ is henselian, we see that $\VF (\dcl \mathbb{S}) = \VF(\mathbb{S})$ if and only if $(\VF(\mathbb{S}), \OO(\mathbb{S}))$ is henselian. By quantifier elimination and routine syntactical inspection, $\RV(\dcl \mathbb{S}) = \RV(\mathbb{S}) $.  The lemma follows.
\end{proof}

We shall impose at a later stage the assumption that $\mathbb{S}$ is $\VF$-generated and (without loss of generality) definably closed,  that is, $\dcl \mathbb{S} = \mathbb{S}$.

\begin{lem}[{\cite[Lemma~4.10]{Yin:QE:ACVF:min}}]\label{finite:VF:project:RV}
For any finite definable set $A$ in $\VF$ there is a definable bijection $f : A \fun U$ for some set $U$ in $\RV$.
\end{lem}

The phenomenon exemplified by this simple lemma is dubbed ``internalizing finite sets.'' Its failure in positive residue characteristics belies the deep obstacles one faces in developing a Hrushovski-Kazhdan integration theory therein; see the discussion around \cite[Lemma~3.9]{hrushovski:kazhdan:integration:vf}.

\begin{lem}[{\cite[Lemma~3.3]{Yin:special:trans}}]\label{dcl:to:ac}
For $a, b \in \VF$ and $t \in \RV$, if $b$ is algebraic over $(a, t)$ then it is algebraic over $a$.
\end{lem}

\begin{lem}[{\cite[Lemma~4.3]{Yin:QE:ACVF:min}}]\label{exchange}
The exchange principle holds for single elements in both sorts, with all algebraic closures taken over the same fixed parameter base:
\begin{itemize}
 \item For any $a$, $b \in \VF$, if $a \in \acl(b) \mi \acl(\0)$ then $b \in \acl(a)$.
 \item For any $t$, $s \in \RV$, if $t \in \acl(s) \mi \acl(\0)$ then $s \in \acl(t)$.
\end{itemize}
\end{lem}


\begin{lem}\label{finite:fib:def}
Let $f : A \fun \mdl P(B)$ be a definable finitary function, where $A$, $B$ are subsets of $\VF$. Then there is a finite definable subset $B' \sub B$ such that, for every $b \in B$, the fiber $A_b = \set{a \in A \given b \in f(a)}$ is infinite if and only if  $b \in B'$. The same holds if $A$, $B$ are subsets of $\RV$.
\end{lem}
\begin{proof}
Let $b \in B$ such that $A_b$ is infinite. Then, by compactness, there is an element $a \in A_b \mi \acl(b)$. Since $f$ is finitary, which means $b \in \acl(a)$, we must have, by Lemma~\ref{exchange},  $b \notin \acl(a) \mi \acl(\0)$ and hence $b \in \acl(\0)$. It follows that there is an algebraic subset $B' \sub B$ such that if $A_b$ is infinite then $b \in B'$, for otherwise, by compactness again, there would be a $b \in B$ that lies outside every algebraic subset of $B$ and yet the size of $A_b$ is not bounded by any natural number. Clearly we can adjust $B'$ so that it contains exactly those $b \in  B$ with $A_b$ infinite. The second case is similar.
\end{proof}


\begin{defn}
We set $\dim_{\VF}(\varnothing)=-\infty$. For a nonempty definable set $A$, the \memph{$\VF$-dimension},  denoted by $\dim_{\VF}(A)$, is the largest natural number $k$ such that, possibly after re-indexing of the $\VF$-coordinates, $\pr_{\leq k}(A_t)$ has nonempty interior for some $t \in A_{\RV}$.
\end{defn}

It is a fact that if $A \sub \VF^n$ is definable then $\dim_{\VF}(A)$ equals the Zariski dimension of the Zariski closure of $A$, see for example \cite[\S~3.8]{hrushovski:kazhdan:integration:vf} for more  on dimensions.

\begin{lem}[{\cite[Lemma~4.6]{Yin:special:trans}}]\label{full:dim:open:poly}
Let $A$ be a definable set with $A_{\VF} \sub \VF^n$. Then $\dim_{\VF} (A) = n$ if and only if there is a $t \in A_{\RV}$ such that $A_t$ contains an open polydisc.
\end{lem}

\begin{lem}\label{ball:to:ac}
Let $b \in \VF$ and $\ga \sub \VF^n$ be a polydisc with $\dim_{\VF}(\ga) = n$. If $b$ is algebraic over $\code \ga$ then it is indeed algebraic.
\end{lem}
\begin{proof}
Let $\alpha = \rad(\ga)$. Then, for any $a \in \ga$, $b$ is algebraic over $(a, \alpha)$ and hence, by Lemma~\ref{dcl:to:ac}, is algebraic over $a$. Since there is an $a \in \ga$ such that $a_i \notin \acl(a_{\tilde i}, b)$ for all $i \in [n]$, by Lemma~\ref{exchange},  $b$ must be algebraic.
\end{proof}


\begin{lem}[{\cite[Lemma~4.15]{Yin:QE:ACVF:min}}]\label{effectiveness}
If $\mathbb{S}$ is $(\VF, \RV,\Gamma)$-generated then every algebraic (respectively, definable) closed disc  contains an algebraic (respectively, definable) point.
\end{lem}

\begin{lem}[{\cite[Lemma~4.17]{Yin:QE:ACVF:min}}]\label{algebraic:balls:definable:centers}
If $\mathbb{S}$ is $(\VF, \Gamma)$-generated then every algebraic (respectively, definable) disc  contains an algebraic (respectively, definable) point.
\end{lem}

A pillar of the structure of definable sets in $\bb U$ is \emph{\cmin-minimality}, meaning that every definable subset of $\VF$ is a finite boolean combination of valuative discs, allowing additional parameters to name the individual discs.


\begin{cor}\label{open:disc:def:point}
Suppose that $\mathbb{S}$ is $(\VF, \RV,\Gamma)$-generated.
Let $\ga \sub \VF$ be a definable disc and $A$ a definable subset of $\VF$. If $\ga \cap A$ is a nonempty proper subset of $\ga$ then $\ga$ contains a definable point.
\end{cor}
\begin{proof}
It is not hard to see that, by \cmin-minimality, if $\ga \cap A$ is a nonempty proper subset of $\ga$ then $\ga$ contains a definable closed disc and hence the claim is immediate by Lemma~\ref{effectiveness}.
\end{proof}

\begin{defn}
Let $D$ be a set of parameters. We say that a (not necessarily definable) nonempty set $A$ \memph{generates a (complete) $D$-type} if, for every $D$-definable set $B$, either $A \sub B$ or $A \cap B = \0$. In that case, $A$ is \memph{$D$-type-definable} if no set properly contains $A$ and also generates a $D$-type. If $A$ is $D$-definable and generates a $D$-type, or equivalently, if $A$ is both $D$-definable and $D$-type-definable then we say that $A$ is \memph{$D$-atomic} or \memph{atomic over $D$}.
\end{defn}

We simply say ``atomic'' when $D =\0$.

In the literature, a type could be a partial type and hence a type-definable set may have nontrivial intersection with a definable set. In this paper, since partial types do not play a role, we shall not carry the superfluous qualifier ``complete'' in our terminology.

\begin{lem}[{\cite[Lemma~3.3]{Yin:int:acvf}}]\label{atomic:base:change}
Suppose that $\Gamma(\mathbb S)$ is nontrivial. Let $A$ be an atomic set. Then $A$ is $\gamma$-atomic for all $\gamma \in \Gamma$.
\end{lem}


\begin{lem}\label{at:ball:t}
Let $\ga \sub \VF^n$ be an atomic open polydisc. Then $\ga$ is $t$-atomic for all $t \in \RV$.
\end{lem}
\begin{proof}
If disc parameters have previously been named, the unary assertion follows directly from \cite[Lemma~3.9]{Yin:int:acvf}, in its stated $\bb U^{\bullet}$ context: for any formula $\varphi(x,t)$ over this base, the definable map $x\mapsto\{s\in\RV:\varphi(x,s)\}$ is constant on the atomic open disc. Thus naming $t$ does not split its type. Without disc parameters, the following argument applies. For the case $n=1$, if $\ga$ were not $t$-atomic then, by Corollary~\ref{open:disc:def:point} and Lemma~\ref{dcl:to:ac}, it would contain a point algebraic over $\mathbb S$. This point belongs to a finite definable set. Atomicity would then force the whole infinite open disc $\ga$ to lie in that finite set, a contradiction. For the case $n>1$, by induction on $n$, $\ga_1$ is $t$-atomic and for every $a \in \ga_1$, $\ga_{>1}$ is $(a, t)$-atomic. So $\ga$ cannot have a nonempty proper $t$-definable subset.
\end{proof}


\begin{lem}\label{atom:self}
Suppose that $\Gamma(\mathbb S)$ is nontrivial. Let $A \sub \VF^n$ be a set that generates a type and $\gb \sub A$ an open (or closed) polydisc. Then $\gb$ is $\code \gb$-atomic.
\end{lem}
\begin{proof}
We do induction on $n$. The base case $n=1$ is just \cite[Lemma~3.4]{Yin:int:acvf}. For the case $n>1$, write $\gb = \prod_i \gb_i$, $\gb' = \prod_{i>1} \gb_i$, and $\gamma=\rad(\gb_1) $. For any $a \in \gb_1$, $A_a$ generates a type. So, by Lemma~\ref{atomic:base:change} and the inductive hypothesis, $\gb'$ is $(a, \gamma, \code{\gb'})$-atomic and hence $\code \gb$-atomic. Similarly $\gb_1$ is $\code \gb$-atomic. It follows that $\gb$ must be $\code \gb$-atomic.
\end{proof}


%there is a definable set $\{\gb_1, \ldots, \gb_k\}$ of open discs of the same radius such that $f(\ga) = \bigcup_i \gb_i$ and for every $\gb_i$,    $\gb_i \cap f(a)$ is a singleton.

\begin{lem}\label{at:rv:split}
Let $\ga \sub \VF^n$ be an atomic open polydisc and $f : \ga \fun \mdl P(\VF)$ a definable finitary function such that $\rv \rest f(a)$ is injective for every $a \in \ga$.  Then  $\rv \circ f$ is constant.
\end{lem}
\begin{proof}
This is clear since, by Lemma~\ref{at:ball:t}, for every $t \in \rv(\bigcup f(\ga))$, $\ga$ is $t$-atomic and hence $f$ induces a $t$-definable function $\ga \fun t^\sharp$. The lemma follows.
\end{proof}

A \memph{thin annulus} is a set of the form $\gc(a,\gamma)\setminus\go(a,\gamma)=\{x\in\VF:\vv(x-a)=\gamma\}$, where $a\in\VF$ and $\gamma\in\Gamma$. Equivalently, it is a closed disc with one maximal open subdisc removed.

\begin{lem}[{\cite[Lemma~3.8]{Yin:int:acvf}}]\label{hae:component:11}
Let $A \sub \VF$ be an atomic open disc or an atomic closed disc or an atomic thin annulus. Let $f : A \fun \VF$ be a definable function. Then $f(A)$ is also atomic of one of these three forms.
\end{lem}


\begin{lem}[{\cite[Lemma~3.13]{Yin:int:acvf}}]\label{atomic:open:ball:open}
Let $A \sub \VF$ be an atomic open disc and $f : A \fun \VF$ a definable function. Then either $f$ is constant or $f(A)$ is also an atomic open disc.
\end{lem}

\begin{lem}\label{atomic:open:multidisc}
Let $\ga$, $f$ be as in Lemma~\ref{at:rv:split} and suppose that $f$ is not constant. Then there are $t_1, \ldots, t_k \in \RV_\infty$ and for each $i$, a $t_i$-definable function $f_i : \ga \fun t^\sharp_i$ such that $f = \bigcup_i f_i$ and each $f_i(\ga)$ is either a point or a $t_i$-atomic open disc.
\end{lem}
\begin{proof}
By Lemma~\ref{at:rv:split} and injectivity of $\rv$ on each finite fiber, $f$ decomposes into finitely many $t$-definable single-valued branches $g:\ga\fun t^\sharp$. A branch with $t=\infty$ is the constant zero branch. Any constant branch already has the required form, so fix a nonconstant branch, for which $t\in\RV$. Lemma~\ref{at:ball:t} makes $\ga$ atomic over $t$.

Every coordinate line of $\ga$ is atomic over its fixed coordinates and $t$. By Lemma~\ref{atomic:open:ball:open}, its image under $g$ is a point or an open disc. Any two points of $\ga$ can be joined by changing one coordinate at a time. Starting at a point with image $p$ and ending at a point with different image, the first nonconstant line image is an open neighborhood of $p$ contained in $g(\ga)$. Thus $g(\ga)$ is open. It is a proper definable subset of $\VF$, since it lies in $t^\sharp$.

By \cmin-minimality, every point of a proper definable open subset of $\VF$ lies in a maximal open disc contained in that subset. Choose such a disc $D\sub g(\ga)$ containing one value of $g$. An open line image that meets $D$ is contained in $D$: intersecting open discs are nested, and maximality excludes a larger one inside $g(\ga)$. A point image also stays in $D$ when its initial value is in $D$. The coordinate chains therefore keep every value of $g$ in $D$, so $g(\ga)=D$. The definable image of an atomic set is atomic over the same base, hence this open disc is $t$-atomic, as required.

\end{proof}


\begin{cor}\label{part:rv:cons}
Let $A \sub \VF^n$ be a definable set and $f : A \fun \mdl P(\VF)$ a definable finitary function as in Lemma~\ref{at:rv:split}. Then there are a definable finite partition $(A_i)_i$ of $A$ and, for each $i$ and all open polydiscs $\ga \sub A_i$, finitely many $(\code\ga,t_{\ga j})$-definable functions $f_{\ga j} : \ga \fun t^\sharp_{\ga j}$, where $t_{\ga j} \in \RV_\infty$, such that  $f \rest \ga = \bigcup_j f_{\ga j}$  and each $f_{\ga j}(\ga)$ is either a point or an open disc.
\end{cor}
\begin{proof}
By compactness, there is a uniform finite bound $N$ on the cardinalities of the fibers $f(a)$. For an open-polydisc code $d$, write $\ga_d$ for its polydisc and let $P(d)$ say that $\rv\circ(f\rest\ga_d)$ is constant and that the image of each of its single-valued $\rv$-tagged branches is a point or an open disc. This is a first-order condition: use the finite disjunction over the possible numbers $0\leq k\leq N$ of distinct tags, and express a point or an open disc by its usual uniform codes. Injectivity of $\rv$ on each fiber gives uniqueness of each tagged branch.

For $a\in A$, let $D_a\sub A$ be the complete type-definable set containing $a$. By Lemma~\ref{atom:self}, every open polydisc $\ga_d\sub D_a$ is $\code\ga_d$-atomic. Lemma~\ref{atomic:open:multidisc}, over $\mathbb S$ together with this disc parameter, therefore gives $P(d)$. If no finite part of the type of $a$ had this property for every polydisc contained in it, compactness would produce an open-polydisc code $d$ with $\ga_d\sub D_a$ and $\neg P(d)$, a contradiction. Thus there is a definable set $A_a\sub A$ containing $a$ such that $P(d)$ holds whenever $\ga_d\sub A_a$. A second application of compactness gives finitely many such sets covering $A$; take the partition they generate.

For a polydisc $\ga$ in a partition piece, choose its finitely many distinct tags $t_{\ga j}\in\RV_\infty$. The full graph of its $j$th branch is
\[
\{(x,y):x\in\ga,\ y\in f(x),\ \rv(y)=t_{\ga j}\}.
\]
It is $(\code\ga,t_{\ga j})$-definable, the branches reconstruct $f\rest\ga$, and $P(d)$ gives the asserted images.
\end{proof}


\begin{defn}\label{defn:dtdp}
Let $f : A \fun B$ be a bijection between two sets $A$ and $B$, each with exactly one $\VF$-coordinate. We say that $f$ is \memph{concentric} if, for every $\VF$-fiber $f_{t}$ of $f$ and all open polydiscs $\ga \sub \dom(f_{t})$, $f_{t}(\ga)$ is also an open polydisc; if both $f$ and $f^{-1}$ are concentric then $f$ has the \memph{disc-to-disc property}.
\end{defn}

\begin{lem}[{\cite[Proposition~3.19]{Yin:int:acvf}}]\label{open:pro}
Let $f : A \fun B$ be a definable bijection between two sets $A$ and $B$, each with exactly one $\VF$-coordinate. Then there is a definable finite partition $(A_i)_i$ of $A$ such that every restriction $f \rest A_i$ has the disc-to-disc property.
\end{lem}


\subsection{Continuity and fiberwise properties}\label{sec:cont:fib:prop}

\begin{defn}
We set $\dim_{\RV}(\varnothing)=-\infty$. For a nonempty definable subset $U \sub \RV^m_\infty$, the \memph{$\RV$-dimension}, denoted by $\dim_{\RV}(U)$, is the smallest number $k$ such that there is a definable finite-to-one function $f: U \fun \RV^k$ ($\RV^0$ is taken to be the singleton $\{\infty\}$).
\end{defn}


\begin{defn}
Let $A$ be a definable subset of $\VF^n$. The \memph{$\RV$-boundary} of $A$, denoted by $\partial_{\RV}A$, is the definable subset of $\rv(A)$ such that $t \in \partial_{\RV} A$ if and only if $t^\sharp \cap A$ is a proper nonempty subset of $t^\sharp$. The definable set $\rv(A) \mi \partial_{\RV}A$, denoted by $\ito_{\RV}(A)$, is called the \memph{$\RV$-interior} of $A$.
\end{defn}


Obviously, $A \sub \VF^n$ is an $\RV$-pullback if and only if $\partial_{\RV} A$ is empty.


\begin{lem}\label{RV:bou:dim}
Let $A$ be a definable subset of $\VF^n$. Then $\dim_{\RV}(\partial_{\RV} A) < n$.
\end{lem}
\begin{proof}
For $n=0$, the boundary is empty and has dimension $-\infty$. For $n=1$, take a finite Boolean disc decomposition of $A$. A disc containing $0$ is an $\RV$-pullback, whereas a disc avoiding $0$ lies in one $\RV$-disc. A partially met $\RV$-fiber therefore contains a member of this finite disc family that avoids $0$, so there are only finitely many such fibers.

For the inductive step, fix $i\in[n]$ and put $X_i=\pr_{\tilde i}(A)$ and $T_a=\partial_{\RV}A_a$ for $a\in X_i$. Each $T_a$ is finite. For each $t\in T_a$, Corollary~\ref{open:disc:def:point}, over the parameters $(a,t)$, gives a definable point in $t^\sharp$. By compactness there are finitely many definable partial choice functions covering all such pairs $(a,t)$; order this list and use the first available choice to obtain a definable function $c(a,t)\in t^\sharp$. Set $f_i(a)=\{c(a,t):t\in T_a\}$ if $T_a\ne\varnothing$, and $f_i(a)=\{0\}$ otherwise. Thus every fiber is finite and contains exactly one representative of each boundary tag when that boundary is nonempty; in particular $\rv\rest f_i(a)$ is injective.

First split $X_i$ by the pattern of its zero coordinates, and apply Corollary~\ref{part:rv:cons} on the remaining coordinates of each piece. This gives a definable finite partition $(A_{ij})_j$ of $X_i$ such that, whenever $t^\sharp\sub A_{ij}$, the finite set $\rv(f_i(a))$ is independent of $a\in t^\sharp$. The zero-pattern split allows $t$ to have $\infty$ coordinates: after deleting them, $t^\sharp$ is an open polydisc in the remaining coordinates. Write $f_{ij}=f_i\rest A_{ij}$ and let $C_{ij}(t)$ be this common finite set for $t\in\ito_{\RV}(A_{ij})$.

Regard each $C_{ij}$ as its graph in $\RV_\infty^n$, retaining the input tags and placing its output tag in coordinate $i$. Put
\[
 C=\bigcup_{i,j}C_{ij},\qquad
 B=\bigcup_{i,j}\{s\in\rv(A):\pr_{\tilde i}(s)\in\partial_{\RV}A_{ij}\}.
\]
Each graph $C_{ij}$ is finite over its $n-1$ input coordinates, so $\dim_{\RV}(C)<n$. By induction, $\dim_{\RV}(\partial_{\RV}A_{ij})<n-1$; with only one remaining coordinate, this gives $\dim_{\RV}(B)<n$. Finally, for a partially met $\RV$-polydisc, start in $A$ and change coordinates one at a time to reach a point outside $A$. The first exit identifies a boundary tag of a unary section. Its other tags either lie in some $\partial_{\RV}A_{ij}$, placing the full tuple in $B$, or in $\ito_{\RV}(A_{ij})$, placing it in $C$. Hence $\partial_{\RV}A\sub B\cup C$, as required.
\end{proof}


Let $f : \VF^n \fun \VF^m$ be a definable function.

\begin{lem}[{\cite[Lemma~8.7]{Yin:special:trans}}]\label{fun:almost:cont}
There is a definable closed set $A \sub \VF^n$ with $\dim_{\VF}(A) < n$ such that $f \rest (\VF^n \mi A)$ is continuous.
\end{lem}

Per the usual terminology, this lemma says that $f$ is continuous \memph{almost everywhere}.

\begin{defn}\label{defn:diff}
In this definition, for a nonzero tuple $x\in\VF^k$ we set $\vv(x)=\min_i\vv(x_i)$, and we set $\vv(0)=\infty$. This minimum convention is separate from the coordinatewise valuation image used in Terminology~\ref{bounded}.


For any $a \in \VF^n$, we say that $f$ is \memph{differentiable at $a$} if there is a linear map $\lambda : \VF^n \fun \VF^m$ (of $\VF$-vector spaces) such that, for any $\epsilon \in \Gamma$, if $b \in \VF^n\setminus\{0\}$ and $\vv(b)$ is sufficiently large then
\[
\vv(f(a + b) - f(a) - \lambda(b)) - \vv(b) > \epsilon.
\]
It is straightforward to check that if such a linear function $\lambda$ (represented by a matrix with entries in $\VF$) exists then it is unique and hence is called the \memph{derivative of $f$ at $a$}.

Write $f = (f_1, \ldots, f_m)$. For $a = (a_i, a_{\tilde i}) \in \VF^n$, if the derivative of the function $f_j \rest (\VF \times a_{\tilde i})$ at $a_i$ exists then we call it the \memph{$ij$th partial derivative of $f$ at $a$}. Whenever $f$ is differentiable at $a$, the $(j,i)$-entry of its derivative matrix is this $ij$th partial derivative.
\end{defn}


We differentiate functions between arbitrary definable sets by ``forgetting'' the $\RV$-coordinates. More precisely, let
\[
f : \VF^n \times \RV^m \fun \VF^{n'} \times \RV^{m'}
\]
be a definable function. For each fixed input $\RV$-tuple $t$, generic local constancy of definable $\RV$-valued maps \cite[Lemma~8.9]{Yin:special:trans} implies that, outside a definable set of $\VF$-dimension less than $n$, the output $\RV$-tuple is constant on a neighborhood of each input $a$. If this tuple is $s$, then $\dom(f_{(t,s)})$ contains a neighborhood of $a$. At such a point $(a,t)$, we define the partial derivatives of $f$ to be those of this local $\VF$-map $f_{(t,s)}$. If $n = n'$ and all the partial derivatives exist, the \memph{Jacobian of $f$ at $(a,t)$} is the determinant of the matrix whose $(j,i)$-entry is the $ij$th partial derivative, and is denoted by $\jcb_{\VF} f(a,t)$.

\begin{lem}\label{diff:almost:every}
Assume $n=n'$. The Jacobian of $f$ exists and is a continuous function almost everywhere.
\end{lem}
\begin{proof}
First remove the definable locus where the output $\RV$-tuple is not locally constant in the input $\VF$-coordinates. On each resulting open fiber with both input and output $\RV$-tuples fixed, extend the $\VF$-map by zero outside its domain and apply \cite[Corollary~9.9]{Yin:special:trans} to its partial derivatives. Their nonexistence and discontinuity loci are definable and have $\VF$-dimension less than $n$ in every such fiber. Retaining or projecting away $\RV$-coordinates does not increase $\VF$-dimension, so the combined exceptional locus still has $\VF$-dimension less than $n$. Off this locus all partial derivatives, and hence their determinant, are continuous. In arity zero the Jacobian is the empty determinant $1$.
\end{proof}


\begin{ter}\label{bounded}
We say that a set $I \sub \Gamma_{\infty}^n$ is \memph{$\gamma$-bounded}, where $\gamma \in \Gamma$, if it is contained in the box $[\gamma, \infty]^n$, and is \memph{doubly $\gamma$-bounded} if it is contained in the box $[-\gamma, \gamma]^n$. More generally, let $A$ be a subset of $\VF^n \times \RV_\infty^m \times \Gamma_\infty^l$ and
\[
\operatorname{valim}(A) = \set{(\vv(a), \vrv(t), \gamma) \given (a, t, \gamma) \in A} \sub \Gamma_{\infty}^{n+m+l},
\]
where $\vv(a)$ and $\vrv(t)$ are taken coordinatewise. We say that $A$ is \memph{$\gamma$-bounded} if $\operatorname{valim}(A)$ is, and likewise for \memph{doubly $\gamma$-bounded}. A set is \memph{bounded} if it is $\gamma$-bounded for some $\gamma\in\Gamma$, and \memph{doubly bounded} if it is doubly $\gamma$-bounded for some $\gamma\in\Gamma$. The notation $A_{\Gamma}$ continues to mean projection to existing $\Gamma$-coordinates.
\end{ter}


\begin{defn}[Contractions]\label{defn:corr:cont}
A function $f : A \fun B$ is \memph{$\rv$-contractible} if there is a (necessarily unique) function $f_{\downarrow} : \rv(A) \fun \rv(B)$, called the \memph{$\rv$-contraction} of $f$, such that
\[
(\rv \rest B) \circ f = f_{\downarrow} \circ (\rv \rest A).
\]
Similarly, it is \memph{$\res$-contractible} (respectively, \memph{$\vv$-contractible}) if the same holds in terms of $\res$ (respectively, $\vv$ or $\vrv$, depending on the coordinates) instead of $\rv$.
\end{defn}


\begin{lem}\label{db:to:db}
Let $U \sub \RV^k$ be a doubly bounded set and $f: U \fun \RV$ a definable function. Then $f(U)$ is also doubly bounded. The same holds if the codomain of $f$ is $\Gamma$.
\end{lem}
\begin{proof}
For an $\RV$-valued $f$, put $h=\vrv\circ f$; for a $\Gamma$-valued $f$, put $h=f$. We retain the induction on $k$, with the case $k=0$ immediate. For $k>1$ and $t\in\pr_1(U)$, the inductive hypothesis makes $h(U_t)$ doubly bounded. Its finite bound
\[
\gamma_t=\sup\set{|h(t,x)|\given x\in U_t}
\]
exists by definable completeness of $\Gamma$, and the least-upper-bound formula makes $t\efun\gamma_t$ definable. The projection $\pr_1(U)$ is doubly bounded, so the unary $\Gamma$-valued case bounds all $\gamma_t$ and hence bounds $h(U)$. Thus both claims reduce to the case $k=1$ and a definable function $g:U\fun\Gamma$. Then, by \cmin-minimality, we may assume that $g$  $\vrv$-contracts to a function $g_{\downarrow} : \vrv(U) \fun \Gamma$. Since definable functions in $\Gamma$ are piecewise $\Q$-linear, the range of $g_{\downarrow}$ must be doubly bounded.
\end{proof}


% the following lemma has been moved to the section on proper covariants bijecetions, since it is not used before and is a corollary of the sujectivity of the lifting morphism
%Recall that, by \cmin-minimality, every definable set $A \sub \VF$ may be written as a unique boolean combination of discs $\ga_i$. Any atomic subset of these discs is called an \memph{atomic band} of $A$. If $\gB$ is an atomic band of $A$ then every disc in it has the same attributes: radius, open or closed, a positive or a negative boolean component, etc. So they may be considered as attributes of $\gB$. The totality of these attributes is called the \memph{signature} of $\gB$.
%
%\begin{lem}\label{invar:db}
%Let $A \sub \VF^n \times \RV^m$ be a definable set such that $A_{\VF}$ is bounded and $A_{\RV}$ is doubly bounded. Let $f : A \fun \Gamma$ be a definable function  that is constant on $A \cap \ga$ for all open polydiscs $\ga$ of radius $\gamma \in \Gamma$. Then $f(A)$ is doubly bounded.
%\end{lem}
%\arthur{I suspect one can get this for free as a corollary of lemma \ref{L:measure:surjective} and \ref{db:to:db}, if we can postpone the proof. It seems to be the case, since the lemma is only used in the last section.}
%\begin{proof}
%By Lemma~\ref{db:to:db} and \omin-minimality in the $\Gamma$-sort, it is enough to consider the case $A \sub \VF^n$. It is easy to see that, by induction on $n$, this may be further reduced to the case $n=1$. By the assumption, we may assume that, for every open disc $\ga$ of radius $\gamma$, either $\ga \cap A = \0$ or $\ga \sub A$, and moreover $A \cap \MM_\gamma = \0$, that is, $A$ is doubly bounded.
%
%Suppose for contradiction that $f(A)$ contains a definable interval $I$ of the form $(\beta, \infty)$. By \cmin-minimality, for  each $\alpha \in I$, if $\delta^{\dsh} \cap f^{-1}(\alpha)$ is a nonempty proper subset of $\delta^{\dsh}$ for some  $\delta \in \Gamma$ then some atomic band $\gB_\alpha$ of $f^{-1}(\alpha)$ is contained in $\delta^{\dsh}$. Therefore, shrinking $I$ if necessary, we can construct a definable function $g$ on $I$ such that
%\begin{itemize}
%  \item either $g(\alpha)  = \gB_\alpha$ of the same signature for all $\alpha \in I$ or
%  \item $g(\alpha) \in \Gamma$ and $g(\alpha)^{\dsh} \sub f^{-1}(\alpha)$ for all $\alpha \in I$.
%\end{itemize}
%In the latter case, we may assume that $g$ is $\Q$-linear, but then its image cannot be doubly bounded since $g$ is injective, contradicting the choice of $I$. The former case requires further argument.
%
%Consider the function $\bar g$ on $I$ given by $\alpha \efun \rad(\gB_\alpha)$. Since $\rad(\gB_\alpha) \leq \gamma$ for all $\alpha \in I$ and $A$ is bounded,  $\bar g(I)$ must be doubly bounded. By \omin-minimality in the $\Gamma$-sort, shrinking $I$ if necessary, we may assume that $\bar g$ is constant, say $\bar g(I) = \delta$, and either every $\gB_\alpha$ is closed or every $\gB_\alpha$ is open. In the former case, by Lemmas~\ref{effectiveness} and \ref{dcl:to:ac}, $g(I)$ must be finite; but this is impossible since either every $\gB_\alpha$ is a positive boolean component or every $\gB_\alpha$ is a negative boolean component. In the latter case, for the same reason, we may assume that $g(I)$ is infinite. Then a moment of reflection shows that there would be a closed disc $\gb$ of radius $\delta$ and a $\code \gd$-definable function, induced by the ``inverse'' of $g$, from the set of maximal open subdiscs of $\gd$ onto an infinite subset of $\Gamma$, which is easily seen to contradict \cmin-minimality.
%\end{proof}
%


\begin{defn}
Let $A \sub \VF^n$ and $f : A \fun \mdl P(\VF^m)$ be a definable function. For $a \in \VF^n$ and $L \sub \VF^m$, we say that $L$ is a \memph{limit set of $f$ at $a$}, written as $\lim_{A \rightarrow a} f \sub L$, if for every $\epsilon \in \Gamma$ there is a $\delta \in \Gamma$ such that if $c \in \go(a,\delta) \cap (A\setminus\{a\})$ then
\[
f(c)\sub\bigcup_{b\in L}\go(b,\epsilon).
\]
This is a local definition and imposes no global boundedness hypothesis. A finite limit set forces the union of the values $f(c)$ to be bounded on some punctured neighborhood of $a$. Whenever boundedness of the range is assumed, it means boundedness of $\bigcup_{c\in A}f(c)$.

A limit set $L$ of $f$ at $a$ is \memph{minimal} if no proper subset of $L$ is a limit set of $f$ at $a$.
\end{defn}


Observe that if $\lim_{A \rightarrow a} f \sub L$ and $b \in L$ is not isolated in $L$ then actually $\lim_{ A \rightarrow  a} f \sub L \mi b$. So in a minimal limit set every element is isolated. Moreover, if a minimal limit set $L$ exists then its topological closure is unique:

\begin{lem}[{\cite[Lemma~9.2]{Yin:special:trans}}]
Let $L_1, L_2 \sub \VF^m$ be two minimal limit sets of $f$ at $a$. Then the closures of $L_1$, $L_2$ coincide.
\end{lem}


This lemma justifies writing $\lim_{ A \rightarrow a} f = L$ when $L$ is a closed (hence the unique) minimal limit set of $f$ at $a$, in particular, when $L$ is finite and minimal.

Let $A \sub \VF^n$ be a definable set with $\dim_{\VF}(A) = n$ and $f : A \fun \mdl P(\VF^m)$ a definable finitary function. We say that $f$ is \memph{continuous} if $\lim_{ A \rightarrow a} f = f(a)$ for all $a \in A$.

\begin{lem}\label{fini:conti}
Suppose $n>0$. Then $f$ is continuous almost everywhere.
\end{lem}
\begin{proof}
This can be reduced to the familiar case that $f$ is indeed a function into $\VF^m$ as follows. By Lemma~\ref{finite:VF:project:RV} and compactness, there are a definable set $U$ in $\RV$ and a definable function $g: A \times \VF^m \fun A \times U$ such that, for all $a \in A$, $\pr_A(g(a, b)) = a$  and $g \rest (a \times f(a))$ is injective. For each $t \in U$, let
\[
A_t = \set{a \in A \given (a,t) \in g(a \times f(a))}
\]
and $f_t \sub f$ be the obvious $t$-definable function $A_t \fun \VF^m$. The discontinuity locus $E$ of $f$ is definable: the limit condition is first-order, and minimality of the finite set $f(a)$ can be tested by removing one of its elements at a time. Suppose for contradiction that $\dim_{\VF}(E)=n$. By Lemma~\ref{full:dim:open:poly}, $E$ contains a nonempty open polydisc. Apply generic local constancy \cite[Lemma~8.9]{Yin:special:trans} to the finite label-set map
\[
a\efun g(a\times f(a))_{\RV}.
\]
After shrinking to a nonempty open polydisc $P\sub E$, this set is a fixed finite set $T$. If $T$ is empty, then $f$ is the constant empty correspondence on $P$ and is continuous there, a contradiction.

For each $t\in T$, the inverse branch $f_t$ is defined on all of $P$. Extend it by zero outside $A_t$ and apply Lemma~\ref{fun:almost:cont}. The finitely many branch-exception sets have union of $\VF$-dimension less than $n$, so choose $a\in P$ outside their union. All branches are continuous at $a$. Their finite union has limit set $f(a)$; this limit set is minimal because the distinct values at $a$ can be separated by sufficiently small discs and, since $n>0$, the punctured polydisc still accumulates at $a$. Thus $f$ is continuous at $a$, contradicting $P\sub E$. Hence $\dim_{\VF}(E)<n$.
\end{proof}


\begin{lem}[{\cite[Lemma~9.5]{Yin:special:trans}}]\label{lim:exists}
Suppose that $f : \MM \mi 0 \fun \mdl P (A)$ is a definable finitary function, where $A$ is a closed and bounded set in $\VF$. Then there is a definable finite set $L \sub A$ such that $\lim_{\MM \mi 0 \rightarrow 0} f = L$.
\end{lem}

We remark that the finite-limit theorem also holds in the coordinatewise projective line $(\mathbb P^1(\VF))^m$, with the valuation topology: if $A\subseteq\VF^m$ is not bounded, take its closure there and allow infinite coordinates in $L$. This is the projective version of the same cited theorem; no topological local compactness is assumed.


\begin{lem}[Curve selection]\label{curve:select}
Suppose that $\mathbb S$ is $\VF$-generated. Let $A\sub\VF^n$ be definable and let $b$ be a point in its closure in the coordinatewise projective line $(\mathbb P^1(\VF))^n$, with the valuation topology. Then there is a $b$-definable continuous finitary function $g:\MM\setminus\{0\}\fun\mdl P(A)$ whose values are nonempty and such that
\[
\lim_{\MM\setminus\{0\}\to0}g=\{b\}.
\]
\end{lem}
\begin{proof}
Use a coordinatewise projective chart around $b$ to reduce to the case $b\in\VF^n$: at an infinite coordinate use inversion and restrict $A$ to the chart. These changes are definable over $b$ and preserve convergence and continuity locally.

Naming the finite tuple $b$ and a nonzero $c\in\MM$ preserves $\VF$-generation and the standing nontriviality of the ground valuation. Hence $\acl(\mathbb S,b,c)$ is a model of $\ACVF(0,0)$. Since $b\in\overline A$, the definable set $A\cap\go(b,\vv(c))$ is nonempty, so elementarity gives a point of that intersection algebraic over $(\mathbb S,b,c)$. Compactness supplies a $b$-definable finitary function $g_0$ on $\MM\setminus\{0\}$ such that
\[
\varnothing\ne g_0(c)\sub A\cap\go(b,\vv(c))
\]
for every $c$. Thus $g_0$ converges to $\{b\}$, and nonempty values make this limit minimal.

By Lemma~\ref{fini:conti}, with its one-dimensional source, $g_0$ is continuous outside a finite set. Choose a $b$-definable clopen disc $D\sub\MM$ around $0$ such that $g_0$ is continuous on $D\setminus\{0\}$. Also $\acl(\mathbb S,b)$ is a model, so $A$ contains a finite nonempty $b$-definable set $E$. Define
\[
g(c)=\begin{cases}
g_0(c),&c\in D\setminus\{0\},\\
E,&c\in\MM\setminus D.
\end{cases}
\]
The two pieces are relatively clopen in $\MM\setminus\{0\}$, so $g$ is continuous, finite-valued and nonempty everywhere. Its limit at $0$ is still $\{b\}$. Finally transport the construction back through the projective chart. This avoids assuming that a definable value-group radius is the valuation of a definable scalar.
\end{proof}


This is an analogue of the highly useful \omin-minimal curve selection.

\begin{rem}
Although the $\Gamma$-sort is essentially an \omin-minimal group, curve selection fails due to the presence of the top element $\infty$. For example, the set $\Gamma \times (0, \infty) \sub  \Gamma^2_\infty$ contains curves approaching the boundary point $(0, 0)$ or $(\infty, \infty)$ but no curves approaching the boundary point $(\infty, 0)$, because every definable one-parameter curve is piecewise rational-affine. Near a one-sided endpoint, one affine piece remains, and a positive affine coordinate cannot tend to $0$ while another coordinate tends to $\infty$.
\end{rem}


%In addition, we can make this ``curve'' $f \circ g$ approach $b$ as fast as we wish. More precisely, if $l : \Gamma^+ \fun \Gamma^+$ is a definable increasing function then $g$ can be so constructed that $\vv(f(g(c))) > l(\vv(c))$ for all $c \in \MM \mi 0$.


\begin{lem}\label{cb:to:cb}
Let $A\subseteq\VF^n$ be definable, closed and bounded, and let $f:A\fun\VF^m$ be definable and continuous. Then $f(A)$ is also closed and bounded.
\end{lem}
\begin{proof}
The conclusion is geometric, so we may temporarily enlarge the parameter base to a small $\VF$-generated base over which all the data are definable. The standing nontriviality of the ground valuation is preserved. Over this base, and after naming any additional finite $\VF$ tuple, algebraic closure is a model of $\ACVF(0,0)$. Elementarity and compactness therefore give finite nonempty definable choices in the nonempty families below. No choice over $\mathbb S$ itself is asserted.

First let $b\in\overline{f(A)}\setminus f(A)$, where this closure is affine, and name $b$. For every $c\in\MM\setminus\{0\}$ the set
\[
\{x\in A:f(x)\in\go(b,\vv(c))\}
\]
is nonempty. Applying the model argument over the enlarged base together with $(b,c)$, and then compactness, gives a definable finitary function $F:\MM\setminus\{0\}\fun\mdl P(A)$ such that every $F(c)$ is nonempty and $f(F(c))\subseteq\go(b,\vv(c))$. Lemma~\ref{lim:exists}, in its finite-dimensional form, gives a finite minimal limit set $L\subseteq A$. This set is nonempty, since $F$ has nonempty values arbitrarily close to $0$. For every $a\in L$, minimality implies that values of $F$ meet every neighborhood of $a$ for inputs arbitrarily close to $0$: otherwise $a$ could be removed from $L$. Continuity at $a$ and the displayed containment force $f(a)=b$, contradicting $b\notin f(A)$. Thus $f(A)$ is closed.

If $f(A)$ were unbounded, one coordinate $f_i$ would be unbounded. For every $c\in\MM\setminus\{0\}$ the definable set
\[
\{x\in A:\vv(f_i(x))<-\vv(c)\}
\]
would be nonempty. The same model argument and compactness give finite nonempty choices $F(c)$ inside these sets. Again Lemma~\ref{lim:exists} supplies a nonempty finite minimal limit set $L\subseteq A$. Choose $a\in L$. Continuity makes $f_i$ bounded on a neighborhood of $a$ in $A$, while minimality makes $F(c)$ meet that neighborhood for arbitrarily small nonzero $c$. The inequality $\vv(f_i(x))<-\vv(c)$ for all $x\in F(c)$ contradicts this local bound. Hence $f(A)$ is bounded as well.
\end{proof}


\begin{cor}\label{conti:homeo}
If the function $f$ in Lemma~\ref{cb:to:cb} is injective then  it is a homeomorphism from $A$ onto $f(A)$.
\end{cor}

\begin{lem}[The closed graph criterion for continuity]\label{closed:gra}
Let $A\subseteq\VF^n$ and $B\subseteq\VF^m$ be definable, with $B$ closed in affine space. Write $B^{\mathrm p}$ for the closure of $B$ in $(\mathbb P^1(\VF))^m$, with the valuation topology. A definable function $f:A\fun B$ is continuous if and only if its graph is closed in $A\times B^{\mathrm p}$. If $B$ is bounded, then $B^{\mathrm p}=B$, so this is the usual relative closed-graph condition.
\end{lem}
\begin{proof}
The graph of a continuous map into the Hausdorff space $B^{\mathrm p}$ is closed, proving the forward implication.

Conversely, suppose the graph is closed but $f$ is not continuous at $a\in A$. Choose a clopen projective polydisc $U$ around the finite point $f(a)$ such that
\[
D=\{x\in A:f(x)\notin U\}
\]
has $a$ in its closure. As in the preceding proof, temporarily enlarge the parameter base to a small $\VF$-generated base containing the parameters of these data. This changes neither continuity nor graph closure. By Lemma~\ref{curve:select}, there is a definable continuous finitary function $C:\MM\setminus\{0\}\fun\mdl P(D)$ with nonempty values and limit $\{a\}$.

Apply the projective version of Lemma~\ref{lim:exists} to the finite-valued function $c\efun f(C(c))$. It has a finite minimal limit set $L$, which is nonempty because all the values are nonempty. Every point of $L$ belongs to the closed set $B^{\mathrm p}\setminus U$: minimality makes the image meet each neighborhood of such a point for inputs arbitrarily close to $0$. Choose $b'\in L$. This same property, together with $C(c)\to\{a\}$, shows that $(a,b')$ belongs to the closure of the graph of $f$ in $A\times B^{\mathrm p}$. Closedness of the graph gives $b'=f(a)$, contrary to $b'\notin U$. This proves continuity without a global $\VF$-generation hypothesis.
\end{proof}


\begin{lem}\label{lim:weak}
Let $f : A  \fun B$ be a definable bijection between two open subsets of $\VF$. Suppose that $B$ is bounded. Then $f$ is weakly concentric, that is, for every $a \in A$ there is a $b \in \VF$ such that, for every sufficiently large $\gamma \in \Gamma$, $f(\go(a, \gamma) \mi a) = \go(b, \delta) \mi b$ for some $\delta \in \Gamma$.
\end{lem}

Note that, compared to Lemma~\ref{open:pro}, the assumption of this lemma is stronger and its conclusion weaker, but the point is that we do not need a partition to achieve it.

\begin{proof}
Fix $a\in A$. After translating the input and restricting to a small punctured disc about $a$, Lemma~\ref{lim:exists}, applied with the closed bounded target $\overline B$, gives a finite minimal limit. Its cardinality is at most one by \cite[Lemma~9.4]{Yin:special:trans}. It is nonempty because the function has a value at every point of the punctured disc. Write this limit as $\{b\}$.

By shrinking the punctured disc, exclude the at most one input for which $f(x)=b$. The valuation calculation in the proof of \cite[Lemma~8.3]{Yin:special:trans}, applied to the translated germ $f(x)-b$, gives a rational-affine law
\[
\vv(f(x)-b)=\rho(\vv(x-a)),\qquad \rho(\gamma)=q\gamma+c,
\]
for all sufficiently small nonzero $x-a$, where $q\in\mathbb Q_{>0}$ and $c\in\Gamma$. Indeed, unary polynomials have eventually monomial valuation, the definable germ selects one of the finitely many resulting valuation laws, and convergence to $b$ forces positive slope.

By $C$-minimality, $B$ contains a punctured neighborhood of $b$. On this neighborhood the inverse of $f$ has a unique projective limit, using the projective finite-limit theorem and the same cardinality bound. That limit must be $a$, since $x\to a$ implies $f(x)\to b$ and the image avoids $b$ on a sufficiently small punctured disc. Apply the same valuation calculation to the inverse to obtain a positive rational-affine law $\sigma$ for its tail radii. Substituting $y=f(x)$ shows that $\sigma\circ\rho$ is the identity on a tail of $\Gamma$, so the two affine laws are inverse there.

Consequently, for every sufficiently large $\gamma$, with $\delta=\rho(\gamma)$, the image of the whole source thin annulus is exactly the whole target thin annulus:
\[
f\bigl(\gc(a,\gamma)\setminus\go(a,\gamma)\bigr)
 =\gc(b,\delta)\setminus\go(b,\delta).
\]
The forward inclusion is the law for $f$; the reverse inclusion follows by applying the inverse law to every point of the target annulus. Taking the unions over all radii greater than $\gamma$ proves
\[
f\bigl(\go(a,\gamma)\setminus\{a\}\bigr)
 =\go(b,\rho(\gamma))\setminus\{b\},
\]
as required. No assertion that $q=1$ is needed.
\end{proof}


\begin{defn}
Let $p : A \fun \Gamma$ be a definable function. We say that $p$ is an \memph{$\go$-partition} of $A$ if, for every $a \in A$, the function $p$ is constant on $\go(a, p(a)) \cap A$.
\end{defn}


This makes sense even when $A$ has no $\VF$-coordinates, in which case any definable function $A \fun \Gamma$ is an $\go$-partition of $A$.

\begin{lem}\label{vol:par:bounded}
Let $p$ be an $\go$-partition of $A$. Suppose that $A_{\RV}$ is doubly bounded and $A_t$ is closed and bounded for every $t \in A_{\RV}$. Then $p(A)$ is doubly bounded.
\end{lem}
\begin{proof}
First suppose that $A$ has no $\RV$-coordinates. The empty set is immediate. Choose $\eta\in\Gamma$ such that $A\subseteq\OO_\eta^n$. If $p(a)<\eta$ for some $a\in A$, then $A\subseteq\go(a,p(a))$, so the partition property makes $p$ constant on $A$. Otherwise $p(A)\subseteq[\eta,\infty)$. In either case $p(A)$ is bounded below.

Suppose for contradiction that $p(A)$ is unbounded above. We may temporarily enlarge the parameter base to a small $\VF$-generated base, as in Lemma~\ref{cb:to:cb}; the desired bound is a geometric conclusion. For every $c\in\MM\setminus\{0\}$ the definable set
\[
A_{\vv(c)}=\{x\in A:p(x)>\vv(c)\}
\]
is nonempty. Elementarity of the algebraic-closure model over this base together with $c$, followed by compactness, supplies a definable finitary function $F:\MM\setminus\{0\}\fun\mdl P(A)$ with finite nonempty values $F(c)\subseteq A_{\vv(c)}$. Continuity of this choice is not needed.

By Lemma~\ref{lim:exists}, $F$ has a finite minimal limit set $L\subseteq A$, nonempty because its values are nonempty arbitrarily close to $0$. Choose $e\in L$. Minimality implies that, for arbitrarily small nonzero $c$, $F(c)$ meets $\go(e,p(e))$. Choose such a $c$ with $\vv(c)>p(e)$ and $x\in F(c)\cap\go(e,p(e))$. The partition property at $e$ gives $p(x)=p(e)$, whereas the choice gives $p(x)>\vv(c)>p(e)$, a contradiction. This proves the upper bound.

For the general case, each nonempty fiber image $p(A_t)$ is bounded above and below by the case just proved, with $t$ temporarily named. Definable completeness in $\Gamma$ gives finite endpoints
\[
\alpha_t=\inf p(A_t),\qquad \beta_t=\sup p(A_t).
\]
Their defining greatest-lower-bound and least-upper-bound formulas make these functions definable. Thus
\[
t\efun\gamma_t=\max\{|\alpha_t|,|\beta_t|,0\}
\]
is a definable $\Gamma$-valued function on $A_{\RV}$. By Lemma~\ref{db:to:db}, its image is doubly bounded. A common upper bound for the $\gamma_t$ bounds $p(A)$ in both directions, proving the lemma.
\end{proof}


This lemma, in its various incarnations, is crucial for the good behavior of motivic Fourier transform (see \cite[\S~11]{hrushovski:kazhdan:integration:vf} and \cite{yin:hk:part:3}). For essentially the same reason, the  main construction of this paper depends heavily on  it (often in combination with Lemma~\ref{db:to:db}).


\begin{lem}\label{lem:open:fiberwise}
Let $A\subseteq \VF^n$ be a definable set. Suppose that, for each $a \in A_1$, the fiber $A_a \sub \VF^{n-1}$ is open (respectively, closed). Then there is a finite definable set $A' \subseteq A_1$ such that $A\cap ((A_1 \mi A')\times \VF^{n-1})$ is open (respectively, closed) in $(A_1 \mi A')\times \VF^{n-1}$.
\end{lem}
\begin{proof}
We prove the open case over an arbitrary definable unary base $I\subseteq\VF$, allowing empty fibers, with $A\subseteq I\times\VF^{n-1}$. The closed case then follows by taking the relative complement in $I\times\VF^{n-1}$. In the statement take $I=A_1$.

Put $X=I\times\VF^{n-1}$ and let
\[
S=A\cap\Cl_X(X\setminus A).
\]
Thus no point of $S$ is interior to $A$ relative to $X$. It is enough to show that $S_1$ is finite. For each $\gamma\in\Gamma$, set
\[
T_\gamma=\{(a,x)\in S:\go(x,\gamma)\subseteq A_a\}.
\]
Openness of every fiber gives
\[
S_1=\bigcup_{\gamma\in\Gamma}(T_\gamma)_1.
\]
If every $(T_\gamma)_1$ were finite, this definable $\Gamma$-indexed union would be finite as well: lift the index to $\RV$ and use the $\RV$-indexed $\VF$-dimension bound \cite[Lemma~4.3]{Yin:special:trans}. Consequently, if $S_1$ is infinite, some $(T_\gamma)_1$ contains an open disc $D$.

For this geometric argument we may temporarily enlarge the parameter base to a small $\VF$-generated base containing the data, as in Lemma~\ref{cb:to:cb}. For every $a\in D$, the nonempty definable fiber $(T_\gamma)_a$ has an algebraic point by elementarity of the algebraic-closure model over this base together with $a$. Compactness therefore gives a definable finitary function
\[
F:D\fun\mdl P(\VF^{n-1}),\qquad
\varnothing\ne F(a)\subseteq(T_\gamma)_a.
\]
By Lemma~\ref{fini:conti}, choose $a_0\in D$ at which $F$ is continuous, and put $E=F(a_0)$. After shrinking $D$ around $a_0$, continuity gives
\[
F(a)\subseteq\bigcup_{e\in E}\go(e,\gamma)\qquad(a\in D).
\]
Here $E$ is finite and nonempty. The finitely many definable sets
\[
D_e=\{a\in D:F(a)\cap\go(e,\gamma)\ne\varnothing\},\qquad e\in E,
\]
cover $D$. Hence one of them contains an open subdisc $D'$. Fix its index $e$. For each $a\in D'$, choose $x\in F(a)\cap\go(e,\gamma)$. Since $(a,x)\in T_\gamma$ and intersecting open polydiscs of the same radius are equal,
\[
\go(e,\gamma)=\go(x,\gamma)\subseteq A_a.
\]
Thus $D'\times\go(e,\gamma)$ is an open subset of $X$ contained in $A$ and meeting $S$, a contradiction. Therefore $S_1$ is finite. Away from $S_1$, every point of $A$ is interior relative to $X$, which proves the open case and, by relative complementation, the closed case.
\end{proof}


\begin{prop}\label{prop:continuity:fiberwise}
Let $A$ be as in Lemma~\ref{lem:open:fiberwise} and $f : A \fun \VF$  a definable function such that
for each $a \in A_1$, the induced function $f_a : A_a \fun \VF$ is continuous.
Then there is a finite definable set $A' \sub A_1$ such that $f$ is continuous on
$\bigcup_{a\in A_1 \mi A'} a \times A_a$.
\end{prop}

For the closed-graph argument, the projective line $\mathbb P^1(\VF)$ is understood with the valuation topology, as in Lemma~\ref{closed:gra}.

\begin{proof}
Put $I=A_1$ and write $H_a$ for the graph of $f_a$. Let
\[
G_a=\Cl_{\VF^{n-1}\times\mathbb P^1(\VF)}(H_a),\qquad
G=\bigcup_{a\in I}\{a\}\times G_a.
\]
This family is definable: fiberwise closure is expressed uniformly by the valuative neighborhood formula in each affine chart. Every $G_a$ is closed in $\VF^{n-1}\times\mathbb P^1(\VF)$.

Apply the fixed-base version proved in Lemma~\ref{lem:open:fiberwise}, with empty fibers allowed, to $G$ in each of the two standard affine charts of $\mathbb P^1(\VF)$. Outside the union $I'$ of the resulting two finite definable exceptional sets, $G$ is closed in
\[
(I\setminus I')\times\VF^{n-1}\times\mathbb P^1(\VF),
\]
since closedness is local on this open chart cover. Fiberwise continuity gives
\[
H_a=G_a\cap\bigl(A_a\times\mathbb P^1(\VF)\bigr).
\]
Consequently, the graph of the restricted function is closed relative to
\[
\bigl(A\cap((I\setminus I')\times\VF^{n-1})\bigr)\times\mathbb P^1(\VF).
\]
The closed-graph criterion (Lemma~\ref{closed:gra}) proves the asserted continuity.
\end{proof}


This is analogous to \cite[Corollary~6.2.4]{dries:1998} in the \omin-minimal environment and \cite[Lemma~2.38]{Yin:tcon:1.5} in the \T-convex environment.


\section{Hrushovski-Kazhdan style integral}\label{sec:HK:main}

\begin{defn}
\label{def:Kgroup}
The Grothendieck semigroup $\gsk \mdl C$ of a category $\mdl C$ endowed with a binary operation ``$\mi$'' and a binary relation ``$\subseteq$'' (subobject) is the free semigroup generated by the isomorphism classes of $\mdl C$, subject to the usual scissor relation
$[A \mi B] + [B] = [A]$, when $B\subseteq A$,
where $[A]$, $[B]$ denote the isomorphism classes of the objects $A$, $B$, and ``$\mi$'' and ``$\subseteq$'' are compatible with isomorphisms between objects, usually they are just set subtraction and set inclusion. Sometimes $\mdl C$ is also equipped with a binary operation (for example, cartesian product) that induces multiplication in $\gsk \mdl C$, in which case $\gsk \mdl C$ becomes a (commutative) semiring. The formal groupification of  $\gsk \mdl C$, which is then a ring, is denoted by $\ggk \mdl C$. %If a group $G$ acts on the objects of $\mdl C$ and the morphisms of $\mdl C$ are $G$-equivariant, that is, they commute with $G$-actions, then the corresponding $G$-equivariant Grothendieck ring is denoted by $\ggk^G \mdl C$. If $G = \lim_n G_n$ is profinite then we shall always impose the condition that a $G$-action factor through some $G_n$-action.
\end{defn}

\begin{defn}
An \memph{$\RV$-fiber} of a definable set $A$ is a set of the form $A_a$, where $a \in A_{\VF}$. The \memph{$\RV$-fiber dimension} of $A$ is the maximum of the $\RV$-dimensions of its $\RV$-fibers and is denoted by $\dim^{\fib}_{\RV}(A)$.
\end{defn}


\begin{lem}[{\cite[Lemma~4.13]{Yin:special:trans}}]\label{RV:fiber:dim:same}
Let $f : A \fun A'$ be a definable bijection. Then $\dim^{\fib}_{\RV}(A) = \dim^{\fib}_{\RV} (A')$.
\end{lem}

\begin{defn}[$\VF$-categories]\label{defn:VF:cat}
The objects of the category $\VF[k]$ are the definable sets of $\VF$-dimension less than or equal to $k$ and $\RV$-fiber dimension $0$ (that is, all the $\RV$-fibers are finite). Any definable bijection between two such objects is a morphism of $\VF[k]$. Set $\VF_* = \bigcup_k \VF[k]$.
\end{defn}

As soon as one considers adding volume forms to definable sets in $\VF$, the question of ambient dimension arises and, consequently, one has to take ``essential bijections'' as morphisms.

 \begin{defn}
\label{defn:VF:mu:G}
An object of the category $\mgVF[k]$ is a definable pair $(A, \omega)$, where $A \in \VF[k]$, $A_{\VF} \sub \VF^k$, and $\omega : A \fun \Gamma$ is a function, which is understood as a \memph{definable $\Gamma$-volume form} on $A$. A \memph{morphism} between two such objects $(A, \omega)$, $(B, \sigma)$ is a definable \memph{essential bijection} $F : A \fun B$, that is, a bijection that is defined outside definable subsets of $A$, $B$ of $\VF$-dimension less than $k$, such that, for almost all $x \in A$,
\begin{equation}\label{cond:mg}
\omega(x) = \sigma(F(x)) + \vv(\jcb_{\VF} F(x)).
\end{equation}
We also say that such an $F$ is \memph{$\Gamma$-measure-preserving}.
\end{defn}

\begin{nota}
In \cite{hrushovski:kazhdan:integration:vf}, the category $\mgVF[k]$ is denoted by ${\mu_{\Gamma}}{\VF[k]}$ to indicate that the volume forms take values in $\Gamma$ as opposed to $\RV$. Here the subscript ``$\Gamma$'' is dropped since  we will not consider $\RV$-volume forms at all.
\end{nota}

In the definition above and other similar ones below, for the cases $k =0$, the reader should interpret things such as $\VF^0$ and how they interact with other things in a natural way. For instance, $\VF^0$ may be treated as the empty tuple, the only definable set of $\VF$-dimension less than $0$ is the empty set, and $\jcb_{\VF}$ is always $1$ on sets that have no $\VF$-coordinates. So $(A, \omega) \in \mgVF[0]$ if and only if $A$ is a finite definable set in $\RV_\infty$.

Set $\mgVF[{\leq} k] = \coprod_{i \leq k} \mgVF[i]$ and $\mgVF[*] = \coprod_{k} \mVF[k]$; similarly for the other categories below (with or without volume forms).


\begin{rem}\label{mor:equi}
Let $F : (A, \omega) \fun (B, \sigma)$ be a $\mgVF[k]$-morphism. Our intention is that such an $F$ should identify the two objects, that is, be an isomorphism. However, if $F$ is not defined everywhere in $A$ then obviously it does not admit an inverse. We remedy this by introducing the following congruence relation $\sim$ on $\mgVF[k]$. Let $G : (A, \omega) \fun (B, \sigma)$ be another $\mgVF[k]$-morphism. Then $F \sim G$ if $F(a) = G(a)$ for all $a \in A$  outside a definable subset  of $\VF$-dimension $< k$. The morphisms of the quotient category $\mgVF[k] /{\sim}$ have the form $[F]$, where $F$ is a $\mgVF[k]$-morphism. Clearly every $(\mgVF[k] / {\sim})$-morphism is an isomorphism and hence $\mgVF[k] / {\sim}$ is a groupoid. In fact, all the categories of definable sets we shall work with should be and are groupoids.

It is certainly more convenient to work with representatives than equivalence classes. In the discussion below, this quotient category $\mgVF[k] /{\sim}$ will almost never be needed except when it comes to forming the Grothendieck semigroup or, by abuse of terminology, when we speak of two objects of $\mgVF[k]$ being isomorphic.
\end{rem}


\begin{defn}[$\RV$-categories]\label{defn:c:RV:cat}
The objects of the category $\RV[k]$ are the pairs $(U, f)$ with $U$ a set in $\RVV$ and $f : U \fun \RV^k$ a definable finite-to-one function. Given two such objects $(U, f)$, $(V, g)$, any definable bijection $F : U \fun V$ is a \memph{morphism} of $\RV[k]$.
\end{defn}

By Lemma~\ref{finite:VF:project:RV}, the two categories $\VF[0]$, $\RV[0]$ are actually equivalent, similarly for other such categories.

\begin{nota}\label{0coor}
We emphasize that if $(U, f)$ is an object of $\RV[k]$ then $f(U)$ is a subset of $\RV^k$ instead of $\RV^k_\infty$, while $\infty$ can occur in any coordinate of $U$. An object of  $\RV[*]$ of the form $(U, \id)$ is often just written as $U$.

More generally, if $f : U \fun \RV_\infty^k$ is a definable finite-to-one function then $(U, f)$ denotes the obvious object of $\RV[{\leq} k]$. For example, the inclusion $\{\infty\} \fun \RV_{\infty}$ gives rise to an object of $\RV[0]$, the inclusion $\{(1, \infty)\} \fun \RV_{\infty}^2$ gives rise to an object of $\RV[1]$, and so on. Often $f$ will be a coordinate projection (every object in $\RV[*]$ is isomorphic to an object of this form). In that case, $(U, \pr_{\leq k})$ is simply denoted by $U_{\leq k}$ and its class in $\gsk \RV[k]$ by $[U]_{\leq k}$, etc.
\end{nota}

\begin{defn}[$\RES$-categories]\label{defn:RES:cat}
The category $\RES[k]$ is the full subcategory of $\RV[k]$ such that $(U, f) \in \RES[k]$ if and only if $\vrv(U)$ is finite.
\end{defn}

\begin{nota}\label{mor:ftf}
Every $\RV[k]$-morphism $F: (U, f) \fun (V, g)$  induces a definable finite-to-finite correspondence $F^\dag \sub f(U) \times g(V)$, which may also be thought of as a finitary function $f(U) \fun \mdl P(g(V))$. Recall the operator $\jcb_{\Gamma}$ from Notation~\ref{jcbG:op}. For $ u \in U$, we abbreviate $\jcb_{\Gamma}  F^\dag(f(u), g \circ F(u))$ as $\jcb_{\Gamma}  F^\dag(u)$.
\end{nota}


\begin{defn}[$\RV$- and $\RES$-categories with $\Gamma$-volume forms]\label{defn:RV:mu}
An object of the category $\mgRV[k]$ is a definable triple $(U, f, \omega)$, where $(U, f)$ is an object of $\RV[k]$ and $\omega : U \fun \Gamma$ is a function, which is understood as a \memph{definable $\Gamma$-volume form} on $(U, f)$. A \memph{morphism} between two such objects $(U,f, \omega)$, $(V, g, \sigma)$ is an $\RV[k]$-morphism $F: (U, f) \fun (V, g)$ such that for every $ u \in U$,
\begin{equation}\label{cond:mg2}
\omega(u) = (\sigma \circ F)(u) + \jcb_{\Gamma}  F^\dag(u).
\end{equation}

The category $\mgRES[k]$ is the obvious full subcategory of $\mgRV[k]$.
\end{defn}


\begin{defn}[$\Gamma$-categories]\label{def:Ga:cat}
The objects of the category $\Gamma[k]$ are the finite disjoint unions of definable subsets of $\Gamma^k$. A definable bijection between  two such objects is a  morphism of $\Gamma[k]$ if and only if it is the $\vrv$-contraction of a definable bijection.


The category $\Gamma^{\fin}[k]$ is the full subcategory of $\Gamma[k]$ such that $I \in \Gamma^{\fin}[k]$ if and only if $I$ is finite.
\end{defn}

\begin{rem}
If $\gamma \in \Gamma$ is definable then it is in the divisible hull of $\Gamma(\mathbb{S})$. This does not mean that the definable set $\gamma^{\sharp} \sub \RV$ contains a definable point unless $\gamma \in \Gamma(\mathbb{S})$.
\end{rem}


\begin{rem}\label{GLZ:char}
By \cite[Remark~2.28]{Yin:int:expan:acvf}, if a definable function between two sets in $\Gamma$ is a $\vrv$-contraction then it is piecewise $\Z$-affine (with constant terms of the form $\vrv(t)$, where $t \in \RV$ is definable, that is, $t \in \RV(\mathbb{S})$). Moreover,  a definable bijection  between two objects of $\Gamma[k]$ is a  $\Gamma[k]$-morphism if and only if it is definably a piecewise $\mgl_k(\Z)$-transformation. The ``if'' direction is clear. For the ``only if'' direction, see \cite[Lemma~10.1]{hrushovski:kazhdan:integration:vf} or \cite[Lemma~2.29]{Yin:int:expan:acvf}.
\end{rem}


\begin{defn}[$\Gamma$-categories with volume forms]\label{def:Ga:cat:mu}
An object of the category $\mG[k]$ is a definable pair $(I, \omega)$, where $I \in \Gamma[k]$ and $\omega : I \fun \Gamma$ is a function. A $\mG[k]$-morphism between two objects $(I, \omega)$, $(J, \sigma)$ is a $\Gamma[k]$-morphism $F : I \fun J$ such that for all $\alpha  \in I$,
\[
\omega(\alpha) = \sigma(F(\alpha)) + \jcb_{\Gamma} F(\alpha).
\]

The category $\mG^{\fin}[k]$ is the obvious full subcategory of $\mG[k]$.
\end{defn}

\begin{nota}\label{comp:form}
For each $(U, f, \omega) \in \mgRV[k]$, write $\omega_f : U \fun \Gamma$ for the function given by $u \efun \Sigma(\vrv \circ f)(u) + \omega(u)$. Similarly, for $(I, \sigma) \in \mG[k]$, write $\sigma_I : I \fun \Gamma$ for the function given by $\gamma \efun \Sigma \gamma + \sigma(\gamma)$.
\end{nota}

%\begin{defn}[$\Gamma$-categories with counterweights]\label{def:Ga:wgt}
%The category ${\mu_0}{\Gamma[k]}$ is the full subcategory of $\mG[k]$ such that $(I, \omega) \in {\mu_0}{\Gamma[k]}$ if and only if $\omega(\alpha) + \Sigma \alpha = 0$ for all $\alpha \in I$.
%
%The category ${\mu_0}{\Gamma^c[k]}$ is the obvious full subcategory of ${\mu_0}{\Gamma[k]}$.
%\end{defn}

There is a natural map $\Gamma[*] \fun \RV[*]$ given by $I \efun (I^\sharp, \id)$ (see Notation~\ref{gamwhat}). This map induces a commutative diagram in the category of graded semirings:
\[
\bfig
 \Square(0,0)/>->`>->`>->`>->/<400>[{\gsk \Gamma^{\fin}[*]}`{\gsk \RES[*]}`{\gsk \Gamma[*]}`{\gsk \RV[*]}; ```]
\efig
\]
where all the arrows are monomorphisms.

\begin{proof}[Proof of the horizontal monomorphisms]
It suffices to prove that an $\mathbb S$-definable bijection
$F:I^\sharp\fun J^\sharp$ in $\RV[n]$ implies that $I$ and $J$ are
isomorphic in $\Gamma[n]$, retaining all finite disjoint-union tags.
The case $n=0$ is immediate, so suppose $n>0$.

For each $\alpha\in I$, the full fibre $\alpha^\sharp$ is a torsor
for $(\K^\times)^n$ and has a unique residue-generic type. By
residue-field/value-group orthogonality, $\vrv\circ F$ has finite
image on that fibre, with its target tag included. Exactly one target
fibre $\beta^\sharp$ receives its generic part. Its image has residue
dimension $n$, since $F$ is injective, and is therefore generic in
$\beta^\sharp$. The same argument for $F^{-1}$ shows that this
assignment defines a bijection $h:I\fun J$. Definability of residue
dimension in families makes $h$ definable over $\mathbb S$.
After trivializing the two fibres solely for the proof, the generic
restriction of $F$ is a birational map of residue tori: in
characteristic zero, the constructible maps supplied by $F$ and
$F^{-1}$ are rational inverses on dense open subsets.

Fix $\alpha\in I$, put $\beta=h(\alpha)$, and set
\[
 \Delta=\vrv(\RV\cap\dcl(\mathbb S)),\qquad
 L=(\Delta+\textstyle\sum_i\Z\alpha_i+\sum_i\Z\beta_i)/\Delta.
\]
The group $L$ is finitely generated. Let $D$ be the diagonalizable
group over $\K$ with character group $L$. Its two diagonal
representations have weights $\bar\alpha_i$ and $\bar\beta_i$.
Every character $\chi:L\fun\K^\times$ extends to $\Gamma/\Delta$,
since $\K^\times$ is divisible. The resulting RV twist
$u\efun\chi(\vrv(u))u$ fixes the residue field, every $\Gamma$ value,
and every definable RV point. It preserves the RV group operations
and residue addition, so relative RV quantifier elimination makes
it preserve every $\mathbb S$-definable RV relation, including the
graph of $F$. The birational map of the two residue tori is thus
$D$-equivariant. The two actions have the same kernel, by that map
and its inverse, and their joint kernel is trivial because all their
weights generate $L$. Both representations are faithful.

By \cite[Theorem~7.1 and Corollary~2.5]{reichstein:youssin:birational},
there is a matrix $M\in\operatorname{GL}_n(\Z)$ with
$\beta-M\alpha\in\Delta^n$. Choose
$c_i\in\RV\cap\dcl(\mathbb S)$ with these values. Then
\[
 H(u)_i=c_i\prod_j u_j^{M_{ij}}
\]
is an $\mathbb S$-definable bijection of the full fibres at $\alpha$
and $\beta$. This replaces the generic birational map; it need not
extend the given map on its dense open domain. Compactness supplies
finitely many such formulas on a definable partition of $I$.
Their full-preimage maps glue to a bijection contracting to $h$,
so $h$ is a $\Gamma[n]$-morphism. This proves the lower horizontal
positive injection. The upper injection follows from the commuting
square and the two vertical positive injections.
\end{proof}


The map
\[
\gsk \RES[*] \times \gsk \Gamma[*] \fun \gsk \RV[*]
\]
determined by the assignment
\begin{equation}\label{bil:def}
([(U, f)], [I]) \efun [(U \times I^\sharp, f \times \id)]
\end{equation}
is well-defined and is clearly $\gsk \Gamma^{\fin}[*]$-bilinear. Hence it induces a $\gsk \Gamma^{\fin}[*]$-linear map
\begin{equation}\label{tensor:expre}
\Psi: \gsk \RES[*] \otimes_{\gsk \Gamma^{\fin}[*]} \gsk \Gamma[*] \fun \gsk \RV[*],
\end{equation}
which is a homomorphism of graded semirings. By the universal mapping property, groupifying a tensor product in the category of $\gsk \Gamma^{\fin}[*]$-semimodules is, up to isomorphism, the same as taking the corresponding tensor product in the category of $\ggk \Gamma^{\fin}[*]$-modules; the groupification of $\Psi$ is still denoted by $\Psi$. Similarly, there is a
$\gsk \mG^{\fin}[*]$-linear map
\begin{equation}\label{bdd:to:2bdd}
{\mu}\Psi: \gsk \mgRES[*] \otimes_{\gsk \mG^{\fin}[*]} \gsk \mG[*] \fun \gsk \mgRV[*].
\end{equation}
We shall abbreviate $\otimes_{\gsk \Gamma^{\fin}[*]}$, $\otimes_{\gsk \mG^{\fin}[*]}$ both as ``$\otimes$'' below when no confusion can arise.

\begin{prop}[Positive tensor presentation]\label{red:D:iso}
Let $T_+$ and $T_{\mu,+}$ be the sources of $\Psi$ and ${\mu}\Psi$,
respectively. All definitions below use the ground structure $\mathbb S$.
For a $\RES$ object with underlying set $X$, put
\[
 \Lambda_X=\{\delta\in\Gamma\cap\dcl^{\mathrm{eq}}(\mathbb S):
         \exists\ \mathbb S\text{-definable }c:X\to\delta^\sharp\}.
\]
Let $\mathcal C$ be the semiring congruence on $T_+$ generated by
\[
 [\bm X]_r\otimes[U]_m\ \sim\
 [\bm X]_r\otimes[U+\delta e_i]_m,
 \qquad \bm X\in\RES[r],\ U\in\Gamma[m],\ \delta\in\Lambda_X,
\]
where $m\geq1$ and $i$ is a coordinate index. Let $\mathcal C_\mu$
be the semiring congruence on $T_{\mu,+}$ generated by
\[
 [\bm X]_r\otimes[(U,\sigma)]_m\ \sim\
 [\bm X]_r\otimes[(U+\delta e_i,\sigma_\delta)]_m,
 \qquad \sigma_\delta(\eta)=\sigma(\eta-\delta e_i)-\delta,
\]
with $\bm X=(X,f,\omega)\in\mgRES[r]$ and the same admissibility
condition $\delta\in\Lambda_X$.
The homomorphisms $\Psi$ and ${\mu}\Psi$ are surjective, with kernel
congruences $\mathcal C$ and $\mathcal C_\mu$, respectively. They induce
isomorphisms of graded semirings
\[
 T_+/\mathcal C\ \simeq\ \gsk\RV[*],\qquad
 T_{\mu,+}/\mathcal C_\mu\ \simeq\ \gsk\mgRV[*].
\]
\end{prop}
\begin{proof}
The finite product normal form over the ground structure $\mathbb S$ partitions
every $\RV[k]$ object into pieces isomorphic to
$X\times U^\sharp$, where $X$ has finite $\Gamma$ support and has a
finite label of degree $k-m$, and $U\in\Gamma[m]$. For a measured
object the combined weight on a RES factor has finite
$\Gamma$-image. Partition by these values and use the shared finite
measured $\Gamma$ factor in degree zero to transfer each constant weight
to the $\Gamma$ factor. We may therefore normalize every RES combined
weight to $0$ and retain the whole weight as a definable function
$\lambda$ on its $\Gamma$ factor. This gives surjectivity of both positive
tensor maps, including finite and degree-zero factors.

For an admissible $\delta$, choose a witness $c:X\to\delta^\sharp$.
Multiplication of the $i$th free RV coordinate by $c(x)$ gives an
definable bijection between the two product lifts;
division by $c(x)$ is its inverse. Its $\Gamma$-Jacobian is $\delta$,
so the prescribed target form $\sigma_\delta$ makes the map measured.
Thus the displayed generating pairs lie in the respective kernels.
The pair depends on $\delta$, not on the choice of witness.

For the converse, work in one total degree $k$ and represent the two
positive tensors by finite disjoint unions of product objects.
Let $m$ be the largest $\Gamma$ arity on either side of a definable
RV bijection $F$ between the product unions. In total degree $k$, all
maximal-arity factors have RES degree $k-m$. We induct on $m$.

At a full-rank $\Gamma$ type in a source product and a residue type in
its RES factor, the angular normal form of the graph gives, on a
product neighbourhood,
\[
 F(x,u)=\bigl(h(x),c(x)u^M\bigr),\qquad
 M\in\operatorname{GL}_m(\mathbb Z),\qquad
 \vrv(c(x))=\boldsymbol\delta\text{ constant}.
\]
Here the $\Gamma$ rank and the total degree force the target free arity
and the RES degree to agree with those of the source. The residue
output depends only on $x$. After fixing $x$, the generic full-fibre
monomial calculation gives coefficients over $\mathbb S(x)$, which
are retained as definable functions of $x$. Injectivity forces the
square exponent matrix to be unimodular. Finite $\Gamma$ support makes
the values of those coefficients constant after a residue partition,
and their values are definable. Whole-fibre
surjectivity then forces $h$ to be injective; shrinking and taking its
image makes it a RES bijection. Product completeness and compactness
extend these formulas to definable product neighbourhoods. No
definable RV representative of $\boldsymbol\delta$ is chosen.

Compactness supplies finitely many such product neighbourhoods, with
exceptions in finitely many proper definable integral-character
hyperplanes. A branch with smaller target $\Gamma$ arity lies entirely
over such hyperplanes. Apply the same construction to $F^{-1}$ and
reverse those matches too. A common finite Boolean partition of the
source RES and $\Gamma$ factors selects disjoint product cells. The
two-sided cut property gives product images and retains every original
patch image, including the inverse cover. On each matched cell, the
integral monomial part is already a $\Gamma$ morphism and the elementary
pairs handle the coefficient translations one coordinate at a time.

Keep the matched target boxes whole and partition their complements by
a finite product grid. These are positive tensor identities on both
sides. Restricting $F$ to the actual two complements supplies the
residual comparison directly. Every nonempty residual maximal-arity
$\Gamma$ factor is covered by the finite exceptional hyperplanes.
Partition by the first applicable hyperplane. A primitive integral
equation and a unimodular coordinate change put one coordinate at a
fixed definable value. Finite-$\Gamma$ balancing transfers that
coordinate into RES, lowering $\Gamma$ arity by one and preserving total
degree. The old lower-arity factors already satisfy the smaller bound.
Apply induction to the residual tensors and add the matched congruences.
For $m=0$, all $\Gamma$ factors are finite and are absorbed into RES, where
the positive RES embedding into RV gives equality. Finally reassemble
the finitely many total degrees. This proves the unmeasured kernel
assertion without cancelling matched summands.

For the measured kernel, use the preceding normalization on both
sides. A matched rectangle then satisfies
\[
 \lambda(\gamma)=\nu(M\gamma+\boldsymbol\delta).
\]
The unimodular $\Gamma$ morphism and the elementary translated pairs
respect these complete combined weights. Retain each parent's $\Gamma$
weight on every cell of the finite product grid. On the disjoint union
of native charts, the forms glue to a definable form;
the complement inherits this form, and the restricted map $F$ preserves
it. The two residual positive sums therefore present actual measured RV
objects. In the hyperplane reduction, give the fixed finite $\Gamma$
prefix combined weight $0$ and substitute its value into the remaining
$\Gamma$ weight function. Balancing transfers its degree into RES while
retaining the full remaining weight. At rank zero absorb the actual
finite measured $\Gamma$ object with its entire form into measured RES.
The same positive induction and finite-degree reassembly prove the
measured kernel assertion.

Every step uses actual bijections, finite partitions, and positive
semiring operations. Arithmetic in the $\Gamma$-valued forms changes
object data; it does not subtract semiring classes. The quotient
isomorphisms follow from surjectivity and the kernel-congruence
equalities, without additive cancellation or groupification.
\end{proof}


Note that, however, \cite[Proposition~10.10(2)]{hrushovski:kazhdan:integration:vf} does not hold; see Remark~\ref{rem:HK:fail}. One of the main motivations of this paper is to construct an alternative to be used in applications.

\begin{nota}\label{drop:0}
For simplicity, we often drop the constant volume form $0$ from the notation. For instance, if $\bm U$ is an object of $\RV[*]$ then it may also denote the  object $(\bm U, 0)$ of $\mgRV[*]$ with the constant volume form $0$.
\end{nota}

\begin{nota}
Let $[1] \in \gsk \RES[1]$ be the class of the singleton $\{1\}$. The class of the singleton $\{1\}$ in $\gsk \RES[0]$ is the multiplicative identity of $\gsk \RES[*]$ and hence is simply denoted by $1$. We have $[\RV^{\circ\circ}_\infty] = [\RV^{\circ\circ}] + 1$ in $\gsk \RV[{\leq} 1]$.

Let $\isp$ be the (nonhomogeneous) semiring congruence relation on $\gsk \RV[*]$ generated by the pair $([1], [\RV^{\circ\circ}_\infty])$. Let
\[
\bm P = [1]-[\RV^{\ccirc}]  \in \ggk \RV[1].
\]
The corresponding principal ideal of $\ggk \RV[*]$ is thus generated by the element $\bm P - 1$.


Similarly, let $\mgisp$ be the  semiring congruence relation on $\gsk \mgRV[*]$ generated by the pair $([1], [\RV^{\circ\circ}])$, which is homogeneous, and the corresponding principal ideal of $\ggk \mgRV[*]$ is generated by the element $\bm P$.
\end{nota}


\begin{nota}\label{RVcat:tag}
For each  $\bm U = (U, f) \in  \RV[k]$, let $U_f$ be the set $\bigcup \set{f(u)^\sharp \times u \given u \in U}$. Let $\mathbb{L}_{\leq k}: \RV[{\leq}k] \fun \VF[k]$ be the map  given by $\bm U \efun U_f$. Set $\bb L = \bigcup_k \bb L_{\leq k}$.

Let $\mgL_k: \mgRV[k] \fun \mgVF[k]$ be the map given by $(\bm U, \omega) \efun (\bb L \bm U, \bb L \omega)$, where $\bb L \omega$ is the obvious function on $\bb L \bm U$ induced by $\omega$. Set $\mgL = \bigoplus_k \mgL_{k}$.
\end{nota}

Suppose that $\mathbb S$ is $(\VF,\Gamma)$-generated. The map $\bb L$ induces a surjective semiring homomorphism $\gsk \RV[*] \fun \gsk \VF_*$ and the map $\mgL$ induces a surjective graded semiring homomorphism $\gsk \mgRV[*] \fun \gsk \mgVF[*]$, see \cite[\S~4, \S~6]{hrushovski:kazhdan:integration:vf} or \cite[Corollaries~7.7, 10.6]{Yin:special:trans}; we use the same notation for these homomorphisms as well as their groupifications. We have
\[
\bb L ([1]) = [1 + \MM] = [\MM] \dand \bb L ([\RV^{\circ\circ}])=[\MM \mi 0],
\]
and hence $\bb L (\bm P - 1)=0$. Similarly, $\mgL (\bm P)=0$ since, as we have seen above, $\MM$ and $\MM \mi 0$ are in essential bijection. It so happens that these relations are the only ones needed to describe the kernels of $\bb L$ and $\mgL$.


\begin{thm}\label{main:prop:HK}
Suppose that $\mathbb{S}$ is $(\VF, \Gamma)$-generated. For each $k \geq 0$ there is a canonical isomorphism of semigroups
\[
\int_{+} : \gsk  \VF[k] \fun \gsk  \RV[{\leq}k] /  \isp
\]
such that $ \int_{+} [A] = [\bm U]/  \isp$ if and only if $[A] = [{\bb L}{\bm U}]$. These isomorphisms are  compatible with the inductive systems. Passing to  the colimit of the groupifications, we obtain a canonical isomorphism of rings
\[
\int : \ggk \VF_* \fun \ggk  \RV[*] /  (\bm P - 1).
\]
Similarly, for each $k \geq 0$ there is a canonical isomorphism of semigroups
\[
\int_{+}^\mu : \gsk  \mgVF[k] \fun \gsk  \mgRV[k] /  \mgisp
\]
such that $\int_{+}^{\mu} [\bm A] = [\bm U]/  \mgisp$  if and only if $[\bm A] = [{\mgL}{\bm U}]$.
Taking the direct sum of the groupifications yields a canonical isomorphism of graded rings
\[
\int^\mu : \ggk \mgVF[*] \fun \ggk  \mgRV[*] /  (\bm P).
\]
\end{thm}

This is a combination of two main theorems, Theorems~8.8 and 8.29, of \cite{hrushovski:kazhdan:integration:vf}.

\section{Uniform retraction to $\RES$}
\label{sec:retract}


\begin{defn}[Dimension-free $\RES$-categories]\label{df:res}
The objects of the category $\RES$ are obtained from $\RES[*]$ by forgetting the function $f$ in the pair $(U, f)$. Any definable bijection between two such objects is a morphism of $\RES$.
\end{defn}

%\begin{nota}
%Denote the half interval $(0, \infty) \sub \Gamma$ simply by $\bm H$. We can associate two Euler characteristics $\chi_{g}$, $\chi_{b}$ with the $\Gamma$-sort, which are distinguished by $\chi_{g}(\bm H) = -1$ and $\chi_{b}(\bm H) = 0$. They agree on doubly bounded definable sets. Both of them shall be denoted simply by $\chi$ when no distinction is needed. See \cite[Section 9]{hrushovski:kazhdan:integration:vf} for details, where $\chi_g$ is denoted by $\chi$ and $\chi_b$ by $\chi_c$.
%\end{nota}

\begin{rem}
As  the  $\Gamma$-sort is an \omin-minimal group, there is  the \memph{geometric} \omin-minimal Euler characteristic  $\chi_g : \ggk \Gamma[*] \fun \Z$. This is roughly constructed as follows. Choose a finite cell decomposition $X=\coprod_j C_j$ of a definable set $X\sub\Gamma^k$ and set
\[
\chi_g(X)=\sum_j(-1)^{\dim C_j}.
\]
The cell dimensions may vary, including dimension zero.
This does not depend on the chosen partition of $X$ and is invariant under definable bijections. It can also be shown that $\chi_g(X \cap [-\gamma, \gamma]^k)$ stabilizes as $\gamma$ approaches $+\infty$. The induced map, denoted by $\chi_b: \ggk \Gamma[*] \fun \Z$,  is the so-called \memph{bounded} \omin-minimal Euler characteristic. Obviously the two maps $\chi_g$, $\chi_b$  coincide on doubly bounded sets. On the other hand,  $\chi_g((0,+\infty))=-1$ whereas $\chi_b((0,+\infty))=0$. Both of them shall be denoted simply by $\chi$ when no distinction is needed.


See \cite[\S~9]{hrushovski:kazhdan:integration:vf} for more details, where $\chi_g$ is denoted by $\chi$ and $\chi_b$ by $\chi_c$.
\end{rem}


\begin{nota}\label{res:note}
Let $[\A] \in \ggk \RES$ denote the class of the affine line over the residue field. The class of the multiplicative torus $[\A] - 1$ over the residue field is written as $[\G_m]$; note that the multiplicative identity $1$ of $\ggk \RES$ is indeed the class $[1]$, but $[1]$ is not the multiplicative identity of the graded ring $\ggk \RES[*]$.

In the graded ring, $[\A]$ denotes the homogeneous degree-one class $[\G_m]+[1]\in\ggk\RES[1]$, where $[1]$ is the degree-one point class and the multiplicative identity is the degree-zero point class $1$. Every graded localization by $[\A]$ below uses this homogeneous class. After forgetting the grading, the identity is $[\A]=[\G_m]+1$.


Let $!I$ be the ideal of $\ggk \RES$ generated by the elements $[\gamma^\sharp] - [\G_m]$, where $\gamma \in \Gamma$ is definable. The quotient ring $\ggk \RES / !I$ is written as $\sggk \RES$. In the graded ring, $!I[*]$ is generated by the homogeneous degree-one differences $[\gamma^\sharp]-[\G_m]\in\ggk\RES[1]$, and we set $\sggk\RES[*]=\ggk\RES[*]/!I[*]$.

The ideal of $\ggk \mgRES[*]$ generated by the  elements
\[
[(\G_m, \alpha + \beta)] - [(\alpha^\sharp, \beta)] \in \ggk \mgRES[1],
\]
where $\alpha, \beta \in \Gamma$ are definable, is denoted by $!\mu I[*]$, and the (graded) quotient ring $\ggk \mgRES[*] / !\mu I[*]$ by $\sggk \mgRES[*]$.
\end{nota}


The quotient maps from ``$\ggk$'' to ``$\sggk$'' will all be denoted by $\iota$, and will often be omitted when it is clear in the context that the latter one is meant. We will also use the same notation for elements when passing from the former to the latter.

Note that the category $\mgRES[*]$ and the ring $\sggk \mgRES[*]$  will be dismissed after Proposition~\ref{prop:eu:retr:k} as inadequate for our purpose, and a new construction is introduced thereafter.

\subsection{Without volume forms}\label{subs:uni:retr:novol}


The discussion here more or less follows \cite[\S~2.5]{hru:loe:lef} (also see \cite[Theorem~10.5]{hrushovski:kazhdan:integration:vf}).

For $0 \leq l \leq n$, consider the maps
\[
\epsilon^{n,l}_g : \ggk \RES[n{-}l] \otimes_{\Z} \ggk \Gamma[l] \fun \sggk \RES[n]
\]
given by
\begin{equation}\label{kill:gamma}
x \otimes y \efun \chi_g(y) x [\G_m]^l.
\end{equation}
The component maps are first defined on tensor products over $\Z$. Their sum respects the homogeneous balancing relations over $\ggk\Gamma^{\fin}[*]$ after passing to $\sggk\RES[*]$: a finite definable $\Gamma$ set is a disjoint union of its definable singleton values, and moving such a singleton between the two factors changes its contribution from $[\gamma^\sharp]$ to $[\G_m]$, whose difference vanishes in the target.
The resulting tensor assignment also identifies the coefficient-translation pairs of Proposition~\ref{red:D:iso}. Indeed, if a $\RES$ coefficient $R$ admits a definable map $R\fun\delta^\sharp$, the corresponding pair translates a single coordinate of the $\Gamma$ factor by $\delta$. Both $\chi_g$ and $\chi_b$ are invariant under that definable translation. Thus the degree-$n$ sum descends through the kernel congruence of the surjection $\Psi$, giving a homomorphism
\[
\ggk \RV[n] \fun \sggk \RES[n].
\]
For later use, denote the resulting graded homomorphisms, before degree normalization, by
\[
\widetilde{\bb E}_g,\widetilde{\bb E}_b:
\ggk\RV[*]\fun\sggk\RES[*],
\qquad
\widetilde{\bb E}_{\chi}\bigl(\Psi(x\otimes y)\bigr)
=\chi(y)x[\G_m]^l
\quad(y\in\ggk\Gamma[l]).
\]
No further quotient of the $\RES$ target is required for this descent.
Composing this with the homomorphism
\[
\epsilon^{n,m}_g : \sggk \RES[n]  \fun \sggk \RES[m] : x \efun x [\A]^{m-n}
\]
and taking the direct sum of the resulting homomorphisms over all $n$ with $n \leq m$, we obtain a homomorphism
\[
\epsilon^m_g : \ggk \RV[{\leq} m] \fun \sggk \RES[m].
\]
These homomorphisms $\epsilon^m_g$, $m = 0, 1, \ldots$, are compatible in the obvious sense. Consequently, they give rise to a homomorphism
\[
\bb E_g : \ggk \RV[*] \fun \sggk \RES[*][[\A]^{-1}]
\]
whose range is precisely the zeroth piece $(\sggk \RES[*][[\A]^{-1}])_0$ of the $\Z$-graded ring on the righthand side. We have
\begin{equation}\label{eg:isp:van}
\bb E_g([\RV^{\ccirc}] + 1 - [1]) = - [\G_m][\A]^{-1} + 1 - [1][\A]^{-1} = 0.
\end{equation}
All terms have degree zero: the degree-one point maps to $[1][\A]^{-1}$. The cancellation uses the homogeneous identity $[\A]=[\G_m]+[1]$.
Thus $\bb E_g$ induces an eponymous homomorphism
\[
\bb E_g : \ggk \RV[*] / (\bm P - 1) \fun \sggk \RES[*][[\A]^{-1}].
\]

There is a parallel construction, where  we replace $\chi_g$ with $\chi_b$ and $[\A]^{-1}$ with $[1]^{-1}$. The latter replacement is needed so to make $(\bm P - 1)$ vanish as in (\ref{eg:isp:van}). The resulting homomorphism is
\begin{equation}\label{EB}
\bb E_b : \ggk \RV[*] / (\bm P - 1) \fun \sggk \RES[*][[1]^{-1}],
\end{equation}
whose range is precisely the zeroth piece  $(\sggk \RES[*][[1]^{-1}])_0$.

\begin{rem}\label{eb:no:eg}
The zeroth graded piece $(\sggk \RES[*][[\A]^{-1}])_0$ of $\sggk \RES[*][[\A]^{-1}]$ is canonically isomorphic to the colimit  of the groups $\sggk \RES[n]$ with the transition morphisms $\epsilon^{n,m}_g $, which is actually what appears in the construction above. So there is a natural homomorphism from $(\sggk \RES[*][[\A]^{-1}])_0$ to the similar-looking ring $\sggk \RES[[\A]^{-1}]$, induced by the assignment $x \efun x[\A]^{-n}$ for $x \in \sggk \RES[n]$. Note that this homomorphism is not surjective.


Similarly, there is an injective homomorphism $(\sggk \RES[*][[1]^{-1}])_0 \fun \sggk \RES[[1]^{-1}] \cong \sggk \RES$, and this is indeed an isomorphism.

If the codomain of $\bb E_b$ is changed to $\sggk \RES[[\A]^{-1}]$ in the obvious way then it may be compared with $\bb E_g$. They are different since
\begin{equation}\label{eb:neq:eg:counter}
\bb E_b([1]) = 1 \dand \bb E_g([1]) = \bb E_g([\RV^{\circ\circ}]) + 1 = -[\G_m][\A]^{-1} + 1 = [\A]^{-1}.
\end{equation}
We can equalize them by forcing $[\A] = 1$ (hence $\bb E_{g}(x \otimes y) = \bb E_{b}(x \otimes y) = 0$ for homogeneous $y \in \ggk \Gamma[l]$ with $l>0$). The (complex version of the) resulting homomorphism is indeed the one constructed in \cite[(2.5.7)]{hru:loe:lef}.
\end{rem}

The homomorphism $\bb E_b$ will be used in constructing the motivic Milnor fiber below, but not $\bb E_g$. It is not that  $\bb E_g$ would not work; it is just that, conceptually, $\bb E_b$ is a better choice,  see Remarks~\ref{eb:not:eg} and \ref{rem-comp-milnor-fib}.


%\begin{rem}\label{bbE:dim}
%In retrospect, since ``dimensions'' are lost in $\sggk \RES[[\A]^{-1}]$, in computing $\bb E_g([\bm U])$ for $\bm V = (V, g) \in \RV[*]$, we may treat $x$ in (\ref{kill:gamma}) simply as an element in $\ggk \RES$, similarly for $\bb E_b$.
%\end{rem}

\subsection{With volume forms}
Let $\fn(\Gamma, \Z)$ be the set of functions $f: \Gamma \fun \Z$ such that the range of $f$ is finite and $f^{-1}(m)$ is definable for every $m \in \Z$. Addition in $\fn(\Gamma, \Z)$ is defined in the obvious way. Multiplication in $\fn(\Gamma, \Z)$ is given by the convolution product as follows. For $f, g \in \fn(\Gamma, \Z)$ and $\gamma \in \Gamma$, let $h_\gamma$ be the function $\Gamma \fun \Z$ given by $\alpha \efun f(\alpha)g(\gamma - \alpha)$. The range of $h_\gamma$ is obviously finite. For each $m \in \Z$, $h_\gamma^{-1}(m)$ is a finite disjoint union of sets of the form
\[
f^{-1}(m') \cap (\gamma - g^{-1}(m'')),
\]
where $m' m'' = m$, and hence is $\gamma$-definable. Let
\[
 h_\gamma^* = \sum_{m} m \chi(h_\gamma^{-1}(m)) \in \Z
\]
and $f * g : \Gamma \fun \Z$ be the function  given by $\gamma \efun h_\gamma^*$. It is not hard to see that, by \omin-minimality in the $\Gamma$-sort, $f * g = g * f$ and it belongs to $\fn(\Gamma, \Z)$. Thus $\fn(\Gamma, \Z)$ is a commutative ring. Its multiplicative identity is $p_0$. Associativity is checked on the additive generators in Notation~\ref{gam:equ:cons}. Writing $rq_\beta=\eta r$, the mixed $r,q,q$ check reduces to $\eta^2=-\eta$, which holds for both $\eta=-1$ and $\eta=0$; the other generator triples follow directly from the displayed rules.


\begin{nota}\label{gam:equ:cons}
In $\fn(\Gamma, \Z)$, let $r$ be the constant function $1$ and, for each definable element $\gamma \in \Gamma$, let $p_{\gamma}$,  $q_{\gamma}$ be the characteristic functions of $\gamma$, $(\gamma, \infty)$, respectively. By \omin-minimality, $\fn(\Gamma, \Z)$ is generated as a $\Z$-module by elements of the forms $r$, $p_\gamma$, $q_\gamma$. We have the following equalities:
\[
r p_\beta   = r, \quad p_\alpha p_\beta = p_{\alpha+ \beta}, \quad p_\alpha q_\beta  = q_{\alpha+ \beta}, \quad q_\alpha q_\beta = - q_{\alpha+ \beta}.
\]
In addition,
\begin{itemize}
  \item if $\chi = \chi_{g}$ then $r^2 = r q_\beta = - r$,
  \item if $\chi = \chi_{b}$ then $r^2 = r$ and $r q_\beta = 0$.
\end{itemize}
\end{nota}

\begin{rem}\label{int:Gam}
For each $\bm I = (I, \mu) \in \mG[k]$, let $\mu_I : I \fun \Gamma$ be the function given by $\gamma \efun \mu(\gamma) + \Sigma\gamma$ and $\lambda_{\bm I} : \Gamma \fun \Z$  the function given by $\gamma \efun \chi(\mu^{-1}_I(\gamma))$. By \omin-minimality, the range of $\lambda_{\bm I}$ is finite and, for every $m \in \Z$, $\lambda_{\bm I}^{-1}(m)$ is definable. So $\lambda_{\bm I} \in \fn(\Gamma, \Z)$. If $F :(I, \mu) \fun (J, \sigma)$ is a $\mG[k]$-morphism then $F$ restricts to a bijection from $\mu^{-1}_I(\gamma)$ to $\sigma_J^{-1}(\gamma)$ for every $\gamma \in \Gamma$ and hence $\lambda_{\bm I} = \lambda_{\bm J}$ and we may write $\lambda_{\bm I}$ as $\lambda_{[\bm I]}$. Let $\bm J = (J, \sigma) \in \mG[l]$. If $k=l$ then clearly $\lambda_{[\bm I]} + \lambda_{[\bm J]} = \lambda_{[\bm I]+[\bm J]}$. An easy computation shows that, for every $\gamma \in \Gamma$,
\[
(\lambda_{[\bm I]} * \lambda_{[\bm J]})(\gamma) = \chi( (\mu \sigma)^{-1}_{I \times J}(\gamma)),
\]
where $\mu \sigma : I \times J \fun \Gamma$ is the function given by $(\alpha, \beta) \efun \mu(\alpha) + \sigma(\beta)$, and hence $\lambda_{[\bm I]} * \lambda_{[\bm J]} = \lambda_{[\bm I][\bm J]}$. Put $\Gamma_{\mathrm{def}}=\Gamma\cap\dcl^{\mathrm{eq}}(\mathbb S)$ and let $\Z[\Gamma_{\mathrm{def}}]$ be its group ring over $\Z$, viewed as the subring of $\fn(\Gamma,\Z)$ spanned by $p_\gamma$ for all $\gamma\in\Gamma_{\mathrm{def}}$. The degree-zero objects have definable constant weights, so $\ggk\mG[0]\cong\Z[\Gamma_{\mathrm{def}}]$. It follows that the map
\[
\lambda : \ggk \mG[*] \fun \Z[\Gamma_{\mathrm{def}}] \oplus X \fn(\Gamma, \Z)[X]
\]
determined by the assignment
\[
[\bm I] \efun \lambda_{[\bm I]}X^k, \quad \bm I \in \mG[k],
\]
is a graded ring homomorphism. Of course, there are two such homomorphisms $\lambda_g$ and $\lambda_b$, corresponding to the two choices $\chi = \chi_{g}$ and $\chi = \chi_{b}$.
\end{rem}


%recall Notation~\ref{comp:form}


We claim that the assignments for $k\geq1$
\begin{equation}\label{Gam:ass}
p_\gamma X^k \efun  [\G_m]^{k-1}[(\G_m, \gamma)], \quad q_\gamma X^k \efun  [\G_m]^{k-1}[(\{1\}, \gamma)], \quad r X^k \efun 0
\end{equation}
induce a graded ring homomorphism
\[
\psi : \Z[\Gamma_{\mathrm{def}}] \oplus X\fn(\Gamma, \Z)[X] \fun \sggk \mgRES[*] / ([\A]).
\]
Note that $\Z[\Gamma_{\mathrm{def}}]$ may be regarded as a subring of $\ggk \mgRES[0]$ by identifying $p_\gamma$ with $[(\infty, \gamma)]$, which is the intended interpretation of the first assignment for $k = 0$. The $q_\gamma X^k$ and $rX^k$ assignments apply only for $k\geq1$. For the claim, it suffices to check that $\psi$ obeys the equational constraints for the generators above in Notation~\ref{gam:equ:cons}, which are all straightforward (the equality $\psi(q_\alpha X^k)\psi(q_\beta X^l)=\psi(-q_{\alpha+\beta}X^{k+l})\quad(k,l\geq1)$ holds since $[\A] = [\G_m] +[1] = 0$ is forced).


Let $\gh$ denote the element $[(\RV^{\ccirc},  -\vrv)] \in \ggk \mgRV[1]$.

\begin{prop}\label{prop:eu:retr:k}
There are two graded ring homomorphisms
\[
  \mgE_{g}, \mgE_{b}: \ggk \mgRV[*] \fun   \sggk \mgRES[*] / ([\A])
\]
such that
\begin{itemize}
  \item $\bm P \in \ggk \mgRV[1]$ vanishes under both of them,
  \item for all $x \in \ggk \mgRES[*]$, $\mgE_{g}(x) = \mgE_{b}(x) = x/([\A])$,
  \item they are distinguished by $\mgE_{g}(\gh) = +[1]/([\A])$ and $\mgE_{b}(\gh) = 0$.
\end{itemize}
\end{prop}
\begin{proof}
The product map
\[
(\iota /([\A])) \times (\psi \circ \lambda) : \ggk \mgRES[*] \times  \ggk \mG[*] \fun    \sggk \mgRES[*] /([\A])
\]
is $\ggk \mG^{\fin}[*]$-bilinear (because of the ideal $!\mu I[*]$ in the target).
It also identifies the measured coefficient-translation pairs of Proposition~\ref{red:D:iso}. For a translation $\eta=\gamma+\delta e_i$ of $(I,\sigma)$, the target volume form is
\[
\sigma_\delta(\eta)=\sigma(\eta-\delta e_i)-\delta.
\]
Consequently
\[
\sigma_\delta(\eta)+\Sigma\eta=\sigma(\gamma)+\Sigma\gamma.
\]
The corrected-weight fibers are therefore definably bijective, so their two Euler characteristics agree and $\lambda$ has the same value on the pair. The balanced product map thus descends through the kernel congruence of the surjection ${\mu}\Psi$, giving the desired graded homomorphisms. To compute their values, use the explicit tensor preimages
\[
{\mu}\Psi\bigl([1]\otimes1-1\otimes[(0,\infty)]\bigr)=\bm P,
\qquad
{\mu}\Psi\bigl(1\otimes[((0,\infty),-\id)]\bigr)=\gh.
\]
These formulas use preimages under the product map.
For the value on $\gh$, the corrected weight on $((0,\infty),-\id)$ is zero, so
\[
\lambda_g([((0,\infty),-\id)])=-p_0X,
\qquad \mgE_g(\gh)=\psi(-p_0X)=-[\G_m]=+[1]\pmod{([\A])}.
\]
For $\chi_b$ the open ray has characteristic zero, giving $\mgE_b(\gh)=0$.

Finally, $[1]$ is nonzero in the target. After forgetting degree and volume forms, choose auxiliary $\RV$ representatives for the finitely many value fibers of each $\RES$ object. Normalization and quantifier elimination give pure residue-field definable sets, and hence a homomorphism to the Grothendieck ring of pure $\K$-definable sets modulo the affine-line class $\mathbb L$. The fiber relations vanish under this homomorphism, $[\A]$ maps to $\mathbb L$, and $[1]$ maps to $1$. The latter quotient is nonzero: any proposed equality $1=0$ uses only finitely many formulas, parameters and definable bijections. By the Lefschetz principle these data specialize to an algebraic closure of a finite field, with all parameters in some $\mathbb F_q$. A definable bijective relation over $\mathbb F_q$ induces a bijection on its $\mathbb F_q$-points. Counting these points modulo $q$ preserves the defining relations, sends $\mathbb L$ to $0$ and sends $1$ to $1$, a contradiction. Thus the two displayed values on $\gh$ are distinct.
\end{proof}

The discussion above is more or less an expanded version of the proof of \cite[Theorem~10.11]{hrushovski:kazhdan:integration:vf}.

\begin{rem}
Alternatively, for $\chi = \chi_{g}$, we may replace $r X^k \efun 0$ in (\ref{Gam:ass}) with $r X^k \efun [\G_m]^{k-1}[1]$, that is, sending $r X^k$ and $q_0 X^k$ to the same element.  Let $J$ be the enlargement of $!\mu I[*]$ by the elements $[1]-[(\{1\},\gamma)]\in\ggk\mgRES[1]$. Then $\psi$ is still a graded ring homomorphism with this enlarged quotient. Consequently, for instance, we can construct a ring homomorphism
\begin{equation}\label{res:forget}
\ggk \mgRV[*] \fun \ggk \mgRES[*] / (J+([\A])) \to^{\phi} \sggk \RES[*] / ([\A]),
\end{equation}
where the homomorphism $\phi$ is induced by the obvious forgetful functor.


We shall not, however, pursue this further, since taking the quotient by the ideal $([\A])$ is detrimental to our purpose here. To avoid it, we need to work with the categories of doubly bounded objects.
\end{rem}

\begin{defn}[Doubly bounded $\RV$-categories]\label{db:RV}
An object $(U, f) \in \RV[k]$ is \memph{doubly bounded} if $U$ is. Denote the full subcategory of $\RV[k]$ of doubly bounded objects by $\RV^{\db}[k]$, and similarly for the categories $\mgRV^{\db}[k]$, $\mG^{\db}[k]$.
\end{defn}


\begin{rem}\label{tensor:bd}
By Lemma~\ref{db:to:db}, if $(U, f, \omega) \in \mgRV^{\db}[k]$ then $f(U)$, $\omega(U)$ are both doubly bounded as well.
The product construction gives a homomorphism
\[
{\mu}\Psi^{\db}:\gsk\mgRES[*]\otimes\gsk\mG^{\db}[*]
\fun\gsk\mgRV^{\db}[*].
\]
We use it below to compute images of explicit tensor classes. The bounded retraction itself is obtained by restricting the already constructed graded map $\widetilde{\bb E}_b$ after forgetting volume forms.
\end{rem}

\begin{nota}\label{db:generator}
Let $\fn^{\db}(\Gamma,\Z)$ be the subgroup of $\fn(\Gamma,\Z)$ consisting of functions whose nonzero locus is doubly bounded. For each definable $\gamma\in\Gamma$, set
\[
o_\gamma=\begin{cases}
\mathbf1_{(0,\gamma)},&\gamma>0,\\
\mathbf1_{(\gamma,0)},&\gamma<0,\\
-p_0,&\gamma=0.
\end{cases}
\]
As a $\Z$-module, $\fn^{\db}(\Gamma,\Z)$ is generated by $p_\gamma,o_\gamma$. For $\gamma>0$ we have $o_\gamma=q_0-q_\gamma-p_\gamma$, and for $\gamma<0$ we have $o_\gamma=q_\gamma-q_0-p_0$. The convention $o_0=-p_0$ is a convention for functions; it does not identify $o_0$ with the characteristic function of the empty interval $(0,0)$. The convolution products are given by the following cases.
For $\beta\ne0$, put $a=\min(\alpha,\alpha+\beta)$ and $b=\max(\alpha,\alpha+\beta)$. Then $p_\alpha o_\beta$ is the characteristic function of $(a,b)$, and
\[
p_\alpha o_\beta=\begin{cases}
o_b-o_a-p_a,&0\leq a<b,\\
o_a-o_b-p_b,&a<b\leq0,\\
o_a+o_b+p_0,&a<0<b.
\end{cases}
\]
For $\beta=0$, $p_\alpha o_0=-p_\alpha$. Also,
\[
o_\alpha o_\beta=\begin{cases}
-o_{\alpha+\beta},&\alpha,\beta\geq0\text{ or }\alpha,\beta\leq0,\\
-o_\alpha-o_\beta-p_0,&\alpha<0<\beta\text{ or }\beta<0<\alpha.
\end{cases}
\]
\end{nota}


\begin{nota}\label{db:gam:fn}
The homomorphism
\begin{equation*}
\lambda : \ggk \mG^{\db}[*] \fun \Z[\Gamma_{\mathrm{def}}] \oplus X \fn^{\db}(\Gamma, \Z)[X]
\end{equation*}
is constructed as before (there is only one such homomorphism now since  $\chi_{g}$, $\chi_{b}$ agree on doubly bounded sets).
\end{nota}


\begin{nota}\label{pgammabis}
Recall that $\RV^{\ccirc}_\gamma = \rv(\MM_{\gamma} \mi 0)$ (Notation~\ref{defn:disc}). Let $D_{\RV}$ consist of the nonnegative definable $\gamma$ such that $\gamma^\sharp$ contains an $\mathbb S$-definable point. For each $\gamma\in D_{\RV}$, let
\[
\bm P_\gamma = [\RV^{\ccirc} \mi \RV^{\ccirc}_{\gamma}] + [\{t_\gamma\}] - [1] \in \ggk \RV^{\db}[1],
\]
where $t_\gamma \in \gamma^\sharp$ is any definable point.  It also stands for the corresponding element  in $\ggk \mRV^{\db}[1]$ (with the constant volume form $0$). Note that, although $[\{t_\gamma\}] = [1]$ in $\ggk \RV^{\db}[1]$,  $[\{t_\gamma\}] \neq [1]$ in $\ggk \mRV^{\db}[1]$  unless $\gamma = 0$.


Clearly  $\bm P_\gamma$ does not depend on the choice of $t_\gamma \in \gamma^\sharp$. The ideal of $\ggk \mRV^{\db}[*]$   generated by the elements $\bm P_\gamma$ with $\gamma\in D_{\RV}$ is denoted by $(\bm P_\Gamma)$. A definable $\Gamma$ value need not have a definable RV lift, as already noted after Definition~\ref{def:Ga:cat}. All choices are made over the ground structure $\mathbb S$. The images of $(\bm P_\Gamma)$ under the obvious forgetful homomorphisms
\[
\ggk  \mgRV^{\db}[*] \fun \ggk \RV[*], \quad \ggk  \mgRV^{\db}[*] \fun \ggk  \mgRV[*]
\]
are contained in $(\bm P - 1)$, $(\bm P)$, respectively.
\end{nota}

\begin{rem}\label{db:gam:res}
Let $!P$ denote the ideal of $\fn^{\db}(\Gamma, \Z)$ generated by the elements $p_\gamma - p_0$, $o_\gamma + p_0$. By the computation in Notation~\ref{db:generator}, $\fn^{\db}(\Gamma, \Z) / !P \cong \Z$; the quotient homomorphism is denoted by $\chi$ because a simple computation shows that if $\bm I = (I, \mu) \in \mG^{\db}[*]$ then indeed $\lambda_{\bm I} / !P = \chi(I)$. Now it is straightforward to check that the assignments
\begin{equation}\label{Gam:ass:bd}
p_\gamma X^k \efun  [\G_m]^{k}, \quad o_\gamma X^k \efun  - [\G_m]^{k}
\end{equation}
induce a graded ring homomorphism
\[
\psi^{\db} : \Z[\Gamma_{\mathrm{def}}] \oplus X\fn^{\db}(\Gamma, \Z)[X] \fun \sggk \RES[*].
\]
In a nutshell, the  homomorphism $\psi^{\db} \circ \lambda$ is a ``forgetful'' map  given by $[(I, \mu)] \efun \chi(I)[\G_m]^{k}$, $(I, \mu) \in \mG^{\db}[k]$.
\end{rem}


The forgetful homomorphism $\phi : \ggk \mgRES[*] \fun \sggk \RES[*]$ is similar to the one in  (\ref{res:forget}).

\begin{prop}\label{prop:eu:retr:k:db}
There is a graded ring homomorphism
\[
  \mgE^{\db}: \ggk \mgRV^{\db}[*] \fun \sggk \RES[*]
\]
such that $(\bm P_\Gamma)$ vanishes and, for all $x \in \ggk \mgRES[*]$, $\mgE^{\db}(x) = \phi(x)$.
\end{prop}
\begin{proof}
Let $F:\ggk\mgRV^{\db}[*]\fun\ggk\RV[*]$ be the forgetful homomorphism and set
\[
\mgE^{\db}=\widetilde{\bb E}_b\circ F.
\]
This is graded, and its restriction to $\ggk\mgRES[*]$ is $\phi$. On the bounded tensor product its assignment is $\phi\times(\psi^{\db}\circ\lambda)$, by Remark~\ref{db:gam:res}. Thus it remains to compute the radius generators.
If $\gamma=0$, the annulus is empty and the value-zero singleton classes $[\{t_0\}]$ and $[1]$ agree, so $\bm P_0=0$. For an indexed radius $\gamma>0$, we have

\[
{\mu}\Psi^{\db}\bigl(1\otimes[(0,\gamma)]+1\otimes[\{\gamma\}]
+[\{t_\gamma\}]\otimes1-[1]\otimes1\bigr)=\bm P_\gamma.
\]
Its image under $\phi\otimes(\psi^{\db}\circ\lambda)$ is
\[
\chi((0,\gamma))[\G_m]+\chi(\{\gamma\})[\G_m]+[1]-[1]
=-[\G_m]+[\G_m]=0.
\]
The open-interval-plus-endpoint expansion is used only for $\gamma>0$.
\end{proof}

The composition of $\mgE^{\db}$ with the obvious forgetful homomorphism $\sggk \RES[*] \fun \sggk \RES$ is denoted by $\bb E^{\diamond}$.  By Remark~\ref{db:gam:res}, the diagram commutes:
\begin{equation}\label{edb:eb:com}
\bfig
 \Vtriangle(0,0)/->`->`->/<400, 400>[{\ggk \mgRV^{\db}[*]}`{\ggk \RV[*]}`\sggk \RES; `\bb E^{\diamond}`\bb E_b]
\efig
\end{equation}

Here the arrow labelled $\bb E_b$ means the quotient map $\ggk\RV[*]\fun\ggk\RV[*]/(\bm P-1)$, followed by~\eqref{EB} and the degree-zero point-localized identification in Remark~\ref{eb:no:eg}.
This provides an alternative and more direct (but less comprehensive) construction of $\bb E^{\diamond}$.

\begin{rem}\label{eb:not:eg}
Changing $\sggk \RES$ to $\sggk \RES[[\A]^{-1}]$ in (\ref{edb:eb:com}), we see that  $\bb E_g$ does not commute with $\bb E^{\diamond}$: by (\ref{eb:neq:eg:counter}),
\[
\bb E_g([1]) = [\A]^{-1} \quad \text{whereas} \quad \bb E^{\diamond}([1]) = \bb E_b([1]) = 1.
\]
This is but an example of a general phenomenon, see Remark~\ref{eg:eb:factor}.


We can certainly redefine $\bb E^{\diamond}$ as the composition of $\mgE^{\db}$ with  the homomorphism $\sggk \RES[*] \fun \sggk \RES[[\A]^{-1}]$ implicit in Remark~\ref{eb:no:eg} and thereby make it compatible with $\bb E_g$, but as we have mentioned above, $\bb E_b$ is ultimately a better choice; see Remark~\ref{rem-comp-milnor-fib}.
\end{rem}


\begin{nota}[Relative annular ideal]\label{relative:annular:ideal}
Put $B=\ggk\mgRV^{\db}[*]$ and retain the scalar ideal
$I=(\bm P_\Gamma)$ of Notation~\ref{pgammabis}, with the same set
$D_{\RV}$ of radii. For
$\bm U=(U,f,\omega)\in\mRV^{\db}[k]$, $k>0$, suppose that
$\pr_{\widetilde j}\rest f(U)$ is finite-to-one. Let
$\tau:U\to\RV$ be definable with
$0\leq\gamma(u):=\vrv(\tau(u))\leq\Delta$ for one finite named
$\Delta$. Define the relative annular object by the separately tagged
union
\[
 U^{\tau,j}=\{(u,s):0<\vrv(s)\leq\gamma(u)\}
 \ \uplus\ \{(u,\tau(u)):u\in U\}.
\]
Its structural coordinates are $sf_j(u)$ in position $j$ and $f_i(u)$
in every other position; its volume form is $\omega(u)$.
Write $J_{\mathrm{var}}$ for the ideal of $B$ generated by
$[\bm U^{\tau,j}]-[\bm U]$ for all these data. The full aperture graph
is retained, including when its image is infinite. Constant apertures
give $I\subseteq J_{\mathrm{var}}$, but no equality is asserted.
\end{nota}

\begin{lem}[Compatibility with the relative annular ideal]\label{relative:annular:compatibility}
The volume-forgetting map $F:B\to\ggk\RV[*]$ kills
$J_{\mathrm{var}}$. Hence the bounded retractions of this section
factor through $B/J_{\mathrm{var}}$, with the same restrictions
to the residue rings. The ideal is homogeneous and meets the entire
degree-zero subring trivially. Its image in the full measured RV ring
lies in $(\bm P)$.
\end{lem}
\begin{proof}
Retain the actual graph of $\tau$; projection onto $U$ is a definable
isomorphism and preserves the required bounds. Put
\[
 \bm W=(U,f_{\widetilde j},\omega+\vrv(f_j)),\qquad
 \bm W_\gamma=(U,f_{\widetilde j},\omega+\vrv(f_j)+\gamma).
\]
The inactive structural map is finite-to-one, so these are bounded
objects of degree $k-1$. Equality of total weights identifies the
source with $\bm W$ times the structural unit point, and the endpoint
with $\bm W_\gamma$ times that point. The shell has structural map
$(s,f_{\widetilde j})$ and form $\omega+\vrv(f_j)$.

After forgetting volume, the full positive-ray bundle
$C_+=\{(u,s):\vrv(s)>0\}$ is isomorphic to its subobject
$C_{\mathrm{deep}}=\{(u,z):\vrv(z)>\gamma(u)\}$ by
\[
 (u,s)\longmapsto(u,\tau(u)s),
\]
with inverse division by $\tau(u)$. Its complement is the whole
relative shell. Scissors therefore give zero for the shell in the
full ordinary RV ring. The endpoint graph is ordinarily isomorphic
to the source. Thus $F$ kills every generator, and hence the ideal,
before any Euler evaluation. This includes depth zero.

Every generator has positive homogeneous degree, so the degree-zero
projection kills the ideal and proves injectivity on degree zero.
In the measured ring the same shear identifies the positive ray
$C_+(\bm W_\gamma)$ with the deep subray of $C_+(\bm W)$:
the structural shift $\gamma$ equals the shift in the base form.
With $\bm P=[C_+]-[1]$, scissors give the exact identity
\[
 \operatorname{image}([\bm U^{\tau,j}]-[\bm U])
   =([\bm W_\gamma]-[\bm W])\bm P.
\]
This proves the measured ideal containment. Composing ordinary
forgetting with the already constructed Euler maps proves the stated
retraction factorization and leaves their residue restrictions unchanged.
\end{proof}

The pure RV tensor presentation and its intrinsic coefficient-translation
kernel are unchanged. If $q:T\twoheadrightarrow B$ is that presentation,
the relative annular quotient is $T/q^{-1}(J_{\mathrm{var}})$, and
$B/I\to B/J_{\mathrm{var}}$ is the canonical further quotient.
No factorwise scalar-radius presentation of this new quotient is used.

\section{Interlude: RV decomposition and the realization in varieties}\label{csndag}
The results recalled in this section are proved in \cite{FY:Fubini}, together with Fubini theorems for the motivic volume that are not used in this paper. They are used in Section~\ref{sec:spec:pui}.

For an integer matrix $M=(m_{ab})\in\GL_n(\Z)$, write
\[
 h_M:(\RV^\times)^n\longrightarrow(\RV^\times)^n,\qquad
 (h_M(u))_a=\prod_b u_b^{m_{ab}}.
\]
Thus $h_M^{-1}=h_{M^{-1}}$ and $\vrv(h_M(u))=M\vrv(u)$. For $I\subseteq\Gamma^\ell$, the notation $I^\sharp$ means its full inverse image in $(\RV^\times)^\ell$ under $\vrv$.

\begin{lem}[Monomial decomposition {\cite[Lemma~2.1]{FY:Fubini}}]\label{coh:decom}
Let $U\subseteq(\RV^\times)^n$ be definable. There is a finite definable partition $(D_i)_i$ of $\vrv(U)$ such that, writing $U_i=U\cap\vrv^{-1}(D_i)$, there are $M_i\in\GL_n(\Z)$, $0\leq m_i\leq n$, a nonempty definable set $V_i\subseteq(\RV^\times)^{m_i}$, and a nonempty definable set $I_i\subseteq\Gamma^{n-m_i}$ with
\[
 \vrv(V_i)=\{\beta_i\},\qquad
 h_{M_i}(U_i)=V_i\times I_i^\sharp .
\]
All sets and maps are definable over the parameter base $\mathbb S$.
\end{lem}

\begin{nota}[Equivariant residue realization]\label{bipo:gdv}
In this notation we work in the field $\puC$ of complex Puiseux series and take $\mathbb S = \C\dpar t$; the value group is identified with $\Q$. Write $\ggk^{\widehat\mu,\mathrm{raw}}\var_{\C}$ for the Grothendieck ring of complex varieties with good $\widehat\mu$-actions, where \emph{good} means that the action factors through some $\mu_N$ and every orbit is contained in an affine open subset. Write $\gmv$ for its quotient by the ideal generated by
\[
 [X\times(\A^r,\rho)]-[X\times\A^r],
\]
for every variety $X$ with a good $\widehat\mu$-action and every linear action $\rho$ of some $\mu_N$ on $\A^r$, where the action on the second affine factor on the right is trivial. For $q\in\Q$, let $L_q=q^\sharp\sqcup\{0_q\}$ be the one-dimensional residue vector space whose nonzero part is $q^\sharp$; its class in $\sggk\RES$ is $[\G_m]+1=[\A]$. By \cite[Notation~7.3]{FY:Fubini}, there is a canonical isomorphism
\[
 \Theta:\sggk\RES\xrightarrow{\ \sim\ }\gmv,
\]
constructed with the finite Kummer actions of $\widehat\mu$ on the points of $\RES$ objects that are rational over the fields $\C\dpar{t^{1/N}}$. If $h_{M_i}(U_i)=V_i\times I_i^\sharp$ is a monomial decomposition of an object $\bm U$ of $\RV[*]$ as in Lemma~\ref{coh:decom}, with $\ell_i$ free $\RV$ coordinates, then \cite[(7.1)]{FY:Fubini}
\begin{equation}\label{twist:C}
 (\Theta\circ\bb E_b)([\bm U])
   =\sum_i\chi_b(I_i)\,\Theta([V_i])[\G_m]^{\ell_i}.
\end{equation}
This is an evaluation of the retraction $\bb E_b$ of Section~\ref{sec:retract}, so it does not depend on the decomposition.
\end{nota}

\begin{nota}\label{mov:vol}
The motivic volume operator is
\[
 \Vol=\Theta\circ\bb E_b\circ\int:
 \ggk\VF_*\longrightarrow\gmv .
\]
Here $\bb E_b$ is applied after the quotient by $(\bm P-1)$, on which it is defined by Section~\ref{sec:retract}.
\end{nota}


\section{Special covariant bijections}\label{sec-pscb}


From here on, we also assume that the substructure $\mathbb{S}$ of parameters is $\VF$-generated. So every definable disc contains a definable point (Lemma~\ref{algebraic:balls:definable:centers}). Another simple fact we shall occasionally use is that if $\gamma \in \Gamma$ is definable then there is a $\gamma' \geq \gamma$ such that $\gamma'^\sharp$ contains a definable point. Indeed, after replacing $\gamma$ by $\max\{\gamma,0\}$, the punctured ball $\{x\ne0:\vv(x)>\gamma\}$ contains a finite nonempty definable set $E$: elementarity of $\acl(\mathbb S)$ gives an algebraic point, and its finite definable cover may be intersected with this ball. The unordered product $c=\prod_{x\in E}x$ is definable, is nonzero, and has $\vv(c)>\gamma$. Take $\gamma'=\vv(c)$ and the definable point $\rv(c)\in\gamma'^\sharp$.

\subsection{Invariance and covariance}\label{sec:inv:cov}

\begin{defn}\label{defn:binv}
For $\gamma \in \Gamma$, let $\pi_\gamma: \VF \fun \VF / \MM_{\gamma}$ be the quotient  map. If $\gamma = (\gamma_1, \ldots, \gamma_n) \in \Gamma^n$ then $\pi_{\gamma}$ denotes the product of the maps $\pi_{\gamma_i}$.

Let $\alpha \in \Gamma^n$ and $\beta \in \Gamma^m$. We say that a function $f :A \fun B$ with $A \sub \VF^n$ and $B \sub \VF^m$  is \memph{$(\alpha, \beta)$-covariant} if it $(\pi_{\alpha}, \pi_{\beta})$-contracts to a function $f_\downarrow : \pi_{\alpha}(A) \fun \pi_{\beta}(B)$, that is, $\pi_{\beta} \circ f = f_\downarrow \circ \pi_{\alpha}$. For simplicity, we shall often suppress parameters in notation and refer to $(\alpha, \beta)$-covariant functions as $(\alpha, -)$-covariant or $(-,\beta)$-covariant or just  covariant functions.

A set  in $\VF$ is \memph{$\alpha$-invariant} if its characteristic function is $(\alpha, 0)$-covariant.


More generally, for sets $A$, $B$ with $\RV$-coordinates, a function $f : A \fun B$ is \memph{covariant} if every one of its $\VF$-fibers    is covariant (recall Terminology~\ref{rvfiber}). Similarly, a set is \memph{invariant} if every one of its $\VF$-fibers is invariant.
\end{defn}


\begin{defn}\label{prop:pseu}
A definable function $f : A \fun B$ is \memph{proper covariant} if
\begin{itemize}
  \item the sets $A_{\VF}$, $f(A)_{\VF}$ are bounded and the sets $A_{\RV}$, $f(A)_{\RV}$ are  doubly bounded,
  \item for each $\VF$-fiber $f_{t}$ of $f$ there is a $t$-definable tuple $(\alpha_t, \beta_t) \in \Gamma$ such that $f_t$ is $(\alpha_{t}, \beta_{t})$-covariant, $\dom(f_{t})$ is $\alpha_t$-invariant, and $\ran(f_{t})$ is $\beta_t$-invariant.
\end{itemize}
%Accordingly, an invariant set $A$ is \memph{proper} if its projection into the $\RV$-coordinates is proper covariant, in particular, $A_{\VF}$ is bounded and $A_{\RV}$ is doubly bounded.
A definable set $A$ is \memph{proper invariant}   if $A_{\VF}$ is bounded, $A_{\RV}$ is doubly bounded, and for each $\VF$-fiber $A_t$ of $A$ there is a $t$-definable tuple $\alpha_t\in \Gamma$ such that $A_t$ is $\alpha_t$-invariant.
\end{defn}

Covariance also makes sense when the domain or range has no $\VF$-coordinates: the corresponding tuple space is a singleton. In particular, a contraction with no output $\VF$-coordinates is the unique map to the singleton empty tuple; it is empty only when its domain is empty. For a definable map $f:A\fun B\sub\RV^m$, proper covariance is equivalent to proper invariance of all level sets, together with the global conditions that $A_{\VF}$ is bounded and $A_{\RV}$ and $f(A)$ are doubly bounded. For a powerset-valued map, use its definable membership relation $R_f=\{(a,s):s\in f(a)\}$: require these same global bounds on $A$ and on the RV-support of $R_f$, and proper invariance of every membership section $\{a:s\in f(a)\}$. With these boundedness conditions retained, $A$ is proper invariant precisely when its projection to its existing $\RV$-coordinates is proper covariant.


\begin{exam}
A proper invariant set  is not required to be doubly bounded. However, if $\bm U \in \RV^{\db}[*]$ then $\bb L \bm U$ is indeed a doubly bounded proper invariant set.
\end{exam}

\begin{rem}
In defining proper covariance above, if  $A$ is closed and $\infty$ does not occur in the $\RV$-coordinates of $f(A)$ then  there is no need to demand that $f(A)_{\RV}$ be doubly bounded, because it is implied by the other conditions, as follows. There is an $\go$-partition $p: A \fun \Gamma$ such that, for each open polydisc $\gb$ in question, that is, for each open polydisc $\gb \sub A$ of the form $\go(a, p(a))$, $f(\gb)$ lies in the range of some $\VF$-fiber of $f$. By Lemma~\ref{vol:par:bounded}, $p(A)$ is doubly bounded. So the definable surjection $A \fun \vrv(f(A)_{\RV})$ induced by $f$ satisfies the assumption of Lemma~\ref{invar:db} below. So  $f(A)_{\RV}$ is doubly bounded.
\end{rem}

\begin{lem}\label{inv:crit}
Let $A$ be a definable set such that $A_{\VF}$ is bounded and $A_{\RV}$ is doubly bounded. Then $A$ is proper invariant if and only if $A$ is clopen.
\end{lem}
\begin{proof}
The ``only if'' direction follows since each invariant fiber and its complement are unions of open polydiscs. For the converse, on each nonempty fiber put
\[
p(a,t)=\min\{\gamma\in\Gamma:0\leq\gamma,\ \go(a,\gamma)\sub A_t\}.
\]
Openness makes this a nonempty upper ray. It contains its finite endpoint: every point of $\go(a,\gamma)$ belongs to $\go(a,\delta)$ for some $\delta>\gamma$, by density of $\Gamma$ and finiteness of the VF-tuple. Definable completeness therefore gives the stated minimum. Ultrametricity shows that it is constant on its own polydiscs, so it is an $\go$-partition. The lower cutoff also handles full fibers and zero VF-arity. Since the fiber is closed and bounded, Lemma~\ref{vol:par:bounded} bounds the partition radii above. Their supremum is $t$-definable, and is an invariance radius for the fiber. Empty fibers may use radius $0$. Thus all the required radii have the specified parameter base, and $A$ is proper invariant.
\end{proof}


\begin{lem}\label{int:inv}
If $A$, $B$ are proper invariant subsets of $\VF^n \times \RV^m$ then $A \cap B$, $\RVH(A) \mi A$, and $B \mi A$ are also proper invariant.
\end{lem}
\begin{proof}
First, $\RVH(A)$ is proper invariant. Fix an existing RV-tag $t$ and a scalar radius $\delta$ for which $A_t$ is $\delta$-invariant. An RV-polydisc meeting $A_t$ either lies, in each coordinate whose RV-disc has radius at least $\delta$, in the corresponding radius-$\delta$ disc already present in $A_t$, or has an RV-disc of radius less than $\delta$, which is itself $\delta$-invariant. More explicitly, if $x\in\RVH(A)_t$ and $\vv(y-x)>\delta$ coordinatewise, start with $a\in A_t$ with $\rv(a)=\rv(x)$. In coordinates with $\vv(x_i)\leq\delta$, one has $\rv(y_i)=\rv(x_i)$; in the remaining coordinates, both $a_i$ and $y_i$ lie in $\MM_\delta$. Replace $a_i$ by $y_i$ only in these remaining coordinates. The new tuple belongs to $A_t$ by $\delta$-invariance and has the same RV-tuple as $y$. Thus $y\in\RVH(A)_t$. The hull preserves the coordinate values and all existing RV-tags, so it preserves the required bounds. Intersections and differences of clopen fibers are clopen, and the three sets in the statement inherit the bounds from $A$, $B$, or $\RVH(A)$. Lemma~\ref{inv:crit} now applies.
\end{proof}

\begin{lem}\label{conti:covar}
Let $f : A \fun B$ be a definable continuous surjection between two proper invariant sets such that  every $\dom(f_t)$ is open. If every $\ran(f_t)$ is clopen then $f$ is proper covariant. If $f$ is bijective then the same holds under the weaker assumption that   every $\ran(f_t)$ is open.
\end{lem}
\begin{proof}
For every $(r, s) \in f_{\RV}$ and every $b \in \ran(f_{r,s})$, let
\[
q(b,r,s) = \min\set{\beta \in \Gamma \given 0\leq\beta,\ \go((b,s), \beta) \sub \ran(f_{r,s})}.
\]
The definable function $q$ is clearly an $\go$-partition of its domain, which, by the assumption, is closed. By Lemma~\ref{vol:par:bounded}, $\ran(q)$ is doubly bounded by, say, $\beta \in \Gamma^+$. Choose the source partition radii so that their open polydiscs lie in one fixed $\dom(f_{r,s})$ as well as in $A$. This is possible since these domains are open. Continuity on that graph fiber then permits a further shrinking such that its image lies in an open polydisc of radius $\beta$. Thus we obtain another $\go$-partition $p:A\fun\Gamma$ with both requirements. By Lemmas~\ref{inv:crit} and \ref{vol:par:bounded} again,  $p(A)$ is doubly bounded as well. It follows that $f$ is  proper covariant.

For the second claim, since $f$ is bijective, for each $(b, s) \in B$ there is only one  $f_{r,s}$ with $(b,r, s) \in \ran(f_{r,s})$. Thus $r=r(b,s)$ is a definable function on $B$. Since the sets $\ran(f_{r,s})$ are open and form a partition, this function is locally constant on each input VF-fiber. Consequently $q(b,r(b,s),s)$ is an $\go$-partition of the closed set $B$. The rest of the argument is the same.
\end{proof}


\begin{lem}\label{conv:unif}
If $f : A \fun B$ is a proper covariant surjection then both $A$ and $B$ are proper invariant. If, in addition, $f$ is continuous and bijective then $f$ is a proper covariant homeomorphism, which just means that $f^{-1}$ is also proper covariant and continuous. Moreover, in that case, if $A'$ is a proper invariant subset of $A$ then $f \rest A'$ is a proper covariant homeomorphism as well.
\end{lem}
\begin{proof}
Since $f_{\RV}$ is  doubly bounded, by Lemma~\ref{db:to:db}, for all sufficiently large  $\epsilon \in \Gamma$, every $\VF$-fiber $f_{t}$ of $f$ is $(\epsilon, \beta_{t})$-covariant (but not necessarily $(\epsilon, \epsilon)$-covariant) and $\dom(f_{t})$,  $\ran(f_{t})$ are both $\epsilon$-invariant. The first claim follows.

For the second claim, since $\dom(f_{t})$ is clopen, by  Corollary~\ref{conti:homeo}, every $f_t$ is a homeomorphism. Since  $\ran(f_{t})$ is open, it follows from the continuity of $f^{-1}$, as in the proof of Lemma~\ref{conti:covar}, that every  $f^{-1}_t$ is  $(\delta, \epsilon)$-covariant  for all sufficiently large $\delta > \epsilon$.

For the third claim, clearly every $\dom(f_t) \cap A'$ is   clopen and, since $f_t$ is a homeomorphism, every $\ran(f_t) \cap f(A')$ is  clopen too. This is sufficient for the proof of (the first claim of) Lemma~\ref{conti:covar} to go through. By the second claim, then, $f \rest A'$ is a proper covariant homeomorphism.
\end{proof}

\begin{lem}\label{pr:co:home:com}
If $f : A \fun B$ and $g : B  \fun C$ are proper covariant homeomorphisms then $g \circ f$ is also a proper covariant homeomorphism.
\end{lem}
\begin{proof}
By Lemma~\ref{conv:unif}, we just need to show that $g \circ f$ is proper covariant. Then, by Lemma~\ref{conti:covar}, we are reduced to showing that the domain and range of every $\VF$-fiber of $g \circ f$ are open. Arguing as in the proofs of Lemmas~\ref{conti:covar} and \ref{conv:unif}, we see that there are $\epsilon, \epsilon', \epsilon'' \in \Gamma$, each one  sufficiently larger than the next, such that every $\VF$-fiber of $f$ is $(\epsilon, \epsilon')$-covariant and every $\VF$-fiber of $g$ is $(\epsilon', \epsilon'')$-covariant. It follows that the domain of every $\VF$-fiber of $g \circ f$ is $\epsilon$-invariant and hence is clopen. Since the same holds for $(g \circ f)^{-1}$, the range of every $\VF$-fiber of $g \circ f$ is clopen as well.
\end{proof}

\begin{conv}\label{conv:can}
The trivial-looking operation we describe here is actually quite crucial for understanding the discussion below, especially those passages that involve special covariant bijections (defined below). For each definable set $A$, let
\[
\can(A) = \set{(a, \rv(a), t) \given (a, t) \in A \tand a \in A_{\VF}}.
\]
The obvious definable bijection $\can : A \fun \can(A)$ is  referred to as the \memph{regularization} of $A$.

Observe that, for a proper invariant set $A$, $\can(A)$ is proper invariant if and only if $A_{\VF}$ is  doubly bounded. Going forward, when $A$ is doubly bounded and proper invariant, it is often more convenient or even necessary to substitute $\can(A)$ for $A$. Whether such a substitution has been performed (tacitly) or not should be clear in the context (or rather, it is always performed).
\end{conv}
%\begin{defn}
%Let $A$ be a definable set. A \memph{regularization} map is a map of the form
%\[
%x=(x_1, \ldots , x_n) \in A \efun (x,\rv(x_i)) \in A \times \RV,
%\]
%where the $i$th coordinate of $A$ is assumed to be a $\VF$-variable.
%\end{defn}

\subsection{Proper covariant bijections that are special}

\begin{defn}\label{defn:special:bijection}
Let $A$ be a doubly bounded proper invariant set, or a finite disjoint
union of such sets, regularized as in Convention~\ref{conv:can}. We
choose its first VF-coordinate as the coordinate to be changed.
Let $C\subseteq\RVH(A)$ be an RV-pullback and let
\[
 \lambda:\pr_{>1}(C\cap A)\longrightarrow\VF
\]
be a definable continuous focus whose graph lies in $C$. Thus on each
RV-polydisc $\gp\subseteq C$, the value of $\lambda_\gp$ lies in
$\pr_1(\gp)$. Suppose that $\gp\cap A$ is
$(\alpha_\gp,\beta_\gp)$-invariant and $\lambda_\gp$ is
$(\beta_\gp,\alpha_\gp)$-covariant, with the radius tuple
$\rv(\gp)$-definable.

Write $U$ for the full retained RV-label domain on the active source;
regularization makes these labels determine its VF leading terms.
An \memph{aperture} consists of a definable map
$\tau:U\longrightarrow\RV$ and a finite definable bound
$\Delta\in\Gamma_{\geq0}$ such that
\[
 0\leq\gamma(u):=\vrv(\tau(u))\leq\Delta.
\]
Here $\RV$ excludes its absorbing element. If $u_\gp$ is the label of
$\gp$, put
\[
 \gamma_\gp=\rad(\pr_1\gp)+\gamma(u_\gp),\qquad
 t_\gp=\tau(u_\gp)\rv(\pr_1\gp),
 \qquad \gamma_\gp\geq\alpha_\gp.
\]
All these maps and bounds are definable. In particular the
relative coefficient may depend on the whole retained label, whereas
the absolute cutoff still depends on the polydisc as displayed.
The \memph{centripetal transformation} is
\begin{equation*}\tag{CT}\label{centri}
\begin{cases}
 \eta(a,x)=(a-\lambda(x),x),
 & (a,x)\in\gp\cap A,\quad\vv(a-\lambda(x))\leq\gamma_\gp,\\
 \eta(a,x)=(t_\gp,x),
 & (a,x)\in\gp\cap A,\quad\vv(a-\lambda(x))>\gamma_\gp,\\
 \eta=\id,&\text{on }A\setminus C.
\end{cases}
\end{equation*}
The two images carry distinct finite tags, which are retained under
regularization. The finite bound on $\gamma$ makes the image doubly
bounded. The focus, locus, aperture and covariance radii are part of
the transformation's data.

A \memph{bijective lift} additionally includes a definable function
$b:\pr_{>1}(C\cap A)\longrightarrow\VF^\times$, continuous on each
fixed-label fibre, with
\[
 \rv(b(x))=t_\gp\quad\text{on }\pr_{>1}(\gp\cap A).
\]
This is an actual VF map, independent of the changed VF-coordinate;
its existence is part of the lift data. A variable RV aperture alone
is not asserted to have a definable VF section. Replace the contracted
clause by
\[
 \eta^\flat(a,x)=(a-\lambda(x)+b(x),x),
 \qquad \vv(a-\lambda(x))>\gamma_\gp,
\]
retaining its branch tag. Invariance and
$\gamma_\gp\geq\alpha_\gp$ imply that a nonempty core contains the
whole required disc about the focus, and its image is the full RV-disc
$b(x)+\MM_{\gamma_\gp}=t_\gp^\sharp$. The tagged map is injective;
its inverse adds $\lambda$ and subtracts $b$ on the core branch.

A \memph{special covariant transformation} is an alternating finite
composition of these transformations and regularizations. Its
\memph{length} is the number of centripetal components. A choice of
bijective lift data for every component gives a \memph{special covariant
bijection}; its image is denoted by $A^\flat$.

The constant-aperture subclass is included literally. If
$\tau(u)=t=\rv(c)$ is constant with $c$ a definable nonzero VF point,
take $b(x)=c\lambda(x)$; both lifted clauses and all branch tags are
the scalar ones. The existing bounds and VF-generation give
such sufficiently small named coefficients whenever the constant
subclass is used in the existence constructions below. Larger apertures
are permitted when their bounded RV graph and continuous VF endpoint
are included in the data.
\end{defn}

When the contracted clause is absent, the construction is an ordinary
special bijection. The finite apertures above give control of bounds and
proper covariance that is not required for an arbitrary ordinary special
bijection.

\begin{rem}\label{bigger:aper}
Suppose two admissible aperture functions on the same locus satisfy
$\gamma'(u)>\gamma(u)$ everywhere, and keep the full label $u$ and
distinct core tags. On a polydisc $\gp$, increasing the aperture replaces
the old endpoint $t_\gp$ by the relative annulus and new endpoint
\[
 \rv(\MM_{\gamma_\gp}\setminus\MM_{\gamma'_\gp})
 \ \uplus\ \{t'_\gp\}.
\]
Outside the old core the two transformations agree. This comparison is
fibrewise over $u$; no global choice of either endpoint is made.
The released annulus and the retained core are distinguished by their
finite tags. If the old contracted clause occurs, invariance makes the
whole old core present, so the strict aperture increase gives a nonempty
released part. Otherwise the two images agree.
\end{rem}

\begin{rem}\label{spec:addendum}
Suppose that $A$ is proper $\delta$-invariant but is not doubly bounded. Then, by definition, there  are no special covariant transformations on $A$. We can prepare $A$, though, as follows. First increase $\delta$, if necessary, to an attainable definable radius $\delta'\geq\delta$, and write $\delta$ for this larger radius. Invariance is preserved under this increase. Choose a definable point $t_\delta\in\delta^\sharp$ and a definable $c\in\VF$ with $\rv(c)=t_\delta$. For each open polydisc $\gb_1 \times \ldots \times \gb_n \times t \sub A$ of radius $\delta$, if $\gb_i= \MM_\delta$ then replace it with $t_\delta$. The resulting finite disjoint union is the contracted preparation. Its bijective lift replaces $x_i\in\MM_\delta$ by $x_i+c\in t_\delta^\sharp$, carrying a finite disjoint-union tag to distinguish overlapping pieces. This lifted set is doubly bounded and remains $\delta$-invariant. We include this zero-centered preparation among the permitted preliminary operations. It is a translation with derivative $1$; zero-centered pieces are permitted preliminary operations, whereas the focus in Definition~\ref{defn:special:bijection} lies in a nonzero RV-disc.
\end{rem}


\begin{lem}\label{spec:inve}
Special covariant bijections are proper covariant homeomorphisms. If $n$ is the $\VF$-arity, they are totally differentiable with Jacobian $1$ outside a definable set of $\VF$-dimension less than $n$.
\end{lem}
\begin{proof}
The continuity and covariance argument uses the continuity of the focus and endpoint maps, the invariance of their domains, and the bounded relative aperture. The latter bounds the endpoint values and every output VF and RV coordinate. The cutoff is constant on a retained-label fibre, so the two branch domains are clopen. Each branch is an independent-coordinate shear, and its inverse is continuous as well. On each branch the lifted centripetal transformation and its inverse carry sufficiently small open polydiscs to open polydiscs. Their graph fibers therefore have open domains and ranges. Lemmas~\ref{conti:covar}, \ref{conv:unif}, and~\ref{pr:co:home:com} give proper covariance and the homeomorphism property under composition, retaining the intermediate $\RV$-labels as in the standing continuity convention.

We justify the Jacobian assertion on its actual locus. A focus is a definable function of the other $\VF$-coordinates, with its $\RV$-labels fixed. Generic total differentiability in equicharacteristic zero gives a definable exceptional set of dimension less than $n-1$ for each focus and endpoint map. Its cylinder in the current $n$ $\VF$-coordinates has dimension less than $n$. Away from these cylinders the two lifted clauses are
\[
(a,x)\longmapsto(a-\lambda(x),x),\qquad
(a,x)\longmapsto(a-\lambda(x)+b(x),x).
\]
Each has a triangular total derivative with diagonal entries $1$. The branch domains are relatively clopen. Regularization has locally constant $\RV$-labels on the doubly bounded domains, and the preliminary zero-centered translations also have derivative $1$.

Write the finite composite as $T_r\circ\cdots\circ T_1$, and let $E_i$ be a definable exceptional set for $T_i$ in its current source. Remove
\[
E=\bigcup_{i=1}^{r}(T_{i-1}\circ\cdots\circ T_1)^{-1}(E_i).
\]
Definable bijections preserve $\VF$-dimension, so $\dim_{\VF}(E)<n$. At every point outside $E$ all the required total derivatives exist, and the chain rule gives $\jcb_{\VF}(T_r\circ\cdots\circ T_1)=1$. For zero $\VF$-arity the derivative is the empty matrix and its determinant is $1$. No everywhere chain rule for partial derivatives is used.
\end{proof}


So calling them ``special proper covariant bijections'' may seem more appropriate, but the verbosity is somewhat confusing.


\begin{rem}\label{spec:is:iso}
Let $(A,\omega)$, $(B, \sigma)$ be objects of $\mgVF[k]$ and $F : A \fun B$ a special covariant bijection. Then  $F$ represents a $\mgVF[k]$-isomorphism if $\omega(x) = \sigma(F(x))$ for all $x \in A$ outside a definable subset of $\VF$-dimension less than $k$ (recall Remark~\ref{mor:equi}). Include the exceptional set from Lemma~\ref{spec:inve} in this set; the measured morphism convention already permits it.
\end{rem}


\begin{lem}[{\cite[Theorem~5.5]{Yin:special:trans}}]\label{special:bi:term:constant}
Suppose that $A$ is a definable set in $\VF$ and $f : A \fun \mdl P(\RV^m)$ is a definable function. Then there is a special bijection $T$ on $A$ such that $A^\flat$ is an $\RV$-pullback and the function $f \circ T^{-1}$ is $\rv$-contractible.
\end{lem}
\begin{proof}
Apply quantifier elimination to the definable membership relation $R\sub A\times\RV^m$. Simultaneously prepare the finitely many VF-polynomials occurring in a quantifier-free formula for $R$ and for $A$. The cited preparation theorem makes their RV-values constant on each resulting RV-polydisc, with every RV-variable retained as a parameter. Hence every atomic formula, and then the entire Boolean combination, has constant truth value on that polydisc for each output RV-tuple. This is exactly the asserted powerset-valued contraction.
\end{proof}

\begin{rem}\label{special:pseudo}
The proof of Lemma~\ref{special:bi:term:constant} in \cite{Yin:special:trans} actually shows that for every definable set $A$, there is a special bijection $T$ on  $\RVH(A)$  such that $A^{\flat}$ is an $\RV$-pullback. In the present context, for reasons that will become clear, we would like to extend this result, using special covariant bijections  on proper invariant sets $A$. This is not guaranteed by Lemma~\ref{special:bi:term:constant} since the focus maps in $T$ are not required to be (suitably) covariant within each $\RV$-polydisc; except when $A_{\VF} \sub \VF$, in which case  the covariance requirement is half vacuous and it is easy to see how to turn $T$ into a special covariant bijection whose components all have the same aperture. The technique underlying such a modification could be generalized so as to be applicable in other similar situations, see  Terminology~\ref{FMT} below.
\end{rem}

\begin{ter}\label{FMT}
Often in a construction there is obstructing data that is contained in finitely many distinct definable open discs $\ga_i$, all of the same radius, and is presented in the form of a definable function $f : A \fun \VF$ with $A\sub\VF^n\times\RV^m$ and $\ga_i\sub\pr_1(A)$. Choose $a_i\in\ga_i$ only when the current domain satisfies $A\cap(\ga_i\times\VF^{n-1}\times\RV^m)=\ga_i\times A_{a_i}$ (or, for a particular replacement, the required domain inclusion). This condition must be verified again for each current focus-map domain. Then the \memph{FMT procedure} may be applied to circumvent the difficulties, which is simply this: choose a definable point $a_i$ in each $\ga_i$ and replace $f$ with another function $f' : A \fun \VF$ given by $(a,b) \efun f(a_i, b)$ if $a \in \ga_i$ and $(a,b) \efun f(a,b)$ otherwise, now the construction over each $\ga_i$ depends on $a_i$ and hence no longer varies.
\end{ter}


\begin{prop}\label{simplex:with:hole:rvproduct}
Let $H$ be a proper invariant $\RV$-pullback, not necessarily doubly bounded, and $(A_i)_i$ a definable finite partition of $H$ such that each $A_i$ is proper invariant. Then there is a special covariant bijection $T$ on $H$ such that
\begin{itemize}
  \item every $T \rest A_i$ is indeed a special covariant bijection,
  \item every $A_i^{\flat} \sub H^{\flat}$  is a doubly bounded $\RV$-pullback.
%  \item $f \circ T^{-1}$ is $\rv$-contractible.
\end{itemize}
\end{prop}
\begin{proof}
To begin with, by Remark~\ref{spec:addendum}, we may assume that $H$ is doubly bounded. It is equivalent and less cumbersome to construct a special covariant transformation such that its lift is as required (see the last paragraph of this proof). For this existence proposition we use the constant-aperture subclass of Definition~\ref{defn:special:bijection}; thus the induction proves this stronger choice as well. Its components remain admissible in the enlarged class with exactly the same maps. To that end, we proceed by induction on $n$, where $H_{\VF} \sub \VF^n$. The base case $n=1$ follows from Lemma~\ref{special:bi:term:constant} and the discussion in Remark~\ref{special:pseudo}.


For the inductive step, let $A_{i1} = \pr_1(A_i)$ and $H_1 = \pr_1(H)$. Partition $H_1$ into the nonempty Boolean atoms of the finite family $(A_{i1})_i$, and restrict every original $A_i$ over these atoms. The resulting pieces remain proper invariant by Lemma~\ref{int:inv}. Over each atom, group all original pieces with that first-coordinate projection; their fibers, rather than a purported disjoint family of original projections, form the partition used below. For each $a \in H_1$, the inductive hypothesis gives an $a$-definable special covariant transformation $T_a$ on $H_a$ such that every $T_a \rest A_{ia}$ is a special covariant transformation and every $T_a(A_{ia}) \sub T_a(H_a)$ is an $\RV$-pullback. Our goal then is to fuse together these transformations $T_a$ so to obtain one special covariant transformation on $H$ as desired. This is in general not possible without first modifying $H_1$ in a suitable way, which constitutes the bulk of the work below.


Let $U_{ak}$ enumerate the loci of the components of $T_{a}$, $\lambda_{ak}$ the corresponding continuous focus maps, and $(\gamma_{ak}, t_{ak})$ their apertures; for each $\RV$-polydisc $\gq \sub U_{ak}$, the map $\lambda_{ak\gq}$ is $(\alpha_{ak\gq}, \beta_{ak\gq})$-covariant. By compactness, there is a definable set $V \sub \VF \times \RV^l$ such that $\pr_1(V) = H_1$ and, for each $a \in H_1$, the set $V_a$ contains the following $\RV$-data of $T_a$:
\begin{itemize}
  \item $\rv(T_{a} \rest A_{ia})$, $\rv(T_{a})$, and $\rv(U_{ak})$,
  \item the $\VF$-coordinates targeted by the focus maps $\lambda_{ak}$,
  \item the $a$-definable apertures $(\gamma_{ak}, t_{ak})$,
   \item the $(a, \rv(\gq))$-definable tuples $(\alpha_{ak\gq}, \beta_{ak\gq})$;
\end{itemize}
The set $\rv(T_a)$ is redundant but is retained for clarity. Each individual row of these data is doubly bounded: its finite list of apertures is bounded, and the radii indexed by the doubly bounded RV-loci are bounded by Lemma~\ref{db:to:db}. We do not yet claim a uniform bound for $V$. Include in the data the finite formula-choice partition for all focus definitions, and pad shorter lists of components by identities. The following preparation applies to the full membership formula for these data, without requiring a bound for $V$.


Let $\phi(x, y)$ be a quantifier-free formula that defines $V$ and $G_{i}(x)$ enumerate its top terms. By Lemma~\ref{special:bi:term:constant}, there is a special bijection $R : H_1 \fun H_1^\flat$ such that each $A_{i1}^\flat \sub H_1^\flat$ is an $\RV$-pullback and every function $G_{i} \circ R^{-1}$ is $\rv$-contractible. This implies that,  for every $\RV$-polydisc $\gp \sub H_1^\flat$, the $\RV$-data $V_a$ is constant over $a \in R^{-1}(\gp)$. Also, since $H_1$ is an $\RV$-pullback of $\RV$-fiber dimension $0$, by \cmin-minimality and Lemma~\ref{RV:fiber:dim:same}, each focus map in $R$ consists of only finitely many points.

Fix a sufficiently large definable $\delta \in \Gamma$ such that every $A_i$ is proper $\delta$-invariant.  Then there is an algebraic set of open discs $\ga_j \sub H_1$ of radius $\delta$ such that the restriction of $R$ to $H_1 \mi \bigcup_j \ga_j$ is actually a special covariant bijection --- the reason simply being that, after deleting all the discs $\ga_j$, each focus map in $R$ lies outside the set in question. Each $\ga_j$ contains an $\code {\ga_j}$-definable point $a_j$ and, for all $a, a' \in \ga_j$ and every $A_i$, we have $A_{ia} = A_{ia'}$ (because  $A_i$ is $\delta$-invariant). It follows that, over each $\ga_j$, we can use the same  transformation $T_{a_j}$ to achieve the desired effect. This is the first instance where FMT is applied, namely to $T$, and then $R$ is adjusted so that every $\ga_j$ is mapped to the same $\RV$-disc $t^\sharp_{\delta}$, where $t_{\delta} \in \delta^\sharp$ is a chosen definable point.

%\begin{ter}\label{FMT}
%This operation just performed --- modifying a construction over finitely many definable open discs with a chosen definable point in each of them --- will be applied many more times below (even after this proof) and, to reduce verbosity, we shall refer to it by the acronym FMT. Which construction it is being applied to should be clear in context.
%\end{ter}

Therefore, we may assume that $R$ is a special covariant bijection whose components all have the same  aperture $(\delta, t_{\delta})$.


The finite formula-choice partition was included among the sets prepared by $R$. After freezing the finitely many radius-$\delta$ discs split by this partition, its transported pieces are RV-pullbacks. On each such piece use its prescribed focus formula, with the identity padding already chosen. At each stage the union of the selected loci is an RV-pullback, since it is a union of whole RV-polydiscs, and each constituent focus has its graph in that locus. The pieces are clopen; hence the piecewise focus is continuous when the constituent foci are continuous. This supplies one global definition at each stage, with an allowed locus. Taking $R$ to be trivial means replacing $H$, its partition, and the complete fiberwise data by their transported versions along $R$.

We may now establish the postponed uniform bound. On the unexceptional pieces the complete row $V_a$ depends only on $\rv(R(a))$ and the existing RV-tags. On each frozen disc it is the single row chosen at its center, with its finite disc tag carried along. Thus the data are indexed by a doubly bounded RV-base. For each such parameter the row is doubly bounded, as above. The set of its symmetric $\Gamma$ bounds is a nonempty definable upper ray. Definable completeness, followed by increasing an endpoint by a fixed positive definable value if needed, gives a definable bound as a function of that RV-parameter. Lemma~\ref{db:to:db} bounds this function uniformly. In particular the recomputed, frozen family of data is doubly bounded. No bound for the initial arbitrary family $V$ is used.


Over each $t^\sharp \sub H_1$, we can now write $U_{ak}$ as $U_{tk}$, $\alpha_{ak\gq}$ as $\alpha_{k\gq}$ (the first $\RV$-coordinate of $\gq$ is actually $t$), and so on. We are almost ready to fuse together the transformations $T_a$ over $a \in t^\sharp$. The remaining problem  is that, for any $a, a' \in t^\sharp$, although the two focus maps $\lambda_{a1\gq}$, $\lambda_{a'1\gq}$ are both $(\alpha_{1\gq}, \beta_{1\gq})$-covariant,  the images of the same open polydisc of radius $\alpha_{1\gq}$ may lie in two distinct open discs of radius $\beta_{1\gq}$. To solve this problem, consider an open polydisc $\gp$ of radius $\alpha_{1\gq}$ that is contained in $\dom(\lambda_{a1\gq})$ for some (hence all) $a \in t^\sharp$. For each $b \in \gp$, let $\lambda_{1b} : t^\sharp \fun \VF$ be the function defined by $\lambda_1(x, b, z)$. By \cmin-minimality and Corollary~\ref{part:rv:cons}, there are a $b$-definable finite set $C_b \sub \VF$ and, for any $a \in t^\sharp \mi C_b$, an open disc $\ga_a \sub t^\sharp \mi C_b$ around $a$ such that $\lambda_{1b}(\ga_a)$ lies in an open disc of radius $\beta_{1\gq}$. Since for any other $a' \in \ga_a$,  $\lambda_{a'1\gq}(\gp)$ also lies in an open disc of radius $\beta_{1\gq}$, we see that $\lambda_1(\ga_a \times \gp)$ lies in an open disc of radius $\beta_{1\gq}$, where $\lambda_1$ stands for the function defined by $\lambda_1(x, y, z)$. With the remaining RV-labels $z$ fixed, define $E_{\gp,z}$ to be the set of $a$ for which no neighborhood $\ga$ has $\lambda_1(\ga\times\gp)$ contained in a single radius-$\beta_{1\gq}$ disc. Covariance in the $\gp$-variables shows that $E_{\gp,z}\sub C_b$ for every $b\in\gp$: a neighborhood that works at one $b$ works throughout $\gp$. Consequently $E_{\gp,z}$ is finite and is canonically definable from $\code\gp$ and $z$. Lemma~\ref{ball:to:ac} removes the polydisc code from its algebraic dependence, and Lemma~\ref{dcl:to:ac} removes the remaining RV-labels. The union of this uniformly definable family lies in $\acl(\mathbb S)$ and is finite by compactness. Freeze sufficiently small discs around this common finite set. At every stage their product-domain condition is checked for the current focus domain: the current ambient pullback and locus are determined by their RV-tags, and the freezing discs lie within the corresponding first-coordinate RV-discs. Thus replacing the first coordinate by the chosen center leaves all other allowed coordinates in the current focus domain. The earlier product condition for the original partition alone is not being reused without this check. By compactness and FMT, we may now assume that there is an $\go$-partition $p: \dom(\lambda_1) \fun \Gamma$ such that, for each open polydisc $\gb$ in question,  $\lambda_1(\gb)$ lies in a disc of radius $\beta_{1\gq}$, where $\gq$ is related to $\gb$ in the obvious way. By Lemma~\ref{vol:par:bounded}, the image of $p$ is doubly bounded and, by Lemma~\ref{db:to:db}, there is a definable $\gamma_1 \in \Gamma$ with $\gamma_1 \geq \gamma_{t1}$ for all $t \in \rv(H_1)$. We may choose such a $\gamma_1$ together with a definable point $t_1 \in \gamma^\sharp_1$ (see the opening paragraph of this section). All this means that $\lambda_1$ can serve as the focus map of a centripetal transformation $T_1$ on $H$ of aperture $(\gamma_{1}, t_{1})$. For continuity, extend each focus by zero off its current clopen domain. Its fibers in the other VF-coordinates remain continuous. For each fixed tuple of RV-labels, Proposition~\ref{prop:continuity:fiberwise} makes the projection of its actual discontinuity locus to the first VF-coordinate finite. This projection is canonically definable with those labels, so Lemma~\ref{dcl:to:ac} and compactness give one finite definable exceptional set for all labels. Take sufficiently small same-radius discs around this set, refining the invariance radius of the current focus domain. The exceptional set itself meets every such disc; averaging its points in each disc gives centers with a complete definable graph. The domain condition for FMT follows from invariance. On a frozen disc the new focus is a continuous fiber at its center; off their clopen finite union it agrees locally with the original focus. Thus the resulting focus is continuous everywhere on each RV-fiber.


It remains to justify successive aperture changes. We continue in the contracted picture and use a finite induction on the number of components still to be fused, inside the induction on $n$. Consider a component in coordinate $j$ whose old and new absolute aperture radii on the current RV-polydisc $\gq$ are $\rho_\gq\leq\rho'_\gq$. Write $x$ for the other VF-coordinates together with all existing RV-tags, and put $d=a_j-\lambda(x)$. The part newly released from the old contracted clause is
\[
\rho_\gq<\vv(d)\leq\rho'_\gq.
\]
Regularize $d$, retain its RV-tag $s$, and retain the finite tag distinguishing it from the new contracted core. The inequalities on $s$ depend only on the existing RV-code of $\gq$ and on the two apertures. Hence the released block $W$ is an RV-pullback. Let $p$ forget the released VF-coordinate $d$, retaining $s$ and every other tag.

For each original piece $A_i$, the old aperture dominates its invariance radius on this polydisc, because the old transformation restricted to $A_i$ was required to be special covariant. Therefore, throughout the old contracted ball,
\[
(\lambda(x)+d,x)\in A_i
\quad\Longleftrightarrow\quad
(\lambda(x),x)\in A_i.
\]
Thus, in the new image, $W\cap T_1(A_i)=W\cap p^{-1}(B_i)$, where $B_i$ is defined by the condition on the right, together with the RV-inequalities on $s$. This is an equality of full sets, not only of their RV-images. The projected ambient $p(W)$ is an RV-pullback with $n-1$ VF-coordinates. Its bounds are inherited from the current ambient and the bounded apertures. Each $B_i$ is clopen: its field domain is the current projected pullback, and the displayed condition is the preimage of the clopen set $A_i$ under the continuous graph map $x\mapsto(\lambda(x),x)$. Consequently the $B_i$ are proper invariant by Lemma~\ref{inv:crit}. They form the induced partition of $p(W)$.

Apply the $(n-1)$-dimensional induction to this projected partition. Lift every resulting component by carrying $d$ along unchanged. Since $s$ is retained, every fiber of $W$ in the released coordinate is already a whole RV-disc; hence these lifted components are permitted and prepare all $W\cap T_1(A_i)$. The new contracted core has already lost coordinate $j$, so the same lower-dimensional induction applies there directly. After lifting the final construction to a bijection, that core coordinate becomes a full RV-disc and is likewise carried along.

On the complementary branch $\vv(d)\leq\rho_\gq$ the old and new transformations coincide, including their regularization tags. Continue the original remaining components there, restricting their loci to this branch. The annular and core branches are distinguished by the retained tags, so these restrictions and their complements are RV-pullbacks. Each prepared branch therefore extends by the identity to the other branches, as in the finite gluing step above. The upper-dimensional remainder has one fewer component left to fuse, whereas every replacement problem has at most $n-1$ VF-coordinates. This is a finite lexicographic induction and completes the construction. All source and branch tags are retained in its bijective lift.
\end{proof}


This proposition can indeed be applied to any proper invariant set $A$ against its $\RV$-hull $\RVH(A)$, without specifying a partition of $\RVH(A)$ as required, because, by  Lemma~\ref{int:inv},   $\RVH(A) \mi A$ is proper invariant too.

\begin{rem}\label{cons:extra}
As in the statement of Lemma~\ref{special:bi:term:constant} above, if $f : H \fun \mdl P(\RV^m)$ is a proper covariant function then we can demand as a clause of the conclusion of Proposition~\ref{simplex:with:hole:rvproduct} that the function $f \circ T^{-1}$ be $\rv$-contractible, which would follow if it is simply included as a clause of the inductive hypothesis in the proof.
\end{rem}

Let $A$ be a definable set and $\omega : A \fun \Gamma$ a definable function. We say that the definable pair $\bm A = (A, \omega)$ is \memph{proper invariant} if $A$ is proper invariant and $\omega$ is proper covariant, that is, every fiber $A_\gamma$ is proper invariant.

\begin{lem}\label{L:measure:surjective}
Suppose that $A$ is an $\RV$-pullback and  $(A_i)_i$ is a definable finite partition of $A$ such that each $\bm A_i = (A_i, \omega \rest A_i)$ is proper invariant. Then there is a special covariant bijection $T$ as given in Proposition~\ref{simplex:with:hole:rvproduct} such that $\omega \circ T^{-1}$ is $\rv$-contractible. In particular, if $\bm A_i \in \mVF[k]$ then $T$ induces a morphism  $\bm A_i \fun \mL(\bm U_i, \pi_i)$, where $(\bm U_i, \pi_i) \in \mRV^{\db}[k]$.
\end{lem}
\begin{proof}
We view $\omega$ as a definable function $A \fun \mdl P(\RV)$. By Proposition~\ref{simplex:with:hole:rvproduct}, there is a special covariant bijection $T$ on $A$ such that, for each $i$, $T \rest A_{i}$ is a special covariant bijection and $A_{i}^\flat$ is an $\RV$-pullback. Moreover, applying Remark~\ref{cons:extra}, we get that each function $\sigma_i \coloneqq (\omega \circ T^{-1}) \rest A_{i}^\flat$ is constant on every $\RV$-polydisc in $A_{i}^\flat$ and hence induces a function $\pi_i : \rv(A_{i}^\flat) \fun \Gamma$. Now, if $\bm A_i \in \mVF[k]$ then, by Lemma~\ref{RV:fiber:dim:same}, $\bm U_i\coloneqq \rv(A_{i}^\flat)_{\leq k} \in \RV^{\db}[k]$ (recall Notation~\ref{0coor}). We see that $(\bm U_i, \pi_i)$ is as required because, by Remark~\ref{spec:is:iso}, $\bm A_i$ is isomorphic to $(\mathbb{L} \bm U_i, \sigma_i)$.
\end{proof}


As above, this lemma may be applied to any proper invariant object $(A, \omega) \in \mVF[k]$, since the auxiliary pair $(\RVH(A)\mi A,0)$ is also proper invariant. This auxiliary pair need not be an object of $\mVF[k]$: taking an RV-hull can create infinite RV-fibers. Apply the lemma to the partition $(A,\RVH(A)\mi A)$, extending the form by $0$ on the auxiliary piece. Only the selected original piece $(A,\omega)$ is required to belong to $\mVF[k]$ when invoking the final lifting conclusion.


\begin{lem}\label{invar:db}
Let $A$ be a definable set with $A_{\VF}$  bounded and $A_{\RV}$  doubly bounded. Let $f : A \fun \Gamma$ be a definable function  that is constant on $A \cap \ga$ for all open polydiscs $\ga$ of radius $\gamma \in \Gamma$. Then $f(A)$ is doubly bounded.
\end{lem}
\begin{proof}
Observe that, without loss of generality, we can replace $A$ with its $\gamma$-tubular neighborhood and thereby  assume that $(A,f)$ is indeed proper invariant. By Lemma~\ref{L:measure:surjective}, $f$ factors through a doubly bounded definable set in $\RV$. So, by Lemma~\ref{db:to:db},  $f(A)$ must be doubly bounded.
\end{proof}

\begin{cor}
If $(A, \omega)$ is proper invariant then $\omega(A)$ is doubly bounded.
\end{cor}
\begin{proof}
We have seen this many times above: as a consequence of Lemma~\ref{vol:par:bounded}, for all sufficiently large $\beta \in \Gamma$, every fiber $A_\gamma$ of $\omega$ is $\beta$-invariant. So the claim is immediate by Lemma~\ref{invar:db}.
\end{proof}

\subsection{Standard contractions}

\begin{defn}\label{rela:unary}
Let $A \sub \VF^{n} \times \RV^{m}$, $B \sub \VF^{n} \times \RV^{m'}$ be definable sets. A definable bijection $G : A \fun B$ is \memph{relatively unary}, or more precisely, \memph{relatively unary in the $i$th coordinate}, where $i \in [n]$, if $(\pr_{\tilde{i}} \circ G)(x) = \pr_{\tilde{i}}(x)$ for all $x \in A$. If, in addition, $G \rest A_a$ is also a special covariant bijection for every $a \in \pr_{\tilde{i}} (A)$ then $G$ is \memph{relatively special covariant} in the $i$th coordinate.

Let $\bm U = (U, f)$, $\bm V = (V, g)$ be objects of $\RV[k]$. A morphism $F : \bm U \fun \bm V$ is \memph{relatively unary in the $i$th coordinate}, where $i \in [k]$, if $(\pr_{\tilde{i}} \circ g \circ F)(u) = (\pr_{\tilde{i}} \circ f)(u)$ for all $u \in U$.
\end{defn}

Since identity functions are relatively unary in any coordinate, if a  bijection  is piecewise (over a definable finite partition) a composition of relatively unary bijections then it is indeed a composition of relatively unary bijections.

\begin{exam}
Every  special covariant bijection $T$ of length $1$ is relatively special covariant, but clearly not vice versa (to begin with, the relevant data is not necessarily uniform across the fibers in each $\RV$-polydisc).
\end{exam}

Let $A \sub \VF^{n} \times \RV^{m}$ be a proper invariant set. Let $i\in[n]$ and let $T_i:A\fun A^\flat$ be a continuous proper covariant bijection that is
\begin{itemize}
  \item relatively unary in the $i$th coordinate,
  \item totally differentiable with Jacobian $1$ outside a definable subset of $\VF$-dimension less than $n$.
\end{itemize}
Moreover, the $\RV$-coordinates of $A$ are passed on to $A^\flat$ via $T_i$ in the obvious way (the simplest example is that $T_i$ is the regularization of $A$). Adjusting $T_i$ by way of Remark~\ref{spec:addendum} if necessary, we take $A$ and $A^\flat$ to be both doubly bounded. By Lemma~\ref{conv:unif}, $A^\flat$ is proper invariant as well, and $T_i$ is a proper covariant homeomorphism. Suppose that, for all $a \in A_{\tilde i}$ and all $t \in (A_{a})_{\RV}$, the image $T_i(A_{a, t}) \sub A_{a, t}^\flat$ of the fiber $A_{a, t} \sub A$ is an $\RV$-pullback.  Let
\[
 A_{-i} = \bigcup_{a \in A_{\tilde i}} a \times (A_{a}^\flat)_{\RV} \sub \VF^{n-1} \times \RV^{m_i},
\]
which is doubly bounded and proper invariant, and  we think of it as obtained from $A^\flat$ by ``$\rv$-contracting'' the $\VF$-coordinate in question.  Also, under Convention~\ref{conv:can}, $A^\flat$ can be ``fully recovered'' from  $A_{-i}$. Let $\hat T_i : A \fun A_{-i}$ be the function induced by $T_i$.

For any $j \in [n{-}1]$, we repeat the above setup on $A_{-i}$ with respect to the $j$th $\VF$-coordinate, and thereby obtain a doubly bounded proper invariant set $A_{-j} \sub \VF^{n-2} \times \RV^{m_j}$ and a function $\hat T_{j} : A_{-i} \fun A_{-j}$ that is induced by a relatively unary, partially differentiable, proper covariant homeomorphism $T_j : A^\flat \fun A^{\flat\flat}$ (the $\VF$-coordinate in $A^\flat$ that gets $\rv$-contracted in $A_{-i}$ is simply carried along by $T_j$). Continuing thus,  a sequence of such homeomorphisms  $T_{\sigma(1)}, \ldots, T_{\sigma(n)}$ and a corresponding function $\hat T_{\sigma} : A \fun \RV^{l}$ result, where $\sigma$ is the permutation of $[n]$ in question. The composition $T_{\sigma(n)} \circ \ldots \circ T_{\sigma(1)}$, which is referred to as the \memph{lift} of $\hat T_{\sigma}$, is denoted by $T_{\sigma}$. By Lemma~\ref{pr:co:home:com}, $T_{\sigma}$ is a proper covariant homeomorphism. It is totally differentiable with Jacobian $1$ outside a definable lower-dimensional set, by the exceptional-set union in the proof of Lemma~\ref{spec:inve}. For a later stage, carry the already contracted coordinates along and take the corresponding cylinder of its exceptional set; the dimension inequality remains strict. All measured identities associated with a standard contraction are understood on this common good locus. Also note that $T^{-1}_{\sigma}$ is $\rv$-contractible.


\begin{defn}\label{defn:standard:contraction}
The function $\hat T_{\sigma}$, or the set $\hat T_{\sigma}(A)$, is called a \memph{standard contraction} of  $A$.
\end{defn}


\begin{lem}\label{exit:stand}
Every proper invariant set $A \sub \VF^{n} \times \RV^{m}$  admits a standard contraction with respect to any permutation $\sigma$ of $[n]$.
\end{lem}
\begin{proof}
Relabel the coordinates so that $\sigma(n)$ is the base coordinate in the proof of Proposition~\ref{simplex:with:hole:rvproduct}. Strengthen its fiberwise induction by prescribing the order $\sigma(1),\ldots,\sigma(n-1)$ and retaining every label of an already contracted coordinate. The aperture-change induction there is made in the same order: on a newly released block its released coordinate is already a full RV-disc, and the lower-dimensional construction carries this coordinate along. Identity padding allows the finitely many branches to use the same ordered list of coordinate contractions.

Write $a$ for the base coordinate, $x$ for the other VF-coordinates, and $t$ for all old RV-tags. Keep the unary base preparation $R$ deferred. If $\tau(a,t)$ denotes all RV-coordinates of $R(a,t)$, including the finite formula-choice and frozen-disc labels, first apply the grafting map
\[
G(a,x,t)=(a,x,t,\tau(a,t)).
\]
The map $\tau$ is locally constant: it is the RV-coordinate projection of the proper covariant homeomorphism $R$, with its uniformly bounded labels retained. Each of its fibers is invariant, and all these labels are doubly bounded. Thus $G$ and its inverse projection are continuous. For fixed old and new labels, their VF-maps are identities on the clopen fiber of $\tau$. Lemmas~\ref{inv:crit} and~\ref{conti:covar} show that $G$ is a proper covariant homeomorphism; its VF-derivative is the identity on every such fiber. Grafting is allowed for a standard contraction, whose lift need not itself be a special covariant bijection.

At an intermediate fiber stage, let $S_j$ be the relatively unary map constructed over the prepared base $b=R(a,t)_{\VF}$. On the grafted space use the transported map
\[
(a,x,t,\tau)\longmapsto
\bigl(a,S_{j,R(a,t)}(x),t,\tau\bigr),
\]
retaining the other output labels of $S_j$. Its focus is the prepared focus composed with $R$ in the base variable. This is continuous. The permitted loci and the output RV-data depend only on the retained labels, by the preparation of the complete data and the freezing steps. Consequently their pullbacks are clopen, and the same covariance bounds, transported along $R$, give proper covariance. Each map is relatively unary in its prescribed coordinate and has total derivative and Jacobian $1$ off a lower-dimensional set, as in Lemma~\ref{spec:inve}.

After the entire contraction in coordinate $\sigma(j)$, replace that VF-coordinate by its RV-code. Every later $S_j$ is defined on this reduced space; in the lifted picture it therefore carries the contracted coordinate along. Once the first $n-1$ coordinates have been contracted, apply $R$ in the sole remaining VF-coordinate $\sigma(n)$, carrying all accumulated labels. The prepared fiber-image data are constant on its RV-discs, so this last image is an RV-pullback in that coordinate too. The resulting list has precisely the prescribed order and never changes a VF-coordinate after contracting it. It is the required standard contraction.
\end{proof}


\begin{rem}\label{stand:vol}
Suppose that there is a $k \in 0 \cup [m]$ such that  $(A_{a})_{\leq k} \in \RV^{\db}[k]$  for  every $a \in A_{\VF}$. If $k=0$ then this just means $A \in \VF_*$. By Lemma~\ref{RV:fiber:dim:same}, $\hat T_{\sigma}(A)_{\leq n+k}$ is an object of $\RV^{\db}[n{+}k]$. Let $\bm A = (A, \omega)$ be a proper invariant pair such that the proper covariant function $\omega : A \fun \Gamma$   induces a function $\omega_{\hat T_{\sigma}}$ on $\hat T_{\sigma}(A)$ via $T_{\sigma}$, which just means that $\omega \circ T^{-1}_{\sigma}$ is invariant on $\RV$-polydiscs. Then the function $\hat T_{\sigma}$, or the object
\[
\hat T_{\sigma}(\bm A)_{\leq n+k} = (\hat T_{\sigma}(A), \omega_{\hat T_{\sigma}})_{\leq n+k} \in \mRV^{\db}[n{+}k],
\]
is understood as  a \memph{standard contraction} of $\bm A$. Using Lemma~\ref{exit:stand} instead of Proposition~\ref{simplex:with:hole:rvproduct} in the proof of Lemma~\ref{L:measure:surjective}, we see that such a $\hat T_{\sigma}$ does exist. Moreover, if we set $A^\sharp = \bigcup_{a \in A_{\VF}} a \times \bb L (A_a)_{\leq k}$, which is a  proper invariant object of $\VF_*$, and denote the  proper covariant function on $A^\sharp$ induced by  $\omega$ still by $\omega$, then  $T_{\sigma}$  indeed induces a $\mVF[n{+}k]$-morphism
\[
(A^\sharp, \omega) \fun {\mL}\hat T_{\sigma}(\bm A)_{\leq n+k}
\]
that can be written as a composition of relatively unary proper covariant homeomorphisms.

Here $k$ is dubbed the \memph{head start} of $\hat T_{\sigma}$, which is usually implicit in the sequel. In fact,  it never has to be anything other than $0$ except in Lemma~\ref{isp:VF:fiberwise:contract}, and can be circumvented even there. This seemingly needless gadget only serves to make the discussion more streamlined: if $A \in \VF_*$ then the intermediate steps of a standard contraction of $A$ may or may not result in objects of $\VF_*$ and hence the definition cannot be formulated entirely within $\VF_*$.
\end{rem}


%So $\bm A^\sharp \coloneqq (A^\sharp, \omega)$ is an object of $\mVF^{\diamond}[n{+}k]$.

\subsection{Lifting from $\RV$ to $\VF$}


\begin{lem}\label{bijection:partitioned:unary}
Every morphism in $\RV[k]$ can be written as a composition of relatively unary morphisms, similarly in $\mRV[k]$, $\RV^{\db}[k]$, and $\mRV^{\db}[k]$.
\end{lem}
\begin{proof}
It is enough to show this piecewise, and the situation can be  simplified accordingly step by step. Let $F : (U, f) \fun (V, g)$ be an $\RV[k]$-morphism. The idea is to adjust  $f(u)$ towards $g(F(u))$, simultaneously for all $u \in U$, one coordinate at a time, and the issue is how to stay in $\RV[k]$, that is, how to make the adjusted $f$ remain finite-to-one along the way. To that end, we proceed by induction on $n = \dim_{\RV}(U)$ and, without loss of generality, may take $F = \id$. The base case $n=0$ is straightforward and is left to the reader. For the inductive step, we may assume that both $\pr_{\leq n} \rest f(U)$ and $\pr_{\leq n} \rest g(V)$ are finite-to-one. Let $f' = \pr_{\leq n} \circ f$ and $g' = \pr_{\leq n} \circ g$. Since, as in the base case, the two obvious $\RV[k]$-morphisms
\[
(U,f) \fun (U \times 1, f' \times 1), \quad (U \times 1, g' \times 1) \fun (U, g)
\]
are compositions of relatively unary morphisms, it is enough to show that $(U, f') \fun (U, g')$ is such a morphism as well. So we may assume $n=k$.

Write $\dot U = f(U)$ and $\ddot U = g(U)$. For each $t \in \dot U_{< n}$ and each $i \in [n]$, let
$F^\dag_{t, i} : \dot U_t \fun \mdl P(\ddot U_i)$
be the definable finitary function induced by $F^\dag$ (recall Notation~\ref{mor:ftf}). Then, by Lemma~\ref{finite:fib:def}, there is a $t$-algebraic subset $B_{t, i} \sub \ddot U_i \sub \RV$ such that every fiber of $F^\dag_{t, i}$ over an element in $\ddot U_i \mi B_{t, i}$ is finite; denote by $W_{t, i} \sub \dot U$ the union of such fibers. Indeed, by the compactness argument thereof,  $B_{t, i}$ can be defined uniformly. Let $U_{t, i}$ be the set of those $u \in f^{-1}( t \times W_{t, i})$ such that $g(u)_i \notin B_{t, i}$, and $f'_{t, i} : U_{t, i} \fun \RV^n$ the function given by $u \efun (t, g(u)_i)$, which is finite-to-one. Let $U_i = \bigcup_{t}  U_{t, i}$ and $f'_i = \bigcup_{t}  f'_{t, i}$. Then $(U_i, f \rest U_i) \fun (U_i, f'_i)$ is a relatively unary $\RV[n]$-morphism. By applying the inductive hypothesis fiberwise to $(U_i, f'_i) \fun (U_i, g \rest U_i)$ in the obvious sense, we may actually assume $U_i = \0$, for every $i \in [n]$. Now, since $F^\dag(\dot U_t) \sub \prod_{i \in [n]} B_{t, i}$, which is finite, we must have $\dim_{\RV}(U) < n$. So the claim follows from the inductive hypothesis.

For every intermediate finite-to-one coordinate map $h:U\fun\RV^k$, equip $U$ with the form
\[
\omega_h(u)=\omega(u)+\Sigma\vrv(f(u))-\Sigma\vrv(h(u)).
\]
Then $\omega_h+\Sigma\vrv\circ h=\omega+\Sigma\vrv\circ f$, so all intermediate measure identities telescope. If only $n<k$ active coordinates are displayed, pad them by $\mathbf1_{k-n}$ throughout. In the doubly bounded case the intermediate coordinate maps use only coordinates of the original maps and these constant pads, so their supports remain doubly bounded. This gives the measured and doubly bounded versions as well.
\end{proof}


\begin{lem}\label{RVlift}
Let $U$, $V$ be sets in $\RV$ and $f : U \fun \mdl P(V)$ a definable finitary function. Then there is a definable finitary function $f^\sharp : U^\sharp \fun \mdl P(V^\sharp)$ that $\rv$-contracts to $f$. Moreover, for every $(u, v) \in f$, $f^\sharp_{u, v} \coloneqq (u, v)^\sharp \cap f^\sharp$ is a partially differentiable function $u^\sharp \fun v^\sharp$.
\end{lem}

Similar to Definition~\ref{defn:special:bijection}, this is another place where we need the property that every algebraic $\RV$-disc contains an algebraic point, and the easiest way to achieve
this is to assume that the substructure $\bb S$ is $\VF$-generated, which validates the repeated use of Lemma~\ref{algebraic:balls:definable:centers} below.

\begin{proof}
We proceed by induction on $\dim_{\RV}(U) = n$. For the base case $n=0$, $U$ is finite and hence, for every $u \in U$, the $\RV$-polydisc $u^\sharp$ contains a $u$-definable point, and similarly for every element in $f(u) \sub V$. So lifting $f$ as desired is trivial. For the inductive step, we may assume that $\pr_{\leq n} \rest U$ is finite-to-one. Let $U' = \pr_{\leq n}(U)$ and $f' : U' \fun \mdl P(V)$ be the definable finitary function induced by $f$. Observe that if $f'$ can be lifted as desired then  $f$ can be lifted (again trivially) as desired as well. So, without loss of generality,  $U \sub \RV^n$. For a full VF-tuple $a\in U^\sharp$, every RV-tuple $v\in f(\rv(a))$ is algebraic over $a$. Apply Lemma~\ref{algebraic:balls:definable:centers} over this VF-generated extension. Compactness gives a finite $a$-definable family of VF-tuples meeting every such $v^\sharp$, uniformly in $a$. In each $v^\sharp$, average all chosen tuples coordinatewise. The number being averaged is finite and nonzero; residue characteristic zero implies that the average stays in the same RV-polydisc. This defines one output in each prescribed RV-polydisc, without choosing a center over an arbitrary RV-parameter alone.

Remove the noncontinuity and nondifferentiability loci of all these finitely many branches simultaneously. Their union $E$ has VF-dimension less than $n$, by Lemmas~\ref{fun:almost:cont} and~\ref{diff:almost:every}. No full RV-polydisc lies in $E$, so Lemma~\ref{RV:bou:dim} gives $\dim_{\RV}\rv(E)<n$. Apply the inductive hypothesis over the entire parameter set $\rv(E)$ and replace the lift on its full RV-preimage, rather than merely deleting exceptional points. Restrict the inductive relation to the originally prescribed pairs $(u,v)$. The old and replacement domains are unions of whole RV-polydiscs and hence clopen. The resulting branches therefore glue to everywhere continuous, partially differentiable functions with precisely the original RV-contraction. This proves the assertion, together with the strengthened continuous-branch form used below.
\end{proof}


Let $F: \bm U \fun \bm V$ be an $\RV[k]$-morphism, where $\bm U = (U, f)$ and $\bm V = (V, g)$. If $F^{\sharp} : \bb L \bm U \fun \bb L \bm V$ is a definable bijection that $\rv$-contracts to the function $U_f \fun V_g$ induced by $F$ (recall Notation~\ref{RVcat:tag}) then it is called a \memph{lift} of $F$. We shall also think of the finite-to-finite correspondence $F^\dag$ as the $\rv$-contraction of such a lift.

%Assume that the substructure $\mathbb{S}$ is $(\VF, \Gamma)$-generated for the remainder of this section.
%\arthur{We can remove this sentence since $\mathbb{S}$ is assumed $\VF$-generated since the beginning of the section.}

%\begin{lem}[{\cite[Proposition~6.1, Lemma~6.2]{hrushovski:kazhdan:integration:vf}, \cite[Corollary~7.7]{Yin:special:trans}}]\label{L:semigroup:hom}
% $F$ admits a lift.
%\end{lem}

\begin{lem}[{\cite[Lemma~6.3]{hrushovski:kazhdan:integration:vf}}]\label{iso:lifted:vol}
There is a lift $F^{\sharp}$ of $F$ such that
\begin{equation}\label{jac:mat}
\vv(\jcb_{\VF} F^{\sharp}(a,  u)) = \jcb_{\Gamma} F^\dag(u), \quad \text{for all } (a, u) \in \bb L \bm U.
\end{equation}
\end{lem}

See \cite[Corollary~7.7]{Yin:special:trans} for an alternative proof (the weaker qualifier ``for almost all $(a, u) \in \bb L \bm U$'' there may be upgraded to ``for all $(a, u) \in \bb L \bm U$'' here by appealing to Lemma~\ref{RV:bou:dim} at suitable places).

\begin{rem}\label{rvdisc:stre}
Let $t, s \in \RV$ such that the $\RV$-discs $t^\sharp$, $s^\sharp$ contain definable points. It is easy to see that for any definable $c \in \VF^\times$ with $\vv(c) = \vrv(s/t)$, there is a definable bijection $f : t^\sharp \fun s^\sharp$ such that $\ddx f = c$.
\end{rem}


\begin{lem}\label{L:measure:class:lift}
Every $\mRV[k]$-morphism $F : (\bm U, \omega) \fun (\bm V, \sigma)$ admits a lift
\[
F^\sharp: {\mu}\bb{L}(\bm U, \omega) \fun {\mu}\bb{L}(\bm V, \sigma)
\]
that is piecewise a composition of partially differentiable relatively unary $\mVF[k]$-morphisms. Moreover, if $F$ is a $\mRV^{\db}[k]$-morphism then every component of the composition is a  proper covariant homeomorphism and hence so is $F^\sharp$.
\end{lem}
\begin{proof}
We strengthen the induction on $n=\dim_{\RV}(U)$ by requiring every relatively unary branch to be continuous, as well as partially differentiable. In degree zero the coordinate discs have definable centers over their algebraic RV-parameters. Remark~\ref{rvdisc:stre} gives affine lifts; these are continuous, and their slopes give the required volume identity. Use the intermediate forms of Lemma~\ref{bijection:partitioned:unary} throughout.

Suppose $n>0$. First project to $n$ finite-to-one active coordinates, retaining all original labels and padding inactive coordinates by $1$. Lemma~\ref{RVlift}, in the strengthened continuous-branch form proved above, gives the center functions for the triangular corrections that replace the remaining coordinates by these pads. On each fixed RV-branch such a correction is an affine map in its changing coordinate, with continuous coefficients in the other coordinates. Thus it is continuous and relatively unary. Its inverse is of the same form, and Remark~\ref{rvdisc:stre} gives its constant slope. The intermediate forms make these corrections measured morphisms. This reduces the problem to $n=k$.

By Lemma~\ref{bijection:partitioned:unary}, it is enough to lift one relatively unary RV-morphism. Apply Lemma~\ref{iso:lifted:vol} on its fixed RV-fibers. The resulting definable lift has the prescribed Jacobian valuation. Include the noncontinuity loci as well as the nondifferentiability loci of all its branches in one bad set $E$. There are uniformly only finitely many branches over each VF-input, so Lemmas~\ref{fun:almost:cont} and~\ref{diff:almost:every} give VF-dimension less than $k$ for their projected union. The same simultaneous removal is made for every center function used in a triangular correction. By Lemma~\ref{RV:bou:dim}, $\rv(E)$ has RV-dimension less than $k$.

Apply the inductive hypothesis to the original measured RV-morphism restricted over this entire bad parameter set, retaining the corresponding target labels. Replace the old lift on its whole RV-preimage. This changes neither the prescribed RV-morphism nor its original forms. On the complementary RV-preimage all the retained branches are continuous and partially differentiable, and the Jacobian identity makes them measured morphisms. Both domains are unions of full RV-polydiscs, hence clopen, so the two constructions glue with these properties everywhere. The induction is on RV-dimension, even when the ambient degree remains $k$.

In the doubly bounded case all intermediate supports consist of original coordinates and the constant pads, so they retain their bounds. For a fixed input and output RV-label, each relatively unary component and its inverse act between full RV-discs in the changing coordinate, carrying the other coordinates along. These graph-fiber domains and ranges are open. Continuity and Lemmas~\ref{conti:covar} and~\ref{conv:unif} therefore give a proper covariant homeomorphism for each component. Lemma~\ref{pr:co:home:com} gives the same conclusion for their composition.
\end{proof}


\begin{defn}\label{defn:diamond}
The subcategory $\mVF^{\diamond}[k]$ of $\mVF[k]$ consists of the proper invariant objects and the morphisms that are compositions of relatively unary proper covariant homeomorphisms inside $\mVF^{\diamond}[k]$. Explicitly, a proper invariant object is a pair $\bm A = (A, \omega) \in \mVF[k]$ with $A$ proper invariant and $\omega$ proper covariant, and a morphism $F : \bm A \fun \bm B$ of $\mVF^{\diamond}[k]$ is a composition $F = F_l \circ \cdots \circ F_1$ of $\mVF[k]$-morphisms $F_m : \bm A_m \fun \bm A_{m+1}$, where $\bm A_1 = \bm A$, $\bm A_{l+1} = \bm B$, every $\bm A_m$ is a proper invariant object, and every $F_m$ is a relatively unary proper covariant homeomorphism. Thus each intermediate object carries its own proper covariant volume form, and each factor $F_m$ is itself $\Gamma$-measure-preserving, that is, it satisfies (\ref{cond:mg}) outside a definable set of $\VF$-dimension less than $k$.
\end{defn}

\begin{rem}\label{dia:cat:well}
Obviously the composition law holds in $\mgVF^{\diamond}[k]$: two such compositions compose to another one, and the identity is the composition with no factors. Moreover, every morphism in it is an honest bijection, as opposed to merely an essential bijection, and is in effect required to admit an inverse: the inverse of each factor $F_m$ is a relatively unary proper covariant homeomorphism $\bm A_{m+1} \fun \bm A_m$ that is again $\Gamma$-measure-preserving, so $F^{-1} = F_1^{-1} \circ \cdots \circ F_l^{-1}$ is a morphism. Thus, $\mgVF^{\diamond}[k]$ is already a groupoid and there is no need to pass to a quotient category as in Remark~\ref{mor:equi}. On the other hand, Proposition~\ref{simplex:with:hole:rvproduct}, Lemma~\ref{spec:inve}, and Remark~\ref{stand:vol} give various constructions of nontrivial morphisms. For a special covariant bijection $T = T_r \circ \cdots \circ T_1$ on a proper invariant object $(A, \omega)$, the intermediate objects are the images $T_i \circ \cdots \circ T_1(A)$ with the transported forms $\omega \circ (T_i \circ \cdots \circ T_1)^{-1}$. Each transported form is proper covariant, that is, each of its fibers is proper invariant. Indeed, its fibers are the images $T_i \circ \cdots \circ T_1(A_\gamma)$ of the fibers $A_\gamma$ of $\omega$. Each $A_\gamma$ is a proper invariant subset of $A$, so by Lemmas~\ref{spec:inve} and~\ref{conv:unif} the restriction of $T_i \circ \cdots \circ T_1$ to $A_\gamma$ is a proper covariant homeomorphism, and hence its image is proper invariant. Each stage $T_i$ is $\Gamma$-measure-preserving by Lemma~\ref{spec:inve}. Similarly, the lift in Lemma~\ref{L:measure:class:lift} passes through the lifts of the intermediate objects of Lemma~\ref{bijection:partitioned:unary}, whose volume forms factor through their $\RV$-coordinates.

The intermediate objects in Definition~\ref{defn:diamond} cannot be replaced by definable sets without volume forms. Let $t \in \MM$ be definable and nonzero, and let $\kappa : \OO \fun \OO$ be the locally affine homeomorphism that is the identity on $\OO \mi (\MM \cup (1 + \MM))$ and on $\MM \mi t\MM$, and is given by $w \efun tw$ on $t\MM$, by $w \efun t(w - 1)$ on $(1 + \MM) \mi (1 + t\MM)$, and by $w \efun 1 + (w - 1)/t$ on $1 + t\MM$. Then
\[
F(x, z) = \begin{cases} (x, x\kappa(z/x)), & x \neq 0 \text{ and } \vv(z) \geq \vv(x),\\ (x, z), & \text{otherwise},\end{cases}
\]
is a relatively unary proper covariant homeomorphism of $\OO^2$, and $F^{-1} \circ F$ is the identity of $(\OO^2, 0)$. But no proper covariant volume form $\sigma$ on the middle copy of $\OO^2$ makes $F$ $\Gamma$-measure-preserving: by (\ref{cond:mg}), $\sigma \circ F$ would agree almost everywhere with $(x, z) \efun -\vv(\kappa'(z/x))$, which takes the three values $0$ and $\pm \vv(t)$ in every open polydisc around $(0, 0)$, while $F$ is a homeomorphism fixing $(0, 0)$. The proof of Proposition~\ref{kernel:L} applies Corollary~\ref{contraction:same:perm:isp} to each factor $F_m$, and so it needs every $\bm A_m$ to be an object of $\mVF^{\diamond}[k]$.
\end{rem}

\begin{rem}
The main reason that we need the lift in Lemma~\ref{L:measure:class:lift} to be in that particular form is that, as it will become clear, we are forced to work with explicit compositions of relatively unary proper covariant morphisms, due to the failure of generalizing Lemma~\ref{simul:special:dim:1} below to higher dimensions; see \cite[\S~1]{Yin:int:acvf} for further explanation.

Every definable bijection between definable sets with the same number of $\VF$-coordinates is a composition of relatively unary bijections, see \cite[Lemma~5.6]{Yin:int:acvf}.   It is certainly more satisfying to simply take proper covariant homeomorphisms as morphisms in $\mVF^{\diamond}[k]$ and then show that they are all indeed compositions of relatively unary proper covariant homeomorphisms. We suspect this is not true, although it does not seem straightforward to come up with a counterexample.
\end{rem}


\begin{cor}\label{L:sur:c}
The lifting maps ${\mu}\bb L_{k}$ induce surjective homomorphisms, which are often simply denoted by ${\mu}\bb L$, between the Grothendieck semigroups
\[
 \gsk \mRV^{\db}[k] \epi \gsk \mVF^{\diamond}[k].
\]
\end{cor}
\begin{proof}
By Lemma~\ref{L:measure:class:lift}, every $\mRV^{\db}[k]$-morphism can be lifted to a $\mVF^{\diamond}[k]$-morphism. Thus, ${\mu}\bb L_{k}$ induces a map on the isomorphism classes, which is easily seen to be a semigroup homomorphism. By Lemma~\ref{L:measure:surjective} (see the remark thereafter) or Remark~\ref{stand:vol}, this  homomorphism is surjective.
\end{proof}

\begin{lem}\label{simul:special:dim:1}
Let $f : (A, \omega) \fun (B, \sigma)$ be a $\mgVF^{\diamond}[1]$-morphism.  Then there are special covariant bijections $T_A : A \fun A^{\flat}$, $T_B : B \fun B^{\flat}$ such that $A^{\flat}$, $B^{\flat}$ are doubly bounded $\RV$-pullbacks and the functions $f^{\flat}$, $f^{\flat}_{\downarrow}$ are bijective in
the diagram
\[
\bfig
  \square(0,0)/->`->`->`->/<600,400>[A`A^{\flat}`B`B^{\flat};
  T_A`f``T_B]
 \square(600,0)/->`->`->`->/<600,400>[A^{\flat}`\rv(A^{\flat})`B^{\flat} `\rv(B^{\flat});  \rv`f^{\flat}`f^{\flat}_{\downarrow}`\rv]
 \efig
\]
where both squares commute and $f^\flat=T_B\circ f\circ T_A^{-1}$. The special covariant bijections use the bounded variable-aperture data of Definition~\ref{defn:special:bijection}. Moreover, $\omega' = \omega \circ T_A^{-1}$ and $\sigma' = \sigma \circ T_B^{-1}$ are both $\rv$-contractible, and  $f^{\flat}_{\downarrow}$ is a $\mRV^{\db}[1]$-morphism $(\rv(A^{\flat}), \omega')_{1} \fun (\rv(B^{\flat}), \sigma')_1$ (recall Notation~\ref{0coor}).
\end{lem}
\begin{proof}
First use Lemma~\ref{L:measure:surjective} on both sides, retaining the
original labels and the actual transported forms. Thus we may work
with doubly bounded, regularized RV-pullbacks whose forms are constant
on each retained-label fibre. Take the ordinary unary matching
preparations $H_A,H_B$ from \cite[Lemma~5.2]{Yin:int:acvf}. All their
RV labels remain variables over $\mathbb S$.

An ordinary zero is absorbing. Therefore its exceptional centre family
is exactly the inverse image of the terminal zero coordinate, and the
ordinary matching pairs the two families by $f$. Their VF projections
are finite; in the stated category their mixed centre sets are also
finite. The same construction below works for arbitrary bounded
retained RV fibres, using their finite VF projections.

For paired centres $p,q=f(p)$ choose definable offsets
$d_p,e_q$ and whole discs
\[
 D_p=\go(p,\delta_p),\qquad D_q=\go(q,\epsilon_q),\qquad
 f(D_p)=D_q,\qquad \delta_p=\vv(d_p),\quad\epsilon_q=\vv(e_q).
\]
The paired-disc construction follows from Lemma~\ref{lim:weak},
continuity and injectivity: the punctured image disc is centred at
$f(p)$, and adding the centres gives the equality of whole discs.
Choose the discs beyond the separation and form-constancy thresholds
and beyond every nonzero centre value in the two finite ordinary
orbits. These choices are simultaneous and definable; no single
algebraic centre is named. More explicitly, form the definable
relation of bad positive boundary coefficients. Each centre row is
bounded by the local paired-disc assertion. The uniform bound for a
finite VF projection with a bounded RV tail gives one
definable coefficient beyond every bad row. The finite source and target
VF images give a common separation bound. This constructs the actual
offset maps and paired, disjoint whole discs.

Here is the finite one-sided modification. At each stage keep the
actual preparation $P$, its inverse, a live retained-label locus, and
the corresponding ordinary prefix $H$. For every original core $D_p$
maintain the full affine chart
\[
 P(a)_{\VF}=a_{\VF}-p_{\VF}+P(p)_{\VF},\qquad a\in D_p,
\]
and constancy of all current retained labels on that core. A centre
is live exactly while $H(p)_{\VF}\ne0$. After capture its value is
$\delta_p$, and its whole core is frozen. The live locus is the
complement of precisely the cores already captured; on it $P$ and
$H$ agree in the VF-coordinate and every ordinary label. Every
remaining nonzero ordinary centre value is below $\delta_p$.
These assertions hold initially. Regularization preserves them,
since its new leading term is constant on every prescribed core.

At a centripetal instruction on the live locus, the current source
is an RV-pullback, so its actual focus belongs to its source label
fibre. Pull the focus back by $P^{-1}$. The next ordinary instruction
sends this preimage to zero; absorption and the exceptional-family
identity identify it as an original centre $p$. This is an actual
definable centre map, constant on each current retained-label fibre.
Composing it with $p\mapsto d_p$ supplies the VF endpoint. The affine
chart proves that its cutoff pulls back to \emph{exactly} the whole
original $D_p$. Separation makes that branch constant on every other
prescribed core. The captured centre now has value $\delta_p$ and
its labels freeze; frozen fibres are inactive thereafter. On the
remaining live complement the step is exactly the ordinary translation.
Thus all the invariants persist through the finite word.

The relative aperture at that step is the actual label-definable
ratio $\rv(d_p)/\rv(\lambda)$. Its value is nonnegative by the
chosen orbit bound. The full ratio graph is definable on the bounded
RV-label domain, so Lemma~\ref{db:to:db} gives one named upper bound.
No VF section of an arbitrary RV function is invoked: the endpoint
is the previously constructed offset map. Source invariance supplies
the required cutoff radii. These are precisely admissible data in
Definition~\ref{defn:special:bijection}.

At completion each frozen core image is one whole prepared label
fibre and a full nonzero RV-disc. On the live complement equality
of the prepared labels is equivalent to the original ordinary
RV-equivalence. Perform this construction on both sides. The paired
cores and ordinary live matching show that the literal conjugate
$f^\flat=T_BfT_A^{-1}$ takes every whole label fibre bijectively onto
a whole label fibre, in both directions. No finite-point adjustment
is made. This proves both commuting squares and the bijective RV map.

Finally the preparations retain every original label, so the actual
transported forms $\omega',\sigma'$ remain constant on their final
RV-discs. Lemma~\ref{spec:inve} and the chain rule transfer the
original generic measure identity to $f^\flat$. For any matched
pair $\gp,\gq$, put $\kappa=\omega'(\gp)-\sigma'(\gq)$.
The constant-Jacobian-value comparison for a definable bijection of
whole open discs gives
\[
 \omega'(\gp)-\sigma'(\gq)
 =\rad(\gq)-\rad(\gp)
 =\vv\bigl((f^\flat)'(a)\bigr)
\]
on the good locus. Apply the same comparison to $f:D_p\to D_q$
before the core translations; it gives
$\epsilon_q-\delta_p=\omega(D_p)-\sigma(D_q)$.
Thus exceptional core fibres obey the same radius identity as the
live ones. Every whole source disc contains a good point outside
the definable zero-dimensional exception, so the descended volume
identity holds at \emph{every} RV label. The complete form graphs
and bounded RV domains descend through the actual label quotients.
Consequently $f^\flat_\downarrow$ is the required measured RV
bijection. Attach the initial measured preparations by composition;
their point maps and transported forms give the original objects
and the asserted literal conjugacy.
\end{proof}


\section{Integrating proper invariant sets}
\label{section:int}

\subsection{Annular blowups}
Recall that we write $\RV$ to mean the $\RV$-sort without the  element $\infty$, and $\RV_\infty$ otherwise, although quite often the difference is immaterial and it does not matter which set $\RV$ stands for. In the definition below, this difference does matter, in particular, $\RV^{\ccirc}$ denotes the set $\rv(\MM \mi 0)$.

\begin{nota}\label{punc:RVd}
Recall from Notation~\ref{defn:disc} that the set $\rv(\MM_{\gamma} \mi 0)$ is denoted by $\RV^{\ccirc}_\gamma$. For a finite $\gamma \in \Gamma$ with $\gamma \geq 0$ and $t \in \gamma^\sharp$, write
\[
\RV^{\ccirc}(t)_\gamma = (\RV^{\ccirc} \mi \RV^{\ccirc}_\gamma) \uplus t = \set{s \in \RV \given 0 < \vrv(s) \leq \gamma} \uplus t.
\]
The first part is the \memph{shell} and the second the \memph{endpoint}. They carry distinct tags, so for $\gamma > 0$ the endpoint is not identified with the element $t$ of the shell. If $\gamma = 0$ then the shell is empty and $\RV^{\ccirc}(t)_0 = t$ for every $t \in \RV$ with $\vrv(t) = 0$. In all cases $\RV^{\ccirc}(t)_\gamma$ is doubly bounded; in particular, it does not contain $\infty$.
\end{nota}

% be a definable function and $\tau_*: U \fun$ the function given by
%\[
%t \efun \vrv(\tau(t)) - \vrv(f_j(t)).
%\]

\begin{defn}\label{defn:blowup:coa}
Let $\bm U = (U, f, \omega)$ be an object of $\mRV[k]$, where $k > 0$, such that the function $\pr_{\tilde j} \rest f(U)$ is finite-to-one for some $j \in [k]$. Write $f = (f_1, \ldots, f_k)$. An \memph{aperture} on $\bm U$ consists of a definable function $\tau : U \fun \RV$ and a definable bound $\Delta \in \Gamma$ such that
\[
0 \leq \gamma(u) \coloneqq \vrv(\tau(u)) \leq \Delta \quad \text{for every } u \in U.
\]
As in Definition~\ref{defn:special:bijection}, $\RV$ excludes $\infty$ and $\tau$ need not be constant; its whole graph is part of the data, also when its image is infinite. An \memph{elementary annular blowup} of $\bm U$ in the $j$th coordinate of \memph{aperture} $\tau$ is the triple $\bm U^{\flat} = (U^{\flat}, f^{\flat}, \omega^\flat)$, where
\[
U^{\flat} = \bigcup_{u \in U} u \times \RV^{\ccirc}(\tau(u))_{\gamma(u)} = \set{(u, s) \given u \in U,\ 0 < \vrv(s) \leq \gamma(u)} \uplus \set{(u, \tau(u)) \given u \in U},
\]
the shell and the endpoint graph carrying distinct tags, and for every $(u, s) \in U^{\flat}$,
\[
f^{\flat}_{i}(u, s) = f_{i}(u) \text{ for } i \neq j, \quad f^{\flat}_{j}(u, s) = s f_{j}(u), \quad \omega^\flat(u, s) = \omega(u).
\]
If $f^\flat(u, s) = r$ then $f_{\tilde j}(u) = r_{\tilde j}$. Since $\pr_{\tilde j} \rest f(U)$ and $f$ are finite-to-one, only finitely many $u$ qualify, and over each of them $f^\flat$ takes the value $r$ at most once on the shell and at most once on the endpoint graph. Thus $f^\flat : U^\flat \fun \RV^k$ is finite-to-one and $\bm U^\flat$ is an object of $\mRV[k]$. A constant aperture $\tau(u) = t$ gives $U^\flat = U \times \RV^{\ccirc}(t)_{\vrv(t)}$. If $\bm U \in \mRV^{\db}[k]$, then $\bm U^\flat$ is the relative annular object $\bm U^{\tau, j}$ of Notation~\ref{relative:annular:ideal}.


Let $\bm V = (V, g, \sigma)$ be another object of $\mRV[k]$ and $C_i \sub V$ finitely many pairwise disjoint definable sets. Each triple
\[
\bm C_i = (C_i, g \rest C_i, \sigma \rest C_i) \in \mRV[k]
\]
is referred to as a \memph{subobject} of $\bm V$. Suppose that $F_i : \bm U_i \fun \bm C_i$ is a $\mRV[k]$-morphism and $\bm U^{\flat}_i$ is an elementary annular blowup of $\bm U_i$. Let $C = \bigcup_i C_i$, $\bm C = \bigcup_i \bm C_i$,  and $F = \biguplus_i F_i$. Then the object
$
 (\bm V \mi \bm C) \uplus \biguplus_i \bm U_i^{\flat}
$
is  an \memph{annular blowup of $\bm V$ via $F$}, denoted by $\bm V^{\flat}_F$. The subscript $F$ is usually dropped. The object $\bm C$ (or the set $C$) is referred to as the \memph{locus} of $\bm V^{\flat}_F$.

An \memph{annular blowup of length $n$} is a composition of $n$ annular blowups.
\end{defn}

\begin{rem}
Infinite apertures are excluded. Formally, the aperture $\infty$ would replace $\RV^{\ccirc}(t)_\gamma$ by $\RV^{\ccirc}_\infty = \rv(\MM)$, and the construction would become an elementary blowup in the sense of \cite[Definition~6.1]{Yin:int:acvf}. Then $f^\flat_j$ takes the value $\infty$ at the endpoint, so that piece lies in $\RV[{\leq}k]$ rather than in $\RV[k]$ (Notation~\ref{0coor}), and the shell has unbounded valuation. This is the different grading convention of \cite{Yin:int:acvf}; the comparison is only heuristic here. If there is an elementary annular blowup of $\bm U$ then, \textit{a posteriori}, $\dim_{\RV}(U) < k$. For any coordinate of $f(U)$, there is at most one elementary blowup of $\bm U$, whereas there could be many annular elementary blowups of $\bm U$. We should have included the coordinate that is blown up as a part of the data. However, in context, either this is
clear or it does not need to be spelled out, and we shall suppress mention of it for ease of notation.
\end{rem}

\begin{rem}
\label{db:blowup:bound}If $\bm U \in \mRV^{\db}[k]$ then $\bm U^{\flat} \in \mRV^{\db}[k]$ too, because the new coordinate satisfies $0 < \vrv(s) \leq \Delta$. Conversely, for such $\bm U$ the bound $\Delta$ is automatic: by Lemma~\ref{db:to:db}, every definable $\tau : U \fun \RV$ has doubly bounded image, and a bounded definable subset of $\Gamma$ has a definable upper bound. We shall only consider annular blowups  in $\mRV^{\db}[k]$ (then, for simplicity, they will just be referred to as blowups). This notion is dual to that of a  special covariant bijection. As a matter of fact, to make this duality precise is essentially what is left to do for the rest of our main construction.
\end{rem}


For the next few lemmas, let $\bm U = (U, f, \omega)$ and $\bm V = (V, g, \sigma)$ be objects of $\mRV^{\db}[k]$.

%and the aperture $\alpha$ of $\bm U^{\flat}$ is no greater than the aperture $\beta$ of  $\bm V^{\flat}$

\begin{lem}\label{elementary:blowups:preserves:iso:vol}
Suppose that $[\bm U] = [\bm V]$ and $\bm U^{\flat}$, $\bm V^{\flat}$ are elementary blowups of $\bm U$, $\bm V$, in possibly different coordinates. Then there are  blowups $\bm U^{\flat\flat}$, $\bm V^{\flat\flat}$ of $\bm U^{\flat}$, $\bm V^{\flat}$ of length $1$ such that $[\bm U^{\flat\flat}] = [\bm V^{\flat\flat}]$.
\end{lem}
\begin{proof}
Fix an isomorphism $F : \bm U \fun \bm V$. Let $\tau_U$, $\tau_V$ be the apertures of $\bm U^{\flat}$, $\bm V^{\flat}$, in the $j$th and $j'$th coordinates, with depths $\gamma_U = \vrv \circ \tau_U$, $\gamma_V = \vrv \circ \tau_V$ and bounds $\Delta_U$, $\Delta_V$. By \eqref{cond:mg2}, a $\mRV[k]$-morphism is a definable bijection that preserves the \memph{total weight} $\omega + \Sigma \vrv \circ f$; which coordinate carries which factor plays no role, as in \cite[Lemma~6.3]{Yin:int:acvf}. At a point $(u, s)$ of $\bm U^{\flat}$ the total weight is $\omega(u) + \Sigma \vrv(f(u)) + \vrv(s)$, and similarly for $\bm V^{\flat}$.

Put $\tau' = \tau_V \circ F$ and $\gamma' = \gamma_V \circ F$, and partition $U$ into the definable sets
\[
U_{\leq} = \set{u \in U \given \gamma_U(u) \leq \gamma'(u)}, \qquad U_{>} = \set{u \in U \given \gamma_U(u) > \gamma'(u)}.
\]
Let $\bm W$ be the subobject of $\bm U^{\flat}$ formed by the endpoints $(u, \tau_U(u))$ with $u \in U_{\leq}$. Its structural coordinates are $\tau_U(u) f_j(u)$ in position $j$ and $f_i(u)$ elsewhere. Since $f_{\tilde j}$ is finite-to-one on $U$, the projection $\pr_{\tilde j}$ is finite-to-one on the image of $\bm W$, so $\bm W$ admits elementary blowups in the $j$th coordinate. The function $\rho = \tau'/\tau_U$ on $\bm W$ is an aperture, since $0 \leq \vrv \circ \rho = \gamma' - \gamma_U \leq \Delta_V$. Let $\bm U^{\flat\flat}$ be the blowup of $\bm U^{\flat}$ with locus $\bm W$ and aperture $\rho$. Over $u \in U_{\leq}$ it consists of the old shell, the new shell of parameters $s'$ with $0 < \vrv(s') \leq \gamma'(u) - \gamma_U(u)$, and the new endpoint $\rho(u)$. The map
\[
(u, s) \efun (F(u), s), \quad (u, \tau_U(u), s') \efun (F(u), \tau_U(u) s'), \quad (u, \tau_U(u), \rho(u)) \efun (F(u), \tau_V(F(u)))
\]
is a definable bijection onto the part of $\bm V^{\flat}$ over $F(U_{\leq})$: the first two pieces fill the shell $0 < \vrv(s) \leq \gamma'(u)$, and the third goes to the endpoint. It preserves total weights, because $\vrv(\tau_U(u) s') = \gamma_U(u) + \vrv(s')$ and the new endpoint has total depth $\gamma_U(u) + \vrv(\rho(u)) = \gamma'(u)$. When $\gamma_U(u) = \gamma'(u)$ the new shell is empty and $\rho(u)$ is an endpoint of valuation $0$, possibly different from $1$; this covers equal depths with different endpoints.

Symmetrically, blow up $\bm V^{\flat}$ at its endpoints over $F(U_{>})$ with the aperture $(\tau_U \circ F^{-1})/\tau_V$, whose depth lies in $(0, \Delta_U]$. This gives $\bm V^{\flat\flat}$ and an isomorphism from its part over $F(U_{>})$ onto the part of $\bm U^{\flat}$ over $U_{>}$. Elsewhere $\bm U^{\flat\flat}$ and $\bm V^{\flat\flat}$ coincide with $\bm U^{\flat}$ and $\bm V^{\flat}$. Both are blowups of length $1$; if one of the loci is empty, the corresponding blowup is the identity, which Definition~\ref{defn:blowup:coa} permits. Gluing the two isomorphisms gives $[\bm U^{\flat\flat}] = [\bm V^{\flat\flat}]$.
\end{proof}

\begin{cor}\label{eblowup:same:loci:iso}
Let $F : \bm U \fun \bm V$ be a morphism and $\bm U^{\flat}$, $\bm V^{\flat}$ blowups of  $\bm U$, $\bm V$ of length $1$. Then there are blowups $\bm U^{\flat\flat}$, $\bm V^{\flat\flat}$ of $\bm U^{\flat}$, $\bm V^{\flat}$ of length $1$ such that $[\bm U^{\flat\flat}] = [\bm V^{\flat\flat}]$.
\end{cor}
\begin{proof}
Let $C$, $D$ be the loci of $\bm U^{\flat}$, $\bm V^{\flat}$, given by morphisms $F_i : \bm U_i \fun \bm C_i$, $G_l : \bm V_l \fun \bm D_l$ and elementary blowups $\bm U_i^{\flat}$, $\bm V_l^{\flat}$. Refine $C \cap F^{-1}(D)$ and $F(C) \cap D$ into the finitely many pieces $C_i \cap F^{-1}(D_l)$ and $F(C_i) \cap D_l$. On each piece both blowups restrict to elementary blowups of isomorphic objects, with the restricted apertures and possibly different coordinates. Lemma~\ref{elementary:blowups:preserves:iso:vol} gives blowups of length $1$ of the two restricted pieces with isomorphic results; since these pieces are disjoint summands of $\bm U^{\flat}$, $\bm V^{\flat}$, they are blowups of $\bm U^{\flat}$, $\bm V^{\flat}$. In the meantime, blow up $\bm U^{\flat}$ at the locus $F^{-1}(D) \mi C$, which is still a subobject of $\bm U^{\flat}$, via the morphisms $F^{-1} \circ G_l$ and the same elementary blowups $\bm V_l^{\flat}$, restricted. Symmetrically, blow up $\bm V^{\flat}$ at the locus $F(C) \mi D$ via $F \circ F_i$ and $\bm U_i^{\flat}$. These pieces then agree on the two sides, and $F$ matches the parts outside all loci. Thus the resulting compounded blowups of $\bm U^{\flat}$, $\bm V^{\flat}$ of length $1$ are as desired.
\end{proof}


This corollary also holds in $\RV[k]$, which is essentially the only thing that the otherwise formal proof of \cite[Lemma~6.5]{Yin:int:acvf} depends on. Thus, the same proof yields the following analogue in the present context:


\begin{lem}\label{blowup:equi:class:coa}
If $[\bm U] = [\bm V]$ and $\bm U_1$, $\bm V_1$ are blowups of $\bm U$, $\bm V$ of lengths $m$, $n$, respectively, then there are blowups $\bm U_2$, $\bm V_2$ of $\bm U_1$, $\bm V_1$ of lengths $n$, $m$, respectively, such that $[\bm U_2] = [\bm V_2]$.
\end{lem}

\begin{cor}\label{blowup:equi:class}
Suppose that $[\bm U] = [\bm U']$ and $[\bm V] = [\bm V']$. If there are isomorphic blowups of $\bm U$, $\bm V$ then there are isomorphic blowups of $\bm U'$, $\bm V'$.
\end{cor}

\begin{defn}
Let $\mispdb[k]$ be the set of pairs $(\bm U, \bm V)$ of objects of $\mRV^{\db}[k]$ such that there exist isomorphic blowups $\bm U^{\flat}$, $\bm V^{\flat}$. Let $\mispdb[*] = \bigoplus_{k} \mispdb[k]$.
\end{defn}

We will just write $\mispdb$ for all these sets when there is no danger of confusion. By Corollary~\ref{blowup:equi:class}, they may be regarded as binary relations on isomorphism classes.


\begin{lem}\label{isp:congruence:vol}
$\mispdb[k]$ is a semigroup congruence relation and $\mispdb[*]$ is a semiring congruence relation.
\end{lem}
\begin{proof}
This is an easy consequence of Lemma~\ref{blowup:equi:class:coa} (the proof of \cite[Lemma~6.8]{Yin:int:acvf}, which uses \cite[Lemma~6.5]{Yin:int:acvf}, contains more details).
\end{proof}


Let  $T$ be a special covariant bijection on $\bb L (U,f)$. Denote the object $(\mathbb{L} (U,f)^\flat_{\RV})_{\leq k} \in \RV^{\db}[k]$ by $(U_{T})_{\leq k}$ and  the object $(U_{T}, \omega_T)_{\leq k} \in \mRV^{\db}[k]$ by $\bm U_{T}$, where $\omega_T : U_{T} \fun \Gamma$ is the volume form induced by $\omega$.


\begin{lem}\label{special:to:blowup:coa}
The object $\bm U_T$ is isomorphic to a blowup of $\bm U$ of the same length as $T$.
\end{lem}
\begin{proof}
By induction on the length $\lh (T)$ of $T$ and Lemma~\ref{blowup:equi:class:coa}, this is immediately reduced to the case $\lh (T) = 1$; regularizations only add redundant leading-term labels and do not change $\bm U_T$ up to isomorphism. The locus of $T$ is an $\RV$-pullback in $\bb L(U, f)$, hence of the form $\bb L(U_C, f \rest U_C)$ for a subobject $\bm U_C$ of $\bm U$, and $T$ is the identity off it. So $\bm U_T$ is $(\bm U \mi \bm U_C) \uplus (\bm U_C)_T$, and we may assume that the locus is $\bb L(U, f)$.

Let $j$ be the coordinate changed by $T$, $\lambda$ its focus and $(\tau, \Delta)$ its aperture. Since $\bb L(U, f)$ is regularized, the retained label of the polydisc $f(u)^\sharp \times u$ is $u$, so $\tau$ is a function on $U$. We first check that $\pr_{\tilde j} \rest f(U)$ is finite-to-one, as Definition~\ref{defn:blowup:coa} requires. Fix $r \in \pr_{\tilde j}(f(U))$ and a $\VF$-tuple $a \in r^\sharp$. For every $u$ with $f_{\tilde j}(u) = r$, the focus is defined at $(a, u)$ and $\rv(\lambda(a, u)) = f_j(u)$. By Lemma~\ref{dcl:to:ac}, applied to one $\RV$-coordinate of $u$ at a time, each value $\lambda(a, u)$ is algebraic over $a$. So the $a$-definable set of these values is finite, by compactness, and only finitely many values $f_j(u)$ lie over $r$.

Now fix $u \in U$ and $a \in f_{\tilde j}(u)^\sharp$, and let $x$ range over $f_j(u)^\sharp$. Then $\vv(x - \lambda(a, u)) > \vrv(f_j(u))$, and the cutoff of \eqref{centri} is $\vrv(f_j(u)) + \gamma(u)$. On the first branch $\rv(x - \lambda(a, u)) = s f_j(u)$ with $0 < \vrv(s) \leq \gamma(u)$, and every such $s$ occurs; the core goes to the tagged endpoint $\tau(u) f_j(u)$. The other leading terms remain $f_{\tilde j}(u)$ and the label $u$ is retained. Hence $(u, s) \efun (u, s f_j(u))$, with the tags retained, is an isomorphism from the elementary blowup $\bm U^\flat$ of $\bm U$ in the $j$th coordinate of aperture $\tau$ onto $\bm U_T$: the structural coordinates agree, both forms are $\omega(u)$, and the bound $\Delta$ of $T$ is the bound required in Definition~\ref{defn:blowup:coa}.
\end{proof}


\begin{lem}\label{blowup:same:RV:coa}
Let $\bm U^{\flat}$ be a blowup of $\bm U$. Then $\mL \bm U^{\flat}$ is isomorphic to $\mL \bm U$ in $\mVF^{\diamond}[k]$.
\end{lem}
\begin{proof}
By induction on the length of $\bm U^{\flat}$, it suffices to treat a blowup of length $1$. If its locus $\bm C = \bigcup_i \bm C_i$ is given by morphisms $F_i : \bm U_i \fun \bm C_i$ and elementary blowups $\bm U_i^\flat$, then $\mL \bm U^\flat$ is the disjoint union of $\mL(\bm U \mi \bm C)$ and the $\mL \bm U_i^\flat$, while Lemma~\ref{L:measure:class:lift} lifts each $F_i$ to a $\mVF^\diamond[k]$-morphism $\mL \bm U_i \fun \mL \bm C_i$. So it suffices to treat an elementary blowup $\bm U^\flat$, say in the $j$th coordinate with aperture $\tau$.

A point of $\tau(u)^\sharp$ cannot be chosen definably in $u$ alone, so the centres are chosen over the inactive $\VF$-input. The definable function
\[
r \efun \set{(f_j(u), \tau(u) f_j(u)) \given f_{\tilde j}(u) = r}
\]
on $\pr_{\tilde j}(f(U))$ is finitary, because $\pr_{\tilde j} \rest f(U)$ and $f$ are finite-to-one. By Lemma~\ref{RVlift}, in its continuous-branch form, it lifts to continuous, partially differentiable centre functions: for $u \in U$ and $a \in f_{\tilde j}(u)^\sharp$ we obtain $b_0(a, u) \in f_j(u)^\sharp$ and $b_1(a, u) \in (\tau(u) f_j(u))^\sharp$, continuous in $a$. For $x \in f_j(u)^\sharp$ put
\[
\Phi(x, a, u) =
\begin{cases}
(x - b_0(a, u), a, u, s), & \vv(x - b_0(a, u)) \leq \vv(b_1(a, u)),\\
(x - b_0(a, u) + b_1(a, u), a, u, \tau(u)), & \vv(x - b_0(a, u)) > \vv(b_1(a, u)),
\end{cases}
\]
where $s = \rv(x - b_0(a, u)) / f_j(u)$ on the first branch, and the two branches carry the tags of the shell and of the endpoint. Since $\vv(x - b_0(a, u)) > \vrv(f_j(u))$ and $\vv(b_1(a, u)) = \vrv(f_j(u)) + \gamma(u)$, the first branch maps onto the part of $\mL \bm U^\flat$ over the shell, where $0 < \vrv(s) \leq \gamma(u)$, and the second maps onto $b_1(a, u) + \MM_{\vv(b_1(a, u))} = (\tau(u) f_j(u))^\sharp$. At depth zero only the second branch occurs. The inverse adds $b_0$, and also subtracts $b_1$ on the second branch.

On a fixed label the branch condition is clopen, because $\vv(b_1(a, u))$ is constant there, and each branch is a translation in the $j$th coordinate by a continuous function of $a$. Thus $\Phi$ is a relatively unary continuous bijection whose graph fibers have open domains and ranges, and Lemmas~\ref{conti:covar} and~\ref{conv:unif} show that it is a proper covariant homeomorphism. Where $b_0$ and $b_1$ are totally differentiable, its total derivative is triangular with diagonal entries $1$; as in Lemma~\ref{spec:inve}, the exceptional set has $\VF$-dimension less than $k$. Both forms are $\omega(u)$, so $\Phi$ is a $\mVF^\diamond[k]$-morphism $\mL \bm U \fun \mL \bm U^\flat$. For a constant aperture $t = \rv(c)$ one may take $b_1 = c b_0$.
\end{proof}


\begin{lem}\label{kernel:dim:1:coa}
Let $[\bm A] = [\bm B]$ in  $\gsk \mVF^{\diamond}[1]$ and $\bm U, \bm V \in \mRV^{\db}[1]$ be standard contractions of $\bm A$, $\bm B$, respectively. Then $([\bm U], [\bm V]) \in \mispdb$.
\end{lem}
\begin{proof}
This is immediate by Lemmas~\ref{simul:special:dim:1} (applied to $\mL \bm U$, $\mL \bm V$) and \ref{special:to:blowup:coa}.
\end{proof}

\begin{lem}\label{isp:VF:fiberwise:contract}
Let $A'$, $A''$ be definable sets with $A'_{\VF} = A''_{\VF} = A \sub \VF^n$ and $\omega'$, $\omega''$ definable functions from $A'$, $A''$ into $\Gamma$, respectively. Write $\bm A' = (A', \omega')$ and $\bm A'' = (A'', \omega'')$. Suppose that $\bm A'$, $\bm A''$ are proper invariant and there is a $k \in \N$ such that for  every $a \in A$,
\begin{equation}\label{condi:fiber}
([(A'_a, \omega')]_{\leq k}, [(A''_a, \omega'')]_{\leq k}) \in \mispdb.
\end{equation}
Let $\hat T_{\sigma}$, $\hat R_{\sigma}$ be standard contractions of $\bm A'$, $\bm A''$, respectively. Then
\[
([\hat T_{\sigma}(\bm A')]_{\leq n+k}, [\hat R_{\sigma}(\bm A'')]_{\leq n+k}) \in \mispdb.
\]
\end{lem}

Note that condition (\ref{condi:fiber}) makes sense only  over the substructure $\mathbb{S} \la a \ra$.

\begin{proof}
If $n = 0$ there is nothing to prove. Suppose $n > 1$ and put $i = \sigma(1)$. For $b \in \pr_{\tilde i}(A)$, the restrictions of $T_{\sigma(1)}$, $R_{\sigma(1)}$ to the fibers $\bm A'_b$, $\bm A''_b$ are standard contractions of these objects over $\mathbb S \la b \ra$, which is again $\VF$-generated, and (\ref{condi:fiber}) holds for them. So the case $n = 1$, applied over $\mathbb S \la b \ra$, gives
\[
([\hat T_{\sigma(1)}(\bm A')_b]_{\leq 1+k}, [\hat R_{\sigma(1)}(\bm A'')_b]_{\leq 1+k}) \in \mispdb
\]
for every such $b$. The remaining stages of $\hat T_\sigma$, $\hat R_\sigma$ are standard contractions of $\hat T_{\sigma(1)}(\bm A')$, $\hat R_{\sigma(1)}(\bm A'')$, which have $n-1$ $\VF$-coordinates, and the inductive hypothesis with head start $k+1$ gives the claim.

So let $n = 1$ and $A \sub \VF$. We follow the proof of \cite[Lemma~6.14]{Yin:int:acvf} (the proof of \cite[Lemma~5.36]{Yin:tcon:I} is formally the same and has more details), and spell out the uniformity of the witnesses, which is the additional point in the present setting.

\emph{Uniform witnesses.} For $a \in A$, a witness for (\ref{condi:fiber}) consists of two finite chains of blowups of $(A'_a, \omega')_{\leq k}$ and $(A''_a, \omega'')_{\leq k}$, with all their loci, source objects and morphisms, coordinates, aperture graphs and bounds and intermediate objects, together with a final isomorphism and its inverse; all of these are $\mathbb S \la a \ra$-definable. Fix a finite tuple of formulas with parameter $a$ and explicit bounds on the sizes of the fibers in the finite-to-one conditions. Being a witness is then a first-order condition on $a$; $\Gamma$-valued forms and bounds are encoded by their whole $\vrv$-fibers, so no representative is named. By compactness, finitely many such codes cover $A$. Let $D_a$ be the witness given by the first valid code, packed with distinct finite tags into one definable row $D \sub A \times \RV^l$, with shorter chains padded by blowups with empty locus.

\emph{Freezing.} Apply the preparation in the proof of Proposition~\ref{simplex:with:hole:rvproduct} to $A$, against its $\RV$-hull, and to the complete data
\[
a \efun (D_a, A'_a, A''_a, \omega' \rest A'_a, \omega'' \rest A''_a).
\]
Lemma~\ref{special:bi:term:constant} gives a special bijection on $A$ along which the whole membership formula is $\rv$-contractible, no bound on $D$ being required, and FMT on finitely many open discs of a common invariance radius $\delta$ of $\bm A'$, $\bm A''$ and their form levels turns it into a special covariant bijection $F : A \fun A^\sharp$. On a frozen disc $\ga$ with chosen point $a_\ga$, replace the whole witness by $D_{a_\ga}$. For $a \in \ga$, $\delta$-invariance gives $A'_a = A'_{a_\ga}$, $A''_a = A''_{a_\ga}$ and the same forms, so $D_{a_\ga}$ is a witness at $a$ as well; no locus, aperture or comparison map is selected independently of the others. Consequently, for every $\RV$-polydisc $\gp \sub A^\sharp$, the complete row $D_a$, the fibers $A'_a$, $A''_a$ and the forms do not depend on $a \in F^{-1}(\gp)$. In particular, for all $a_1, a_2 \in F^{-1}(\gp)$,
\[
\omega' \rest A'_{a_1} = \omega' \rest A'_{a_2} \dand \omega'' \rest A''_{a_1} = \omega'' \rest A''_{a_2}.
\]

\emph{Descent.} Let $\hat F'$, $\hat F''$ be the standard contractions of $\bm A'$, $\bm A''$ induced by $F$; for each $\RV$-tuple $x$, the set $F(\set{a \given x \in A'_a})$ is a union of polydiscs $\gp$ and hence an $\RV$-pullback, and similarly for $\bm A''$. The row over $\gp$ is definable from the $\RV$-code of $\gp$ and the retained tags, so the loci, aperture graphs, intermediate objects and isomorphisms glue over these codes into $\mathbb S$-definable data. Validity descends: totality, functionality, the inverse identities and the finite-to-one conditions hold fiberwise, and the added coordinate $\rv(F(a))$ is the same on both sides and adds the same weight, so the total weights still match. Every glued intermediate object is a blowup of a doubly bounded object and hence doubly bounded, so each glued aperture function has a definable bound by Remark~\ref{db:blowup:bound}; no bound uniform in the parameter $a$ is assumed. Thus
\[
([\hat F'(\bm A')]_{\leq 1+k}, [\hat F''(\bm A'')]_{\leq 1+k}) \in \mispdb.
\]

\emph{Comparison.} It remains to compare $\hat T_\sigma$ with $\hat F'$, and $\hat R_\sigma$ with $\hat F''$. The bijection $F \circ T_\sigma^{-1}$ between $T_\sigma(A')$ and $F(A')$ is relatively unary in the $\VF$-coordinate and retains the head $\RV$-coordinates as labels. Lemma~\ref{simul:special:dim:1}, whose proof allows arbitrary bounded retained $\RV$-fibers, gives special covariant bijections on both sides whose contractions are matched by a $\mRV^{\db}[1+k]$-morphism. By Lemma~\ref{special:to:blowup:coa}, applied in this $\VF$-coordinate (its proof only uses the coordinate being changed), these contractions are isomorphic to blowups of $\hat T_\sigma(\bm A')_{\leq 1+k}$ and $\hat F'(\bm A')_{\leq 1+k}$. This is the argument of Lemma~\ref{kernel:dim:1:coa}, used in place of \cite[Corollary~6.11]{Yin:int:acvf}. Hence $\hat T_\sigma(\bm A')$ and $\hat F'(\bm A')$ are $\mispdb$-congruent, and similarly on the other side. The lemma follows by Lemma~\ref{isp:congruence:vol}.
\end{proof}


\begin{cor}\label{contraction:same:perm:isp}
Let $\bm A', \bm A'' \in \mVF^{\diamond}[n]$ and suppose that there is a morphism $F : \bm A' \fun \bm A''$ that is relatively unary in the $i$th $\VF$-coordinate. Then for any permutation $\sigma$ of $[n]$ with $\sigma(1) = i$ and any standard contractions $\hat T_{\sigma}$, $\hat R_{\sigma}$ of $\bm A'$, $\bm A''$,
\[
([\hat T_{\sigma}(\bm A')]_{\leq n}, [\hat R_{\sigma}(\bm A'')]_{\leq n}) \in \mispdb.
\]
\end{cor}
\begin{proof}
Since $F$ is relatively unary in the $i$th coordinate, $\pr_{\tilde i}(A') = \pr_{\tilde i}(A'')$; call this common base $A$, and for $a \in A$ let $F_a : A'_a \fun A''_a$ be the induced bijection. If $n = 1$ the claim is Lemma~\ref{kernel:dim:1:coa}. Let $E \sub A'$ be the definable set where $F$ is not totally differentiable or its measure identity fails; then $\dim_{\VF}(E) < n$, so $F$ need not induce a morphism $\bm A'_a \fun \bm A''_a$ for every $a \in A$. Call $a \in A$ \memph{good} if $E_a$ is finite. At a good $a$, $F_a$ is a $\mVF^\diamond[1]$-morphism over $\mathbb S \la a \ra$: it is a fiber of a relatively unary proper covariant homeomorphism, and off the finite set $E_a$ its derivative is the Jacobian of $F$. The restrictions of $\hat T_{\sigma(1)}$, $\hat R_{\sigma(1)}$ to the fibers are standard contractions of $\bm A'_a$, $\bm A''_a$, so Lemma~\ref{kernel:dim:1:coa} gives
\[
([\hat T_{\sigma(1)}(\bm A')_a]_{\leq 1}, [\hat R_{\sigma(1)}(\bm A'')_a]_{\leq 1}) \in \mispdb \quad \text{over } \mathbb S \la a \ra.
\]
The set $B$ of the remaining points of $A$ has $\VF$-dimension at most $n-2$, so it lies in the zero set of a nonzero polynomial $P$ over $\mathbb S$; let $d = \deg P$.

Choose a definable $\delta \in \Gamma$ such that $A'$, $A''$ and the first-stage contractions $\hat T_{\sigma(1)}(\bm A')$, $\hat R_{\sigma(1)}(\bm A'')$, together with their volume forms, are all $\delta$-invariant; invariance of one side alone would not suffice. Since $\Gamma(\mathbb S)$ is nontrivial, there are distinct $c_0, \ldots, c_d \in \VF(\mathbb S)$ with $\vv(c_l) > \delta$. A nonzero polynomial of degree at most $d$ in $n-1$ variables cannot vanish on all of $a + \set{c_0, \ldots, c_d}^{n-1}$. So for each $a \in A$ the lexicographically first point $b$ of this grid outside the zero set of $P$ exists. It is good, lies in $A$ and in the open polydisc of radius $\delta$ about $a$, and belongs to $\dcl(\mathbb S, a)$. By $\delta$-invariance the contracted fibers of both sides at $a$ and at $b$ coincide, forms included. Hence the displayed congruence holds at every $a \in A$, over $\mathbb S \la b \ra \sub \mathbb S \la a \ra$, as condition (\ref{condi:fiber}) requires. Lemma~\ref{isp:VF:fiberwise:contract}, applied to $\hat T_{\sigma(1)}(\bm A')$ and $\hat R_{\sigma(1)}(\bm A'')$ with $n-1$ $\VF$-coordinates and head start $1$, and to the remaining stages of $\hat T_\sigma$, $\hat R_\sigma$, gives the claim.
\end{proof}

\subsection{$2$-cells}

\begin{defn}[$\vv$-affine and $\rv$-affine]\label{rvaffine}
Let $\ga$ be an open disc and $f : \ga \fun \VF$ an injection.
We say that $f$ is \memph{$\vv$-affine} if there is a (necessarily unique) $\gamma \in \Gamma$, called the \memph{shift} of $f$, such that, for all $a, a' \in \ga$,
\[
\vv(f(a) - f(a')) = \gamma + \vv(a - a').
\]
We say that $f$ is \memph{$\rv$-affine} if there is a (necessarily unique) $t \in \RV$, called the \memph{slope} of $f$, such that, for all $a, a' \in \ga$,
\[
\rv(f(a) - f(a')) = t \rv(a - a').
\]
\end{defn}

\begin{defn}\label{defn:balance}
Let $A \sub \VF^2$ be a definable set such that $\ga_1 \coloneqq \pr_1(A)$ and $\ga_2 \coloneqq \pr_2(A)$ are both open discs. Let $f : \ga_1 \fun \ga_2$ be a definable bijection that has the disc-to-disc property.  We say that $f$ is \memph{balanced in $A$} if $f$ is actually $\rv$-affine and there are $t_1, t_2 \in \RV_\infty$, called the \memph{paradigms} of $f$, such that, for every $a \in \ga_1$,
\[
A_a = t_2^\sharp + f(a) \dand f^{-1}(A_a) = a - t_1^\sharp.
\]
\end{defn}

\begin{rem}\label{2cell:linear}
Suppose that $f$ is balanced in $A$ with paradigms $t_1$, $t_2$.
If one of the paradigms is $\infty$ then the other one must be $\infty$. In this case $A$ is just the (graph of the) bijection $f$ itself.

Assume that $t_1$, $t_2$ are not $\infty$. Let $\gB_1$, $\gB_2$ be the sets of closed subdiscs of $\ga_1$, $\ga_2$ of radii $\vrv(t_1)$, $\vrv(t_2)$, respectively. Let $a_1 \in \gb_1 \in \gB_1$ and $\go_1$ be the maximal open subdisc of $\gb_1$ containing $a_1$. Let $\gb_2 \in \gB_2$ be the smallest closed disc containing the open disc $\go_2 \coloneqq A_{a_1}$. Then, for all $a_2 \in \go_2$, we have
\[
\go_2 = t_2^\sharp + f(\go_1) = A_{a_1} \dand A_{a_2} = f^{-1}(\go_2) + t_1^\sharp = \go_1.
\]
This internal symmetry of $A$ is illustrated by the following diagram:
\[
\bfig
  \dtriangle(0,0)|amb|/.``<-/<600,250>[\go_1`f^{-1}(\go_2)`\go_2; \pm t_1^\sharp`\times`f^{-1}]
  \ptriangle(600,0)|amb|/->``./<600,250>[\go_1`f(\go_1)`\go_2; f`` \pm t_2^\sharp]
 \efig
\]
Since $f$ is $\rv$-affine, we see that its slope must be $-t_2/t_1$.

Let $\tor(\code {\gb_1})$, $\tor(\code {\gb_2})$ denote the sets of the maximal open subdiscs of $\gb_1$, $\gb_2$. These may be viewed as $\K$-torsors and as such  are equipped with much of the structure of $\K$. Then the set $A \cap (\gb_1 \times \gb_2)$ may be thought of as the ``line'' in $\tor(\code {\gb_1}) \times \tor(\code {\gb_2})$ given by the equation
\[
x_2 = - \tfrac{t_2}{t_1}(x_1 - \code{\go_1}) + \code{\go_2}.
\]
Thus, by Lemma~\ref{simul:special:dim:1}, the obvious bijection between $A_1 \times t_2^\sharp$ and $t_1^\sharp \times A_2$ is the lift of an $\RV^{\db}[2]$-morphism modulo special covariant bijections; see Lemma~\ref{2:unit:contracted} below for details.
\end{rem}

\begin{defn}\label{def:units}
We say that a set $A$ is a \memph{$1$-cell} if it is either an open disc contained in a single $\RV$-disc or a point in $\VF$. We say
that $A$ is a \memph{$2$-cell} if
\begin{itemize}
 \item $A$ is a subset of $\VF^2$ contained in a single $\RV$-polydisc and $A_1$ is a $1$-cell,
 \item there is a function $\epsilon : A_1  \fun \VF$ and a $t \in \RV$ such that, for every $a \in A_1$, $A_a = t^\sharp + \epsilon(a)$,
 \item one of the following three possibilities occurs:
  \begin{itemize}
   \item $\epsilon$ is constant,
   \item $\epsilon$ is injective, has the disc-to-disc property, and $\rad(\epsilon(A_1)) \geq \vrv(t)$,
   \item $\epsilon$ is balanced in $A$.
  \end{itemize}
\end{itemize}
The function $\epsilon$ is called the \memph{positioning function} of $A$ and the element $t$ the \memph{paradigm} of $A$.

More generally, a set $A$ with exactly one $\VF$-coordinate is a \memph{$1$-cell} if, for each $t \in A_{\RV}$, $A_t$ is a $1$-cell in the above sense; the parameterized version of the notion of a $2$-cell is formulated in the same way.
\end{defn}

\begin{lem}\label{decom:into:2:units}
Let  $A \sub \VF^2$ be a proper invariant set. Then there is a standard contraction $\hat T_\sigma$ of $A$ such that $A_s$ is a $2$-cell for every $s \in \hat T_\sigma(A)$.
\end{lem}
\begin{proof}
This is a variation of \cite[Lemma~4.8]{Yin:int:acvf}. The proof of the latter proceeds by constructing a positioning function $\epsilon_{(t,s)}$ in each $\VF$-fiber, which heavily relies on the use of standard contractions in the preceding auxiliary results, namely \cite[Lemmas~4.2, 4.4]{Yin:int:acvf}. It still goes through here, provided that FMT is applied at suitable places to cut out finitely many exceptional open discs in the first coordinate at which the desired properties of $\epsilon_{(t,s)}$ do not hold, as in the proof of Proposition~\ref{simplex:with:hole:rvproduct}. The positioning functions so obtained are $\vv$-affine. The following refinement of the final unary stage, in the first coordinate, makes them $\rv$-affine.

Fix a positioning function $\epsilon$, with its labels as parameters. For $x \in \dom(\epsilon)$ and $\beta \in \Gamma$, let $R_\epsilon(x, \beta)$ be the set of $t' \in \RV$ such that $\go(x, \beta) \sub \dom(\epsilon)$ and $\rv(\epsilon(z) - \epsilon(y)) = t' \rv(z - y)$ for all $y, z \in \go(x, \beta)$. If $\epsilon$ is $\rv$-affine on a disc then it is so on every subdisc, with the same slope, and two slopes on concentric discs agree. Hence $R_\epsilon(x, \beta)$ has at most one element, and once nonempty it does not change as $\beta$ increases. Moreover, if $\go \sub \dom(\epsilon)$ is an open disc of radius $\rho$ and $\beta > \rho$, then $\go$ contains an open subdisc of radius $\beta$ on which $\epsilon$ is $\rv$-affine: by saturation, choose such a subdisc that is atomic over the parameters together with $\code \go$ and $\beta$, and apply \cite[Lemma~3.18]{Yin:int:acvf}.

Prepare the relation $\set{(x, \beta, t') \given t' \in R_\epsilon(x, \beta)}$, with its $\Gamma$-coordinate encoded by $\vrv$-fibers, together with the cell data, by the method in the proof of Proposition~\ref{simplex:with:hole:rvproduct}: Lemma~\ref{special:bi:term:constant}, followed by FMT on finitely many open discs of an invariance radius $\delta$ of $A$. On a frozen disc $\ga$ with chosen point $a_\ga$, $\delta$-invariance gives $A \cap (\ga \times \VF) = \ga \times A_{a_\ga}$, a $2$-cell with constant positioning function. Let $D$ be any other prepared disc, of radius $\rho_D$. For each $\beta > \rho_D$, $D$ contains a subdisc $\go(c, \beta)$ on which $\epsilon$ is $\rv$-affine, so $R_\epsilon(c, \beta) \neq \0$. By constancy, $R_\epsilon(x, \beta) = R_\epsilon(c, \beta)$ for all $x \in D$, and the slope does not depend on $\beta$. Any two points $y, z \in D$ lie in $\go(y, \beta)$ for some $\beta$ with $\rho_D < \beta < \vv(z - y)$. So $\epsilon \rest D$ is $\rv$-affine.

The fibers $A_a = t_2^\sharp + \epsilon(a)$ and the paradigms are unchanged, and cells of the first two kinds in Definition~\ref{def:units} remain of the same kind on each $D$. For a balanced cell, if $\rho_D < \vrv(t_1)$ then $a - t_1^\sharp \sub D$ for every $a \in D$, so $\epsilon \rest D$ is still balanced in $A \cap (D \times \VF)$. Otherwise $\rad(\epsilon(D)) = \rho_D + \vrv(t_2) - \vrv(t_1) \geq \vrv(t_2)$, which is the second kind. The refined unary stage is a special covariant bijection, so the result is again a standard contraction of $A$.
\end{proof}

We also remark that Lemma~\ref{decom:into:2:units} holds fiberwise for proper invariant sets $A \sub \VF^n$ with $n \geq 2$, that is, there is a standard contraction $\hat T_\sigma$ of $A$ such that, writing $P = \hat T_{\sigma(2)} \circ \hat T_{\sigma(1)}$ for the cumulative contraction of the first two coordinates, for every $(a, s) \in P(A)$, $P^{-1}(a, s)$ is of the form $a \times C$, where $C$ is a $2$-cell. This follows from Lemma~\ref{decom:into:2:units} and the argument in the proof of  Lemma~\ref{exit:stand} (Lemma~\ref{decom:into:2:units} serves as the base case and hence, in the inductive step, we can assume that $P$ is already as desired fiberwise).


For the next two lemmas, let $12$, $21$ denote the permutations of $[2]$ and $\bm A = (A, \omega) \in \mVF^\diamond[2]$.


\begin{lem}\label{2:unit:contracted}
Suppose that $A \sub \VF^2$ is a $2$-cell and $\omega$ is constant. Then there are standard contractions $\hat T_{12}$, $\hat R_{21}$ of $\bm A$ such that $[\hat T_{12}(\bm A)]_{\leq 2} = [\hat R_{21}(\bm A)]_{\leq 2}$.
\end{lem}
\begin{proof}
We check that the maps in the proof of \cite[Lemma~5.7]{Yin:int:acvf} give morphisms in $\mRV^{\db}[2]$. By Definition~\ref{def:units}, either $A$ is a product $A_1 \times (t^\sharp + c)$, or the positioning function $\epsilon$ is balanced in $A$ with paradigms $t_1, t_2 \in \RV$. The first alternative covers a constant positioning function, and also the second kind of $2$-cell, since $\rad(\epsilon(A_1)) \geq \vrv(t)$ makes $t^\sharp + \epsilon(a)$ independent of $a$. In that case contract each factor by the same unary special covariant bijections in either order; the two results differ only by the order of the coordinates, which preserves total weights, so the identity map is the required isomorphism.

In the balanced case put $\gamma = \vrv(t_2) - \vrv(t_1)$, the shift of $\epsilon$. Since $\epsilon^{-1}(A_a) = a - t_1^\sharp$, the fiber of $A$ over $y \in A_2$ in the second coordinate is $\epsilon^{-1}(y) + t_1^\sharp$. For the order $12$ use $(x, y) \efun (x - \epsilon^{-1}(y), y)$, whose image is $t_1^\sharp \times A_2$; for the order $21$ use $(x, y) \efun (x, y - \epsilon(x))$, whose image is $A_1 \times t_2^\sharp$. These are centripetal transformations whose foci lie in the $\RV$-disc of the changed coordinate. The translated coordinate has the fixed leading term $t_1$, respectively $t_2$, so for an aperture whose cutoff exceeds $\vrv(t_1)$, respectively $\vrv(t_2)$, the core clause never occurs. Off a finite set, $\epsilon$ is differentiable with $\vv(\ddx \epsilon) = \gamma$, so $\epsilon : (A_1, 0) \fun (A_2, -\gamma)$ is a $\mVF^\diamond[1]$-morphism. Lemma~\ref{simul:special:dim:1} gives special covariant bijections $T_{A_1}$, $T_{A_2}$ on $A_1$, $A_2$ and a $\mRV^{\db}[1]$-morphism $\psi : \rv(T_{A_1}(A_1)) \fun \rv(T_{A_2}(A_2))$; thus $\jcb_\Gamma \psi = \gamma$. Complete the first translation by $T_{A_2}$ in the second coordinate and the second translation by $T_{A_1}$ in the first coordinate. This gives standard contractions with
\[
\hat T_{12}(A) = t_1 \times \rv(T_{A_2}(A_2)), \qquad \hat R_{21}(A) = \rv(T_{A_1}(A_1)) \times t_2,
\]
both with the constant form. The map $(u, t_2) \efun (t_1, \psi(u))$, keeping all auxiliary tags, has $\Gamma$-Jacobian $\gamma + \vrv(t_1) - \vrv(t_2) = 0$, so it is the required isomorphism. Only the shift $\gamma$ is used; it is the valuation of the slope $-t_2/t_1$ of $\epsilon$ (see the last paragraph of Remark~\ref{2cell:linear}).
\end{proof}

\begin{lem}\label{subset:partitioned:2:unit:contracted}
There are a morphism $\bm A \fun \bm A^*$, relatively unary in both coordinates, and two standard contractions $\hat T_{12}$, $\hat R_{21}$ of $\bm A^*$ such that $[\hat T_{12}(\bm A^*)]_{\leq 2} = [\hat R_{21}(\bm A^*)]_{\leq 2}$.
\end{lem}
\begin{proof}
First prepare the form. By Lemma~\ref{L:measure:surjective} and Remark~\ref{cons:extra}, there is a special covariant bijection $T$ on $A$ such that $\omega \circ T^{-1}$ is $\rv$-contractible; so $\omega = \pi \circ \theta$, where $\theta(a)$ records all the $\RV$-coordinates of $T(a)$. Apply Lemma~\ref{decom:into:2:units} to the graph of $\theta$, with these $\RV$-coordinates retained as labels; its proof is uniform in such labels. We obtain a standard contraction $\hat T_\sigma$ whose fibers are $2$-cells on which $\theta$, and hence $\omega$, is constant. Let $c(a)$ be the $\RV$-tuple consisting of $\theta(a)$ and $\hat T_\sigma(a, \theta(a))$, and put
\[
A^* = \set{(a, c(a)) \given a \in A}, \qquad \omega^*(a, c(a)) = \omega(a).
\]
The map $c$ is the $\RV$-projection of proper covariant homeomorphisms, so it is locally constant, its fibers are proper invariant and its range is doubly bounded. As for the grafting map in the proof of Lemma~\ref{exit:stand}, the map $a \efun (a, c(a))$ and its inverse projection are proper covariant homeomorphisms whose $\VF$-maps are identities on each fiber of $c$. Hence $\bm A \fun \bm A^* = (A^*, \omega^*)$ is a morphism, relatively unary in both coordinates.

For each $s \in c(A)$, the fiber $A^*_s$ is a $2$-cell on which $\omega^*$ is constant. The constructions in the proof of Lemma~\ref{2:unit:contracted}, namely the translations by $\epsilon_s^{-1}$ and $\epsilon_s$ and Lemma~\ref{simul:special:dim:1} applied to the positioning functions with $s$ as a retained label, are uniform in $s$, and the kinds of cells form a definable finite partition of $c(A)$. By compactness, as in the proof of \cite[Corollary~5.8]{Yin:int:acvf}, they glue to standard contractions $\hat T_{12}$, $\hat R_{21}$ of $\bm A^*$ and an isomorphism between $\hat T_{12}(\bm A^*)_{\leq 2}$ and $\hat R_{21}(\bm A^*)_{\leq 2}$.
\end{proof}

\begin{lem}\label{contraction:perm:pair:isp}
Let $\bm A = (A, \omega) \in \mVF^\diamond[n]$. Let $i, j \in [n]$ with $i \neq j$ and $\sigma_1$, $\sigma_2$ be  permutations of $[n]$ such that
\[
\sigma_1(1) = \sigma_2(2) = i, \quad \sigma_1(2) = \sigma_2(1) = j, \quad \sigma_1
\rest \set{3, \ldots, n} = \sigma_2 \rest \set{3, \ldots, n}.
\]
Then, for any standard contractions $\hat T_{\sigma_1}$, $\hat T_{\sigma_2}$ of $\bm A$,
\[
([\hat T_{\sigma_1}(\bm A)]_{\leq n}, [\hat T_{\sigma_2}(\bm A)]_{\leq n}) \in \mispdb.
\]
\end{lem}
\begin{proof}
Let $ij$, $ji$ denote the permutations of $\{i, j\}$ and $E = [n] \mi \{i, j\}$. By compactness and
Lemma~\ref{isp:VF:fiberwise:contract}, it is enough to show
that for  all $a \in A_E$ and all standard
contractions $\hat T_{ij}$, $\hat T_{ji}$ of $\bm A_a = (A_a, \omega \rest A_a)$,
\[
([\hat T_{ij}(\bm A_a)]_{\leq 2}, [\hat T_{ji}(\bm A_a)]_{\leq 2})
\in \mispdb.
\]
By Corollary~\ref{contraction:same:perm:isp} and Lemma~\ref{isp:VF:fiberwise:contract}, it is enough to find a morphism $\bm A_a \fun \bm B$, relatively unary in both coordinates, and  standard contractions $\hat R_{ij}$, $\hat R_{ji}$ of $\bm B$ such that $[\hat R_{ij}(\bm B)]_{\leq 2} = [\hat R_{ji}(\bm B)]_{\leq 2}$. This is just Lemma~\ref{subset:partitioned:2:unit:contracted}.
\end{proof}

\subsection{The kernel of $\mL$ and the bounded integral}

The following proposition is the culmination of the preceding technicalities, it identifies the
congruence relation  $\mispdb$ with that induced by
$\mL$.


\begin{prop}\label{kernel:L}
For $\bm U, \bm V \in \mRV^{\db}[k]$,
\[
[\mu\bb L \bm U] = [\mu\bb L \bm V] \text{ in }\gsk \mVF^\diamond[k] \quad \text{if and only if} \quad ([\bm U], [\bm V]) \in \mispdb.
\]
\end{prop}
\begin{proof}
For $k = 0$ there are no blowups and $\mL$ is the identity on finite sets, so both conditions say that $\bm U \cong \bm V$.

The ``if'' direction simply follows from Lemma~\ref{blowup:same:RV:coa} and Corollary~\ref{L:sur:c}.

For the ``only if'' direction, we first show, by induction on $k \geq 1$, that the class of a standard contraction does not depend on the contraction or its order: if $\bm B \in \mVF^\diamond[k]$ and $\hat T_\sigma$, $\hat T_\tau$ are standard contractions of $\bm B$ for permutations $\sigma$, $\tau$ of $[k]$, then
\begin{equation*}\tag{$\ast$}
([\hat T_\sigma(\bm B)]_{\leq k}, [\hat T_\tau(\bm B)]_{\leq k}) \in \mispdb.
\end{equation*}
If $\sigma = \tau$ this is Corollary~\ref{contraction:same:perm:isp} for the identity of $\bm B$; this settles $k = 1$. Since $\mispdb$ is transitive and every permutation is a product of adjacent transpositions, it remains to treat the case where $\tau$ is $\sigma$ followed by the transposition of the positions $m$ and $m+1$. For $m = 1$ this is Lemma~\ref{contraction:perm:pair:isp}. Let $m \geq 2$ and $h = \sigma(1) = \tau(1)$. By the case $\sigma = \tau$, and Lemma~\ref{exit:stand} applied after the first stage, we may assume that $\hat T_\sigma$ and $\hat T_\tau$ have the same first stage $T_h$. Its lift $T_h(B)$ is an $\RV$-pullback in the coordinate $h$. For a $\VF$-value $a$ of that coordinate, the fiber $T_h(\bm B)_a$ is an object of $\mVF^\diamond[k-1]$ over $\mathbb S \la a \ra$, which is $\VF$-generated, and it depends only on $\rv(a)$. The remaining stages of $\hat T_\sigma$, $\hat T_\tau$ restrict to standard contractions of this fiber whose orders differ by an adjacent transposition. By the inductive hypothesis over $\mathbb S \la a \ra$, their results $U_a$, $V_a$ satisfy $([U_a]_{\leq k-1}, [V_a]_{\leq k-1}) \in \mispdb$. Now apply Lemma~\ref{isp:VF:fiberwise:contract}, with one $\VF$-coordinate and head start $k-1$, to $\bigcup_a a \times U_a$ and $\bigcup_a a \times V_a$. These are $\RV$-pullbacks in their $\VF$-coordinate, so their tautological contractions are standard contractions, and these are $\hat T_\sigma(\bm B)$ and $\hat T_\tau(\bm B)$. This proves $(\ast)$. Note that the base is never enlarged by an $\RV$-parameter.

Now suppose $[\mL \bm U] = [\mL \bm V]$ and let
\[
\mu\bb L \bm U = \bm B_1 \to^{G_1} \bm B_2 \cdots \bm B_l \to^{G_l} \bm B_{l+1} = \mu\bb L \bm V
\]
be a sequence of relatively unary $\mVF^\diamond[k]$-morphisms, which exists by definition. Since $\mL \bm U$ is an $\RV$-pullback, its tautological contraction, which recovers $\bm U$, is a standard contraction with respect to any order, and similarly for $\mL \bm V$. If $l = 0$, then $\bm U$ and $\bm V$ are standard contractions of the same object and $(\ast)$ applies. Otherwise, for each $m \leq l$, let $c_m$ be the coordinate targeted by $G_m$, and choose standard contractions $\bm U_m$ of $\bm B_m$ and $\bm U'_{m+1}$ of $\bm B_{m+1}$ with a common order starting with $c_m$. By Corollary~\ref{contraction:same:perm:isp}, $([\bm U_m]_{\leq k}, [\bm U'_{m+1}]_{\leq k}) \in \mispdb$. By $(\ast)$, $\bm U'_{m+1}$ is congruent to $\bm U_{m+1}$ for $m < l$, $\bm U$ is congruent to $\bm U_1$, and $\bm V$ is congruent to $\bm U'_{l+1}$. Transitivity of $\mispdb$ (Lemma~\ref{isp:congruence:vol}) gives $([\bm U], [\bm V]) \in \mispdb$. This treats every arrow separately, including the last one, chains of every length, and consecutive arrows in the same coordinate.
\end{proof}

\begin{rem}\label{thekernel}
Proposition~\ref{kernel:L} shows that the congruence on $\gsk \mRV^{\db}[k]$ induced by $\mL$ is $\mispdb[k]$, which is generated by the pairs $([\bm U], [\bm U^\flat])$ with $\bm U^\flat$ an elementary blowup of $\bm U$. Consequently, the kernel of the graded ring homomorphism $\ggk \mRV^{\db}[*] \fun \ggk \mVF^\diamond[*]$ induced by $\mL$ is the relative annular ideal $J_{\mathrm{var}}$ of Notation~\ref{relative:annular:ideal}. Indeed, by Definition~\ref{defn:blowup:coa} the generators $[\bm U^{\tau, j}] - [\bm U]$ of $J_{\mathrm{var}}$ are exactly the differences $[\bm U^\flat] - [\bm U]$ for elementary blowups, and these lie in the kernel by Lemma~\ref{blowup:same:RV:coa}. Conversely, a blowup $\bm V^\flat_F$ with locus $\bm C = \bigcup_i \bm C_i$ satisfies $[\bm V^\flat_F] - [\bm V] = \sum_i ([\bm U_i^\flat] - [\bm U_i])$, so a blowup of any length differs from its source by an element of $J_{\mathrm{var}}$; also, the product of a generator with a class $[\bm W]$ is the generator for the pulled-back aperture on the product object. Now let $[\bm U] - [\bm V]$ lie in the kernel, in degree $k$. Then $[\mL \bm U] + b = [\mL \bm V] + b$ in $\gsk \mVF^\diamond[k]$ for some $b$, and $b = [\mL \bm W]$ for some $\bm W$ by Corollary~\ref{L:sur:c}. By Proposition~\ref{kernel:L}, $\bm U \uplus \bm W$ and $\bm V \uplus \bm W$ have isomorphic blowups $(\bm U \uplus \bm W)^\flat \cong (\bm V \uplus \bm W)^\flat$, and therefore
\[
[\bm U] - [\bm V] = ([\bm U \uplus \bm W] - [(\bm U \uplus \bm W)^\flat]) - ([\bm V \uplus \bm W] - [(\bm V \uplus \bm W)^\flat]) \in J_{\mathrm{var}}.
\]
No cancellation in $\gsk \mRV^{\db}[k]$ is used.

For a constant aperture $t \in \gamma^\sharp$, where $\gamma \in D_{\RV}$, put $\bm W = (U, f_{\tilde j}, \omega + \vrv(f_j)) \in \mRV^{\db}[k-1]$. The measure condition gives $\bm U \cong \bm W \times (1, 0)$ and $\bm U^\flat \cong \bm W \times (\RV^{\ccirc}(t)_\gamma, 0)$, hence $[\bm U^\flat] - [\bm U] = [\bm W] \bm P_\gamma$. Thus $(\bm P_\Gamma) \sub J_{\mathrm{var}}$ (Notation~\ref{pgammabis}); equality is not asserted. For a variable aperture, Lemma~\ref{relative:annular:compatibility} gives instead the relative image $([\bm W_\gamma] - [\bm W]) \bm P$ in the full measured ring.
\end{rem}

\begin{thm}\label{main:prop}
For each $k \geq 0$ there is a canonical isomorphism of Grothendieck semigroups
\[
  \int^{\diamond}_+ : \gsk \mVF^\diamond[k] \fun \gsk  \mRV^{\db}[k] /  \mispdb
\]
such that $\int_{+}^{\diamond} [\bm A] = [\bm U]/  \mispdb$ if and only if $[\bm A] = [{\mL}\bm U]$. Taking the direct sum of the groupifications
yields a canonical isomorphism of graded rings
\[
  \int^{\diamond} : \ggk \mVF^\diamond[*] \fun \ggk \mRV^{\db}[*] / J_{\mathrm{var}}.
\]
\end{thm}
\begin{proof}
The semigroup isomorphism is the inverse of the map induced by $\mL$, which is surjective by Corollary~\ref{L:sur:c} and injective by Proposition~\ref{kernel:L}. The groupification of a quotient of a semigroup by a congruence is the quotient of its groupification by the subgroup generated by the differences of congruent pairs; by Lemma~\ref{isp:congruence:vol} and Remark~\ref{thekernel}, in the direct sum over $k$ this subgroup is the homogeneous ideal $J_{\mathrm{var}}$.
\end{proof}


\begin{cor}\label{db:to:novol}
The diagram (\ref{diag-diamond-interpol}) commutes.
\end{cor}
\begin{proof}
Note that, in (\ref{diag-diamond-interpol}),  the first horizontal arrow  is induced by the obvious forgetful functor (it exists
because the morphisms in $\mgVF^\diamond[*]$ are bijections instead of essential bijections, see Remark~\ref{dia:cat:well}), the second horizontal arrow is induced by the subcategory relation, and the bottom row exists by Lemma~\ref{relative:annular:compatibility}. That both squares commute is then simply a consequence of the fact that the image of $[\bm A]$ under any of the three integrals can be represented by a standard contraction of $\bm A$.
\end{proof}


\begin{rem}\label{eg:eb:factor}
We may define the motivic volume operator in Notation~\ref{mov:vol} using $\bb E_g$ instead of $\bb E_b$ (of course $[\A]$ will have to be inverted, see Remark~\ref{eb:no:eg}).
For proper invariant objects, this only changes the image by a factor. To see this, first note that  Theorem~\ref{main:prop} (or rather Corollary~\ref{L:sur:c}) shows that every proper invariant object $A \in \VF_*$ with $\dim_{\VF}(A) = n$ is isomorphic to an object of the form $\bb L \bm U$ with $(\bm U, 0) \in  \mgRV^{\db}[n]$. Therefore, in light of the surjectivity of ${\mu}\Psi$ in (\ref{bdd:to:2bdd}) (Proposition~\ref{red:D:iso}) and the defining conditions of $\bb E_g$, $\bb E_b$ in \S~\ref{subs:uni:retr:novol}, we have
\begin{equation}\label{off:fac:A}
\bb E_g \Big (\int [A] \Big ) = \bb E_b \Big (\int [A] \Big )[\A]^{-n}.
\end{equation}
\end{rem}

\section{Motivic Milnor fiber}\label{sec:spec:pui}

We begin with a very brief discussion on specialization to henselian subfields. Let $\mathbb{M}$ be a $\VF$-generated  substructure. Recall that $\mathbb{S}$ is a part of the language and hence all other substructures contain it. If $X \sub \VF^n  \times \RV^m$ is a definable (and hence quantifier-free definable) set then the \memph{trace} of $X$ in $\mathbb{M}$, denoted by $X(\mathbb{M})$, is the set of $\mathbb{M}$-rational points of $X$, that is,
\[
X(\mathbb{M}) = X \cap (\VF(\mathbb{M})^n  \times \RV(\mathbb{M})^m).
\]
Such a trace is also called a \memph{constructible set} in $\mathbb{M}$ since it is indeed quantifier-free definable in $\mathbb{M}$. Note that, however, if $f : X \fun \Gamma$ is a definable function then the image $f(X(\mathbb{M}))$  is not necessarily a set in $\Gamma(\mathbb{M})$, but rather a set in the divisible hull $\Q \otimes \Gamma(\mathbb{M})$  of $\Gamma(\mathbb{M})$. For instance, if $\mathbb{M} = \C \dpar t$ then $\Gamma(\mathbb{M}) = \Z$ and hence $\gamma \in \Gamma$ is definable if and only if $\gamma \in \Q \otimes \Gamma(\mathbb{M}) = \Q$. On the other hand, if $X$ is a set in $\Gamma$ and $f$ is a piecewise $\mgl_k(\Z)$-transformation on $X$ with constant $\Gamma(\mathbb{M})$-terms  then $f(X(\mathbb{M}))$ is of course a set in $\Gamma(\mathbb{M})$; this is the situation in the $\Gamma$-categories (see Remark~\ref{GLZ:char}).


%\begin{rem}\label{lift:hen:bij}
%Suppose that $X$, $Y$ are definable sets and $g : X(\mathbb{M}) \fun Y(\mathbb{M})$ is a constructible bijection in $\mathbb{M}$. Obviously there may or may not be a definable bijection between $X$ and $Y$, let alone one whose trace equals $g$. But the quantifier-free formula that defines $g$ in $\mathbb{M}$ also gives rise to a definable bijection $f : X' \fun  Y'$ with $X' \sub X$ and $Y' \sub Y$ such that its trace $f'(\mathbb{M}): X'(\mathbb{M}) \fun  Y'(\mathbb{M})$ indeed equals $g$.
%\end{rem}


If $\mathbb{M}$ is definably closed and $\Gamma(\mathbb{M})$ is nontrivial, or equivalently (see Lemma~\ref{cut:to:hensel:substru}), the valued field $(\VF(\mathbb{M}), \OO(\mathbb{M}))$  is henselian, then $\mathbb{M}$ is \memph{functionally closed}, that is, for any definable sets $X$, $Y$ (in the sorts $\VF$ and $\RV$, see Terminology~\ref{nota:sub}) and any definable function $f : X \fun Y$, we have $f(X(\mathbb{M})) \sub Y(\mathbb{M})$. In particular, if $f$ is a bijection then, since $f^{-1}$ is definable too, $f$ restricts to a bijection between $X(\mathbb{M})$ and $Y(\mathbb{M})$. The image $f(X(\mathbb{M}))$ is definable in the structure $\mathbb{M}$, but in general it is neither the trace of $f(X)$ nor constructible in $\mathbb{M}$. For instance, if $\mathbb{M} = \C \dpar t$ and $f(x) = x^2$ on $\VF$ then $f(\VF(\mathbb{M}))$ is the set of squares in $\C \dpar t$, whose nonzero elements are the ones of even valuation, whereas a definable subset of $\VF$ either contains or omits all nonzero points that are sufficiently close to $0$. The closure property is all we need to deduce the results below.


%
%\begin{defn}[$\mathbb{M}$-specialized $\RV$-categories]
%An object of the category $\mgRV[k]_{\mathbb{M}}$ is a triple of the form $\bm U(\mdl M) = (U(\mdl M), f(\mdl M),  \omega \rest U(\mdl M))$, where $\bm U = (U, f, \omega) \in \mgRV[k]$. Any constructible function of the form $F(\mdl M) : \bm U(\mdl M) \fun \bm V(\mdl M)$, where $F : \bm U \fun \bm V$ is a $\mgRV[k]$-morphism, is a morphism of $\mgRV[k]_{\mdl M}$.
%\end{defn}
%
%%The categories $\mgRES[k]_{\mdl M}$, $\mG[k]_{\mdl M}$, etc., are formulated analogously.
%
%We call $\gsk \mgRV[*]_{\mdl M}$,  etc.,  \memph{$\mdl M$-constructible Grothendieck semirings} associated with $\mdl M$.
%
%
%Suppose that $[\bm A] = [\bm C]$, $[\bm B] = [\bm D]$ in $\gsk \mgRV[*]$ and there is a $\gsk \mgRV[*]_{\mdl M}$-morphism $\bm A(\mdl M) \fun \bm B(\mdl M)$. By Remark~\ref{lift:hen:bij}, there is a $\mgRV[*]$-morphism $F : \bm A' \fun \bm B'$ with $\bm A' \sub \bm A$, $\bm B' \sub \bm B$, $\bm A(\mdl M) = \bm A'(\mdl M)$, and $\bm B(\mdl M) = \bm B'(\mdl M)$, and hence there is a $\mgRV[*]$-morphism $G : \bm C' \fun \bm D'$ with $\bm C' \sub \bm C$, $\bm D' \sub \bm D$, $\bm C(\mdl M) = \bm C'(\mdl M)$, and $\bm D(\mdl M) = \bm D'(\mdl M)$. Thus the following binary relation is well-defined and is indeed a semiring congruence relation:
%\[
%\mathfrak{R}_{\mdl M} = \{ ([\bm A], [\bm B]) \in (\gsk \mgRV[*])^2 : [\bm A(\mdl M)] = [\bm B(\mdl M)] \text{ in } \gsk \mgRV[*]_{\mdl M}\}.
%\]
%Obviously $\gsk \mgRV[*]_{\mdl M} \cong \gsk \mgRV[*] / \gR_{\mdl M}$.


%It is also routine to check that the diagram commutes:
%\begin{equation}\label{dia:tensor:hen}
%\bfig
% \Square(0,0)/->`->`->`->/<400>[{\gsk \mgRES[*] \otimes \gsk \mG[*]}`{\gsk \mgRES[*]_{\mdl M} \otimes \gsk \mG[*]_{\mdl M}}`{\gsk \mgRV[*]}`{\gsk \mgRV[*]_{\mdl M}}; -/{\gR_{\mdl M}} \otimes - / {\gG_{\mdl M}}`\mu \Psi`{\mu \Psi_{\mdl M}}`-/{\gR_{\mdl M}}]
%\efig
%\end{equation}
%where the tensor product on the righthand side is taken over $\gsk \mG^{\fin}[*]_{\mdl M}$ and the isomorphism .
%Similar constructions are available for the other categories and their doubly bounded versions.To simplify the notation, we shall write $\mgRV[*]_{\C \dpar{ t^{1/m} }}$ as $\mgRV[*]_m$.


\subsection{Piecewise retraction to $\RES$}\label{piece:res}

From here on, we work in the field $\puC$ of complex Puiseux series and take  $\mathbb{S} = \C \dpar{ t }$. The value group $\Gamma$ is identified with $\Q$. We also consider the family of henselian subfields $\C \dpar{ t^{1/m} }$,  $m \in \Z^+$, regarded as substructures. The value group $\Gamma(\C \dpar{ t^{1/m} })$ is identified with $m^{-1} \Z$.


\begin{nota}\label{loc:rings}
For any ring $R$, let $R[T^{\Q}]$ denote the ring of Puiseux polynomials over $R$, that is, the group ring of $\Q$ over $R$.

Let $\gmv[[\A]^{-1}][T,T^{-1}]$ be the obvious subring of $\gmv[[\A]^{-1}][T^{\Q}]$. The canonical image of $\gmv [T^{\Q}]$ in  $\gmv[[\A]^{-1}][T^{\Q}]$ is just denoted as such. The assignment $T \efun [\A]$ determines a ring homomorphism
\begin{equation}\label{pui:laurent}
  \bm \eta : \gmv[[\A]^{-1}][T,T^{-1}] \fun \gmv[[\A]^{-1}].
\end{equation}
\end{nota}

%(recall Terminology~\ref{twist:rv:ob})

Recall Notation~\ref{comp:form}, Lemma~\ref{coh:decom} and the isomorphism $\Theta$ of Notation~\ref{bipo:gdv}. Let $\bm U = (U, f, \omega) \in \mgRV^{\db}[k]$, where $U \sub \RV^n$. In particular, if $f=(f_1,\dots,f_k)$, we have introduced the shorthand
\[
\omega_{f}: x\in U \efun \omega(x)+\sum_{i=1}^k\vrv(f_i(x)).
\]
For $\gamma \in \vrv(U)$ write $U_\gamma = U \cap \vrv^{-1}(\gamma)$. The image $\omega_f(U_\gamma)$ is finite for every $\gamma \in \vrv(U)$, because the $\K$-sort and the $\Gamma$-sort are orthogonal to each other. Its size is bounded uniformly in $\gamma$, so, listing the values of $\omega_f$ on $U_\gamma$ in increasing order, we get a finite definable partition of $U$ on each piece of which $\omega_f$ $\vrv$-contracts to a function. Applying Lemma~\ref{coh:decom} to each piece, it follows that there is a monomial decomposition $(U_i)_i$ of $U$, with
\[
h_{M_i}(U_i) = V_i \times I_i^\sharp, \qquad \vrv(V_i) = \{\beta_i\}, \qquad I_i \sub \Gamma^{\ell_i},
\]
such that each restriction $\omega_f \rest U_i$ $\vrv$-contracts to a function $\sigma_i : D_i = \vrv(U_i) \fun \Gamma$. Write $D_{i,m}$ for the set $D_i \cap (m^{-1}\Z)^n$ of  $m^{-1}\Z$-rational points of $D_i$; it is finite because $D_i$ is doubly bounded. We assign to $\bm U$ the expression
\begin{equation}\label{hm:ass}
  h_m(\bm U) = \sum_i \Theta([V_i]) [\G_m]^{\ell_i} \sum_{\gamma \in D_{i,m}} T^{-m \sigma_i(\gamma)},
\end{equation}
which is  a finite sum and hence  belongs to $\gmv[[\A]^{-1}][T^{\Q}]$. Note that $M_i \in \GL_n(\Z)$ preserves the lattice $(m^{-1}\Z)^n$, and hence $D_{i,m}$ is empty unless $\beta_i$ is a tuple in $m^{-1}\Z$.

If $\bm U$ is already of the form $(V, g, \omega') \times (I, \sigma) \in \mgRES[i] \times \mG^{\db}[j]$, where $\vrv(V)$ is a singleton $\{\delta\}$ and $\omega'_g$ is constant, then clearly
\[
h_m(\bm U) = \Theta([V])T^{-m\omega'_g} \cdot [\G_m]^{j} \sum_{\gamma \in I_{m}} T^{-m (\sigma(\gamma) + \Sigma \gamma )}
\]
if $\delta$ is a tuple in $m^{-1}\Z$, where $I_m = I \cap (m^{-1}\Z)^j$, and $h_m(\bm U) = 0$ otherwise. Every object of $\mgRES[i]$ is a finite disjoint union of such objects $(V, g, \omega')$. On the image of the homomorphism ${\mu} \Psi^{\db}$ in Remark~\ref{tensor:bd}, this is essentially how the assignment is given in \cite[\S~8.2]{hru:loe:lef}. We do not take it as the definition of $h_m$, since that would require its compatibility with the relations among the tensor classes. Computationally, too,  (\ref{hm:ass}) is more effective.


The ``piecewise retraction'' formula (\ref{hm:ass}) is indeed quite similar to the ``uniform retraction'' formula (\ref{twist:C}). The difference is just that the coefficient $\chi_b(I_i)$ in the latter is replaced in the former by a formal ``counting'' sum $\sum_{\gamma \in D_{i,m}} T^{-m \sigma_i(\gamma)}$ that also takes the volume form into account. As noted after (\ref{twist:C}), that the ``uniform retraction'' formula does not depend on the choice of the monomial decomposition is simply a consequence of the construction of $\bb E_b$ itself. The same holds for the ``piecewise retraction'' formula, but we do need an argument.

For $w \in \Gamma$, write
\[
U^{(m)}_w = \set{u \in U \given \vrv(u) \in (m^{-1}\Z)^n \text{ and } \omega_f(u) = w}.
\]
Since $\vrv(U)$ is doubly bounded, the set $\vrv(U) \cap (m^{-1}\Z)^n$ is finite, and each of its points is definable. Hence $U^{(m)}_w$ is a definable set with finite value support, that is, an object of $\RES$, and it is nonempty for only finitely many $w$.

\begin{lem}\label{muhat:section}
For every $\bm U = (U, f, \omega) \in \mgRV^{\db}[k]$,
\[
h_m(\bm U) = \sum_{w} \Theta([U^{(m)}_w])\, T^{-mw}.
\]
In particular, the  assignment (\ref{hm:ass}) does not depend on the choice of the monomial decomposition and is invariant on isomorphism classes.
\end{lem}
\begin{proof}
Let $(U_i)_i$ be a monomial decomposition as above. For $\gamma \in D_{i,m}$, write $M_i \gamma = (\beta_i, \delta)$. Then $h_{M_i}$ restricts to a definable bijection between $U_i \cap \vrv^{-1}(\gamma)$ and $V_i \times \delta^\sharp$, and $\Theta([\delta^\sharp]) = [\G_m]^{\ell_i}$ by the defining relations of $\sggk \RES$. Since $\omega_f$ is constant on $U_i \cap \vrv^{-1}(\gamma)$, with value $\sigma_i(\gamma)$, the set $U^{(m)}_w$ is the disjoint union of the sets $U_i \cap \vrv^{-1}(\gamma)$ with $\gamma \in D_{i,m}$ and $\sigma_i(\gamma) = w$, each of which has class $\Theta([V_i]) [\G_m]^{\ell_i}$ under $\Theta$. Summing over $w$ gives the displayed equality. Its right-hand side does not involve the decomposition.

Let $\bm V = (V, g, \pi) \in \mgRV^{\db}[k]$, where $V \sub \RV^{n'}$, and let $F : \bm U \fun \bm V$ be a $\mgRV^{\db}[k]$-morphism. We need to show that $h_m(\bm U)$ and $h_m(\bm V)$ are  equal.
By (\ref{cond:mg2}) and Notation~\ref{jcbG:op}, $\omega_f = \pi_g \circ F$. So it is enough to show that, for every $u \in U$, $\vrv(u)$ is a tuple in $m^{-1}\Z$ if and only if $\vrv(F(u))$ is. Both conditions define subsets of $U$ over $\mathbb{S}$, so this may be checked for the points of $U$ in $\puC$. The residue fields of $\puC$ and of $\C \dpar{ t^{1/m} }$ are both $\C$, and hence an element of $\RV(\puC)$ lies in $\RV(\C \dpar{ t^{1/m} })$ if and only if its value is in $m^{-1}\Z$. The claim follows because $\C \dpar{ t^{1/m} }$ is functionally closed and $F$, $F^{-1}$ are definable. Therefore $F$ restricts to a definable bijection between $U^{(m)}_w$ and $V^{(m)}_w$ for every $w$, and these two objects of $\RES$ have the same image under $\Theta$.
\end{proof}


Thus, $h_m$ may be viewed as a map on $\gsk \mgRV^{\db}[k]$. It is additive, since $(U \uplus U')^{(m)}_w = U^{(m)}_w \uplus U'^{(m)}_w$, and multiplicative, since the weight of a product is the sum of the weights, so that $(U \times U')^{(m)}_w$ is the disjoint union of the sets $U^{(m)}_{w_1} \times U'^{(m)}_{w_2}$ with $w_1 + w_2 = w$. So we have in effect constructed a ring homomorphism
\[
\bm h_m : \ggk \mgRV^{\db}[*] \fun  \ggk^{\hat \mu} \var_{\C}[[\A]^{-1}][T^{\Q}].
\]

Let $\ggk^\natural_m \mgRV^{\db}[*]$ denote the subring $(\bm h_m)^{-1}(\gmv[[\A]^{-1}][T, T^{-1}])$ of $\ggk \mgRV^{\db}[*]$.

\begin{lem}\label{PG:van}
The homomorphism
\[
\bm \eta \circ \bm h_m : \ggk^\natural_m \mgRV^{\db}[*] \fun \gmv[[\A]^{-1}]
\]
vanishes on $J_{\mathrm{var}} \cap \ggk^\natural_m \mgRV^{\db}[*]$, where $J_{\mathrm{var}}$ is the relative annular ideal of Notation~\ref{relative:annular:ideal}.
\end{lem}
Recall that $(\bm P_\Gamma) \sub J_{\mathrm{var}}$ (Remark~\ref{thekernel}). Note that the ideal $(\bm P_\Gamma)$ of $\ggk \mgRV^{\db}[*]$ in Notation~\ref{pgammabis} is now generated by the elements $\bm P_\gamma$ with $\gamma \in \Z_{\geq 0}$ (because the point $t_\gamma$ there needs to be definable, which is possible only if $\gamma \in \Z$ in the current setting).


\begin{proof}
Let $[\bm U^{\tau, j}] - [\bm U]$ be a generator of $J_{\mathrm{var}}$, where $\bm U = (U, f, \omega) \in \mgRV^{\db}[k]$ and $\tau : U \fun \RV$ is a definable function with $0 \leq \gamma(u) = \vrv(\tau(u)) \leq \Delta$. We first show that $\bm h_m([\bm U^{\tau, j}] - [\bm U])$ is a multiple of $[\A] T^{-1} - 1$ in $\gmv[[\A]^{-1}][T^{\Q}]$.

Let $U^{(m)}$ be the set of those $u \in U$ such that $\vrv(u)$ is a tuple in $m^{-1}\Z$. As in the proof of Lemma~\ref{muhat:section}, since $\tau$ is definable and $\C \dpar{ t^{1/m} }$ is functionally closed, $\gamma(u) \in m^{-1}\Z$ for every $u \in U^{(m)}$. So $m \gamma(u)$ is an integer with $0 \leq m\gamma(u) \leq m \Delta$, and there is a finite definable partition of $U^{(m)}$ such that, on each piece $R$, the functions $\omega_f$ and $\gamma$ are constant, say with values $w$ and $M/m$. Put $a = \Theta([R])$.

The weight of $\bm U^{\tau, j}$ at a point $(u, s)$ of the shell is $\omega_f(u) + \vrv(s)$, and at a point $(u, \tau(u))$ of the endpoint graph it is $\omega_f(u) + \gamma(u)$. The points of $U^{\tau, j}$ over $R$ whose values lie in the lattice are therefore the following. First, the endpoint graph $\set{(u, \tau(u)) \given u \in R}$, which is in definable bijection with $R$ and has weight $w + M/m$. Secondly, for each $i = 1, \ldots, M$, the shell slice $R \times (i/m)^\sharp$, which has weight $w + i/m$ and whose image under $\Theta$ is $a([\A] - 1)$. So the contribution of $R$ to $\bm h_m([\bm U^{\tau, j}] - [\bm U])$ is
\[
a T^{-mw} \Big( ([\A] - 1) \sum_{i = 1}^{M} T^{-i} + T^{-M} - 1 \Big) = a T^{-mw} ([\A] T^{-1} - 1) \sum_{i = 0}^{M-1} T^{-i}.
\]
Summing over the pieces $R$ proves the claim.

Now let $x \in J_{\mathrm{var}} \cap \ggk^\natural_m \mgRV^{\db}[*]$. By the claim, $\bm h_m(x) = ([\A] T^{-1} - 1) Q$ for some $Q \in \gmv[[\A]^{-1}][T^{\Q}]$. Let $\operatorname{pr}_0$ be the additive map on $\gmv[[\A]^{-1}][T^{\Q}]$ that retains the terms with integer exponents. It satisfies $\operatorname{pr}_0(p Q') = p \operatorname{pr}_0(Q')$ for every Laurent polynomial $p$. Since the exponents of $\bm h_m(x)$ are integers,
\[
\bm h_m(x) = \operatorname{pr}_0(\bm h_m(x)) = ([\A] T^{-1} - 1) \operatorname{pr}_0(Q),
\]
which vanishes after passing to $\gmv[[\A]^{-1}]$ via $\bm \eta$.
\end{proof}

For instance, for $\gamma \in \Z_{\geq 0}$, the image of $[\RV^{\circ \circ} \mi \RV^{\circ \circ}_{\gamma}] + [\{t_\gamma\}]$ under $\bm h_m$ in $\gmv[[\A]^{-1}][T, T^{-1}]$ is
\[
 ([\A] - 1) \sum_{i = 1}^{m\gamma} T^{-i} + T^{-m\gamma},
\]
which, after passing to $\gmv[[\A]^{-1}]$ via $\bm \eta$, becomes $1 = \bm \eta (\bm h_m([1]))$.


\begin{rem}\label{dia:to:db:dag}
If $\bm U = (U, f, l) \in \mgRV^{\db}[*]$ with $l \in m^{-1}\Z$ (constant volume form) then the exponents in (\ref{hm:ass}) are all integers and hence $[\bm U] \in \ggk^\natural_m \mgRV^{\db}[*]$. Actually we shall only need the case $l = 0$.

The ring $\bigcap_{m \in \Z^+} \ggk^\natural_m \mgRV^{\db}[*]$ is denoted by $\ggk^\natural \mgRV^{\db}[*]$.

If $\bm A = (A, l) \in \mgVF^\diamond[*]$ with $l$  constant then $\int^\diamond [\bm A]$ may be expressed as $[(U, f, l)]/ J_{\mathrm{var}}$, and hence if $l \in \Z$ then $\int^\diamond [\bm A]$ belongs to the image of $\ggk^\natural \mgRV^{\db}[*]$ in $\ggk \mgRV^{\db}[*] / J_{\mathrm{var}}$, which we write as $\ggk^\natural \mgRV^{\db}[*]/ J_{\mathrm{var}}$. Two elements of $\ggk^\natural_m \mgRV^{\db}[*]$ that represent $\int^\diamond [\bm A]$ differ by an element of $J_{\mathrm{var}} \cap \ggk^\natural_m \mgRV^{\db}[*]$. So, by Lemma~\ref{PG:van}, for every $m \in \Z^+$, the expression $(\bm \eta \circ \bm h_m \circ \int^\diamond)([\bm A])$ designates a unique element  in $\gmv[[\A]^{-1}]$.
Let $\ggk^{\Z} \mgRV^{\db}[*]$ be the subring of $\ggk \mgRV^{\db}[*]$ generated by the classes $[(U, f, l)]$ with $l \in \Z$, and $\ggk^{\Z} \mgVF^\diamond[*]$ the subring of $\ggk \mgVF^\diamond[*]$ generated by the classes $[(A, l)]$ with $l \in \Z$. Since a product of two such classes is again of this form, their elements are the finite sums of such classes and their negatives. The first ring is contained in $\ggk^\natural \mgRV^{\db}[*]$, and, by the above, $\int^\diamond$ maps the second into the image of the first in $\ggk \mgRV^{\db}[*]/J_{\mathrm{var}}$, which we write as $\ggk^{\Z} \mgRV^{\db}[*]/J_{\mathrm{var}}$.
\end{rem}


Denote by $\RES_{m}$ the full subcategory of $\RES$ such that $U \in \RES_{m}$ if and only if every coordinate of every element of $U$ is either $\infty$ or has value in $m^{-1} \Z$, or equivalently, $U \in \RES_{m}$ if and only if the action on $U$ of the kernel of the canonical projection $\hat \mu \fun \mu_m$ is trivial.

For $\gamma \in \Q$, let $L_\gamma = \gamma^\sharp \cup \{0_\gamma\}$ be the one-dimensional residue vector space of Notation~\ref{bipo:gdv}. Its class in $\sggk \RES$ is $[\G_m] + 1 = [\A]$.

Let $\beta = (\beta_1, \ldots, \beta_n) \in (m^{-1} \Z)^n$ with $\beta_i \geq 0$ and $A \sub \OO^n \times \RV^l$ be a proper $\beta$-invariant definable set such that $\pvf \rest A$ is finite-to-one. Then there is a set
\[
A[m;\beta] \sub \prod_i \C[t^{1/m}] / t^{\beta_i + 1/m} \times \RV(\C \dpar{t^{1/m}})^l
\]
such that, for every $t \in \RV(\C \dpar{t^{1/m}})^l$, the $\VF$-fiber $A_t(\C \dpar{t^{1/m}})$ is, under the obvious quotient map, the preimage of the fiber $A[m;\beta]_t$. Note that $\pvf \rest A[m;\beta]$ is still finite-to-one, for otherwise, by the $\beta$-invariance of $A$, $\pvf \rest A$ would fail to be finite-to-one. We can think of $\C[t^{1/m}] / t^{\beta_i + 1/m}$ as the object $\prod_{\gamma} L_\gamma$ of $\RES_{m}$, where $\gamma$ runs through the elements of $m^{-1}\Z$ with $0 \leq \gamma \leq \beta_i$ and the coefficient of $t^\gamma$ is regarded as an element of $L_\gamma$. Then, since every element in the value group $\Q$ is definable and the values of the $\RV$-coordinates of $A[m;\beta]$ are doubly bounded,  it follows that  $A[m;\beta]$ is indeed a finite disjoint union of definable sets in $\RV_\infty$ (no extra parameters other than those in $\mathbb{S} = \C \dpar{ t }$ are needed) and hence, as such, is an object of $\RES_{m}$ (see \cite[\S~4.2]{hru:loe:lef} for more details). The same applies to graphs: if $B \sub \OO^n \times \RV^{l'}$ is a proper $\beta$-invariant definable set and $F : A \fun B$ is a definable bijection that induces a bijection between $A[m;\beta]$ and $B[m;\beta]$, then the induced bijection is definable, and hence $[A[m;\beta]] = [B[m;\beta]]$ in $\ggk \RES_m$.


\begin{lem}[{\cite[Lemma~4.2.1]{hru:loe:lef}}]
Let $\beta' \in (m^{-1} \Z)^n$ with $\beta_i \leq \beta'_i$ for all $i$. Then
\begin{equation}\label{arc:com}
  [A[m;\beta']] = [A[m;\beta]][\A]^{m \Sigma (\beta' - \beta)} \in \sggk \RES.
\end{equation}
\end{lem}
\begin{proof}
Since $A$ is $\beta$-invariant, a point of $A$ that is rational over $\C \dpar{t^{1/m}}$ remains in $A$ when the coefficients of $t^\gamma$ in its $i$th $\VF$-coordinate, $\gamma > \beta_i$, are changed. So
\[
A[m;\beta'] = A[m;\beta] \times \prod_i \prod_{\gamma} L_\gamma,
\]
where, for each $i$, $\gamma$ runs through the $m(\beta'_i - \beta_i)$ elements of $m^{-1}\Z$ with $\beta_i < \gamma \leq \beta'_i$. The equality follows since $[L_\gamma] = [\A]$ in $\sggk \RES$.
\end{proof}
The equality (\ref{arc:com}) does not hold in $\ggk \RES_{m}$ in general, because $L_\gamma$ is not isomorphic to $L_0$ unless $\gamma \in \Z$: under $\Theta$ it becomes an affine line on which $\mu_m$ acts through the character $\zeta \efun \zeta^{m\gamma}$.


Thus the element $[A[m;\beta]][\A]^{- m\sum \beta }$ in $\sggk \RES[[\A]^{-1}]$ does not depend on $\beta$; we denote it by $\tilde A[m]$.

The localization of $\Theta$ at $[\A]$ is still denoted by $\Theta$.

\begin{lem}\label{inv:dir:com}
Suppose that $(A, 0) \in \mgVF^\diamond[n]$ and $A \sub \OO^n \times \RV^l$. Then, in $\gmv[[\A]^{-1}]$,
\[
(\bm \eta \circ \bm h_m) \Big( \int^\diamond [(A, 0)] \Big) = \Theta(\tilde A[m]).
\]
\end{lem}
\begin{proof}
We first show that $\tilde A[m]$ is not changed by the operations that make up a special covariant bijection. So let $A \sub \OO^n \times \RV^l$ be doubly bounded, proper invariant and regularized, and let $\eta^\flat : A \fun A^\flat$ be a lifted centripetal transformation as in Definition~\ref{defn:special:bijection}, followed by regularization. Write a point of $A$ as $(a, x)$, where $a$ is the first $\VF$-coordinate, and suppress the $\RV$-labels. On its two branches $\eta^\flat$ is given by
\[
(a, x) \efun (a - \lambda(x), x), \qquad (a, x) \efun (a - \lambda(x) + b(x), x),
\]
according to whether $\vv(a - \lambda(x)) \leq \gamma_\gp$ or not, and the image carries the $\RV$-label of its first coordinate and the tag of the branch. Since $a$ and $\lambda(x)$ lie in the same $\RV$-disc and $\vv(b(x)) = \gamma_\gp \geq 0$, we have $A^\flat \sub \OO^n \times \RV^{l'}$.

Choose $\beta = (\beta_1, \beta') \in (m^{-1}\Z)^n$ as follows. Both $A$ and $A^\flat$ are proper invariant (Lemmas~\ref{spec:inve} and~\ref{conv:unif}), and hence $\epsilon$-invariant for every sufficiently large $\epsilon \in \Gamma$. First take $\beta_1$ so large that $\beta_1 \geq \gamma_\gp$ for every $\gp$ and that $A$ and $A^\flat$ are $\beta_1$-invariant. The functions $\lambda$ and $b$ are continuous and, over each label, their domains are closed and bounded. So, as in the proof of Lemma~\ref{conti:covar}, there is a $\delta \in \Gamma$ such that every $\VF$-fiber of $\lambda$ and of $b$ is $(\delta, \beta_1)$-covariant. Then take every coordinate of $\beta'$ at least $\delta$ and at least $\beta_1$, and so large that the labels of $A$ and of $A^\flat$ are constant on their open polydiscs of radius $\beta$. Then $A$ and $A^\flat$ are $\beta$-invariant.

Let $(a, x)$, $(a', x')$ be two points of $A$ with the same labels such that $a - a' \in \MM_{\beta_1}$ and $x - x' \in \MM_{\beta'}$. Then $\lambda(x) - \lambda(x') \in \MM_{\beta_1}$ and $b(x) - b(x') \in \MM_{\beta_1}$. Since $\beta_1 \geq \gamma_\gp$, the two points lie on the same branch, the first coordinates of their images differ by an element of $\MM_{\beta_1}$, and these have the same $\RV$-label, because their values are at most $\gamma_\gp$. The inverse of $\eta^\flat$ adds $\lambda(x)$, and subtracts $b(x)$ on the second branch, so it has the same property. Since $\eta^\flat$ and its inverse are definable and $\C \dpar{t^{1/m}}$ is functionally closed, $\eta^\flat$ induces a bijection between $A[m;\beta]$ and $A^\flat[m;\beta]$. Hence $\tilde A[m] = \widetilde{A^\flat}[m]$. Note that, for a given $\beta_1$, the other cutoffs $\beta'$ have to be large enough: in general $\eta^\flat$ does not induce a map between the truncations for an arbitrary large $\beta$. This is harmless, since $\tilde A[m]$ does not depend on $\beta$.

The same argument, with no focus, applies to regularization and to the preliminary translations in Remark~\ref{spec:addendum}.

By Lemma~\ref{L:measure:surjective}, there exists a special covariant bijection $F : (A, 0) \fun (\bb L \bm U, 0)$, where $\bm U = (U, f) \in \RV^{\db}[n]$, and then $\int^\diamond [(A, 0)] = [(U, f, 0)] / J_{\mathrm{var}}$. By the above, $\tilde A[m] = \widetilde{\bb L \bm U}[m]$. Thus, we may assume that $A$ is already of the form $\bb L \bm U$.

Let $U^{(m)}$ be the set of those $u \in U$ such that $\vrv(u)$ is a tuple in $m^{-1}\Z$. The points of $\bb L \bm U$ that are rational over $\C \dpar{t^{1/m}}$ lie over $U^{(m)}$. There is a finite definable partition of $U^{(m)}$ such that, on each piece $R$, the tuple $\vrv \circ f$ is constant, say with value $\alpha = (\alpha_1, \ldots, \alpha_n) \in (m^{-1}\Z)^n$, where $\alpha_i \geq 0$. Take $\beta_i \geq \alpha_i$ for all $i$ and all pieces. Over a point $u \in R$, the $i$th $\VF$-coordinate of a rational point of $\bb L \bm U$ has no terms below $t^{\alpha_i}$, its coefficient of $t^{\alpha_i}$ is $f_i(u)$, and its coefficients of $t^\gamma$, $\alpha_i < \gamma \leq \beta_i$, are unrestricted. So the part of $\bb L \bm U[m;\beta]$ over $R$ is in definable bijection with $R \times \prod_i \prod_\gamma L_\gamma$, where $\gamma$ runs through the $m(\beta_i - \alpha_i)$ elements of $m^{-1}\Z$ with $\alpha_i < \gamma \leq \beta_i$, and its class in $\sggk \RES$ is $[R] [\A]^{m \Sigma (\beta - \alpha)}$. Therefore
\[
\Theta(\widetilde{\bb L \bm U}[m]) = \sum_R \Theta([R]) [\A]^{-m \Sigma \alpha}.
\]
On the other hand, the weight of $(U, f, 0)$ on $R$ is $\Sigma \alpha$, and so, by Lemma~\ref{muhat:section}, the right-hand side is $(\bm \eta \circ \bm h_m)([(U, f, 0)])$.
\end{proof}


\subsection{Constructing motivic Milnor fiber}\label{sec:milnor}

Finally, we proceed to replicate the construction in \cite[\S~8]{hru:loe:lef} so as to recover the motivic zeta function with coefficients in $\gmv[[\A]^{-1}]$ (this is the purpose of Lemma~\ref{inv:dir:com}) and the corresponding motivic Milnor fiber.


Let $X$ be a smooth connected complex algebraic variety of dimension $d$ and $f : X \fun \A^1$ a nonconstant regular function. Let $0$ denote a point of $X(\C)$ with $f(0) = 0$; the case of interest is that of a critical point of $f$. For $m \geq 1$, we consider the set of truncated arcs
\begin{equation}\label{f:m:arc}
\mdl X_{m} = \set{ \varphi \in X(\C[t] / (t^{m+1})) \given \varphi(0) = 0 \text{ and } f(\varphi) = t^m \mod t^{m+1}},
\end{equation}
which may be viewed in a natural way as the set of closed points of an algebraic variety over $\C$ that carries a natural $\mu_{m}$-action, and the so-called \memph{nonarchimedean Milnor fiber}
\[
\mdl X_f = \set{ x \in X(\MM) \given \rv(f(x)) = \rv(t)}.
\]
Here $X(\MM)$ stands for the set of the points of $X(\OO)$ that specialize to $0$. To regard $\mdl X_f$ as a definable set, choose an \'etale map $q$ from a neighborhood of $0$ in $X$ to $\A^d$ with $q(0) = 0$. By Hensel's lemma, $q$ induces a bijection between $X(\MM)$ and $\MM^d$, and we identify $\mdl X_f$ with its image in $\MM^d$. In these coordinates $f$ is given on $\MM^d$ by a power series with coefficients in $\C$ and no constant term, so that $f(x') - f(x) \in \MM_1$ whenever $x' - x \in \MM_1^d$. Since $\vv(f(x)) = 1$ for $x \in \mdl X_f$, it follows that $\mdl X_f$ is a (quantifier-free) definable subset of $\MM^d$ that is $\beta$-invariant for every $\beta \geq 1$. Therefore, $\mdl X_f$ may be regarded as an object of $\mgVF^\diamond[d]$ equipped with the constant volume form $0$. The chart $q$ is fixed; the coefficients computed in (\ref{z:coeff:change}) below do not depend on it. The  \memph{motivic zeta function} of $f$ is then the power series
\begin{equation}\label{def:zeta}
Z_f(T) = \sum_{m \in \Z^+} \Big(\bm \eta \circ \bm h_m \circ \int^{\diamond} \Big)([\mdl X_f]) T^m \in \gmv[[\A]^{-1}] \dbra T.
\end{equation}
By Remark~\ref{dia:to:db:dag} and Lemma~\ref{inv:dir:com}, we have, in $\gmv[[\A]^{-1}]$,
\begin{equation}\label{z:coeff:change}
 \Big(\bm \eta \circ \bm h_m \circ \int^{\diamond} \Big)([\mdl X_f]) = \Theta(\tilde{\mdl X_f}[m]) = [\mdl X_{m}][\A]^{-md},
\end{equation}
where the second equality follows if  we change the variable $t$ to $t^{1/m}$ in (\ref{f:m:arc}) (see the discussion in \cite[\S~6.2]{hru:loe:lef} for  detail). More precisely, put $s = t^{1/m}$ and take $\beta = (1, \ldots, 1)$. The truncation $\mdl X_f[m;\beta]$ is taken modulo $s^{m+1}$, and for a point $x$ of $\mdl X_f$ that is rational over $\C \dpar{s}$ the condition $\rv(f(x)) = \rv(t)$ reads $f(x) = s^m \mod s^{m+1}$. The chart $q$ identifies the truncated arcs $\varphi \in X(\C[s]/(s^{m+1}))$ with $\varphi(0) = 0$ with the elements of $(s\C[s]/(s^{m+1}))^d$. Hence $\mdl X_f[m;\beta]$ is $\mdl X_m$, written in the variable $s$, the action of $\mu_m$ through $s \efun \zeta s$ is the usual one on truncated arcs, and the normalizing factor is $[\A]^{-m \Sigma \beta} = [\A]^{-md}$. This shows that the coefficients of $Z_f(T)$ may also be written as $[\mdl X_{m}][\A]^{-md}$, which is how they are usually defined.


\begin{rem}\label{rem:HK:fail}
One of the key steps of the main construction in \cite[\S~8]{hru:loe:lef} is to show (\ref{z:coeff:change}), see \cite[Proposition~8.3.1]{hru:loe:lef}, but with $\int^\diamond$  replaced by $\int$. This, and a few other statements in \cite[\S~8]{hru:loe:lef}, rely on a tensor decomposition of the ring $\ggk \mathrm{vol}\RV[*]$, as claimed in \cite[Proposition~10.10(2)]{hrushovski:kazhdan:integration:vf}. Unfortunately, this result does not hold, and neither do the surjectivity claims for the morphisms (8.1.4), (8.1.6) and (8.1.7) of \cite{hru:loe:lef}. Here is an explicit object outside their images.

Recall from \cite[\S~8.1]{hru:loe:lef} that the objects of $\mathrm{vol}\RV[n]$ are the objects $(X, f)$ of $\RV[n]$, and that a morphism $(X, f) \fun (X', f')$ is a definable bijection $h : X \fun X'$ such that $w(f(x)) = w(f'(h(x)))$ for all $x \in X$, where $w(y_1, \ldots, y_n) = \vrv(y_1) + \cdots + \vrv(y_n)$. By Notation~\ref{jcbG:op}, these are exactly the morphisms of $\mgRV[n]$ between objects with the constant volume form $0$. The category $\mathrm{vol}\RES[n]$ is the full subcategory of the objects in $\RES[n]$, and the morphisms in question are induced by $((W, g), \Delta) \efun (W \times \Delta^\sharp, g \times \id)$, where $\Delta \sub \Gamma^l$ is definable. So every class in their images is the class of a finite disjoint union $Y$ of such products. The categories in \cite[Proposition~10.10(2)]{hrushovski:kazhdan:integration:vf} have the same morphisms \cite[Definition~5.21]{hrushovski:kazhdan:integration:vf}.

Let $X = \set{u \in \RV \given 0 < \vrv(u) < 1}$ and $f(u) = u^2$, so that $(X, f)$ is an object of $\mathrm{vol}\RV[1]$ whose image in $\Gamma$ is doubly bounded. Suppose that $h : (X, f) \fun Y$ is an isomorphism in $\mathrm{vol}\RV[1]$, with $Y$ as above. Since the degree is $1$, every product in $Y$ is of one of two kinds: either $W \in \RES[1]$ and $l = 0$, and then the weight $w \circ (g \times \id)$ takes only finitely many values on it, all in $\Q$; or $W \in \RES[0]$, that is, $W$ is finite, and $\Delta \sub \Gamma$, and then the point $(p, y)$ has weight $\vrv(y)$. The points of the finite sets $W$ lie in the model $\puC$ of $\ACVF(0,0)$, so they all lie in $\RV(\C\dpar{t^{1/N}})$ for some $N \in \Z^+$. Each $\Delta$ is a finite union of points and intervals with endpoints in $\Q \cup \{\pm\infty\}$. Choose a multiple $m$ of $N$ such that $2/m$ is smaller than every positive endpoint of the sets $\Delta$ and every positive weight on the products of the first kind.

Since $\C\dpar{t^{1/m}}$ is functionally closed, $h$ restricts to a weight-preserving bijection between the points of $X$ and of $Y$ that are rational over $\C\dpar{t^{1/m}}$. The weights of the rational points of $X$ are the numbers $2i/m$ with $0 < i < m$; in particular, no rational point of $X$ has weight $1/m$. The rational point $\rv(t^{1/m})$ of $X$ has weight $2/m$, so its image lies in a product $W \times \Delta^\sharp$ of the second kind with $2/m \in \Delta$. By the choice of $m$, $(0, 2/m] \sub \Delta$, hence $1/m \in \Delta$, and for any $p \in W$ the point $(p, \rv(t^{1/m}))$ of $Y$ is rational over $\C\dpar{t^{1/m}}$ and has weight $1/m$. This is a contradiction. Therefore $[(X, f)]$ is not in the image of (8.1.4), (8.1.6) or (8.1.7) of \cite{hru:loe:lef}, nor in the image of the map in \cite[Proposition~10.10(2)]{hrushovski:kazhdan:integration:vf}.

In the category $\mgRV[1]$ of this paper, by contrast, the identity of $X$ is an isomorphism from $(X, f, 0)$ onto $(X, \id, \vrv)$, which lies in the image of ${\mu}\Psi$ (Proposition~\ref{red:D:iso}): the volume form absorbs the difference of the weights, whereas in $\mathrm{vol}\RV[1]$ the weight is tied to the structural map.
\end{rem}

\begin{nota}\label{dag:loc}
Let $\gmv[[\A]^{-1}][T]_{\dag}$ be the localization of $\gmv[[\A]^{-1}][T]$ by the multiplicative family generated by the elements $1 - [\A]^{a}T^b$, where $a \in \Z$ and $b \in \Z^+$; it may be regarded as a subring of  $\gmv[[\A]^{-1}] \dbra T$.
\end{nota}

It is known that $Z_f(T)$ belongs to $\gmv[[\A]^{-1}][T]_{\dag}$ and, letting ``$T$ go to infinity'' as described in \cite[\S~8.4]{hru:loe:lef}, we get a limit
\[
\mathscr S_f \coloneqq - \lim_{T \limplies \infty} Z_f(T) \in \gmv[[\A]^{-1}],
\]
which is understood as the \memph{motivic Milnor fiber} attached to $f$.

We recall how the limit is taken. Let $R$ be a commutative ring, $A \in R$ a unit, and $R[T]_{\dag}$ the localization of $R[T]$ by the multiplicative family generated by the elements $1 - A^{a}T^b$, where $a \in \Z$ and $b \in \Z^+$, regarded as a subring of $R \dbra T$. For $h = P / Q \in R[T]_{\dag}$ set $\deg h = \deg P - \deg Q$; this does not depend on the choice of $P$ and $Q$, since the leading coefficient of $Q$ is a unit. If $\deg h \leq 0$ then $\lim_{T \limplies \infty} h$ is $0$ if $\deg h < 0$, and is the quotient of the leading coefficients of $P$ and $Q$ if $\deg h = 0$. Equivalently, $\lim_{T \limplies \infty} h$ is the constant term of the expansion of $h$ in powers of $T^{-1}$. Hence the elements of degree at most $0$ form a subgroup of $R[T]_{\dag}$ on which $\lim_{T \limplies \infty}$ is additive.

The next lemma is \cite[Lemma~8.5.2]{hru:loe:lef}. The proof given there goes through a description of the Grothendieck semiring of bounded definable sets in $\Gamma$ from \cite[\S~2]{hru:kaz:value:ring} and through Hadamard products of rational series. We include a direct proof. Recall that $\Gamma$ is identified with $\Q$, so that $\Gamma(\mathbb{S}) = \Z$, and that a definable set in $\Gamma$ is a Boolean combination of sets of the form $g > 0$ or $g = 0$, where $g$ is an affine function with coefficients and constant term in $\Q$. For a set $I \sub \Gamma^n$ and $m \in \Z^+$ we write $I_m = I \cap (m^{-1}\Z)^n$. By an \memph{open simplex} of dimension $k$ in $\Gamma^n$ we mean a set of the form
\[
\set{ \textstyle\sum_{i=0}^k t_i v_i \given t_i > 0, \ \sum_{i=0}^k t_i = 1 },
\]
where the \memph{vertices} $v_0, \ldots, v_k \in \Q^n$ are affinely independent.

\begin{lem}\label{lat:zeta}
Let $I \sub \Gamma^n$ be a bounded definable set and $\sigma : I \fun \Gamma$ a function that is, on each piece of a finite definable partition of $I$, of the form $\gamma \efun \sum_i a_i \gamma_i + b$ with $a_i$, $b$ in $\Z$. Let
\[
Z(\sigma)(T) = \sum_{m \in \Z^+} \sum_{\gamma \in I_{m}} A^{-m \sigma(\gamma)} T^m \in R \dbra T.
\]
Then $Z(\sigma)(T)$ belongs to $R[T]_{\dag}$, $\deg Z(\sigma)(T) \leq 0$, and $\lim_{T \limplies \infty}Z(\sigma)(T) = -\chi(I)$.
\end{lem}
\begin{proof}
\emph{Step 1: bands.} Let $D \sub \Gamma^{n}$ be an open simplex with vertices $v_0, \ldots, v_k$ and let $f$, $g$ be affine functions with rational coefficients and constant terms such that $f < g$ on $D$. We claim that the band $\set{(x, y) \given x \in D, \ f(x) < y < g(x)}$ is a finite disjoint union of open simplices. Put $d_i = g(v_i) - f(v_i)$. Since $g - f$ is affine and positive on $D$, we have $d_i \geq 0$ for all $i$, and not all $d_i$ vanish. For $-1 \leq j \leq k$ let $h_j$ be the affine function on $D$ with $h_j(v_i) = f(v_i)$ for $i \leq j$ and $h_j(v_i) = g(v_i)$ for $i > j$. Thus $h_j(x) = f(x) + \sum_{i > j} t_i d_i$ for $x = \sum_i t_i v_i$, and hence
\[
g = h_{-1} \geq h_0 \geq \cdots \geq h_k = f, \qquad h_{j-1} - h_j = t_j d_j .
\]
On $D$ we have $h_{j-1} > h_j$ if $d_j > 0$ and $h_{j-1} = h_j$ otherwise. So the band is the disjoint union of the bands $B_j$ between $h_j$ and $h_{j-1}$ for the $j$ with $d_j > 0$, and of the graphs over $D$ of those $h_j$ that lie strictly between $f$ and $g$. The graph of an affine function over $D$ is an open simplex of dimension $k$. If $d_j > 0$ then $B_j$ is the open simplex of dimension $k+1$ with vertices $p_i = (v_i, h_j(v_i))$, $0 \leq i \leq k$, and $q = (v_j, g(v_j))$. Indeed, the $p_i$ are affinely independent and $q - p_j = (0, d_j)$; and a point $\sum_i s_i p_i + s q$ with positive coefficients of sum $1$ is $(x, h_j(x) + s d_j)$, where $x \in D$ has the coordinates $t_i = s_i$ for $i \neq j$ and $t_j = s_j + s$, so that $s$ ranges over $0 < s < t_j$, that is, $h_j(x) < y < h_{j-1}(x)$.

\emph{Step 2: triangulation.} We claim that if $g_1, \ldots, g_r$ are affine functions on $\Gamma^n$ with rational coefficients and constant terms and $D \sub \Gamma^n$ is a bounded Boolean combination of the sets $g_i > 0$, $g_i = 0$, then $D$ is a finite disjoint union of open simplices on each of which every $g_i$ has constant sign. This is proved by induction on $n$, the case $n = 0$ being trivial. For $n \geq 1$, write the coordinates as $(x, y)$ with $x \in \Gamma^{n-1}$ and choose $N \in \Z^+$ with $D \sub [-N, N]^n$. We may add $y - N$ and $y + N$ to the list. Write $g_i = c_i y + a_i(x)$ and, for $c_i \neq 0$, put $\varphi_i = - a_i / c_i$, so that the sign of $g_i$ is determined by the sign of $y - \varphi_i(x)$. Apply the inductive hypothesis to the box $[-N, N]^{n-1}$ and the functions $a_i$ with $c_i = 0$, the differences $\varphi_i - \varphi_{i'}$, and $x_j \pm N$. On each open simplex $D'$ so obtained the functions $\varphi_i$ are ordered: their distinct restrictions are $\psi_1 < \cdots < \psi_s$, with $\psi_1 \leq -N$ and $\psi_s \geq N$. The cylinder over $D'$ is the disjoint union of the graphs of the $\psi_j$, the bands between consecutive $\psi_j$, and two rays, and every $g_i$ has constant sign on each of these sets. So $D$ is a union of some of the graphs and bands. The graphs are open simplices, and the bands are disjoint unions of open simplices by Step~1.

\emph{Step 3: one simplex.} Let $D \sub \Gamma^n$ be an open simplex of dimension $k$ with vertices $w_0, \ldots, w_k$, and suppose that $\sigma(\gamma) = a \cdot \gamma + b$ on $D$, with $a \in \Z^n$ and $b \in \Z$. Put $L(x, m) = a \cdot x + b m$ for $(x, m) \in \Z^{n+1}$. The map $\gamma \efun (m\gamma, m)$ identifies the index set of the double sum $Z(\sigma \rest D)(T)$ with $C = \set{(x, m) \in \Z^n \times \Z^+ \given x/m \in D}$, and
\[
Z(\sigma \rest D)(T) = \sum_{(x,m) \in C} A^{-L(x,m)} T^m .
\]
Choose $m_j \in \Z^+$ with $m_j w_j \in \Z^n$ and put $u_j = (m_j w_j, m_j) \in \Z^{n+1}$. The $u_j$ are linearly independent, and $C$ is the set of the points $\sum_j \lambda_j u_j$ in $\Z^{n+1}$ with all $\lambda_j \in \Q$ positive: if $x/m = \sum_j t_j w_j$ then $(x, m) = \sum_j (m t_j / m_j) u_j$, and conversely $m = \sum_j \lambda_j m_j$ is a positive integer. Each $\lambda_j > 0$ is uniquely of the form $n_j + \mu_j$ with $n_j \in \N$ and $0 < \mu_j \leq 1$. So every point of $C$ is uniquely of the form $p + \sum_j n_j u_j$ with $n_j \in \N$ and $p$ in the finite set $\Pi$ of the points $\sum_j \mu_j u_j$ in $\Z^{n+1}$ with $0 < \mu_j \leq 1$, and every such sum is in $C$. Since $L$ and $(x, m) \efun m$ are additive, summing the geometric series gives
\[
Z(\sigma \rest D)(T) = \frac{\sum_{p \in \Pi} A^{-L(p)}\, T^{m(p)}}{\prod_{j=0}^{k} \bigl(1 - A^{-L(u_j)}\, T^{m_j}\bigr)},
\]
where $m(p)$ stands for the last coordinate of $p$. The denominator has degree $M = \sum_j m_j$ and leading coefficient $(-1)^{k+1} A^{-L(\sum_j u_j)}$. For $p = \sum_j \mu_j u_j \in \Pi$ we have $m(p) = \sum_j \mu_j m_j \leq M$, with equality only for $p = \sum_j u_j$; so the numerator has degree $M$ and leading coefficient $A^{-L(\sum_j u_j)}$. Therefore $Z(\sigma \rest D)(T)$ belongs to $R[T]_{\dag}$, has degree $0$, and $\lim_{T \limplies \infty}Z(\sigma \rest D)(T) = (-1)^{k+1}$.

\emph{Conclusion.} Apply Step~2 to $I$ and to the affine functions that define $I$ and the pieces of the given partition. We get a partition of $I$ into open simplices $D$ on each of which $\sigma$ is of the form considered in Step~3. The simplices are obtained as graphs and bands over simplices, so they are cells, and $\chi(D) = (-1)^{\dim D}$. Since $Z(\sigma)(T) = \sum_D Z(\sigma \rest D)(T)$, the additivity of the limit and of $\chi$ gives $\lim_{T \limplies \infty}Z(\sigma)(T) = \sum_D (-1)^{\dim D + 1} = -\chi(I)$.
\end{proof}

\begin{rem}
Recall the homomorphism $\bb E^{\diamond}$ from (\ref{edb:eb:com}).
By Lemma~\ref{coh:decom} and Notation~\ref{bipo:gdv}, we can construct the homomorphism
\[
\Theta \circ \bb E^{\diamond} : \ggk \mgRV^{\db}[*] \fun \gmv
\]
simply using the expression in (\ref{twist:C}) with respect to  monomial decompositions.
\end{rem}

\begin{thm}\label{direct:mil}
$\mathscr S_f = (\Theta \circ \bb E^{\diamond} \circ \int^{\diamond})([\mdl X_f])$.
\end{thm}
\begin{proof}
Let $[\bm U] = [(U, f, l)] \in \ggk^\natural \mgRV^{\db}[*]$ with $l \in \Z$  and consider the zeta function
\begin{equation}\label{U:zeta}
Z([\bm U])(T) = \sum_{m \in \Z^+} (\bm \eta \circ \bm h_m)([\bm U]) T^m \in \gmv[[\A]^{-1}]\dbra T.
\end{equation}
If $[\bm U] / J_{\mathrm{var}} = \int^{\diamond}([\mdl X_f])$ then this is $Z_f(T)$. Thus it is enough to show that  $Z([\bm U])(T)$ belongs to $\gmv[[\A]^{-1}][T]_{\dag}$, and $\lim_{T \limplies \infty} Z([\bm U])(T)$ exists and equals $- (\Theta \circ \bb E^{\diamond})([\bm U])$.

Both sides of the desired equality are additive. So, taking a monomial decomposition as in (\ref{hm:ass}), we may assume that $h_M(U) = V \times I^\sharp$ with $M \in \GL_n(\Z)$, $\vrv(V) = \{\beta\}$ and $I \sub \Gamma^\ell$ bounded (since $\bm U$ is doubly bounded), and that the function $l_f: U \fun \Gamma$ (Notation~\ref{comp:form}) $\vrv$-contracts to a function $\sigma : D = \vrv(U) \fun \Gamma$.  Set $v = \Theta([V]) [\G_m]^\ell$. Then, by (\ref{twist:C}),
\[
(\Theta \circ \bb E^{\diamond})([\bm U]) = \chi(I) v.
\]
On the other hand, by Lemma~\ref{muhat:section},
\[
(\bm \eta \circ \bm h_m)([\bm U]) = v \sum_{\gamma \in D_m} [\A]^{-m \sigma(\gamma)}, \qquad D_m = D \cap (m^{-1}\Z)^n.
\]
Let $N$ be the least positive integer such that $N \beta$ is a tuple in $\Z$. Since $M$ preserves the lattice $(m^{-1}\Z)^n$ and $M D = \{\beta\} \times I$, the set $D_m$ is empty unless $N$ divides $m$, and if $m = N m'$ then $\gamma \efun N \delta$, where $M \gamma = (\beta, \delta)$, is a bijection between $D_m$ and $(N I)_{m'} = N I \cap (m'^{-1}\Z)^\ell$. The function $\sigma$ is the $\vrv$-contraction of the definable function $u \efun \rv(t)^l \prod_j f_j(u)$ on $U$. So, by Remark~\ref{GLZ:char}, it is piecewise $\Z$-affine with constant terms in $\Z$. It follows that the function $\sigma'$ on $N I$ given by $\sigma'(N \delta) = N \sigma(\gamma)$ is piecewise of the form $\delta' \efun \sum_i a_i \delta'_i + b$ with $a_i$, $b$ in $\Z$. So Lemma~\ref{lat:zeta} applies to $\sigma'$, the ring $\gmv[[\A]^{-1}]$ and its unit $[\A]$; writing $S$ for the variable, we get
\[
Z(\sigma')(S) \in \gmv[[\A]^{-1}][S]_{\dag} \quad \text{and} \quad \lim_{S \limplies \infty}Z(\sigma')(S) = -\chi(N I) = - \chi(I).
\]
Since $m \sigma(\gamma) = m' \sigma'(N \delta)$ for $m = N m'$, the formula for $(\bm \eta \circ \bm h_m)([\bm U])$ above gives
\[
Z([\bm U])(T) = v \sum_{m' \in \Z^+} \sum_{\delta' \in (N I)_{m'}} [\A]^{-m' \sigma'(\delta')} T^{N m'} = v Z(\sigma')(T^N),
\]
and we conclude that $Z([\bm U])(T)$ belongs to $\gmv[[\A]^{-1}][T]_{\dag}$ and $\lim_{T \limplies \infty}Z([\bm U])(T) = - \chi(I) v$.
\end{proof}


\begin{rem}
The construction above no longer needs to go through this additional localization process ``$\loc$'' employed throughout \cite[\S~8]{hru:loe:lef}.
\end{rem}


\begin{rem}
\label{sec:hm:bounded}
The homomorphisms $\bm \eta \circ \bm h_m$ are defined on doubly bounded objects, and Theorem~\ref{direct:mil} does not extend to bounded objects in the obvious way. Consider the object $\bm U = (\RV^{\circ\circ}, \id, 0)$, which is bounded but not doubly bounded. The expression (\ref{hm:ass}) still makes sense for it, as a series in $T^{-1}$: the weight $\omega_f = \vrv$ takes each value $i/m$, $i \in \Z^+$, on a set of class $[\A] - 1$, and we get
\[
([\A] - 1) \sum_{i > 0} T^{-i} = \frac{[\A] - 1}{T-1}.
\]
Substituting $[\A]$ for $T$ in the partial sums gives $([\A] - 1) \sum_{0 < i \leq M} [\A]^{-i} = 1 - [\A]^{-M}$, so the value to assign to $\bm U$ is $1$, for every $m$; this is the value of the right-hand side at $T = [\A]$ once $1 - [\A]^{-1}$ is inverted, as in \cite[\S~8.2]{hru:loe:lef}. The zeta function of $\bm U$ is then $\sum_{m \in \Z^+} T^m = T / (1 - T)$ and $- \lim_{T \limplies \infty}$ of it is $1$, whereas $(\Theta \circ \bb E_b)([\bm U]) = 0$. Such a failure explains the need to introduce an interpolant, such as the composite map $- \lim \circ Z \circ \int^\diamond$, in (\ref{diag-diamond-zeta-function}), for which Theorem \ref{direct:mil} may be formulated and proven.

We do not pursue the extension of $\bm h_m$ and of the zeta function to bounded objects. The expression (\ref{hm:ass}) is not summable for every bounded object: for $(\RV^{\circ\circ}, \id, -\vrv)$ the weight $\omega_f$ vanishes identically and the expression is $([\A] - 1) \sum_{i > 0} 1$.
\end{rem}


\begin{rem}\label{rem-comp-milnor-fib}
The expression $(\Theta \circ \bb E^{\diamond})([\bm U])$, as an element in $\gmv[[\A]^{-1}]$, does not actually involve  $[\A]^{-1}$. Still, since the coefficients of $Z([\bm U])(T)$ do involve $[\A]^{-1}$ and the natural homomorphism $\gmv \fun \gmv[[\A]^{-1}]$ is not injective (by \cite{borisov:zero:divisor}, $[\A]$ is a zero divisor in $\ggk \var_{\C}$, and equipping varieties with the trivial action embeds $\ggk \var_{\C}$ into $\gmv$, since forgetting the action is a left inverse), we cannot really take  the motivic Milnor fiber $\mathscr S_f$ of $f$ in $\gmv$, at least not if $\mathscr S_f$ is viewed as something obtained through  $Z_f(T)$. It is this point of view that forces us to work with an integral whose target only involves doubly bounded sets in $\RV$, namely $\int^{\diamond}$, instead of $\int$, so as to facilitate the computation of the coefficients of $Z_f(T)$, and consequently with the nonarchimedean Milnor fiber $\mdl X_f$, which is proper invariant, instead of, perhaps, the more obvious definable set
\[
\mdl X_t = \set{x \in X(\MM) \given f(x) = t},
\]
which is not proper invariant. This set $\mdl X_t$ may be understood as containing the $\puC$-rational points of  the analytic Milnor fiber introduced in \cite{nicaise:sebag:motivic:serre} (see Remark~\ref{rem-comp-milnor-fib-bis} below for more) and does play a role in \cite{hru:loe:lef} (but not in this paper).


On the other hand, in light of Theorem~\ref{direct:mil}, we can forego the zeta function point of view and recover $\mathscr S_f$ directly as $(\Theta \circ \bb E^{\diamond} \circ \int^{\diamond})([\mdl X_f])$. In that case there is truly no need to invert $[\A]$. In fact, we can also recover $\mathscr S_f$  directly as  $\Vol([\mdl X_f])$ (Notation~\ref{mov:vol}), the result is the same because the left half of the diagram (\ref{diag-diamond-interpol}) commutes (Corollary~\ref{db:to:novol}).
\end{rem}


\begin{rem}\label{rem-comp-milnor-fib-bis}
For $t' \in \rv(t)^\sharp$ write $\mdl X_{t'} = \set{x \in X(\MM) \given f(x) = t'}$, so that $\mdl X_f = \bigcup_{t'\in \rv(t)^\sharp} \mdl X_{t'}$. All $t' \in \rv(t)^\sharp$ have the same type over $\C \cup \RV$. Indeed, by quantifier elimination, a formula in one $\VF$-variable $x$ with parameters in $\C \cup \RV$ is a Boolean combination of equalities $P(x) = 0$ and of conditions on $\rv(P_1(x)), \ldots, \rv(P_r(x))$ and the parameters in $\RV$, where $P, P_i \in \C[x]$; and if $a_j x^j$ is the term of lowest degree of a nonzero $P \in \C[x]$ then $\rv(P(t')) = \rv(a_j) \rv(t)^j$. (On the points of $\puC$ this is the observation that there is an automorphism $\sigma$ of $\puC$ over $\C$ with $\sigma(t') = t$ that fixes $\RV(\puC)$ pointwise.)

Now let $\mathbb S_0 \sub \mathbb S$ be the substructure generated by $\C$ and $t$. It is $\VF$-generated and $\Gamma(\mathbb S_0) = \Z$, so the results of the previous sections hold over $\mathbb S_0$, and $\mdl X_t$ is definable over $\mathbb S_0$. Since $\bb L$ is surjective on isomorphism classes over $\mathbb S_0$ (see the discussion before Theorem~\ref{main:prop:HK}), there are an object $\bm U$ of $\RV[*]$ and a bijection $F_t : \mdl X_t \fun \bb L \bm U$ that are definable over $\mathbb S_0$, that is, defined by formulas over $\C$ in which $t$ is the only other parameter. Replace $t$ by $t' \in \rv(t)^\sharp$ in these formulas. Since $t'$ and $t$ have the same type over $\C \cup \RV$, the object so defined is again $\bm U$, and we get a bijection $F_{t'} : \mdl X_{t'} \fun \bb L \bm U$. Therefore
\[
x \efun (f(x), F_{f(x)}(x))
\]
is a bijection from $\mdl X_f$ onto $\rv(t)^\sharp \times \bb L \bm U$. It is definable over $\C$ and $\rv(t)$, hence over $\mathbb S_0$ and over $\mathbb S$, and so is $F_t$. As $\rv(t)^\sharp \times \bb L \bm U$ is in definable bijection with the lift of the object $\{\rv(t)\} \times \bm U$, whose class is $[1] [\bm U]$, we conclude that
\begin{equation*}
\int [\mdl X_f] = \int [\rv(t)^\sharp] \int [\mdl X_t] =  [1]\int [\mdl X_t].
\end{equation*}
Since $\bb E_b([1]) = 1 $, this and Remark \ref{rem-comp-milnor-fib} show that
\[
\mathscr S_f  = \Vol([\mdl X_f]) = \Vol([\mdl X_t]).
\]


This is first shown in \cite{nicaise_ps_tropical}, but without taking the $\hat \mu$-actions into account, and also in \cite{Nic:Pay:trop:fub, forey_virtual_2017}. The arguments there all rely on resolution of singularities as well as other algebro-geometric machinery.
\end{rem}


\providecommand{\bysame}{\leavevmode\hbox to3em{\hrulefill}\thinspace}
\providecommand{\MR}{\relax\ifhmode\unskip\space\fi MR }
% \MRhref is called by the amsart/book/proc definition of \MR.
\providecommand{\MRhref}[2]{%
  \href{http://www.ams.org/mathscinet-getitem?mr=#1}{#2}
}
\providecommand{\href}[2]{#2}
\begin{thebibliography}{1}

\bibitem{borisov:zero:divisor}
Lev A. Borisov, {The class of the affine line is a zero divisor in the
  {G}rothendieck ring}, J. Algebraic Geom. \textbf{27} (2018), no.~2, 203--209.

\bibitem{denef_loser_motivic_igusa}
Jan~Denef and Fran\c cois~Loeser,
\newblock Motivic {{Igusa}} zeta functions,
\newblock {\em Journal of Algebraic Geometry}, \textbf{7} (1998), no.~3, 505--537.

\bibitem{denefloeser:arc}
\bysame, \emph{{Geometry on arc spaces of algebraic varieties}}, {3rd European
  congress of mathematics (ECM), Barcelona, Spain, July 10--14, 2000. Volume
  I}, Basel: Birkh\"auser, 2001, pp.~327--348.

\bibitem{den:loe:2002}
\bysame, \emph{Lefschetz numbers of iterates of the monodromy and truncated
  arcs}, Topology \textbf{41} (2002), no.~5, 1031--1040.

\bibitem{dries:1998}
Lou van~den Dries, \emph{Tame topology and \omin-minimal structures}, LMS Lecture Note
  Series, vol. 248, Cambridge University Press, Cambridge, UK, 1998.

\bibitem{engler:prestel:2005}
Antonio~J. Engler and Alexander Prestel, \emph{Valued fields}, Springer-Verlag,
  Berlin, 2005.

\bibitem{FY:Mot-Int-Mil-Fib}
Goulwen Fichou and Yimu Yin, {Motivic integration and Milnor fiber}, {\em Journal of the European Mathematical Society} \textbf{24} (2022), no.~5, 1617--1678.


\bibitem{forey_virtual_2017}
Arthur Forey, {Virtual rigid motives of semi-algebraic sets}, {\em Selecta Mathematica (New Series)}, \textbf{25} (2019),  no.~1, Art. 6.


\bibitem{FY:Fubini}
Arthur Forey and Yimu Yin, {Fubini theorems for motivic volume}, manuscript in the repository \url{https://github.com/daviddhilbert/Fubini-theorems}, file \texttt{paper/fubini-theorems.tex}.

\bibitem{hrushovski:kazhdan:integration:vf}
Ehud Hrushovski and David Kazhdan, {Integration in valued fields}, in
  {\em Algebraic geometry and number theory}, Progr. Math., vol. 253, Birkh\"{a}user,
  Boston, MA, 2006, pp.~261--405.

\bibitem{hru:kaz:value:ring}
Ehud Hrushovski and David Kazhdan, {The value ring of geometric motivic integration, and the
  {I}wahori {H}ecke algebra of {${\rm SL}_2$}}, Geom. Funct. Anal. \textbf{17} (2008),
  no.~6, 1924--1967, With an appendix by Nir Avni.

\bibitem{hru:loe:lef}
Ehud Hrushovski and Fran{\c{c}}ois Loeser, {Monodromy and the {L}efschetz
  fixed point formula}, {\em Annales Scientifiques de l'\'Ecole Normale Sup\'erieure} (4) \textbf{48} (2015),
  no.~2, 313--349.


\bibitem{Nic:Pay:trop:fub}
Johannes Nicaise and Sam Payne, {A tropical motivic {F}ubini theorem with
  applications to {D}onaldson-{T}homas theory}, {\em Duke Mathematical Journal} \textbf{168} (2019), no.~10, 1843–1886.

\bibitem{nicaise_ps_tropical}
Johannes Nicaise, Sam Payne, and Franziska Schroeter,
{Tropical refined curve counting via motivic integration}, {\em Geometry and Topology} \textbf{22} (2018),  no.~6, 3175--3234.

\bibitem{nicaise:sebag:motivic:serre}
Johannes Nicaise and Julien Sebag,
{Motivic Serre invariants, ramification, and the analytic Milnor fiber}, {\em Inventiones Mathematicae} \textbf{168}  (2007), no.~1, 133--173.

\bibitem{NS-deg-stab-rat}
Johannes Nicaise and Evgeny Shinder,
\newblock The motivic nearby fiber and degeneration of stable rationality,
\newblock {\em Inventiones Mathematicae} \textbf{217} (2019), no.~2, 377--413.


\bibitem{thuong-int-conj}
L{\^e}~Quy Thuong,
\newblock Proofs of the integral identity conjecture over algebraically closed
  fields,
\newblock {\em Duke Mathematical Journal} \textbf{164} (2015), no.~1, 157--194.

\bibitem{Yin:QE:ACVF:min}
Yimu Yin, {Quantifier elimination and minimality conditions in
  algebraically closed valued fields}, arXiv:1006.1393v1, 2010.


\bibitem{Yin:special:trans}
\bysame, {Special transformations in algebraically closed valued fields},
  {\em Annals of Pure and Applied Logic} \textbf{161} (2010), no.~12, 1541--1564.

\bibitem{Yin:int:acvf}
\bysame, {Integration in algebraically closed valued fields}, {\em Annals of
  Pure and Applied Logic} \textbf{162} (2011), no.~5, 384--408.

\bibitem{Yin:int:expan:acvf}
\bysame, {Integration in algebraically closed valued fields with
  sections}, {\em Annals of Pure and Applied Logic} \textbf{164} (2013), no.~1,
  1--29.

\bibitem{yin:hk:part:3}
\bysame, {{F}ourier transform of the additive group in algebraically
  closed valued fields}, {\em Selecta Mathematica (New Series)} \textbf{20} (2014), no.~4,
  1111--1157.

\bibitem{Yin:tcon:I}
\bysame, {Generalized {E}uler characteristic in power-bounded {$T$}-convex
  valued fields}, {\em Compositio Mathematica} \textbf{153} (2017), no.~12,
  2591--2642.

\bibitem{Yin:tcon:1.5}
\bysame, \emph{Integration in power-bounded {$T$}-convex valued fields},
  (2018), arXiv:1808.07129 [math.LO].


\bibitem{reichstein:youssin:birational}
Zinovy Reichstein and Boris Youssin,
{A birational invariant for algebraic group actions},
{\em Pacific Journal of Mathematics} \textbf{204} (2002), no.~1,
223--246, DOI \href{https://doi.org/10.2140/pjm.2002.204.223}{10.2140/pjm.2002.204.223}.

\end{thebibliography}


\end{document} 
